\documentclass[11pt,oneside,openany]{book}

% BEGIN preamble
\usepackage[T1]{fontenc}
\usepackage{lmodern,amsmath,amssymb,amsthm,mathtools}
\usepackage[a4paper,margin=25mm,headheight=15pt]{geometry}
\usepackage{microtype,enumitem,xcolor,longtable,booktabs,array,needspace}
\usepackage[colorlinks=true,linkcolor=blue!40!black,urlcolor=blue!40!black]{hyperref}
\usepackage[nameinlink,noabbrev]{cleveref}
\usepackage{fancyhdr}
\hypersetup{pdftitle={Measure, Integration and Real Analysis: Unified Notes},pdfauthor={MathGenome preparation series}}
\pagestyle{fancy}\fancyhf{}
\fancyhead[L]{\small MathGenome: Measure, Integration and Real Analysis}
\fancyhead[R]{\small\thepage}
\renewcommand{\headrulewidth}{0.3pt}
\setlength{\parindent}{0pt}\setlength{\parskip}{5pt}
\setlength{\emergencystretch}{3em}\allowdisplaybreaks[2]\raggedbottom
\setcounter{tocdepth}{1}
\newtheorem{theorem}{Theorem}[chapter]
\newtheorem{lemma}[theorem]{Lemma}
\newtheorem{proposition}[theorem]{Proposition}
\newtheorem{corollary}[theorem]{Corollary}
\theoremstyle{definition}
\newtheorem{definition}[theorem]{Definition}
\newtheorem{construction}[theorem]{Construction}
\crefname{construction}{construction}{constructions}\Crefname{construction}{Construction}{Constructions}
\newtheorem{example}[theorem]{Example}
\theoremstyle{remark}
\newtheorem{remark}[theorem]{Remark}
\newcounter{foundation}\renewcommand{\thefoundation}{F\arabic{foundation}}
\crefname{foundation}{foundation}{foundations}\Crefname{foundation}{Foundation}{Foundations}
\newcommand{\R}{\mathbb R}\newcommand{\C}{\mathbb C}
\newcommand{\N}{\mathbb N}\newcommand{\Z}{\mathbb Z}\newcommand{\Q}{\mathbb Q}
\newcommand{\eps}{\varepsilon}\newcommand{\ee}{\varepsilon}
\newcommand{\abs}[1]{\lvert#1\rvert}\newcommand{\norm}[1]{\lVert#1\rVert}
\newcommand{\ip}[2]{\langle#1,#2\rangle}
\newcommand{\BB}{\mathcal B}\newcommand{\SSS}{\mathcal S}
\newcommand{\calA}{\mathcal A}\newcommand{\LL}{\mathcal L}
\newcommand{\one}{\mathbf1}\newcommand{\ind}[1]{\mathbf1_{#1}}
\newcommand{\aev}{\text{a.e.}}\newcommand{\dd}{\,d}
\newcommand{\Eplus}{[0,\infty]}
\DeclareMathOperator{\id}{id}\DeclareMathOperator{\supp}{supp}
\DeclareMathOperator{\diam}{diam}\DeclareMathOperator{\dist}{dist}
\DeclareMathOperator{\esssup}{ess\,sup}\DeclareMathOperator{\Var}{Var}
\DeclareMathOperator{\sgn}{sgn}\DeclareMathOperator{\Lip}{Lip}
\newcommand{\family}[1]{}
\newcommand{\source}[1]{\par\begingroup\small\sffamily\color{black!65}Source comparison: #1\par\endgroup}

% END preamble

\begin{document}
\frontmatter

% BEGIN frontmatter
\hypersetup{pageanchor=false}
\begin{titlepage}
\centering\vspace*{24mm}
{\Huge\bfseries Measure, Integration\par and Real Analysis\par}
\vspace{12mm}{\Large Unified Notes for MathGenome\par}
\vspace{10mm}{\large D11: twenty-three mathematical topic families\par}
\vfill
An original synthesis of the adopted results from\par
Axler; Royden--Fitzpatrick; Folland; Stein--Shakarchi; Cohn\par
\vspace{8mm}English manuscript prepared from 2 November 2026\par
\end{titlepage}
\hypersetup{pageanchor=true}
\chapter*{Purpose and reading conventions}
\addcontentsline{toc}{chapter}{Purpose and reading conventions}
This volume develops the selected measure and integration core, its principal
real-variable applications, and explicitly selected representation and geometric
bridges. It synthesizes mathematical content rather than concatenating textbook
chapters. Different notations and genuinely different hypotheses are reconciled
by statements and proofs, not by calling unlike objects by the same name.
The adopted scope is the twenty-three-family framework in the accompanying
result ledger. It does not purport to include every exercise or every advanced
topic from the five sources. General operator theory, probability and ergodic
theory, Fourier and distribution theory, and advanced geometric measure theory
remain subjects for separate volumes.

Each independently adopted assertion has a stable label. Statements may be
proved by any valid equivalent route, and established earlier MathGenome
mathematics may be used with its actual hypotheses. Internal steps of one
particular textbook proof are not additional independent acceptance conditions.
This permission does not remove a theorem's quantifiers, hypotheses, mathematical
objects, or conclusions. A finite-dimensional special case does not replace
a stated general result, and a representation does not replace the bridge
identifying it with the original object.

Source references give comparison locations; the exposition and arguments are
newly written here. Physical PDF pages are distinguished from printed book
pages in the source register. A manuscript proof is not a Lean certificate.
The later formalization must still pass the normal specification, compilation,
kernel and independent fidelity checks.

% END frontmatter

\tableofcontents

% BEGIN foundations
\chapter*{Relative foundations and canonical conventions}
\addcontentsline{toc}{chapter}{Relative foundations and canonical conventions}
\begin{description}[style=nextline,leftmargin=13mm]
\item[F\_SET]\refstepcounter{foundation}\label{base:d11-set}Classical logic, sets, functions, countable
families and quotients, choice, well-ordering and Zorn's lemma are background.
Choice-dependent examples are stated as such.
\item[F\_SCALAR]\refstepcounter{foundation}\label{base:d11-scalar}The complete ordered real field, the
complex field, elementary scalar limits, series and finite sums are background.
No measure construction or integral convergence theorem is assumed here.
\item[F\_LINEAR]\refstepcounter{foundation}\label{base:d11-linear}Finite-dimensional algebra,
determinants, Euclidean inner products and the vocabulary of normed spaces
are prior mathematics. The infinite-dimensional projection and representation
results actually used in this volume are proved here.
\item[F\_CALCULUS]\refstepcounter{foundation}\label{base:d11-calculus}Established one-variable
calculus and the chain-rule and inverse-function machinery of the earlier
single- and multivariable MathGenome components are available. A classical
Riemann integral is distinguished from the Lebesgue integral constructed here.
The nonlinear Lebesgue Jacobian formula is a theorem of this volume, not
an imported foundation.
\item[F\_TOPO]\refstepcounter{foundation}\label{base:d11-topology}The published point-set topology
component in MathGenome v0.5.0 supplies compactness, metric completeness,
locally compact Hausdorff theory, Urysohn--Tietze and the relevant partition
and product compactness theorems. Their exact separation, compactness and
local-finiteness hypotheses are retained. Compactness alone does not include
Hausdorffness, and a bounded closed ball need not be compact.
\end{description}

The notes are self-contained relative to these declared prior foundations,
not a reconstruction of all set theory or the real numbers. Every new
measure-theoretic or representation result adopted below is proved, or reduced
to explicitly identified earlier results. References between chapters express
mathematical prerequisites; chapter order is a navigation convention, not a
restriction against first proving an independently stated later lemma.

Positive measures take values in $[0,\infty]$. Empty spaces and zero measures
are allowed unless excluded. Completeness of a measure means measurability of
all subsets of null sets. It has no automatic connection to metric completeness.
Euclidean dimensions are positive integers unless another range is stated.
$\lambda_d$ denotes completed Euclidean Lebesgue measure, and
$\lambda_d^B$ its Borel restriction. Abstract products and their completions
are distinguished. Scalar integrability means $\int|f|<\infty$; equalities of
linear integrals are used only with their integrability hypotheses. A library's
totalized integral on nonintegrable functions is not this convention.

$L^p$ uses the ordinary almost-everywhere quotient. Cohn's local-null
$L^\infty$ is separately named and is compared by a proved bridge. A Radon
measure in this volume is Borel, finite on compact sets, outer regular on Borel
sets and inner regular on open sets; full compact inner regularity on all Borel
sets is stated separately. All density uniqueness is almost everywhere.
Stieltjes generators are right-continuous and use $(a,b]$. Hausdorff measures
use coefficient one in diameter covers, with counting at exponent zero.
Daniell--Stone uniqueness is asserted on its precise generated sigma-ring.

% END foundations

\mainmatter

% BEGIN chapters/01
\chapter{Measurable Spaces and Measures}\label{ch:d11-01}

A measure begins with a choice of sets that may be measured. No topology is
part of that choice unless it is explicitly supplied. Throughout, sums of
nonnegative extended real numbers mean suprema of their finite partial sums;
in particular, rearranging such a sum does not change it. We use
$0\cdot\infty=0$ and allow empty spaces and zero measures.

\section{The measurable structure}

\begin{definition}[Algebras and measurable spaces]
\family{D11.01.measurable_space}\label{d11:01:measurable-space}
An algebra $\calA$ on a set $X$ contains $X$ and is closed under complements
relative to $X$ and finite unions. A sigma-algebra $\Sigma$ is an algebra
closed under countable unions. The pair $(X,\Sigma)$ is a measurable space;
its measurable sets are the members of $\Sigma$.
\end{definition}

De Morgan's identities give closure under the corresponding intersections.
In particular $\varnothing\in\calA$, and differences of members of an
algebra belong to it. For a sequence of measurable sets we write
$\limsup E_n=\bigcap_N\bigcup_{n\geq N}E_n$ and
$\liminf E_n=\bigcup_N\bigcap_{n\geq N}E_n$.

\begin{proposition}[Generation]
\family{D11.01.generated_sigma}\label{d11:01:generated-sigma}
For every $\mathcal C\subseteq\mathcal P(X)$ there is a smallest sigma-algebra
$\sigma(\mathcal C)$ containing $\mathcal C$. It is the intersection of all
sigma-algebras containing $\mathcal C$. If $\mathcal C\subseteq\mathcal D$,
then $\sigma(\mathcal C)\subseteq\sigma(\mathcal D)$, and
$\sigma(\sigma(\mathcal C))=\sigma(\mathcal C)$.
\end{proposition}
\begin{proof}
The class being intersected is nonempty because it contains $\mathcal P(X)$.
An intersection of sigma-algebras contains $X$ and is closed under complements
and countable unions: each operation can be checked in every member of the
class. Thus the intersection is itself a sigma-algebra. Its definition gives
both minimality and the stated monotonicity. Applying minimality to a
sigma-algebra already containing its generators gives idempotence.
\end{proof}

\begin{proposition}[Trace structure]
\family{D11.01.trace_sigma}\label{d11:01:trace-sigma}
If $E\in\Sigma$, then
$\Sigma|_E=\{E\cap A:A\in\Sigma\}$ is a sigma-algebra on $E$; its members
are exactly the $\Sigma$-measurable subsets of $E$.
\end{proposition}
\begin{proof}
It contains $E=E\cap X$. The complement in $E$ of $E\cap A$ is
$E\cap A^c$, and a countable union of traces is the trace of the corresponding
union. Since $E$ is measurable, every trace is a measurable subset of $E$.
Conversely a measurable $B\subseteq E$ satisfies $B=E\cap B$.
\end{proof}

\section{Positive measures}

\begin{definition}[Measure]
\family{D11.01.positive_measure}\label{d11:01:positive-measure}
A positive measure on $(X,\Sigma)$ is a function
$\mu:\Sigma\to\Eplus$ such that $\mu(\varnothing)=0$ and
\[
 \mu\Bigl(\bigcup_{n=1}^{\infty}E_n\Bigr)=\sum_{n=1}^{\infty}\mu(E_n)
\]
for every pairwise disjoint measurable sequence. A finitely additive measure
on an algebra satisfies the same requirement for finite disjoint families.
\end{definition}

\begin{theorem}[Elementary bounds]
\family{D11.01.measure_bounds}\label{d11:01:measure-bounds}
If $A\subseteq B$ are measurable, then $\mu(A)\leq\mu(B)$. If moreover
$\mu(A)<\infty$, then $\mu(B\setminus A)=\mu(B)-\mu(A)$. For every
measurable sequence,
$\mu(\bigcup_n E_n)\leq\sum_n\mu(E_n)$.
\end{theorem}
\begin{proof}
Countable additivity implies finite additivity by appending empty sets.
The disjoint decomposition $B=A\sqcup(B\setminus A)$ gives
$\mu(B)=\mu(A)+\mu(B\setminus A)$, proving monotonicity and, when the
subtracted term is finite, the difference formula. Put
$D_1=E_1$ and $D_n=E_n\setminus\bigcup_{j<n}E_j$. These sets are disjoint,
their union is $\bigcup_nE_n$, and $D_n\subseteq E_n$. Therefore
$\mu(\bigcup_nE_n)=\sum_n\mu(D_n)\leq\sum_n\mu(E_n)$.
\end{proof}

\begin{theorem}[Continuity of a measure]
\family{D11.01.continuity_sets}\label{d11:01:continuity-sets}
If $E_n\uparrow E$, then $\mu(E_n)\uparrow\mu(E)$. If $E_n\downarrow E$
and $\mu(E_N)<\infty$ for some $N$, then $\mu(E_n)\downarrow\mu(E)$.
\end{theorem}
\begin{proof}
In the increasing case use the disjoint sets $D_1=E_1$,
$D_n=E_n\setminus E_{n-1}$: the finite partial sums of
$\sum_n\mu(D_n)$ are $\mu(E_n)$. In the decreasing case discard the first
$N-1$ terms. The sets $E_N\setminus E_n$ increase to $E_N\setminus E$.
The first part and the finite difference formula give
$\mu(E_N)-\mu(E_n)\to\mu(E_N)-\mu(E)$, proving the assertion.
\end{proof}

\begin{theorem}[Continuity criteria]
\family{D11.01.countable_additivity_criterion}\label{d11:01:countable-additivity-criterion}
For a finitely additive $p:\Sigma\to\Eplus$, countable additivity is
equivalent to continuity from below. Continuity at the empty set from above
is also sufficient, and is necessary if $p(X)<\infty$.
\end{theorem}
\begin{proof}
Necessity follows from the preceding theorem. If $B_j$ are disjoint, apply
continuity from below to $S_n=\bigcup_{j\leq n}B_j$ to obtain
$p(\bigcup_jB_j)=\lim_n\sum_{j\leq n}p(B_j)$.
For the other sufficient condition, let $B=\bigcup_jB_j$ and
$T_n=\bigcup_{j>n}B_j\downarrow\varnothing$. Finite additivity gives
$p(B)=\sum_{j\leq n}p(B_j)+p(T_n)$. If the sum of the $p(B_j)$ is infinite,
monotonicity forces $p(B)=\infty$. Otherwise $p(T_n)\to0$ implies that some
tail value is finite; the displayed equality then makes $p(B)$ finite and
its limit gives countable additivity. For a finite measure continuity at
the empty set is exactly the decreasing part of the preceding theorem.
\end{proof}

\begin{definition}[Size assumptions]
\family{D11.01.finite_sigma_finite}\label{d11:01:finite-sigma-finite}
A measure is finite if $\mu(X)<\infty$, sigma-finite if
$X=\bigcup_n A_n$ for measurable $A_n$ with $\mu(A_n)<\infty$, and a
probability measure if $\mu(X)=1$. A measurable subset is sigma-finite when
it has such a cover by measurable subsets of itself.
\end{definition}

\begin{lemma}[Finite-measure exhaustions]
\family{D11.01.sigma_finite_exhaustion}\label{d11:01:sigma-finite-exhaustion}
A sigma-finite space admits a disjoint measurable partition $X=\bigsqcup_nD_n$
with finite measures, and an increasing finite-measure exhaustion.
\end{lemma}
\begin{proof}
Given a cover $A_n$, set $D_n=A_n\setminus\bigcup_{j<n}A_j$.
Every point belongs to the set with its first covering index. The finite unions
$F_n=\bigcup_{j\leq n}D_j$ increase to $X$ and have finite measure by finite
subadditivity. Either representation conversely supplies the defining cover.
This includes $X=\varnothing$, for which all the sets may be empty.
\end{proof}

\begin{proposition}[Elementary measures]
\family{D11.01.elementary_measures}\label{d11:01:elementary-measures}
Counting measure assigns a finite measurable set its cardinality and an infinite
measurable set infinity. For $x\in X$, $\delta_x(E)=\one_E(x)$ is a measure.
For $x_n\in X$ and $a_n\in\Eplus$,
$\mu(E)=\sum_n a_n\one_E(x_n)$ is a measure; locations may repeat.
\end{proposition}
\begin{proof}
The cardinality of a disjoint union is the sum of the cardinalities when the
union is finite. If it is infinite, either a member is infinite or the finite
cardinalities have unbounded partial sums. This proves countable additivity
of counting measure. A point belongs to at most one member of a disjoint
family, proving the assertion for $\delta_x$, including the empty family.
For the last assertion, disjointness gives
$\one_{\cup_j E_j}(x_n)=\sum_j\one_{E_j}(x_n)$. Multiplication by $a_n$
preserves this identity even when $a_n=\infty$, because at most one of the
indicators is one. Interchange the two nonnegative sums. This interchange
is elementary: each iterated sum is the supremum of finite rectangular
partial sums, and every finite set of pairs is contained in a finite rectangle.
The empty set has value zero in every construction.
\end{proof}

\begin{definition}[Nullity and completeness]
\family{D11.01.null_complete}\label{d11:01:null-complete}
A measurable $N$ is null if $\mu(N)=0$. An arbitrary set is negligible if it
is contained in a measurable null set. The measure space is complete if every
negligible set is measurable, hence null by monotonicity. This is not metric
completeness.
\end{definition}

\begin{proposition}[Restriction]
\family{D11.01.restriction_measure}\label{d11:01:restriction-measure}
For $E\in\Sigma$, restriction defines a measure on $(E,\Sigma|_E)$ by
$\mu_E(A)=\mu(A)$ for measurable $A\subseteq E$. On $X$ the formula
$\mu^E(A)=\mu(A\cap E)$ also defines a measure. Its mass is $\mu(E)$;
finiteness and sigma-finiteness are inherited from $\mu$.
\end{proposition}
\begin{proof}
Intersections with $E$ preserve empty sets and disjoint unions, so countable
additivity follows immediately from that of $\mu$. For a finite-measure
cover $A_n$ of $X$, the sets $A_n\cap E$ cover $E$ and have finite measure.
For $\mu^E$ they also give a finite-measure cover of $X$, since
$\mu^E(A_n)\leq\mu(A_n)$. The finite case follows by evaluating $X$ or $E$.
\end{proof}

\begin{proposition}[Sums and increasing limits]
\family{D11.01.measure_sums_limits}\label{d11:01:measure-sums-limits}
If $\mu_n$ are measures and $a_n\geq0$ are finite, then
$\sum_n a_n\mu_n$ is a measure. If $\mu_n(A)$ increases with $n$ for every
measurable $A$, then $\mu(A)=\lim_n\mu_n(A)$ is a measure.
\end{proposition}
\begin{proof}
For sums use the nonnegative double-sum identity proved in the preceding
elementary-measure argument. For increasing limits and disjoint $E_j$, put
$b_{n,j}=\mu_n(E_j)$. Then
\[
 \sup_n\sum_j b_{n,j}=\sup_n\sup_k\sum_{j\leq k}b_{n,j}
 =\sup_k\sum_{j\leq k}\sup_n b_{n,j}
 =\sum_j\sup_n b_{n,j}.
\]
The finite-sum equality holds by increasing scalar convergence, including
infinite limits. This is countable additivity; the value at the empty set is zero.
\end{proof}

\begin{corollary}[First Borel--Cantelli lemma]
\family{D11.01.borel_cantelli_one}\label{d11:01:borel-cantelli-one}
If $\sum_n\mu(E_n)<\infty$, then $\mu(\limsup E_n)=0$.
\end{corollary}
\begin{proof}
The set $\limsup E_n$ is contained in $\bigcup_{n\geq N}E_n$ for every $N$.
Its measure is therefore at most $\sum_{n\geq N}\mu(E_n)$, which tends to
zero. Independence and probability normalization are not involved.
\end{proof}

\section{Generating classes and uniqueness}

\begin{theorem}[Pi--lambda]
\family{D11.01.pi_lambda}\label{d11:01:pi-lambda}
A pi-system is closed under finite intersections. A lambda-system contains
$X$ and is closed under complements and countable disjoint unions.
For every pi-system $\mathcal P$, the lambda-system it generates equals
$\sigma(\mathcal P)$.
\end{theorem}
\begin{proof}
A lambda-system is closed under differences $B\setminus A$ when $A\subseteq B$
are its members: $A$ and $B^c$ are disjoint, and complement their union.
It is consequently closed under increasing unions by disjointifying the
sequence. Conversely $X$, nested differences and increasing unions imply
the lambda-system axioms: complements are differences from $X$; if $A,B$
are disjoint, complement $(X\setminus B)\setminus A$ to obtain $A\cup B$;
then use increasing finite unions for a disjoint sequence.

Let $\mathcal D$ be the smallest lambda-system containing $\mathcal P$.
For $A\in\mathcal P$, the family
$\mathcal D_A=\{B:A\cap B\in\mathcal D\}$ is a lambda-system in the
nested-difference formulation: it contains $X$ because $A\in\mathcal D$,
and intersections with $A$ preserve nested differences and increasing unions.
It contains $\mathcal P$, so it contains $\mathcal D$. Fix now $B\in\mathcal D$.
The same argument says that $\{A:A\cap B\in\mathcal D\}$ is a lambda-system;
the first step shows it contains $\mathcal P$, so it contains $\mathcal D$.
Thus $\mathcal D$ is closed under finite intersections. With complements it
is an algebra, and its increasing-union closure makes it a sigma-algebra.
Minimality gives $\sigma(\mathcal P)\subseteq\mathcal D$; the reverse
inclusion holds because every sigma-algebra is a lambda-system.
\end{proof}

\begin{theorem}[Uniqueness on generators]
\family{D11.01.measure_uniqueness_generators}\label{d11:01:measure-uniqueness-generators}
Let $\sigma(\mathcal P)=\Sigma$, with $\mathcal P$ a pi-system. Finite
measures agreeing on $\mathcal P$ and on $X$ agree on $\Sigma$.
The same conclusion holds if they agree on $\mathcal P$ and there are
$P_n\in\mathcal P$ increasing to $X$ with common finite measures.
\end{theorem}
\begin{proof}
In the finite case the equality class
$\mathcal D=\{A\in\Sigma:\mu(A)=\nu(A)\}$ contains $X$. Complements
preserve equality by subtraction from the common finite mass, and disjoint
unions preserve it by countable additivity. Pi--lambda therefore gives
$\Sigma\subseteq\mathcal D$.
In the second case apply the finite case to
$\mu_n(A)=\mu(A\cap P_n)$ and $\nu_n(A)=\nu(A\cap P_n)$.
They agree on $X$ and on $\mathcal P$, since intersections of generators
are generators. For every $A$, continuity from below then gives
$\mu(A)=\lim_n\mu_n(A)=\lim_n\nu_n(A)=\nu(A)$.
\end{proof}

\begin{theorem}[Set monotone classes]
\family{D11.01.monotone_class}\label{d11:01:monotone-class}
A monotone class is closed under increasing countable unions and decreasing
countable intersections. The smallest monotone class containing an algebra
$\calA$ is $\sigma(\calA)$.
\end{theorem}
\begin{proof}
Let $\mathcal M$ be that smallest class. The class
$\{E:E^c\in\mathcal M\}$ is a monotone class containing $\calA$, hence
contains $\mathcal M$. Thus $\mathcal M$ is closed under complements.
For $A\in\calA$, the class $\{E:A\cap E\in\mathcal M\}$ is a monotone
class containing $\calA$, so $A\cap E\in\mathcal M$ for $E\in\mathcal M$.
For fixed $E\in\mathcal M$, repeat this argument with
$\{F:E\cap F\in\mathcal M\}$; it contains $\calA$ by the first step.
Therefore $\mathcal M$ is closed under finite intersections and complements,
and thus finite unions. An arbitrary countable union is an increasing union
of finite unions, so $\mathcal M$ is a sigma-algebra. The two minimality
properties give the result.
\end{proof}

\begin{corollary}[Bounded functional monotone classes]
\family{D11.01.functional_monotone_class}\label{d11:01:functional-monotone-class}
Let $H$ be a real vector space of bounded real functions on $X$, containing
$1$ and $\one_A$ for every member $A$ of a pi-system $\mathcal P$.
Suppose $H$ is closed under pointwise increasing limits of uniformly bounded
sequences whose limit is bounded. Then $H$ contains every bounded
$\sigma(\mathcal P)$-measurable real function.
\end{corollary}
\begin{proof}
Here measurability means that inverse images of Borel subsets of $\R$ belong
to $\sigma(\mathcal P)$; it will be developed in Chapter~\ref{ch:d11-04}.
The class $\mathcal D=\{A:\one_A\in H\}$ contains $X$, is closed under
complements because $1-\one_A\in H$, and is closed under disjoint countable
unions because their finite partial indicator sums increase and lie between
zero and one. Pi--lambda gives $\sigma(\mathcal P)\subseteq\mathcal D$.
For a bounded nonnegative measurable $f$, the dyadic lower approximations
$s_n=2^{-n}\lfloor2^nf\rfloor$ take finitely many values, have measurable
level sets, increase to $f$, and are bounded by the bound for $f$.
They are finite linear combinations of indicators in $H$, so $f\in H$.
For an arbitrary bounded real $f$, add a sufficiently large nonnegative
constant and then subtract it. This proof uses the explicit dyadic construction,
not a forward integration theorem.
\end{proof}

\section{Boundaries of the axioms}

\begin{example}[A union need not be a sigma-algebra]
\family{D11.01.union_sigma_failure}\label{d11:01:union-sigma-failure}
On $X=\{1,2,3\}$, the sigma-algebras
$\{\varnothing,X,\{1\},\{2,3\}\}$ and
$\{\varnothing,X,\{2\},\{1,3\}\}$ have a union that is not an algebra:
it contains $\{1\}$ and $\{2\}$ but not their union.
\end{example}

\begin{example}[Finite additivity is insufficient]
\family{D11.01.finite_additivity_failure}\label{d11:01:finite-additivity-failure}
On the finite/cofinite algebra of $\N$, let $p(A)=0$ for finite $A$ and
$p(A)=1$ for cofinite $A$. It is finitely additive: in a finite disjoint family
at most one set is cofinite, and the union is cofinite exactly when such a
member occurs. But $p(\N)=1\ne\sum_n p(\{n\})=0$. Thus it admits no
countably additive extension to $\mathcal P(\N)$.
\end{example}

\begin{example}[Decreasing continuity needs a finite term]
\family{D11.01.continuity_above_failure}\label{d11:01:continuity-above-failure}
With the Lebesgue measure constructed in Chapter~\ref{ch:d11-03},
$E_n=(n,\infty)$ decrease to $\varnothing$ but all have infinite measure.
The limit of their measures is therefore not the measure of their intersection.
This example is not used in the construction of Lebesgue measure.
\end{example}
\source{COH 1.1--1.2 and 1.6; AXL 5A; ROY 9.1.}

% END chapters/01


% BEGIN chapters/02
\chapter{Outer Measures, Extension and Completion}\label{ch:d11-02}

There are two distinct tasks in extending an elementary notion of size.
First one constructs a countably additive measure on a suitable domain.
One then asks whether the extension is unique. Neither task should be
confused with adjoining all subsets of null sets.

\section{Outer measures and the splitting test}

\begin{definition}[Outer measure]
\family{D11.02.outer_measure}\label{d11:02:outer-measure}
An outer measure is a function $m^*:\mathcal P(X)\to\Eplus$ that is zero
on the empty set, monotone, and countably subadditive. Its domain is every
subset of $X$; countable additivity on that whole domain is not required.
\end{definition}

\begin{proposition}[Covering costs]
\family{D11.02.cover_outer_measure}\label{d11:02:cover-outer-measure}
Let $\mathcal C\subseteq\mathcal P(X)$ contain $\varnothing$ and $X$, and
let $\rho:\mathcal C\to\Eplus$ satisfy $\rho(\varnothing)=0$. Then
\[
 m^*(E)=\inf\left\{\sum_n\rho(C_n):C_n\in\mathcal C,
                  \ E\subseteq\bigcup_n C_n\right\}
\]
is an outer measure. The infimum of an empty class is interpreted as infinity.
\end{proposition}
\begin{proof}
The cover by empty sets gives $m^*(\varnothing)=0$, and every cover of a
larger set covers a smaller one, proving monotonicity. For subadditivity the
case $\sum_jm^*(E_j)=\infty$ needs no argument. Otherwise, for $\eps>0$
choose for each $j$ a cover $C_{j,n}$ of $E_j$ with total cost at most
$m^*(E_j)+\eps2^{-j}$. These choices exist by the definition of an infimum;
all the relevant infima are finite. Enumerate the countable family indexed
by pairs $(j,n)$ as one sequence. It covers $\bigcup_jE_j$ and its cost is
at most $\sum_jm^*(E_j)+\eps$. Let $\eps\downarrow0$.
\end{proof}

\begin{definition}[Caratheodory measurability]
\family{D11.02.caratheodory_measurable}\label{d11:02:caratheodory-measurable}
A set $A\subseteq X$ is $m^*$-measurable if
\[
 m^*(E)=m^*(E\cap A)+m^*(E\setminus A)\qquad(E\subseteq X).
\]
The quantifier tests all subsets $E$. Subadditivity already gives the
$\leq$ direction. To prove the reverse direction it is enough to test
$m^*(E)<\infty$, since an infinite left side makes it automatic.
\end{definition}

\begin{theorem}[Caratheodory construction]
\family{D11.02.caratheodory_theorem}\label{d11:02:caratheodory-theorem}
The $m^*$-measurable sets form a sigma-algebra $\mathcal M(m^*)$.
The restriction of $m^*$ to it is a complete measure.
\end{theorem}
\begin{proof}
The empty set and $X$ pass the splitting test, which is symmetric in a set
and its complement. If $A,B$ pass the test, split $E$ first by $A$ and then
split $E\cap A$ by $B$. This gives
\[
 m^*(E)=m^*(E\cap A\cap B)+m^*(E\cap A\setminus B)+m^*(E\setminus A).
\]
Subadditivity bounds the last two terms below by
$m^*(E\setminus(A\cap B))$. Thus $A\cap B$ passes the test.
The class is consequently an algebra.

For disjoint $A_j$ in this algebra, repeated splitting gives, for every $n$,
\[
 m^*(E)=\sum_{j\leq n}m^*(E\cap A_j)
          +m^*\left(E\setminus\bigcup_{j\leq n}A_j\right)
 \geq\sum_{j\leq n}m^*(E\cap A_j)
          +m^*\left(E\setminus\bigcup_jA_j\right).
\]
Take the supremum in $n$, then use subadditivity on $E\cap\bigcup_jA_j$.
The result is the reverse splitting inequality for $\bigcup_jA_j$.
Hence that union is measurable. General countable unions are disjointified
inside the algebra, so the class is a sigma-algebra. Taking
$E=\bigcup_jA_j$ in the same inequality proves the reverse countable
additivity inequality; subadditivity proves the other one.

Finally, if $m^*(N)=0$ and $Z\subseteq N$, monotonicity gives $m^*(Z)=0$.
For any $E$, subadditivity and monotonicity imply
\[
 m^*(E)\leq m^*(E\setminus Z)+m^*(E\cap Z)
           =m^*(E\setminus Z)\leq m^*(E).
\]
Thus $Z$ is measurable and null. This proves completeness without any
sigma-finiteness assumption.
\end{proof}

\section{Premeasures and their extensions}

\begin{definition}[Premeasure]
\family{D11.02.premeasure}\label{d11:02:premeasure}
A premeasure on an algebra $\calA$ is a function $p:\calA\to\Eplus$ with
$p(\varnothing)=0$ such that
$p(\bigcup_n A_n)=\sum_n p(A_n)$ whenever the $A_n$ are disjoint members
of $\calA$ and their union belongs to $\calA$. It is sigma-finite when $X$
has a countable cover by members of $\calA$ of finite $p$-value.
\end{definition}

\begin{theorem}[Agreement of the covering construction]
\family{D11.02.premeasure_outer_agreement}\label{d11:02:premeasure-outer-agreement}
For a premeasure $p$ define
\[
 p^*(E)=\inf\left\{\sum_n p(A_n): A_n\in\calA,
                               \ E\subseteq\bigcup_nA_n\right\}.
\]
Then $p^*$ is an outer measure, $\calA\subseteq\mathcal M(p^*)$, and
$p^*(A)=p(A)$ for $A\in\calA$.
\end{theorem}
\begin{proof}
The preceding covering proposition applies. A premeasure is finitely additive,
and hence monotone on its algebra by the same disjoint-difference argument as
for measures. If $A\subseteq\bigcup_n A_n$ with $A,A_n\in\calA$, put
$B_n=A\cap(A_n\setminus\bigcup_{j<n}A_j)$. These are disjoint algebra
members with union $A$, so
$p(A)=\sum_np(B_n)\leq\sum_np(A_n)$. Taking the infimum gives
$p(A)\leq p^*(A)$; covering $A$ by itself and empty sets gives the reverse.

For fixed $A\in\calA$, split every covering set $C_n$ of an arbitrary $E$
into $C_n\cap A$ and $C_n\setminus A$. Finite additivity gives
\[
 \sum_n p(C_n)=\sum_n p(C_n\cap A)+\sum_n p(C_n\setminus A)
 \geq p^*(E\cap A)+p^*(E\setminus A).
\]
The inequality remains valid for infinite costs. Its infimum proves the
Caratheodory test for $A$.
\end{proof}

\begin{theorem}[Existence and the precise uniqueness assertion]
\family{D11.02.extension_existence_uniqueness}\label{d11:02:extension-existence-uniqueness}
The measure $\mu=p^*|_{\sigma(\calA)}$ extends $p$. If $\nu$ is another
extension on $\sigma(\calA)$, then $\nu(E)\leq\mu(E)$ for every measurable
$E$, with equality whenever $\mu(E)<\infty$. If $p$ is sigma-finite, the
extension on $\sigma(\calA)$ is unique.
\end{theorem}
\begin{proof}
Existence follows from Caratheodory's theorem because the measurable class
contains $\calA$ and hence $\sigma(\calA)$. For any algebra cover of $E$,
$\nu(E)\leq\sum_np(A_n)$; take its infimum.

Suppose $\mu(E)<\infty$ and choose a cover with cost less than
$\mu(E)+\eps$. Let $B=\bigcup_nA_n$. Disjointify this union within the
algebra. Since both measures extend $p$, countable additivity shows
$\nu(B)=\mu(B)<\infty$. Also $\mu(B\setminus E)<\eps$. Therefore
\[
 \mu(E)\leq\mu(B)=\nu(B)=\nu(E)+\nu(B\setminus E)
 \leq\nu(E)+\mu(B\setminus E)<\nu(E)+\eps.
\]
Let $\eps\downarrow0$ and combine with $\nu(E)\leq\mu(E)$.
For a sigma-finite premeasure, finite unions of a finite-valued algebra cover
give an increasing finite-valued exhaustion in $\calA$. Apply
\cref{d11:01:measure-uniqueness-generators} to the pi-system $\calA$.
No sigma-finiteness was used for existence.
\end{proof}

\section{Completion is a separate operation}

\begin{proposition}[Completion]
\family{D11.02.completion}\label{d11:02:completion}
For a measure $\mu$ on $\Sigma$, let
\[
 \overline\Sigma=\{E\subseteq X:\ A\subseteq E\subseteq B
       \text{ for some }A,B\in\Sigma,\ \mu(B\setminus A)=0\}.
\]
Then $\overline\mu(E)=\mu(A)=\mu(B)$ is well-defined and is a complete
measure extending $\mu$. Its sigma-algebra is generated by $\Sigma$ and all
subsets of measurable $\mu$-null sets. Every complete extension contains
this sigma-algebra and agrees with $\overline\mu$ there.
\end{proposition}
\begin{proof}
Finite additivity with a null difference gives $\mu(A)=\mu(B)$ without
subtracting infinite quantities. For two sandwiches, $A_1\subseteq B_2$
and $A_2\subseteq B_1$, so these common values agree by monotonicity.
Complementing a sandwich gives
$B^c\subseteq E^c\subseteq A^c$ with null difference. Taking countable
unions of sandwiches gives
\[
 \bigcup_n A_n\subseteq\bigcup_n E_n\subseteq\bigcup_n B_n,
 \qquad
 (\bigcup_nB_n)\setminus(\bigcup_nA_n)
       \subseteq\bigcup_n(B_n\setminus A_n),
\]
whose last set is null. Thus $\overline\Sigma$ is a sigma-algebra.
For disjoint $E_n$ their inner sets $A_n$ are disjoint, and the union
sandwich shows
$\overline\mu(\bigcup_nE_n)=\mu(\bigcup_nA_n)=\sum_n\overline\mu(E_n)$.
If $\overline\mu(E)=0$, an outer sandwich set $B$ is $\mu$-null;
every subset of $E$ is sandwiched between $\varnothing$ and $B$.
Hence the extension is complete.

Every sandwich set is $E=A\cup(E\setminus A)$, where the second piece is a
subset of the measurable null set $B\setminus A$. This proves the claimed
generation, including the reverse inclusion. A complete extension must
contain those generators. On a sandwich its value is trapped between the
equal values of $A$ and $B$, proving agreement and minimality.
\end{proof}

\begin{theorem}[Induced outer measure and sigma-finite completion]
\family{D11.02.induced_outer_completion}\label{d11:02:induced-outer-completion}
For a sigma-finite measure $\mu$ on $\Sigma$, put
$\mu^*(E)=\inf\{\mu(A):E\subseteq A,\ A\in\Sigma\}$.
This is an outer measure, $\mathcal M(\mu^*)=\overline\Sigma$, and its
measure on that class is $\overline\mu$.
\end{theorem}
\begin{proof}
Countable subadditivity follows by choosing measurable supersets with
summable errors, just as for covering costs. On $\Sigma$ the value equals
$\mu$. Every $A\in\Sigma$ is $\mu^*$-measurable: split each measurable
superset of an arbitrary $E$ by $A$ and take its infimum. The resulting
Caratheodory measure is complete, so it contains and agrees with the
completion by the preceding proposition.

We prove the reverse inclusion. Whenever $\mu^*(H)<\infty$, choose
measurable $B_j\supseteq H$ with
$\mu(B_j)<\mu^*(H)+1/j$. Their intersection $B$ contains $H$ and has
$\mu(B)=\mu^*(H)$ by monotonicity and the defining infimum.
In particular every outer-null set has a measurable null envelope.
Let $F_n\in\Sigma$ be a finite-measure exhaustion and let $E$ pass the
Caratheodory test. The set $H_n=E\cap F_n$ also passes that test and has
finite outer measure. Choose its measurable envelope $B_n\subseteq F_n$
as above. Splitting $B_n$ by $H_n$ gives
\[
 \mu(B_n)=\mu^*(H_n)+\mu^*(B_n\setminus H_n).
\]
The first term equals the finite left side, so the remainder is outer-null
and has a measurable null envelope. Thus $H_n$ is a completion-measurable
set. Their union is $E$, proving $E\in\overline\Sigma$.
The equality of measures follows from extension agreement on each sandwich.
The exhaustion, and therefore the sigma-finite hypothesis, was used exactly
in this reverse inclusion.
\end{proof}

\section{Two failures worth retaining}

\begin{example}[Nonunique extensions]
\family{D11.02.extension_nonuniqueness}\label{d11:02:extension-nonuniqueness}
Let $\calA$ be the algebra of finite unions of intervals $(a,b]$, $(-\infty,b]$
and $(a,\infty)$, together with $\varnothing$ and $\R$; all displayed finite
endpoints are real and $a<b$. Every nonempty member contains a nondegenerate
interval. Set $p(\varnothing)=0$ and $p(A)=\infty$ otherwise.
It is a premeasure: a disjoint union in the algebra is nonempty exactly when
at least one of its members is nonempty, and both sides of the premeasure
identity are then infinity. Its generated sigma-algebra is $\BB(\R)$.
Counting measure and the measure that is infinity on every nonempty Borel
set both extend $p$, but have values $1$ and infinity at a singleton.
Thus extension existence does not imply uniqueness.
\end{example}

\begin{example}[An outer measure that is not additive]
\family{D11.02.outer_not_additive}\label{d11:02:outer-not-additive}
On $X=\{a,b\}$ let $m^*(E)=1$ for nonempty $E$, and $m^*(\varnothing)=0$.
It is monotone; if a countable union is nonempty some member is nonempty,
so it is subadditive. But $m^*(X)=1\ne2=m^*(\{a\})+m^*(\{b\})$.
The empty set and $X$ pass the splitting test; each singleton fails it when
tested against $E=X$. Thus its Caratheodory sigma-algebra is precisely
$\{\varnothing,X\}$.
\end{example}
\source{COH 1.3 and 1.5; FOL 1.4, including its extension and completion exercises.}

% END chapters/02


% BEGIN chapters/03
\chapter{Lebesgue Measure in Euclidean Space}\label{ch:d11-03}

We now construct volume rather than assume it as a measure. The elementary
product of side lengths belongs to finite-dimensional geometry; countable
additivity, regularity and completion will be consequences of the construction.
Throughout this chapter $d\geq1$. The notation $\lambda_d$ denotes completed
Lebesgue measure, and $\lambda_d^B$ its Borel restriction.

\section{Borel sets and outer volume}

\begin{definition}[Euclidean Borel sets]
\family{D11.03.borel_euclidean}\label{d11:03:borel-euclidean}
The Borel sigma-algebra $\BB(\R^d)$ is generated by the Euclidean open sets.
A Borel set means a member of this sigma-algebra, not an arbitrary set to
which a completed measure will eventually assign a value.
\end{definition}

\begin{theorem}[Generators]
\family{D11.03.borel_generators}\label{d11:03:borel-generators}
The Borel sigma-algebra is generated by rational open boxes, by bounded
half-open boxes, or by the coordinate half-spaces $\{x:x_j\leq q\}$ with
$q\in\Q$. On $\R$, each of the open, closed and half-open interval families,
and each of the open or closed ray families, also generates it.
\end{theorem}
\begin{proof}
Every listed generator is Borel. Rational open boxes form a countable basis,
so every open set is the union of the basis members it contains and hence a
countable union. Rays with arbitrary endpoints are countable unions or
intersections of rays with rational endpoints. Differences of coordinate rays
give strips; finite intersections give half-open boxes. Conversely an open
box is a countable union of half-open boxes lying inside it. These observations
prove both inclusions for each multidimensional family. In one dimension,
complements interchange open and closed rays; differences of rays give
half-open intervals, while open intervals are countable unions of closed
intervals lying inside them. Rays are countable unions of bounded intervals.
Thus each one-dimensional family gives all the others.
\end{proof}

\begin{proposition}[Outer volume]
\family{D11.03.lebesgue_outer}\label{d11:03:lebesgue-outer}
For $E\subseteq\R^d$ define
\[
 \lambda_d^*(E)=\inf\left\{\sum_n\prod_{j=1}^d(b_{nj}-a_{nj}):
 E\subseteq\bigcup_n\prod_{j=1}^d(a_{nj},b_{nj})\right\},
\]
using bounded open boxes. This is an outer measure and the outer measure of
every bounded box, with any endpoint convention, is its product of side lengths.
\end{proposition}
\begin{proof}
One may adjoin the empty set at cost zero and $\R^d$ at cost infinity to the
cover class. These additions do not change the infimum: empty sets contribute
nothing, and a cover containing a set of infinite cost cannot lower a finite
infimum. If the original infimum is infinite it is unchanged as well.
The covering-cost construction therefore proves the outer-measure axioms.

Here is the geometric lower bound needed for normalization. If a closed
box $Q=\prod_j[a_j,b_j]$ is covered by finitely many open boxes $R_i$, then
\begin{equation}\label{eq:d11-finite-box-cover}
 \operatorname{vol}(Q)\leq\sum_i\operatorname{vol}(R_i).
\end{equation}
If a side of $Q$ has length zero the inequality is immediate. Otherwise cut
each coordinate axis at the finitely many endpoints of $Q$ and the $R_i$,
inside a large box containing them. The resulting rectangular cells have
additive elementary volumes: this is repeated distributivity of the finite
products of side-length sums. The interior of every positive-volume cell of
$Q$ lies in one of the covering boxes. Indeed choose a point of that interior
and a box covering it; no boundary of this covering box crosses the cell's
interior, because all such boundaries are cuts of the grid. Assign the cell
to one such box. The cells assigned to $R_i$ have disjoint interiors and total
elementary volume at most $\operatorname{vol}(R_i)$, again by the same grid
identity. Summing proves \eqref{eq:d11-finite-box-cover}; grid faces do not
contribute a positive side-length product.

Every countable open-box cover of the compact $Q$ has a finite subcover,
so \eqref{eq:d11-finite-box-cover} gives
$\lambda_d^*(Q)\geq\operatorname{vol}(Q)$. Enlarging each side by $\delta$
gives an open cover whose cost tends to $\operatorname{vol}(Q)$ as
$\delta\downarrow0$, proving the opposite inequality, including degenerate
boxes. A bounded box with other endpoint conventions lies between closed
inner boxes and open outer boxes with volumes arbitrarily close to its
side-length product. Monotonicity completes the normalization.
\end{proof}

\begin{proposition}[The Lebesgue measures]
\family{D11.03.lebesgue_measure}\label{d11:03:lebesgue-measure}
Every Borel set is Caratheodory measurable for $\lambda_d^*$. Let $\LL_d$ be
the full Caratheodory sigma-algebra. The restrictions $\lambda_d^B$ to
$\BB(\R^d)$ and $\lambda_d$ to $\LL_d$ are sigma-finite measures assigning
box volume; $\lambda_d$ is complete.
\end{proposition}
\begin{proof}
Fix a coordinate half-space $H=\{x:x_j\leq t\}$. Split each box in a cover
of an arbitrary $E$ by $H$. The two resulting boxes may have mixed endpoint
conventions, but the preceding proposition gives their outer measures, and
their volumes sum to the original box volume. Thus
\[
 \lambda_d^*(E\cap H)+\lambda_d^*(E\setminus H)
 \leq\sum_n\operatorname{vol}(R_n).
\]
Take the infimum over covers. This is the reverse Caratheodory inequality;
subadditivity gives the other direction. Hence all coordinate half-spaces
are measurable, so the generator theorem gives all Borel sets. The
Caratheodory theorem now supplies countable additivity and completeness.
Finally $[-n,n]^d$ are measurable of finite volume and cover $\R^d$.
\end{proof}

\begin{corollary}[Countable null sets]
\family{D11.03.null_countable}\label{d11:03:null-countable}
Singletons and countable subsets of $\R^d$ have outer measure zero. Every
subset of a Lebesgue-null set is Lebesgue measurable and null.
\end{corollary}
\begin{proof}
A singleton is contained in boxes of arbitrarily small volume. Countable
subadditivity handles countable sets. The final assertion is the completeness
part of the Caratheodory construction, or follows directly from the outer-null
splitting test.
\end{proof}

\section{Regularity before completion}

\begin{theorem}[Euclidean regularity]
\family{D11.03.euclidean_regularity}\label{d11:03:euclidean-regularity}
For $E\in\LL_d$,
\[
 \lambda_d(E)=\inf_{E\subseteq U,\ U\ \text{open}}\lambda_d(U)
            =\sup_{K\subseteq E,\ K\ \text{compact}}\lambda_d(K).
\]
If $\lambda_d(E)<\infty$, then for every $\eps>0$ there are compact
$K\subseteq E\subseteq U$ with $U$ open and $\lambda_d(U\setminus K)<\eps$.
\end{theorem}
\begin{proof}
If $\lambda_d(E)<\infty$, choose a box cover of total volume less than
$\lambda_d(E)+\eps$ and take its open union $U$. Then
$\lambda_d(U\setminus E)<\eps$. For infinite measure the stated outer
infimum equality follows from monotonicity.

If $E$ is bounded, take a closed bounded box $Q$ containing it. Apply the
finite outer approximation to $Q\setminus E$, obtaining open $O$ with
$\lambda_d(O)<\lambda_d(Q\setminus E)+\eps$. The set $K=Q\setminus O$
is compact and contained in $E$. Since $Q\setminus E\subseteq O\cap Q$,
finite subtraction gives
\[
 \lambda_d(E\setminus K)
 =\lambda_d(O\cap Q)-\lambda_d(Q\setminus E)<\eps.
\]
For arbitrary $E$, its intersections with $[-n,n]^d$ increase to it.
If $\lambda_d(E)$ is finite, choose $n$ leaving mass less than $\eps$,
and then apply the bounded case. If it is infinite, these bounded
intersections have arbitrarily large measure, and the bounded case gives
compact subsets of arbitrarily large measure. This proves the inner supremum
in all cases. In the finite case split the requested error between the outer
and inner approximations to obtain $\lambda_d(U\setminus K)<\eps$.
\end{proof}

\begin{corollary}[Global small errors]
\family{D11.03.global_closed_open_approximation}\label{d11:03:global-closed-open-approximation}
For every $E\in\LL_d$ and $\eps>0$ there are a closed $F$ and an open $G$
with $F\subseteq E\subseteq G$ and
$\lambda_d(E\setminus F)<\eps$, $\lambda_d(G\setminus E)<\eps$.
There is no finite-measure assumption on $E$; $F$ need not be compact.
\end{corollary}
\begin{proof}
Put $E_n=E\cap[-n,n]^d$. By finite outer approximation choose open
$G_n\supseteq E_n$ with $\lambda_d(G_n\setminus E_n)<\eps2^{-n}$.
Then $G=\bigcup_nG_n$ contains $E$, and
$G\setminus E\subseteq\bigcup_n(G_n\setminus E_n)$ has measure less than
$\eps$. Apply this already established outer approximation to $E^c$ and
complement the resulting open set. This gives the closed $F$. The argument
does not try to subtract two infinite measures.
\end{proof}

\begin{theorem}[Borel and box approximation]
\family{D11.03.borel_approximation}\label{d11:03:borel-approximation}
Every $E\in\LL_d$ lies between an $F_\sigma$ set $F$ and a $G_\delta$ set
$G$ with both $E\setminus F$ and $G\setminus E$ null. If $\lambda_d(E)$
is finite, it can be approximated in symmetric-difference measure by a finite
disjoint union of bounded boxes.
\end{theorem}
\begin{proof}
Choose closed $F_n\subseteq E\subseteq G_n$ open with both errors less than
$2^{-n}$. Let $F=\bigcup_nF_n$ and $G=\bigcap_nG_n$. The error for $F$
is at most the error for every $F_n$, hence zero; the same argument works
for $G$. These are Borel sets by their displayed forms.

For finite measure, first choose $n$ such that
$\lambda_d(E\setminus[-n,n]^d)<\eps/3$. Cover the bounded part by open
bounded boxes whose total volume is less than its measure plus $\eps/3$.
Their union $O$ has finite measure and has excess less than $\eps/3$ over
that bounded part. Finite partial unions increase to $O$, so one such union
$V$ leaves measure less than $\eps/3$. Then
$\lambda_d(E\mathbin\triangle V)<\eps$. To disjointify a finite box union,
cut at all its coordinate endpoints. The resulting products of open
coordinate intervals and singleton endpoints are finitely many pairwise
disjoint boxes partitioning the union. Degenerate boxes have zero volume;
they may instead be discarded if only the symmetric-difference assertion
is wanted. This proves the claimed finite approximation.
\end{proof}

\begin{theorem}[Identification of the completion]
\family{D11.03.lebesgue_completion}\label{d11:03:lebesgue-completion}
The space $(\R^d,\LL_d,\lambda_d)$ is the completion of
$(\R^d,\BB(\R^d),\lambda_d^B)$. In particular, a Lebesgue set is a Borel
set modulo a subset of a Borel null set.
\end{theorem}
\begin{proof}
The preceding theorem sandwiches $E\in\LL_d$ between Borel $F,G$ with
$\lambda_d^B(G\setminus F)=0$. Thus $E$ belongs to the completion of the
Borel space. Conversely $\lambda_d$ is a complete extension of the Borel
measure, so it contains and agrees with that completion by
\cref{d11:02:completion}. This proves equality of both sigma-algebras and
measures. The decomposition $E=F\cup(E\setminus F)$ gives the last assertion;
conversely adjoining a subset of a Borel null set is permitted in the completion.
\end{proof}

\section{Symmetry and uniqueness}

\begin{theorem}[Translations and dilations]
\family{D11.03.translation_dilation}\label{d11:03:translation-dilation}
For $a\in\R^d$, $r\ne0$ and any $E\subseteq\R^d$,
$\lambda_d^*(a+rE)=|r|^d\lambda_d^*(E)$. This invertible affine map
preserves Borel and Lebesgue measurability in both directions, with the same
formula for measurable sets. For $r=0$ the image is empty or a singleton;
the numerical outer-measure formula still holds with $0\cdot\infty=0$,
but there is no reverse measurability implication.
\end{theorem}
\begin{proof}
Transform an open-box cover by $x\mapsto a+rx$. Its costs are multiplied
by $|r|^d$, so $\lambda_d^*(a+rE)\leq|r|^d\lambda_d^*(E)$.
Apply the inverse transformation to obtain the opposite inequality.
These inequalities also handle infinite outer measure, since the scalar
factor is strictly positive and finite. Borel measurability is preserved
because the map and its inverse are continuous. For Caratheodory measurability,
write an arbitrary test set as the image of its inverse image, split there by
$E$, and multiply the splitting identity by $|r|^d$. Thus the image passes
the test exactly when $E$ does. The zero-dilation conclusion follows from
the singleton calculation and is independent of measurability of $E$.
\end{proof}

\begin{theorem}[Uniqueness of volume]
\family{D11.03.euclidean_measure_uniqueness}\label{d11:03:euclidean-measure-uniqueness}
A Borel measure assigning all bounded half-open boxes their volumes is
$\lambda_d^B$. Every nonzero translation-invariant Borel measure finite on
bounded Borel sets is $c\lambda_d^B$ for a unique $c>0$.
\end{theorem}
\begin{proof}
Bounded half-open boxes, with the empty set, form a generating pi-system.
The boxes $(-n,n]^d$ give a common finite exhaustion for any two measures
with the asserted values. Generator uniqueness proves the first assertion.

For a translation-invariant measure $\mu$, put $c=\mu([0,1)^d)<\infty$.
Partition the unit cube into $2^{kd}$ translates of $[0,2^{-k})^d$.
Finite additivity gives each such small cube measure $c2^{-kd}$.
Consequently a half-open box with dyadic rational endpoints, partitioned
into a sufficiently fine finite grid of such cubes, has measure $c$ times
its volume. If $c=0$, the countable cover of $\R^d$ by integer translates
of the unit cube would make $\mu$ zero. Thus $c>0$. The dyadic-endpoint
boxes form a generating pi-system and have increasing finite-measure
exhaustions $[-n,n)^d$. Applying generator uniqueness to $\mu$ and
$c\lambda_d^B$ proves equality. Evaluation on the unit cube gives uniqueness
of $c$.
\end{proof}

\section{Null, large and nonmeasurable sets}

\begin{proposition}[Cantor sets]
\family{D11.03.cantor_null_uncountable}\label{d11:03:cantor-null-uncountable}
The middle-thirds Cantor set $C$ is compact, perfect, nowhere dense,
of cardinality continuum, and null. For every $0<\alpha<1$ there is a
closed nowhere-dense subset of $[0,1]$ of measure $\alpha$.
\end{proposition}
\begin{proof}
After $n$ middle-third deletions the retained set $C_n$ is a union of
$2^n$ closed intervals of length $3^{-n}$. Their intersection $C$ is compact
and has measure at most $(2/3)^n$ for every $n$, hence zero. It has no
nondegenerate interval because every such interval has positive measure.
Its points are exactly the numbers admitting a ternary expansion using
only $0$ and $2$. The map
$(s_n)\mapsto\sum_n2s_n3^{-n}$ from binary sequences is injective: at the
first differing digit the difference $2\cdot3^{-k}$ exceeds the maximum
remaining tail $3^{-k}$. Its range is $C$. This proves its cardinality.
Changing a digit arbitrarily far out gives a distinct point arbitrarily
near any prescribed point of $C$, so no point is isolated. Closedness and
empty interior give nowhere density.

For the second construction, remove total length
$\delta_n=(1-\alpha)2^{-n}$ at stage $n$, equally from the middle of the
$2^{n-1}$ surviving intervals. At every stage the surviving components
have equal length, and the proposed deletion is smaller than each component:
the total length before the deletion exceeds $\alpha>0$ plus that deletion.
Thus the construction is possible. The nested compact retained sets have
measures $1-\sum_{j\leq n}\delta_j$, tending to $\alpha$; continuity from
above gives that measure for their intersection. Each component has length
at most $2^{-n}$, so the intersection contains no nondegenerate interval.
It is therefore closed and nowhere dense.
\end{proof}

\begin{example}[Vitali's obstruction]
\family{D11.03.vitali_nonmeasurable}\label{d11:03:vitali-nonmeasurable}
Assume choice and select $V\subseteq[0,1]$ containing exactly one
representative of every equivalence class for $x\sim y$ when $x-y\in\Q$.
Then $V$ is not Lebesgue measurable.
\end{example}
\begin{proof}
The sets $V+q$, $q\in\Q\cap[-1,1]$, are disjoint: equality
$v+q=w+r$ implies $v\sim w$, hence $v=w$ and $q=r$.
Their union $U$ contains $[0,1]$ and is contained in $[-1,2]$.
If $V$ were measurable and null, countable additivity would give
$\lambda_1(U)=0$. If it had positive measure, the infinitely many disjoint
translates would give $\lambda_1(U)=\infty$. Both conclusions contradict
$1\leq\lambda_1(U)\leq3$.
\end{proof}

\begin{corollary}[A nonmeasurable subset in every positive outer-measure set]
\family{D11.03.positive_set_nonmeasurable_subset}\label{d11:03:positive-set-nonmeasurable-subset}
Assuming choice, every $B\subseteq\R$ with $\lambda_1^*(B)>0$ contains
a non-Lebesgue-measurable subset.
\end{corollary}
\begin{proof}
Some bounded part $E=B\cap[-m,m]$ has positive outer measure; otherwise
countable subadditivity would give $\lambda_1^*(B)=0$. Select one
representative $v\in E$ from each rational-translation class meeting $E$.
The rational translates of this representative set $V$ with
$q\in[-2m,2m]$ cover $E$, are disjoint, and lie in $[-3m,3m]$.
If $V$ were measurable and null the covering would contradict positive
outer measure of $E$. If it had positive measure, its infinitely many
bounded disjoint translates would contradict finite measure of $[-3m,3m]$.
Thus $V\subseteq B$ is the required nonmeasurable subset.
\end{proof}

\begin{proposition}[The Cantor-function bridge]
\family{D11.03.singular_function_route}\label{d11:03:singular-function-route}
The Cantor set supports a continuous nondecreasing map
$\Gamma:[0,1]\to[0,1]$ with $\Gamma(0)=0$, $\Gamma(1)=1$, constant on
each component of $[0,1]\setminus C$. Its derivative, absolute-continuity
and Stieltjes consequences will be treated in
\cref{d11:12:singular-continuous-counterexample}.
\end{proposition}
\begin{proof}
On $C$, replace the digits $0,2$ of a ternary expansion by binary digits
$0,1$. This is nondecreasing: at the first differing digit the later binary
tail cannot reverse the ordering. Every binary sequence comes from a ternary
sequence in $C$, so the map is onto $[0,1]$. On a component of $C_n$ its
oscillation is at most $2^{-n}$. Since distinct components of $C_n$ are
separated, this proves continuity on $C$: a convergent sequence in $C$
eventually lies in the same $C_n$ component as its limit.
The two endpoints of a deleted interval have the same value, by the binary
identity ending in infinitely many ones versus the next terminating binary
expansion. Fill that interval with the common value. The extension remains
nondecreasing and onto. A jump of a monotone function would omit the open
interval between a one-sided limit and its value; surjectivity rules out
all such gaps, including at the endpoints. Hence the extension is continuous.
This constructs the object used later without assuming its later derivative
or measure representation.
\end{proof}

\begin{example}[A Lebesgue set that is not Borel]
\family{D11.03.lebesgue_not_borel}\label{d11:03:lebesgue-not-borel}
There is a non-Borel subset $N$ of the Cantor set. It is Lebesgue measurable
and null, so Borel Lebesgue measure is not complete.
\end{example}
\begin{proof}
There are at most continuum many Borel sets. To verify the bound, start with
the countable rational-box basis. A countable well-founded tree, with leaves
labelled by basis sets and internal vertices labelled by complement or
countable union, codes a Borel set by recursive evaluation. The class of
sets so coded contains the basis and is closed under complement and
countable union: adjoin a new root, and in the union case attach the countably
many coding trees beneath it. The resulting tree is countable and
well-founded. Thus every Borel set has such a code. Trees can be represented
as subsets of the countable set of finite sequences of natural numbers;
their countably many labels can also be encoded by a subset of a countable
set. Hence the total number of codes is at most continuum.
But $C$ has cardinality continuum, so Cantor's theorem says that
$\mathcal P(C)$ has strictly larger cardinality. Some $N\subseteq C$ is
therefore not Borel. It is nevertheless measurable and null because $C$ is
a Borel null set and $\lambda_1$ is complete.
\end{proof}

\begin{theorem}[Steinhaus]
\family{D11.03.steinhaus}\label{d11:03:steinhaus}
If $E\in\LL_1$ has positive measure, then $E-E$ contains an open interval
about zero.
\end{theorem}
\begin{proof}
Inner regularity supplies compact $K\subseteq E$ with
$0<\lambda_1(K)<\infty$. Choose open $U\supseteq K$ with
$\lambda_1(U)<2\lambda_1(K)$. Compactness gives $\delta>0$ such that
$K+h\subseteq U$ for $|h|<\delta$: choose finitely many balls with half
radii covering $K$, whose full-radius balls lie in $U$, and take the minimum
half radius. Then $K$ and $K+h$ are both contained in $U$. Translation
invariance and the finite union formula give
\[
 \lambda_1(K\cap(K+h))=2\lambda_1(K)-\lambda_1(K\cup(K+h))>0.
\]
A point $x=y+h$ in that intersection has $x,y\in K$, so $h=x-y\in E-E$.
Thus $(-\delta,\delta)\subseteq E-E$.
\end{proof}
\source{COH 1.1, 1.3--1.5 and 2.1; ROY 2.4 and 2.6--2.7; FOL 2.6; AXL 2D.}

% END chapters/03


% BEGIN chapters/04
\chapter{Measurable Functions and Almost-Everywhere Statements}\label{ch:d11-04}

We distinguish three objects: a function with a specified value at every
point, its class modulo null changes, and a representative chosen from that
class. Incomplete measures make this distinction visible even before norms
are introduced. Scalar functions take finite real or complex values unless
an extended codomain is displayed.

\section{Tests and closure operations}

\begin{definition}[Measurable maps]
\family{D11.04.measurable_map}\label{d11:04:measurable-map}
A map $f:(X,\Sigma)\to(Y,\mathcal T)$ is measurable if
$f^{-1}(B)\in\Sigma$ for every $B\in\mathcal T$.
Scalar measurability uses the Borel sigma-algebra on $\R$, $\C$, or
$\overline\R=[-\infty,\infty]$. A function on a measurable subset $E$
uses the trace sigma-algebra on $E$.
\end{definition}

\begin{theorem}[Generator tests]
\family{D11.04.generator_test}\label{d11:04:generator-test}
If $\mathcal T=\sigma(\mathcal C)$, then $f$ is measurable if and only if
$f^{-1}(C)\in\Sigma$ for every $C\in\mathcal C$.
Real or extended-real measurability is equivalent to measurability of its
strict lower level sets, or of its non-strict lower level sets; rational
thresholds suffice. Upper level tests are equivalent as well.
\end{theorem}
\begin{proof}
The family $\{B\subseteq Y:f^{-1}(B)\in\Sigma\}$ is a sigma-algebra,
because inverse images preserve complements and countable unions.
Containing the generators is therefore sufficient and plainly necessary.
For scalar maps apply the generating-ray description of Borel sets.
For example
$\{f<a\}=\bigcup_{q\in\Q,\ q<a}\{f\leq q\}$ and
$\{f\leq a\}=\bigcap_{q\in\Q,\ q>a}\{f<q\}$.
Complements give upper levels. These identities also hold at points where
$f$ is infinite; the singleton endpoints of $\overline\R$ are obtained
by countable intersections of rays.
\end{proof}

\begin{theorem}[Composition]
\family{D11.04.composition}\label{d11:04:composition}
If $f:(X,\Sigma)\to(Y,\mathcal T)$ and
$g:(Y,\mathcal T)\to(Z,\mathcal U)$ are measurable, then $g\circ f$ is
measurable. In particular continuous scalar maps may be composed after
measurable scalar maps, with the appropriate Borel structures.
\end{theorem}
\begin{proof}
For $C\in\mathcal U$,
$(g\circ f)^{-1}(C)=f^{-1}(g^{-1}(C))\in\Sigma$.
A continuous map has open inverse images of open sets, and hence is Borel
measurable by the generator test. The middle measurable structure in this
calculation is a single specified $\mathcal T$; replacing it by a completion
on only one side is not justified by the formula.
\end{proof}

\begin{theorem}[Coordinate tests]
\family{D11.04.coordinates_complex}\label{d11:04:coordinates-complex}
A map $f:X\to\R^d$ is Borel measurable exactly when each coordinate is
measurable. A map to $\C$ is measurable exactly when its real and imaginary
parts are measurable. Finite tuples of measurable scalar maps are measurable;
modulus and complex conjugation preserve measurability.
\end{theorem}
\begin{proof}
Necessity follows by composing with continuous coordinate projections.
For sufficiency, an inverse image of a rational box is the finite
intersection of coordinate inverse images of intervals. These boxes generate
the target Borel sigma-algebra. Identify $\C$ with $\R^2$.
The maps $(u,v)\mapsto\sqrt{u^2+v^2}$ and $(u,v)\mapsto(u,-v)$ are
continuous, proving the final assertions by composition.
\end{proof}

\begin{theorem}[Algebra and comparisons]
\family{D11.04.scalar_operations}\label{d11:04:scalar-operations}
For measurable finite scalar $f,g$, the functions $af+bg$ and $fg$ are
measurable, and $f/g$ is measurable on $\{g\ne0\}$.
For extended-real measurable $f,g$, the sets $\{f<g\},\{f\leq g\},\{f=g\}$
and the functions $\max(f,g),\min(f,g)$ are measurable. Nonnegative extended
functions are closed under addition, multiplication with $0\cdot\infty=0$,
and multiplication by finite nonnegative constants.
\end{theorem}
\begin{proof}
Finite scalar operations are continuous functions of a finite tuple;
division is continuous on the open set of pairs with nonzero denominator,
and that denominator condition defines a measurable domain. For comparisons,
\[
 \{f<g\}=\bigcup_{q\in\Q}\bigl(\{f<q\}\cap\{g>q\}\bigr).
\]
The formula is valid for extended values too. Complementing $\{g<f\}$ gives
$\{f\leq g\}$, and intersecting the two non-strict comparisons gives equality.
The level sets of a maximum or minimum are finite intersections or unions
of the corresponding level sets.
For nonnegative extended maps use finite truncations $f\wedge n,g\wedge n$.
Their sums converge to $f+g$. Their products converge to the prescribed
extended product: if one value is zero every truncated product is zero;
if both are positive and at least one is infinite the products tend to
infinity; otherwise they converge as finite scalar products. Measurability
of these limits follows from the countable level formulas in the next proof,
which apply independently to any sequence of measurable functions.
No undefined subtraction of two infinities is being used.
\end{proof}

\begin{theorem}[Countable limits]
\family{D11.04.countable_limits}\label{d11:04:countable-limits}
For measurable $f_n:X\to\overline\R$, their pointwise supremum, infimum,
limsup and liminf are measurable. Their finite-limit set is measurable;
the limit there, extended by zero elsewhere, is measurable. Everywhere
finite real or complex pointwise limits of measurable maps are measurable.
\end{theorem}
\begin{proof}
The identity $\{\sup_n f_n>a\}=\bigcup_n\{f_n>a\}$ proves the assertion
for suprema; lower levels prove it for infima. Then
$\limsup_n f_n=\inf_N\sup_{n\geq N}f_n$ and the analogous liminf formula
give their measurability. The finite-limit set is
\[
 D=\{\limsup f_n=\liminf f_n\}\cap\{-\infty<\limsup f_n<\infty\}.
\]
Its measurability uses only the comparison-set formula already proved,
not any closure under products. Restrict the measurable limsup to $D$ and
assign zero on $D^c$; inverse images are unions of the corresponding
measurable pieces. Real and imaginary coordinate limits give the complex case.
These level-set arguments also justify the limit step in the preceding theorem.
\end{proof}

\section{Simple functions and explicit approximations}

\begin{definition}[Simple functions]
\family{D11.04.simple_function}\label{d11:04:simple-function}
A simple scalar function has finite range. It is measurable if and only if
each nonzero level set is measurable, equivalently if it has a representation
$f=\sum_{j=1}^n c_j\one_{E_j}$ with measurable disjoint $E_j$.
A nonnegative simple function has finite nonnegative real values. There is
no finite-measure-support condition in this definition.
\end{definition}
Indeed the zero level is the complement of the finite union of nonzero
levels; every inverse image is a finite union of levels. Conversely singleton
inverse images of a measurable scalar function are measurable. Empty
representations and zero coefficients cause no difficulty.

\begin{theorem}[Simple approximation]
\family{D11.04.simple_approximation}\label{d11:04:simple-approximation}
Every measurable $f:X\to\Eplus$ is an increasing pointwise limit of finite
nonnegative measurable simple functions $s_n\leq f$.
For measurable $f:X\to\overline\R$ one can choose real simple $s_n\to f$
with $|s_n|\leq|s_{n+1}|\leq|f|$; bounded real $f$ has uniform such
approximations. Complex measurable functions have pointwise simple
approximations, uniform when bounded.
\end{theorem}
\begin{proof}
For $f\geq0$ put
\[
 s_n=2^{-n}\left\lfloor 2^n(f\wedge n)\right\rfloor.
\]
At $f=\infty$ this expression means $s_n=n$. Its finitely many level sets
are intersections of measurable rays. Dyadic refinement never lowers the
rounded-down value, and increasing the truncation level cannot lower it,
so $s_n\uparrow f$. For finite $f(x)$, once $n>f(x)$ the error is less than
$2^{-n}$; for infinite $f(x)$ the values tend to infinity.
Apply this construction to $f^+=\max(f,0)$ and $f^-=\max(-f,0)$ and subtract
the two finite approximations. The supports of these two parts are disjoint,
so absolute values increase and stay below $|f|$. For bounded $f$, truncation
eventually has no effect anywhere and the error bound is uniform.
Approximate the two coordinates of a complex function. The Euclidean norm
of the error is at most the sum of the coordinate errors, giving uniform
approximation in the bounded case as well.
\end{proof}

\section{Null exceptions and versions}

\begin{definition}[Almost everywhere]
\family{D11.04.almost_everywhere}\label{d11:04:almost-everywhere}
A property holds $\mu$-almost everywhere if it holds outside some
$N\in\Sigma$ with $\mu(N)=0$. Equality almost everywhere is an equivalence
relation on scalar functions: reflexivity uses the empty exceptional set,
symmetry uses the same set, and transitivity uses the union of two null sets.
Its equivalence classes are not actual functions. On an incomplete space the
failure set itself need not be measurable, although it has a measurable null
envelope.
\end{definition}

\begin{lemma}[A common exceptional set]
\family{D11.04.common_exceptional_set}\label{d11:04:common-exceptional-set}
Countably many a.e. assertions hold simultaneously outside one measurable
null set. A.e. equality is consequently preserved by finite scalar operations
and countable pointwise limits whenever the stated operations are defined.
\end{lemma}
\begin{proof}
Choose a null exceptional set $N_n$ for each assertion and take
$N=\bigcup_nN_n$. Countable subadditivity gives $\mu(N)=0$.
Outside $N$ the assertions are ordinary pointwise identities, where the
operations and limits preserve equality. This proves a.e. equality of the
results; it does not claim measurability of arbitrary values assigned on $N$.
\end{proof}

\begin{theorem}[A.e. changes and measurability]
\family{D11.04.ae_measurability_complete}\label{d11:04:ae-measurability-complete}
If $\mu$ is complete, a function equal a.e. to a measurable extended-real
function is measurable. An a.e. limit of measurable extended-real functions
has a measurable version on any measure space, and the given limit itself
is measurable when the space is complete.
\end{theorem}
\begin{proof}
If $g=f$ off a measurable null $N$, then
$\{g<a\}=(\{f<a\}\setminus N)\cup(\{g<a\}\cap N)$.
Completeness makes the second piece measurable and the generator test applies.
For an a.e. limit choose a measurable null $N$ outside which convergence
holds. Replace each $f_n$ by zero on $N$. The resulting measurable sequence
has an everywhere extended-real limit $h$, equal to the given limit off $N$
and zero on $N$. The countable-limit theorem makes $h$ measurable.
In a complete space the first part applies to the given limit. Coordinate
arguments give the corresponding finite real and complex version statements.
\end{proof}

\begin{theorem}[Versions for a completion]
\family{D11.04.completion_representative}\label{d11:04:completion-representative}
An extended-real function $f$ is $\overline\Sigma$-measurable if and only if
there are $\Sigma$-measurable extended-real $f_0\leq f\leq f_1$ everywhere
with $f_0=f_1$ a.e. In particular it has a $\Sigma$-measurable version.
For finite real or complex $f$ this version can be chosen finite-valued.
\end{theorem}
\begin{proof}
Approximate a completion-measurable extended-real $f$ by real simple $s_n$
as above. For each of their countably many level sets $E_{nj}$ choose
$A_{nj},B_{nj}\in\Sigma$ with
$A_{nj}\subseteq E_{nj}\subseteq B_{nj}$ and
$\mu(B_{nj}\setminus A_{nj})=0$. Replacing the level sets in $s_n$ by
$A_{nj}$ gives a $\Sigma$-measurable simple $t_n$ equal to $s_n$ off the
single measurable null set $N=\bigcup_{n,j}(B_{nj}\setminus A_{nj})$.
Define $h=\limsup_n t_n$ on $N^c$ and $h=0$ on $N$. This is
$\Sigma$-measurable and equals $f$ on $N^c$. Define $f_0=h$ on $N^c$
and $-\infty$ on $N$, and $f_1=h$ on $N^c$ and $+\infty$ on $N$.
These functions give the required bounds and agreement.

Conversely if such bounds exist, choose a measurable null set $N$ outside
which $f_0=f_1$. Then $f=f_0$ there. The completion is complete, so the
preceding a.e.-modification theorem applied in the completed space proves
completion measurability of $f$.
If $f$ is finite real-valued, the constructed $h$ is already finite everywhere:
it equals $f$ off $N$ and is zero on $N$. Apply this to both coordinates
for complex $f$. The bounding pair may need infinite values even though a
finite unbounded $f$ has a finite version; bounding and agreeing a.e. are
different requirements.
\end{proof}

\begin{definition}[Essential bounds]
\family{D11.04.essential_bounds}\label{d11:04:essential-bounds}
For measurable finite scalar $f$ set
\[
 \esssup|f|=\inf\{M\in[0,\infty): |f|\leq M\ \mu\text{-a.e.}\},
 \qquad \inf\varnothing=\infty.
\]
The function is essentially bounded if this number is finite. Changing $f$
on a measurable null envelope does not change the set of admissible bounds.
This is the ordinary a.e. convention, not the locally-null convention used
in a separate Cohn variant later.
\end{definition}

\begin{theorem}[Attainment of the essential bound]
\family{D11.04.essential_bound_attained}\label{d11:04:essential-bound-attained}
If $M=\esssup|f|<\infty$, then $|f|\leq M$ a.e.
Also $\esssup|f|=0$ iff $f=0$ a.e., and $|f|\leq|g|$ a.e. implies
$\esssup|f|\leq\esssup|g|$.
\end{theorem}
\begin{proof}
Choose admissible bounds $M_n<M+1/n$. Outside the union of their null
exceptional sets, $|f|\leq M_n$ for all $n$, hence $|f|\leq M$.
This proves the zero assertion in one direction, and $f=0$ a.e. makes zero
an admissible bound in the other. If $|f|\leq|g|$ a.e., every admissible
bound for $g$ is one for $f$ after adjoining the comparison's null exception.
Taking infima proves the inequality, including infinite bounds.
\end{proof}

\section{Why the qualifications matter}

\begin{example}[An incomplete-space modification]
\family{D11.04.incomplete_modification_failure}\label{d11:04:incomplete-modification-failure}
Let $N\subseteq C$ be the non-Borel Cantor subset from
\cref{d11:03:lebesgue-not-borel}. On the Borel Lebesgue space, $\one_N$
equals zero outside the Borel null set $C$ but is not Borel measurable:
the inverse image of $(1/2,\infty)$ is $N$. Thus an arbitrary null modification
need not be measurable before completion.
\end{example}

\begin{example}[Uncountable suprema]
\family{D11.04.uncountable_sup_failure}\label{d11:04:uncountable-sup-failure}
Each $\one_{\{x\}}$, $x\in N$, is Borel measurable, but their supremum
is $\one_N$, which is not. Thus the countability in the limit theorem is
a real hypothesis, even for nonnegative functions bounded by one.
\end{example}

\begin{example}[A measurable modulus is insufficient]
\family{D11.04.measurable_modulus_insufficient}\label{d11:04:measurable-modulus-insufficient}
Let $V\subseteq[0,1]$ be non-Lebesgue-measurable and put
$f=2\one_V-1$ on $[0,1]$. Then $|f|=1$ is measurable and has finite
integral by the simple-integral definition of the next chapter. But $f$
is not measurable, since $\{f>0\}=V$. A comparison with an integrable
modulus never licenses dropping measurability of the scalar function itself.
\end{example}

\begin{remark}[The middle sigma-algebra is part of composition]
\family{D11.04.lebesgue_composition_boundary}\label{d11:04:lebesgue-composition-boundary}
Continuity of an inner map ensures measurable inverse images of Borel sets,
not automatically of every subset of a target null set. Testing a composition
with $g=\one_A$ makes this precise: measurability of all such compositions
requires that every admitted target measurable $A$ have measurable inverse
image. For a continuous strictly increasing $h:[a,b]\to[c,d]$ the exact
additional criterion, proved in \cref{d11:09:von-neumann-composition}, is
absolute continuity of $h^{-1}$. That independent later theorem and
\cref{d11:09:composition-measurability-failure} supply the refined criterion
and counterexample; neither is being assumed to prove the general composition
theorem here. This is a forward interface marker, not a second copy of that proof.
\end{remark}
\source{COH 2.1--2.2 and 2.6; AXL 2E and 7A; ROY 6.6 for the later composition boundary.}

% END chapters/04


% BEGIN chapters/05
\chapter{The nonnegative integral}\label{ch:d11-05}
Throughout this chapter $(X,\Sigma,\mu)$ is an arbitrary positive measure
space. Neither finiteness, sigma-finiteness nor completeness is implicit.
All sums and integrals of nonnegative quantities take values in $\Eplus$;
in particular $0\cdot\infty=\infty\cdot0=0$.

\begin{definition}[Integral of a simple function]
\family{D11.05.simple_integral}\label{d11:05:simple-integral}
If $s=\sum_{j=1}^m a_j\one_{E_j}$, where the $E_j\in\Sigma$ are pairwise
disjoint and $0\leq a_j<\infty$, set
\[
 I_\mu(s)=\sum_{j=1}^m a_j\mu(E_j).
\]
One may include the complement of the union among the sets with coefficient
zero. Thus this definition applies to every nonnegative measurable simple
function, without any restriction on the measure of its support.
\end{definition}

\begin{lemma}[Independence and elementary algebra]
\family{D11.05.simple_integral_well_defined}\label{d11:05:simple-integral-well-defined}
The value $I_\mu(s)$ is independent of the disjoint representation.
For nonnegative measurable simple $s,t$ and finite $a\geq0$,
\[
 I_\mu(as)=aI_\mu(s),\qquad I_\mu(s+t)=I_\mu(s)+I_\mu(t).
\]
If $s\leq t$, then $I_\mu(s)\leq I_\mu(t)$.
\end{lemma}
\begin{proof}
Add the zero levels so that each representation partitions $X$, and refine
two representations by all intersections of their level sets. Finite
additivity of $\mu$ shows that replacing any level by this finite disjoint
partition preserves its contribution. On every nonempty intersection the two
coefficients representing the same function agree; empty intersections have
measure zero. This proves independence. On a common finite measurable
partition for $s$ and $t$, the three assertions reduce respectively to
$a(b\mu A)=(ab)\mu A$, $(b+c)\mu A=b\mu A+c\mu A$, and
$b\leq c\Rightarrow b\mu A\leq c\mu A$. These identities remain true when
$\mu A=\infty$, with the stipulated zero convention. Finite summation gives
the conclusions.
\end{proof}

\begin{definition}[Nonnegative integral and its approximation property]
\family{D11.05.nonnegative_integral}\label{d11:05:nonnegative-integral}
For a measurable $f:X\to\Eplus$ define
\[
 \int f\,d\mu=\sup\{I_\mu(s):s\text{ measurable simple},\ 0\leq s\leq f\}.
\]
This agrees with $I_\mu$ on simple functions. If $s_n$ are any nonnegative
measurable simple functions with $s_n\uparrow f$, then
\[
 \int f\,d\mu=\lim_n I_\mu(s_n).
\]
\end{definition}
\begin{proof}
For simple $f$, the supremum is at most $I_\mu(f)$ by the preceding lemma
and attains it at $f$. For a general increasing simple approximation let
$L=\sup_n I_\mu(s_n)\leq\int f\,d\mu$. Fix a simple
$s=\sum_{i=1}^m a_i\one_{A_i}\leq f$ with disjoint levels, and $0<c<1$.
Ignore its zero coefficients and put
$A_{i,n}=A_i\cap\{s_n\geq ca_i\}$. These sets increase to $A_i$:
at a point of $A_i$, the limit $f\geq a_i>ca_i$.
Since $s_n\geq\sum_i ca_i\one_{A_{i,n}}$, continuity from below yields
\[
 L\geq c\lim_n\sum_i a_i\mu(A_{i,n})=cI_\mu(s).
\]
The equality of the limit and the finite sum follows by taking limits of
each nonnegative term; if one limit is infinite both sides are infinite.
Letting $c\uparrow1$ and then taking the supremum over $s$ gives
$L\geq\int f\,d\mu$. This proves the property directly from the definition;
it does not presuppose the monotone convergence theorem below.
\end{proof}

\begin{theorem}[Algebra and restriction]
\family{D11.05.nonnegative_algebra}\label{d11:05:nonnegative-algebra}
For measurable $f,g:X\to\Eplus$ and finite $a\geq0$,
\[
 \int(f+g)\,d\mu=\int f\,d\mu+\int g\,d\mu,
 \qquad \int af\,d\mu=a\int f\,d\mu.
\]
Pointwise order implies order of integrals. For $E\in\Sigma$ write
$\int_E f\,d\mu=\int\one_Ef\,d\mu$; this is also the integral of $f|_E$
on the trace space with the restricted measure.
\end{theorem}
\begin{proof}
Order follows because every simple minorant of $f$ is a minorant of $g$
when $f\leq g$. Choose increasing simple approximations $s_n\uparrow f$
and $t_n\uparrow g$. Then $s_n+t_n\uparrow f+g$, so the approximation
property and simple additivity give
$\int(f+g)=\lim_n(I_\mu(s_n)+I_\mu(t_n))=\int f+\int g$.
For $a>0$ use $as_n\uparrow af$; for $a=0$ both sides are zero by convention.
For restriction, the defining simple sums on $E$ are exactly the sums of the
extensions by zero to $X$. They give equality first for simple functions and
then, using increasing approximations, for every nonnegative $f$.
\end{proof}

\begin{theorem}[Null sets and order almost everywhere]
\family{D11.05.null_invariance}\label{d11:05:null-invariance}
If measurable $f,g:X\to\Eplus$ agree $\mu$-almost everywhere, their
integrals agree. For every measurable null set $N$, $\int_N f\,d\mu=0$,
even if $f$ is infinite there. If $f\leq g$ almost everywhere, then
$\int f\,d\mu\leq\int g\,d\mu$.
\end{theorem}
\begin{proof}
A simple function supported on $N$ has integral zero. Every simple
minorant of $\one_Nf$ is supported there except possibly on zero levels,
so the supremum defining its integral is zero. Decompose
$f=\one_Nf+\one_{X\setminus N}f$ and apply additivity. If $f=g$ off $N$,
the second terms coincide; if $f\leq g$ there, their second terms are
ordered. This proves both assertions. All functions used are measurable;
no subsets of $N$ have been declared measurable by assumption.
\end{proof}

\begin{theorem}[Monotone convergence]
\family{D11.05.monotone_convergence}\label{d11:05:monotone-convergence}
Let $f_n,f:X\to\Eplus$ be measurable. If, outside one measurable null
set, $f_n\leq f_{n+1}$ for every $n$ and $f_n\to f$, then
\[
 \int f\,d\mu=\lim_n\int f_n\,d\mu=\sup_n\int f_n\,d\mu.
\]
\end{theorem}
\begin{proof}
Reset all functions to zero on the stated null set. Their integrals are
unchanged, so assume pointwise increase everywhere. Put
$L=\sup_n\int f_n\leq\int f$. Given a disjoint simple
$s=\sum_i a_i\one_{A_i}\leq f$ and $0<c<1$, put
$A_{i,n}=A_i\cap\{f_n\geq ca_i\}$ for $a_i>0$.
They increase to $A_i$, and
\[
 \int f_n\,d\mu\geq c\sum_i a_i\mu(A_{i,n}).
\]
Continuity from below gives $L\geq cI_\mu(s)$. Taking $c\uparrow1$ and
the supremum over simple $s$ proves the reverse inequality. Since the
integrals are increasing even before the reset, their supremum is their
limit. This proof requires the measurability of the stated limit $f$;
on an incomplete space arbitrary changes of its null-set values are not
automatically measurable.
\end{proof}

\begin{theorem}[Nonnegative series]
\family{D11.05.beppo_levi}\label{d11:05:beppo-levi}
For any measurable $f_n:X\to\Eplus$,
\[
 \int\sum_{n=1}^\infty f_n\,d\mu
   =\sum_{n=1}^\infty\int f_n\,d\mu.
\]
Consequently, if $E_n$ are pairwise disjoint measurable sets, then
$\int_{\bigcup_nE_n}f\,d\mu=\sum_n\int_{E_n}f\,d\mu$ for every
nonnegative measurable $f$.
\end{theorem}
\begin{proof}
The partial sums are nonnegative measurable and increase to the series.
Finite additivity of the integral and monotone convergence prove the first
formula. For the second apply it to $f_n=\one_{E_n}f$; their sum is
$\one_{\bigcup_nE_n}f$, also at points where $f=\infty$ by the zero convention.
\end{proof}

\begin{theorem}[Fatou's inequality]
\family{D11.05.fatou}\label{d11:05:fatou}
For every sequence of nonnegative measurable $f_n$,
\[
 \int\liminf_n f_n\,d\mu\leq\liminf_n\int f_n\,d\mu.
\]
If a measurable $f\geq0$ satisfies $f\leq\liminf_n f_n$ almost everywhere,
then $\int f\,d\mu\leq\liminf_n\int f_n\,d\mu$ as well.
\end{theorem}
\begin{proof}
The measurable functions $h_n=\inf_{k\geq n}f_k$ increase to
$\liminf f_n$. Since $h_n\leq f_k$ for $k\geq n$,
$\int h_n\leq\inf_{k\geq n}\int f_k$. Monotone convergence and passage
to the increasing limits prove the first inequality. The almost-everywhere
order property proves the second.
\end{proof}

\begin{theorem}[Markov's inequality]
\family{D11.05.markov}\label{d11:05:markov}
If $f:X\to\Eplus$ is measurable and $t>0$, then
\[
 t\mu\{f\geq t\}\leq\int_{\{f\geq t\}}f\,d\mu
 \leq\int f\,d\mu.
\]
Thus $\mu\{f\geq t\}\leq t^{-1}\int f\,d\mu$, as an extended inequality.
\end{theorem}
\begin{proof}
The pointwise inequalities
$t\one_{\{f\geq t\}}\leq\one_{\{f\geq t\}}f\leq f$ imply the result
by order and the simple integral formula. Multiplication by the positive
finite number $t^{-1}$ gives the last form, including infinite values.
\end{proof}

\begin{theorem}[What zero and finite integrals imply]
\family{D11.05.zero_integral}\label{d11:05:zero-integral}
For measurable $f:X\to\Eplus$, $\int f\,d\mu=0$ if and only if $f=0$
almost everywhere. If $\int f\,d\mu<\infty$, then $f<\infty$ almost
everywhere and $\{f>0\}$ is a countable union of measurable sets of finite
measure.
\end{theorem}
\begin{proof}
If the integral is zero, Markov's inequality gives
$\mu\{f\geq1/n\}=0$ for every $n$. Their union is $\{f>0\}$, hence is
null. The converse is null invariance. Write $I=\int f<\infty$.
For every integer $n\geq1$, $\{f=\infty\}\subseteq\{f\geq n\}$ and
$\mu\{f\geq n\}\leq I/n$, so $\mu\{f=\infty\}=0$.
Moreover $\mu\{f\geq1/n\}\leq nI<\infty$ and these sets cover $\{f>0\}$.
\end{proof}

\begin{construction}[A measure with density]
\family{D11.05.weighted_measure}\label{d11:05:weighted-measure}
For measurable $h:X\to\Eplus$, the prescription
\[
 \nu(E)=\int_E h\,d\mu\qquad(E\in\Sigma)
\]
defines a positive measure. It is finite when $\int h\,d\mu<\infty$ and
vanishes on every measurable $\mu$-null set. No integrability assumption
is needed to construct it.
\end{construction}
\begin{proof}
The integral over the empty set is zero. On a countable disjoint family,
the second assertion of the nonnegative-series theorem is exactly countable
additivity of $\nu$. Its total mass is $\int h$. The null-set assertion is
\cref{d11:05:null-invariance}.
\end{proof}

\begin{theorem}[Atomic integrals]
\family{D11.05.atomic_integrals}\label{d11:05:atomic-integrals}
For nonnegative measurable $f$, $\int f\,d\delta_x=f(x)$.
On a countable set with its full power-set sigma-algebra and counting
measure, $\int f=\sum_x f(x)$. More generally, if
$\mu=\sum_n a_n\delta_{x_n}$ with $a_n\in\Eplus$, then
\[
 \int f\,d\mu=\sum_n a_nf(x_n).
\]
Repeated points and infinite weights are allowed.
\end{theorem}
\begin{proof}
For a disjoint simple $s=\sum_jb_j\one_{E_j}$ its integral with respect
to the last measure is
\[
 \sum_j b_j\sum_n a_n\one_{E_j}(x_n)
 =\sum_n a_n\sum_j b_j\one_{E_j}(x_n)=\sum_n a_ns(x_n).
\]
The finite interchange is valid in nonnegative extended arithmetic; when
$b_j=0$ all the corresponding summands are zero. Take $s_k\uparrow f$.
For each $n$, $a_ns_k(x_n)\uparrow a_nf(x_n)$: if $a_n=\infty$ and
$f(x_n)>0$, a sufficiently late value is positive, and otherwise every
term is zero. The increasing limit commutes with the nonnegative sum,
because both orders give the supremum of sums over finite sets of indices
and finite approximation levels. The approximation property proves the
formula. The Dirac and counting assertions are its special cases.
\end{proof}

\begin{theorem}[Continuity with respect to the domain]
\family{D11.05.continuity_domains}\label{d11:05:continuity-domains}
If $E_n\uparrow E$ are measurable and $f\geq0$ is measurable, then
$\int_{E_n}f\uparrow\int_Ef$. If $E_n\downarrow E$ and
$\int_{E_N}f<\infty$ for some $N$, then $\int_{E_n}f\downarrow\int_Ef$.
\end{theorem}
\begin{proof}
Apply continuity from below or continuity from above to the positive
measure $\nu(A)=\int_Af\,d\mu$. Its finite-mass hypothesis in the second
case is exactly the hypothesis displayed here.
\end{proof}

\begin{example}[Strict Fatou inequality through escape of mass]
\family{D11.05.fatou_strict_mass_escape}\label{d11:05:fatou-strict-mass-escape}
On $\R$ let $f_n=\one_{[n,n+1)}$. Every fixed point is eventually outside
these intervals, so $f_n\to0$ pointwise. Each integral is the interval
length $1$. Therefore
\[
 \int\liminf_n f_n\,d\lambda=0<1=\liminf_n\int f_n\,d\lambda.
\]
Nonnegativity and pointwise convergence by themselves do not justify
interchanging a limit and an integral.
\end{example}
\source{COH 2.3--2.4; FOL 2.2; AXL 4B--4C; ROY 4.2 and 18.1--18.2.}

% END chapters/05


% BEGIN chapters/06
\chapter{Absolutely integrable scalar functions}\label{ch:d11-06}
Nonnegative integrals are allowed to be infinite. A real or complex integral
in this chapter is a finite scalar, and is defined only under absolute
integrability. This distinction prevents undefined subtractions of infinities.

\begin{definition}[Absolute integrability]
\family{D11.06.absolute_integrability}\label{d11:06:absolute-integrability}
A measurable $f:X\to\R$ or $\C$ is integrable if
$\int|f|\,d\mu<\infty$. For real $f$ write
$f^+=\max(f,0)$ and $f^-=\max(-f,0)$, so that
$|f|=f^++f^-$ and $f=f^+-f^-$. For an extended-real measurable $f$ with
$\int|f|<\infty$, replace $f$ by zero on the measurable null set
$\{|f|=\infty\}$ before assigning it a scalar integral. The resulting
finite version is unique almost everywhere.
\end{definition}

\begin{construction}[Real and complex integrals]
\family{D11.06.scalar_integral}\label{d11:06:scalar-integral}
For real integrable $f$ set
\[
 \int f\,d\mu=\int f^+\,d\mu-\int f^-\,d\mu.
\]
Real integrability is equivalent to finiteness of both terms on the right.
For complex integrable $f$ set
\[
 \int f\,d\mu=\int\operatorname{Re}f\,d\mu
             +i\int\operatorname{Im}f\,d\mu.
\]
Complex integrability is equivalent to integrability of both coordinates.
These definitions give the usual finite scalar sum for an integrable
measurable simple function.
\end{construction}
\begin{proof}
Nonnegative additivity gives $\int|f|=\int f^++\int f^-$, proving the
real equivalence and ensuring that the displayed difference is finite.
The inequalities
$|\operatorname{Re}f|,|\operatorname{Im}f|\leq|f|
\leq|\operatorname{Re}f|+|\operatorname{Im}f|$
prove the complex equivalence. On each disjoint level of a simple function,
the positive and negative parts split its real coefficient, and the two
coordinates split a complex coefficient. Finite summation therefore gives
the claimed scalar sum. A nonzero coefficient on a level of infinite
measure would violate integrability, so no infinite scalar product occurs.
\end{proof}

\begin{theorem}[Linearity, order and absolute-value bound]
\family{D11.06.linearity_order_norm}\label{d11:06:linearity-order-norm}
If scalar functions $f,g$ are integrable and $a,b$ are scalars in their
field, then $af+bg$ is integrable and
\[
 \int(af+bg)\,d\mu=a\int f\,d\mu+b\int g\,d\mu.
\]
For real $f\leq g$ almost everywhere, $\int f\leq\int g$. In both fields,
$\big|\int f\,d\mu\big|\leq\int|f|\,d\mu$.
\end{theorem}
\begin{proof}
The bound $|af+bg|\leq|a||f|+|b||g|$ proves integrability. For real
$f,g$, the pointwise identity
\[
 (f+g)^++f^-+g^-=(f+g)^-+f^++g^+
\]
has integrable nonnegative terms. Integrate and rearrange the finite
numbers to obtain additivity. Positive real homogeneity follows by taking
parts, negative homogeneity by interchanging the parts, and zero
homogeneity directly. A real nonnegative function has nonnegative integral;
applying this to $g-f$ off a measurable null exceptional set proves order,
since the integral of each part over that null set is zero.
Coordinate definitions give complex additivity and real homogeneity;
$\int if=i\int f$ follows by interchanging the two coordinates with a
sign. Combining real and imaginary scalar parts proves complex linearity.
For real $f$, $-|f|\leq f\leq|f|$ proves the absolute-value bound.
For complex $f$ with nonzero integral choose $z\in\C$, $|z|=1$, such that
$z\int f=|\int f|$. Then
\[
 \left|\int f\right|=\int\operatorname{Re}(zf)\,d\mu
 \leq\int|f|\,d\mu.
\]
When the integral is zero the bound is immediate.
\end{proof}

\begin{theorem}[Integrable domination]
\family{D11.06.integrable_domination}\label{d11:06:integrable-domination}
If $f$ is measurable and $|f|\leq g$ almost everywhere for an integrable
measurable $g\geq0$, then $f$ is integrable. In particular, if $f=0$ off
a measurable $E$ of finite measure and $|f|\leq M$ almost everywhere,
where $M<\infty$, then $\int|f|\leq M\mu(E)$.
If $f$ is integrable and measurable $h$ is essentially bounded, then
$hf$ is integrable and
\[
 \int|hf|\,d\mu\leq(\esssup|h|)\int|f|\,d\mu.
\]
\end{theorem}
\begin{proof}
Order almost everywhere for nonnegative integrals gives
$\int|f|\leq\int g<\infty$. The finite-support assertion uses the
dominator $M\one_E$. For the last assertion the essential bound is attained
as an almost-everywhere bound by \cref{d11:04:essential-bound-attained};
use $|hf|\leq(\esssup|h|)|f|$ off its null exceptional set and nonnegative
homogeneity. This also covers essential bound zero.
\end{proof}

\begin{theorem}[Almost-everywhere invariance and restriction]
\family{D11.06.scalar_ae_invariance}\label{d11:06:scalar-ae-invariance}
If measurable scalar $f,g$ agree almost everywhere, one is integrable if
and only if the other is, and their integrals then agree. An integrable
function remains integrable on every measurable subset. Removing a
measurable null subset changes neither its integral nor its absolute integral.
\end{theorem}
\begin{proof}
The absolute values agree almost everywhere, so their nonnegative
integrals coincide. In the real case the positive and negative parts also
agree almost everywhere, and in the complex case so do the two
coordinates. Their defining integrals therefore agree. Finally
$|\one_Ef|\leq|f|$, proving restriction, and the null-set assertion follows
by applying the first assertion to $f$ and $\one_{X\setminus N}f$.
\end{proof}

\begin{theorem}[Small sets and finite-measure tails of one integral]
\family{D11.06.absolute_continuity_integral}\label{d11:06:absolute-continuity-integral}
If $f$ is integrable, then for every $\eps>0$ there is $\delta>0$ such
that $E\in\Sigma$, $\mu(E)<\delta$ imply $\int_E|f|<\eps$.
Also, for every $\eps>0$ there is $A\in\Sigma$, $\mu(A)<\infty$, with
$\int_{X\setminus A}|f|<\eps$. These assertions hold on every positive
measure space.
\end{theorem}
\begin{proof}
The supremum definition and finiteness of $\int|f|$ provide a nonnegative
simple $s\leq|f|$ for which
$\int|f|-\int s<\eps/2$. Since $|f|=s+(|f|-s)$, nonnegative additivity
gives $\int(|f|-s)<\eps/2$. Choose finite $M\geq0$ bounding $s$ and set
$\delta=\eps/(2(M+1))$. Then
\[
 \int_E|f|\leq M\mu(E)+\int(|f|-s)<\eps
 \quad\hbox{if }\mu(E)<\delta.
\]
Every positive level of $s$ has finite measure, since its positive
coefficient times that measure is bounded by $\int|f|$. Thus
$A=\{s>0\}$ is a finite union of finite-measure sets and
$\int_{X\setminus A}|f|\leq\int(|f|-s)<\eps$.
The argument uses no convergence theorem for signed functions.
\end{proof}

\begin{theorem}[Dominated convergence]
\family{D11.06.dominated_convergence}\label{d11:06:dominated-convergence}
Let measurable $f_n,f:X\to\R$ or $\C$ satisfy $f_n\to f$ almost
everywhere. Suppose $g\geq0$ is integrable and $|f_n|\leq g$ almost
everywhere for every $n$. Then $f_n,f$ are integrable and
\[
 \int|f_n-f|\,d\mu\longrightarrow0,
 \qquad \int f_n\,d\mu\longrightarrow\int f\,d\mu.
\]
No finiteness, completeness or sigma-finiteness assumption on $\mu$ is needed.
\end{theorem}
\begin{proof}
Take one measurable null set containing the convergence failures, every
domination failure and $\{g=\infty\}$, and reset the functions and $g$ to
zero there. This preserves their integrals and yields finite $g$, pointwise
convergence, and $|f_n|,|f|\leq g$ everywhere. Integrability follows by
domination. For real functions Fatou applied to $g+f_n$ and $g-f_n$ gives
\[
 \int g+\int f\leq\int g+\liminf_n\int f_n,
 \qquad
 \int g-\int f\leq\int g-\limsup_n\int f_n.
\]
All scalar integrals here lie between $-\int g$ and $\int g$, so finite
subtraction is legitimate. These two inequalities prove integral
convergence. Apply that real conclusion to $|f_n-f|$, which converges to
zero and is bounded by $2g$, to obtain the asserted absolute-integral
convergence. For complex functions first use the real conclusion for both
coordinates; the same real conclusion for the real nonnegative functions
$|f_n-f|\leq2g$ again proves convergence in absolute integral. Alternatively
the scalar integral convergence then follows from the norm bound.
\end{proof}

\begin{corollary}[Bounded convergence on a fixed finite domain]
\family{D11.06.bounded_convergence}\label{d11:06:bounded-convergence}
If $\mu(X)<\infty$, scalar $f_n,f$ are measurable, $f_n\to f$ almost everywhere, and
$|f_n|\leq M$ almost everywhere for one finite $M$, then
$\int|f_n-f|\to0$ and $\int f_n\to\int f$. The same conclusion holds
when the $f_n$ are all supported in one fixed measurable finite-measure set,
with the same measurable limit and the same uniform almost-everywhere bound.
\end{corollary}
\begin{proof}
Use dominated convergence with $g=M\one_X$ in the first case and
$g=M\one_E$ in the second. In the second case the limit is zero off $E$
almost everywhere, which is all the theorem needs.
\end{proof}

\begin{theorem}[Absolutely summable series under the integral]
\family{D11.06.absolute_series_interchange}\label{d11:06:absolute-series-interchange}
If measurable scalar $f_n$ satisfy $\sum_n\int|f_n|\,d\mu<\infty$, then
$\sum_n|f_n(x)|<\infty$ almost everywhere. Define $f=\sum_nf_n$ where
the series is absolutely convergent and $f=0$ elsewhere. Then $f$ is
measurable and integrable, its partial sums converge to $f$ in absolute
integral, and
\[
 \int f\,d\mu=\sum_n\int f_n\,d\mu,
\]
where the scalar series on the right converges absolutely.
\end{theorem}
\begin{proof}
By the nonnegative-series theorem, $h=\sum_n|f_n|$ has finite integral.
Its infinite-value set is measurable and null. Off that set the scalar
series is Cauchy and converges by completeness of $\R$ or $\C$.
The measurable-limit construction, with value zero on the exceptional set,
makes $f$ measurable. Its modulus and every partial-sum modulus are bounded
by $h$ almost everywhere, so dominated convergence proves integrability
and convergence in absolute integral. Finite linearity identifies the
partial-sum integrals, and the norm bound gives
$\sum_n|\int f_n|\leq\sum_n\int|f_n|<\infty$.
One can also read off the error bound
$\int|f-\sum_{n\leq N}f_n|\leq\sum_{n>N}\int|f_n|$ directly.
\end{proof}

\begin{theorem}[Additivity and continuity in the domain]
\family{D11.06.scalar_domain_additivity}\label{d11:06:scalar-domain-additivity}
For integrable $f$ and pairwise disjoint measurable $E_n$,
\[
 \int_{\bigcup_nE_n}f\,d\mu=\sum_n\int_{E_n}f\,d\mu,
\]
with an absolutely convergent scalar series. If measurable $E_n$ increase
or decrease to $E$, then $\int_{E_n}f\to\int_Ef$, without an additional
finite-measure hypothesis.
\end{theorem}
\begin{proof}
Nonnegative domain additivity gives
$\sum_n\int|\one_{E_n}f|=\int_{\bigcup_nE_n}|f|\leq\int|f|<\infty$.
Apply the preceding series theorem to these functions. For monotone sets,
their indicators converge pointwise to $\one_E$; the functions
$\one_{E_n}f$ are all bounded in modulus by the integrable $|f|$.
Dominated convergence proves the second assertion in both directions.
\end{proof}

\begin{theorem}[Detecting order and equality by integrals]
\family{D11.06.integral_order_detection}\label{d11:06:integral-order-detection}
If real integrable $f$ satisfies $\int_Ef\,d\mu\geq0$ for every
$E\in\Sigma$, then $f\geq0$ almost everywhere. Consequently, if scalar
integrable $f,g$ have $\int_Ef=\int_Eg$ for every measurable $E$, then
$f=g$ almost everywhere.
\end{theorem}
\begin{proof}
On $A=\{f<0\}$, $f=-f^-$. The hypothesis gives
$0\leq\int_Af=-\int f^-\leq0$. Thus $\int f^-=0$, whence $f^-=0$
almost everywhere by the zero-integral criterion. For real $f,g$ apply
this assertion to both $f-g$ and $g-f$. For complex $f,g$, the equality
of integrals implies equality of the real and imaginary coordinate
integrals on every set. The real conclusion applies to each coordinate;
the union of the two exceptional null sets is null.
\end{proof}

\begin{theorem}[Differentiating a parameter integral]
\family{D11.06.parameter_differentiation}\label{d11:06:parameter-differentiation}
Let $I\subseteq\R$ be an open interval and $F:X\times I\to\R$.
Assume each $F(\cdot,t)$ is measurable and integrable, each
$F(x,\cdot)$ is differentiable on $I$, and there is integrable $g\geq0$
such that, for every $x\in X$ and all distinct $s,t\in I$,
\[
 \left|\frac{F(x,t)-F(x,s)}{t-s}\right|\leq g(x).
\]
Then $H(t)=\int F(x,t)\,d\mu(x)$ is differentiable on $I$,
$\partial_tF(\cdot,t)$ is measurable and integrable, and
\[
 H'(t)=\int\partial_tF(x,t)\,d\mu(x).
\]
\end{theorem}
\begin{proof}
Fix $t\in I$ and choose nonzero $h_k\to0$ with $t+h_k\in I$.
The difference quotients are measurable and converge pointwise to
$\partial_tF(\cdot,t)$, so this derivative is measurable. It is bounded
in modulus by $g$, hence integrable. For an arbitrary sequence of nonzero
$h_k\to0$ remaining in $I$, dominated convergence and linearity give
\[
 \frac{H(t+h_k)-H(t)}{h_k}
 =\int\frac{F(x,t+h_k)-F(x,t)}{h_k}\,d\mu(x)
 \longrightarrow\int\partial_tF(x,t)\,d\mu(x).
\]
The sequential criterion for a real limit now proves the derivative
identity. Indeed, failure of that limit would supply a sequence of
increments tending to zero on which a fixed positive deviation persists,
contrary to what was just proved. No joint measurability of $F$ is used.
\end{proof}

\begin{theorem}[Continuity of a parameter integral]
\family{D11.06.parameter_continuity}\label{d11:06:parameter-continuity}
Let $I=[a,b]$, $a<b$, and let $F:X\times I\to\C$ have measurable
integrable sections $F(\cdot,t)$. Suppose $|F(x,t)|\leq g(x)$ for every
$x,t$, where $g\geq0$ is integrable. If $F(x,t)\to F(x,t_0)$ as
$t\to t_0$ within $I$ for every $x$, then
$\int F(\cdot,t)\to\int F(\cdot,t_0)$. In particular continuous
parameter sections yield a continuous parameter integral on $I$.
\end{theorem}
\begin{proof}
For every sequence $t_n\to t_0$ in $I$, dominated convergence gives the
claimed convergence along that sequence. If continuity at $t_0$ failed,
points at distance less than $1/n$ from $t_0$ with a fixed positive
integral deviation would form a contradicting sequence. This also proves
the one-sided assertion at an endpoint.
\end{proof}

\begin{example}[Two failures without an integrable dominator]
\family{D11.06.domination_counterexamples}\label{d11:06:domination-counterexamples}
On $\R$, $f_n=n^{-1}\one_{(0,n)}$ satisfies $\sup_x|f_n(x)|=1/n$,
yet $\int f_n=1$ for every $n$. Uniform convergence on an infinite-measure
space need not give integral convergence. On $(0,1)$,
$g_n=n\one_{(0,1/n)}$ tends to zero at every point and again has integral
$1$. Thus finiteness of the domain and pointwise convergence alone are
insufficient. Neither sequence has a common integrable nonnegative
dominator, since such a dominator would make dominated convergence
contradict its constant integrals.
\end{example}

\begin{remark}[The interface with uniform integrability]
\family{D11.06.vitali_route}\label{d11:06:vitali-route}
One integrable dominator supplies two uniform controls. If $|f_n|\leq g$
almost everywhere, then
\[
 \sup_n\int_E|f_n|\leq\int_Eg,
 \qquad
 \sup_n\int_{X\setminus A}|f_n|\leq\int_{X\setminus A}g.
\]
The small-set and finite-measure-tail assertions of
\cref{d11:06:absolute-continuity-integral} make the right sides arbitrarily
small with one choice of $\delta$ and one choice of $A$, respectively,
independent of $n$. Chapter~\ref{ch:d11-17} distinguishes the definitions
and proves the corresponding Vitali theorems. Those theorems are not
prerequisites for the dominated-convergence proof above.
\end{remark}
\source{COH 2.4--2.5; FOL 2.3; AXL 4C--4D; ROY 4.4--4.5 and 18.2.}

% END chapters/06


% BEGIN chapters/07
\chapter{Modes of convergence}\label{ch:d11-07}
All scalar functions in this chapter are finite valued and measurable, unless
a different range is explicitly stated. Global convergence in measure is
used literally; on an infinite space it must not be silently replaced by
convergence on each finite-measure subset.

\begin{definition}[Pointwise, uniform and measure convergence]
\family{D11.07.convergence_modes}\label{d11:07:convergence-modes}
For $f_n,f:X\to\R$ or $\C$, pointwise convergence means convergence at
each $x$, and almost-everywhere convergence means convergence outside a
measurable null set. Uniform convergence means that for every $\eps>0$
there is $N$ such that $|f_n(x)-f(x)|<\eps$ for every $n\geq N$ and every
$x\in X$. Almost-uniform convergence means that for every $\delta>0$
there is $E\in\Sigma$ with $\mu(E)<\delta$ such that convergence is
uniform on $X\setminus E$. Convergence in measure means
\[
 \mu\{|f_n-f|>\eps\}\longrightarrow0\quad(\eps>0).
\]
The sequence is Cauchy in measure if, for every $\eps,\delta>0$, some $N$
satisfies $\mu\{|f_n-f_m|>\eps\}<\delta$ for all $m,n\geq N$.
\end{definition}

\begin{definition}[$L^p$ convergence on representatives]
\family{D11.07.lp_convergence}\label{d11:07:lp-convergence}
For $1\leq p<\infty$, $f_n\to f$ in $L^p$ means that all these functions
have finite integrals of their $p$th absolute powers and
$\int|f_n-f|^p\,d\mu\to0$. For $p=\infty$ they are essentially bounded
and $\esssup|f_n-f|\to0$. The quotient spaces and their norms will be
constructed in Chapter~\ref{ch:d11-10}; the convergence assertions here
can already be stated using representatives.
\end{definition}

\begin{theorem}[Norm convergence implies convergence in measure]
\family{D11.07.norm_to_measure}\label{d11:07:norm-to-measure}
For $1\leq p<\infty$,
\[
 \mu\{|f_n-f|>\eps\}\leq\eps^{-p}\int|f_n-f|^p\,d\mu
 \qquad(\eps>0).
\]
Consequently $L^p$ convergence, including $p=\infty$, implies global
convergence in measure on any positive measure space.
\end{theorem}
\begin{proof}
On the displayed set $|f_n-f|^p>\eps^p$, so Markov's inequality applied
to $|f_n-f|^p$ gives the bound. In the essential-supremum case, for each
$\eps>0$ and all sufficiently large $n$ the attained almost-everywhere
bound on $|f_n-f|$ is smaller than $\eps$. The deviation set is then null.
\end{proof}

\begin{theorem}[Almost-everywhere convergence on a finite space]
\family{D11.07.finite_ae_to_measure}\label{d11:07:finite-ae-to-measure}
If $\mu(X)<\infty$ and $f_n\to f$ almost everywhere, then $f_n\to f$ in
measure. On an arbitrary space the conclusion holds after restriction to
each measurable subset of finite measure.
\end{theorem}
\begin{proof}
For fixed $\eps>0$ the measurable functions
$\one_{\{|f_n-f|>\eps\}}$ tend to zero almost everywhere and are bounded
by the integrable $\one_X$. Bounded convergence says their integrals,
which are the deviation measures, tend to zero. Apply the same argument
to the restricted space for the second assertion.
\end{proof}

\begin{theorem}[Subsequences and uniqueness]
\family{D11.07.measure_subsequence_uniqueness}\label{d11:07:measure-subsequence-uniqueness}
On any positive measure space, convergence in measure to $f$ supplies a
subsequence converging to $f$ almost everywhere. If the same sequence
converges in measure to $f$ and $g$, then $f=g$ almost everywhere.
\end{theorem}
\begin{proof}
Choose strictly increasing $n_k$ such that
$\mu\{|f_{n_k}-f|>2^{-k}\}<2^{-k}$. The sum of these measures is finite,
so the first Borel--Cantelli lemma says that almost every point belongs to
only finitely many such sets. At those points $|f_{n_k}-f|\leq2^{-k}$
eventually, proving convergence. For uniqueness use
\[
 \{|f-g|>\eps\}\subseteq
 \{|f-f_n|>\eps/2\}\cup\{|f_n-g|>\eps/2\}.
\]
The right side has measure tending to zero, so the left side is null.
Taking the union over $\eps=1/j$ gives $f=g$ almost everywhere.
\end{proof}

\begin{theorem}[Subsequence characterization on a finite space]
\family{D11.07.finite_subsequence_characterization}\label{d11:07:finite-subsequence-characterization}
If $\mu(X)<\infty$, convergence in measure to $f$ is equivalent to the
assertion that every subsequence has a further subsequence converging to
$f$ almost everywhere. The forward implication holds without finiteness.
\end{theorem}
\begin{proof}
Every subsequence of a sequence convergent in measure still converges in
measure, so the preceding theorem proves the forward implication.
If convergence in measure fails, there are $\eps,\delta>0$ and a
subsequence with $\mu\{|f_{n_k}-f|>\eps\}\geq\delta$ for every $k$.
An almost-everywhere convergent further subsequence would converge in
measure by the finite-space theorem, contradicting these lower bounds.
\end{proof}

\begin{theorem}[Summable integral errors]
\family{D11.07.summable_errors_ae}\label{d11:07:summable-errors-ae}
If $\sum_n\int|f_n-f|\,d\mu<\infty$, then $f_n\to f$ almost everywhere,
without any finiteness or completeness assumption on $\mu$.
\end{theorem}
\begin{proof}
The nonnegative-series theorem gives
$\int\sum_n|f_n-f|=\sum_n\int|f_n-f|<\infty$.
Thus this series is finite almost everywhere. At every such point its
terms tend to zero.
\end{proof}

\begin{theorem}[Completeness for convergence in measure]
\family{D11.07.cauchy_measure_complete}\label{d11:07:cauchy-measure-complete}
Every globally Cauchy-in-measure sequence has a finite measurable scalar
limit in measure. A subsequence converges to this limit almost everywhere,
and the limit is unique modulo ordinary almost-everywhere equality.
The measure space may have arbitrary total mass.
\end{theorem}
\begin{proof}
Choose increasing $n_j$ such that, with $u_j=f_{n_j}$,
\[
 \mu(B_j)<2^{-j},\qquad B_j=\{|u_{j+1}-u_j|>2^{-j}\}.
\]
For example first choose Cauchy thresholds for the pairs
$(\eps,\delta)=(2^{-j},2^{-j})$ and take $n_j$ increasing beyond all
thresholds through $j$. Put $F_k=\bigcup_{j\geq k}B_j$.
Then $\mu(F_k)\leq2^{1-k}$ and $N=\bigcap_kF_k$ is a measurable null
set. For $x\notin F_k$ and $i>j\geq k$,
\[
 |u_i(x)-u_j(x)|\leq\sum_{r=j}^{i-1}2^{-r}\leq2^{1-j}.
\]
Thus $u_j(x)$ is Cauchy for $x\notin N$ and has a finite scalar limit.
Set $f$ to this limit off $N$ and to zero on $N$. The measurable-limit
construction in Chapter~\ref{ch:d11-04} makes this specified version
measurable; no completeness of $\Sigma$ is needed. Passing $i\to\infty$
above gives $|u_j-f|\leq2^{1-j}$ outside $F_k$ for $j\geq k$.
For any $\eps,\delta>0$, choose $k$ with $2^{1-k}<\delta$ and then $j$
large enough that $j\geq k$ and $2^{1-j}\leq\eps$.
The deviation set for $u_j$ is contained in $F_k$, so $u_j\to f$ globally
in measure. For the full sequence use
\[
 \{|f_n-f|>\eps\}\subseteq
 \{|f_n-u_j|>\eps/2\}\cup\{|u_j-f|>\eps/2\}.
\]
The Cauchy property bounds the first measure by $\delta/2$ when both
indices are beyond its threshold; fix a sufficiently late such $j$ so
that the second measure is below $\delta/2$. This proves $f_n\to f$ in
measure. The constructed subsequence converges almost everywhere by its
definition, and uniqueness was proved above.
\end{proof}

\begin{theorem}[Egorov]
\family{D11.07.egorov}\label{d11:07:egorov}
Suppose $\mu(X)<\infty$ and $f_n\to f$ almost everywhere. For every
$\eps>0$ there is measurable $E$ with $\mu(E)<\eps$ such that convergence
is uniform on $X\setminus E$. On a Lebesgue measurable $E_0\subseteq\R$
of finite measure, a retained subset can be chosen closed in $\R$:
there is closed $F\subseteq E_0$ with $\lambda(E_0\setminus F)<\eps$
and uniform convergence on $F$.
\end{theorem}
\begin{proof}
Remove a measurable null set $N$ off which convergence holds. For
$k\geq1$, let
$A_{k,m}=\bigcup_{n\geq m}\{|f_n-f|>1/k\}$. These sets decrease in $m$,
and their intersection is contained in $N$. Continuity from above,
using $\mu(X)<\infty$, gives $\mu(A_{k,m})\to0$.
Choose $m_k$ with $\mu(A_{k,m_k})<\eps 2^{-k-1}$. Then
$E=N\cup\bigcup_kA_{k,m_k}$ has measure less than $\eps$, and on its
complement $|f_n-f|\leq1/k$ for $n\geq m_k$. Given any positive
tolerance choose $k$ with $1/k$ smaller than it, proving uniformity.
For the real-line assertion first spend $\eps/2$ on the removed set.
Global closed approximation from \cref{d11:03:global-closed-open-approximation}
gives closed $F$ inside the remaining measurable subset, losing measure
less than $\eps/2$. Restriction preserves uniform convergence, and the
two losses sum to less than $\eps$.
\end{proof}

\begin{theorem}[Consequences of almost-uniform convergence]
\family{D11.07.almost_uniform_implications}\label{d11:07:almost-uniform-implications}
On every positive measure space, almost-uniform convergence implies both
almost-everywhere convergence and global convergence in measure.
\end{theorem}
\begin{proof}
Choose measurable removed sets $E_k$ with $\mu(E_k)<2^{-k}$ and uniform
convergence on their complements. A point at which convergence fails must
belong to every $E_k$, so the measurable set of convergence failures is
contained in the null set $\bigcap_kE_k$. For measure convergence, fix
$\eps,\delta>0$, choose a removed set of measure below $\delta$, and use
uniform convergence on its complement. Eventually the $\eps$-deviation
set is contained in that removed set, hence has measure below $\delta$.
\end{proof}

\begin{theorem}[Lusin on the real line, without a finite-domain assumption]
\family{D11.07.lusin_real_line}\label{d11:07:lusin-real-line}
Let $E\subseteq\R$ be Lebesgue measurable and $f:E\to\R$ finite and
measurable for the trace sigma-algebra. For every $\eps>0$ there are a
closed $F\subseteq E$ and a continuous $g:\R\to\R$ such that
\[
 \lambda(E\setminus F)<\eps,\qquad g|_F=f|_F.
\]
The same assertion holds for complex-valued $f$ and $g$.
\end{theorem}
\begin{proof}
Enumerate a countable basis $(I_j)_{j\geq1}$ of open intervals for $\R$
and put $A_j=f^{-1}(I_j)\subseteq E$. All these sets, and $E\setminus A_j$,
are Lebesgue measurable in $\R$. The global closed approximation theorem,
which allows sets of infinite measure, gives closed
$P_j\subseteq A_j$ and $Q_j\subseteq E\setminus A_j$ with each of
$\lambda(A_j\setminus P_j)$ and
$\lambda((E\setminus A_j)\setminus Q_j)$ below $\eps2^{-j-3}$.
Choose also closed $F_0\subseteq E$ with
$\lambda(E\setminus F_0)<\eps/2$, and set
\[
 F=F_0\cap\bigcap_{j=1}^\infty(P_j\cup Q_j).
\]
This is closed, contained in $E$, and its omitted measure is at most the
sum of the stated losses, strictly less than $\eps$. Within $F$ we have
$f^{-1}(I_j)\cap F=P_j\cap F$, whose complement is $Q_j\cap F$.
Both are relatively closed, so each inverse image of a basis interval is
relatively open. Every open set is a countable union of basis intervals;
hence $f|_F$ is continuous.

For completeness an explicit real-line extension now avoids any boundedness
condition. If $F$ is empty take $g=0$. Otherwise put $g=f$ on $F$.
Every bounded component $(a,b)$ of the open complement has $a,b\in F$;
on it define
\[
 g(t)=\frac{b-t}{b-a}f(a)+\frac{t-a}{b-a}f(b).
\]
On an unbounded component use the constant value of its finite endpoint.
These prescriptions cover the complement: for a point outside a nonempty
closed subset, the largest interval about it avoiding that subset has
endpoints given by the supremum of the subset to its left and the infimum
to its right, with the absent side unbounded. The function is continuous
inside each component. To check it at $x\in F$, take any sequence
$t_m\to x$. Terms in $F$ have $g(t_m)\to f(x)$ by continuity on $F$.
For terms in complement intervals there are two possibilities after taking
a subsequence. If one interval recurs infinitely often, $x$ is its finite
endpoint and its affine or constant formula has limit $f(x)$. Otherwise
take pairwise distinct intervals. Their lengths tend to zero: for any
$r>0$, only finitely many disjoint intervals intersecting
$(x-r/4,x+r/4)$ can have length at least $r$, since each occupies a segment
of length at least $r/2$ in the bounded interval $(x-2r,x+2r)$.
The at most two unbounded intervals cannot occur in this distinct sequence.
Both endpoints therefore tend to $x$. Their $f$-values tend to $f(x)$,
and each $g(t_m)$ is a convex combination of those endpoint values.
Thus this subsequence also converges to $f(x)$. If the original sequence
failed to converge in value, a subsequence bounded away from $f(x)$ would
contradict these alternatives. This proves continuity everywhere.

For complex $f$, apply the real assertion to both coordinates with loss
$\eps/2$ each. Intersect the retained closed sets and combine the two
continuous real extensions as real and imaginary parts. The loss is below
$\eps$, and agreement holds on the intersection.
\end{proof}

\begin{corollary}[Continuous functions converging almost everywhere]
\family{D11.07.continuous_ae_sequence}\label{d11:07:continuous-ae-sequence}
For measurable $E\subseteq\R$ and finite real measurable $f:E\to\R$
there are continuous $g_n:\R\to\R$ with $g_n(x)\to f(x)$ for almost
every $x\in E$.
\end{corollary}
\begin{proof}
Use Lusin with losses $2^{-n}$ to obtain closed $F_n\subseteq E$ and
continuous $g_n$ agreeing with $f$ on $F_n$. Since
$\sum_n\lambda(E\setminus F_n)<\infty$, Borel--Cantelli says almost
every $x\in E$ lies in only finitely many omitted sets. At such an $x$,
$g_n(x)=f(x)$ for all sufficiently large $n$, which proves the assertion.
\end{proof}

\begin{theorem}[Dominated convergence in measure]
\family{D11.07.dominated_measure_lone}\label{d11:07:dominated-measure-lone}
Suppose $f_n\to f$ in measure and $|f_n|\leq g$ almost everywhere for
one integrable $g\geq0$. Then $f$ is integrable,
$\int|f_n-f|\to0$, and $\int f_n\to\int f$ on an arbitrary measure space.
\end{theorem}
\begin{proof}
An almost-everywhere convergent subsequence exists. Remove the common
measurable null set of its convergence and all domination failures;
passing to its pointwise limit gives $|f|\leq g$ there. Thus $f$ is
integrable. If $\int|f_n-f|$ did not tend to zero, a subsequence would
have these integrals bounded below by a fixed positive number. It still
converges in measure to $f$, and so has a further subsequence converging
almost everywhere. Dominated convergence for that further subsequence,
with dominator $2g$ for its errors, contradicts the lower bound.
The scalar integral conclusion follows from
$|\int f_n-\int f|\leq\int|f_n-f|$.
\end{proof}

\begin{theorem}[Fatou for convergence in measure]
\family{D11.07.fatou_measure}\label{d11:07:fatou-measure}
If finite nonnegative measurable $f_n,f$ satisfy $f_n\to f$ globally in
measure, then $\int f\leq\liminf_n\int f_n$, allowing infinite integrals.
\end{theorem}
\begin{proof}
If the right side is infinite there is nothing to prove. Otherwise choose
a subsequence whose integrals converge to that finite lower limit.
It has a further subsequence converging almost everywhere to $f$.
The ordinary nonnegative Fatou inequality along that further subsequence
gives the desired bound.
\end{proof}

\begin{definition}[Local convergence in measure]
\family{D11.07.local_measure}\label{d11:07:local-measure}
The sequence converges locally in measure to $f$ if, for every $A\in\Sigma$
with $\mu(A)<\infty$ and every $\eps>0$,
$\mu(A\cap\{|f_n-f|>\eps\})\to0$.
Global convergence implies this property by monotonicity. Almost-everywhere
convergence implies it by applying the finite-space theorem on each $A$.
The example below shows that local convergence need not be global.
\end{definition}

\begin{theorem}[Operations on limits in measure]
\family{D11.07.measure_algebra_operations}\label{d11:07:measure-algebra-operations}
If $f_n\to f$ and $g_n\to g$ in measure, then
$af_n+bg_n\to af+bg$ in measure for fixed scalars $a,b$ on any measure
space. If $\mu(X)<\infty$, also $f_ng_n\to fg$ in measure.
For a uniformly continuous $\phi:\R\to\R$ or $\phi:\C\to\C$ on the
chosen scalar field, $\phi\circ f_n\to\phi\circ f$ in measure.
\end{theorem}
\begin{proof}
The inequality
$|a(f_n-f)+b(g_n-g)|\leq|a||f_n-f|+|b||g_n-g|$
places every positive deviation set inside the union of two error sets
with fixed positive thresholds, ignoring a term if its coefficient is
zero. Their measures tend to zero. For products assume $\mu(X)<\infty$.
Since $f,g$ are finite, continuity from above gives
$\mu(\{|f|>M\}\cup\{|g|>M\})\to0$ as $M\to\infty$.
Outside this union and $\{|f_n-f|>1\}$ we have $|f_n|\leq M+1$, and
\[
 |f_ng_n-fg|\leq(M+1)|g_n-g|+M|f_n-f|.
\]
For fixed $M$, the set where the right side exceeds $\eps$ has measure
tending to zero, by using thresholds $\eps/(2(M+1))$ for both errors.
The limsup of the product deviation measures is therefore at most the
displayed tail measure. Let $M\to\infty$.
Finally, given $\eps>0$, uniform continuity supplies $r>0$ such that
$|u-v|<r$ implies $|\phi(u)-\phi(v)|<\eps$.
The $\eps$-deviation set after composition is contained in
$\{|f_n-f|>r/2\}$, whose measure tends to zero. Continuity also ensures
that the composed functions are measurable.
\end{proof}

\begin{example}[The typewriter sequence]
\family{D11.07.typewriter_counterexample}\label{d11:07:typewriter-counterexample}
On $[0,1)$ enumerate
\[
 f_{2^k+j}=\one_{[j2^{-k},(j+1)2^{-k})},
 \qquad k\geq0,\quad0\leq j<2^k.
\]
For $1\leq p<\infty$, each function in block $k$ has
$\int|f_n|^p=2^{-k}$; these tend to zero, proving $L^p$ and measure
convergence to zero. Every point belongs to exactly one interval in every
block, so it sees infinitely many ones. For every block with $k\geq1$
it also lies outside at least one of that block's intervals, so it sees
infinitely many zeros. Thus pointwise convergence fails everywhere.
\end{example}

\begin{example}[Failures on an infinite space]
\family{D11.07.infinite_measure_failures}\label{d11:07:infinite-measure-failures}
The sequence $f_n=\one_{[n,n+1)}$ on $\R$ tends to zero everywhere.
For any finite-measure measurable $A$,
\[
 \lambda(A\cap[n,n+1))\leq\lambda(A\cap[n,\infty))\longrightarrow0
\]
by continuity from above; hence it converges locally in measure.
The global $1/2$-deviation measure is always $1$, so global measure
convergence fails. Almost-uniform convergence fails as well: if the
convergence were uniform off a removed set $E$ of measure less than
$\eps<1$, eventually every interval $[n,n+1)$ would be contained in $E$,
because a value $1$ outside $E$ would violate a uniform tolerance $1/2$.
Even one such interval contradicts $\lambda(E)<1$.
\end{example}

\begin{theorem}[Which modes descend to almost-everywhere classes]
\family{D11.07.ae_invariance_modes}\label{d11:07:ae-invariance-modes}
Suppose measurable scalar functions satisfy $f_n=g_n$ almost everywhere
for every $n$ and $f=g$ almost everywhere. Almost-everywhere, global and
local in-measure, almost-uniform, and $L^p$ convergence ($1\leq p\leq\infty$,
when defined) are invariant under these replacements. Pointwise and
uniform convergence on the whole space need not be invariant.
\end{theorem}
\begin{proof}
Place all the equality failures in one measurable null set $N$. Off $N$
the two error sequences agree pointwise. This immediately transfers
almost-everywhere convergence. Their deviation sets have symmetric
difference contained in $N$, so their measures, and their measures after
intersection with a finite-measure set, agree. Their absolute-power
integrals and essential bounds also agree, as do the membership conditions
for the relevant $L^p$ classes. For almost-uniform convergence adjoin $N$
to any removed set; its measure is unchanged and the two sequences then
coincide on the retained set. Apply this in both directions.
On $[0,1]$ the sequence $n\one_{\{0\}}$ is almost everywhere equal to the
zero sequence, but diverges at zero and has uniform error $n$.
\end{proof}
\source{COH 3.1 and 3.3; FOL 2.4; AXL 3B and 4D; ROY 3.2--3.3 and 18.3.}

% END chapters/07


% BEGIN chapters/08
\chapter{Product measures and iterated integrals}\label{ch:d11-08}
The ordinary product sigma-algebra, its measure, and the completion of that
measure are three separate constructions. Sections of ordinary-product
measurable functions are measurable at every parameter. After completion
that assertion generally holds only almost everywhere, with the factor
measures completed as well.

\begin{definition}[Product sigma-algebra]
\family{D11.08.product_sigma}\label{d11:08:product-sigma}
For measurable spaces $(X,\Sigma)$ and $(Y,\mathcal T)$ put
\[
 \Sigma\otimes\mathcal T
 =\sigma\{A\times B:A\in\Sigma,\ B\in\mathcal T\}.
\]
It is the smallest sigma-algebra on $X\times Y$ making both coordinate
projections measurable. Indeed their inverse images are $A\times Y$ and
$X\times B$, whose intersections are all generating rectangles; conversely
these inverse images belong to the generated sigma-algebra. This is a
sigma-algebra of subsets of the product, not a set of pairs of subsets,
and no completion is included in the notation.
\end{definition}

\begin{theorem}[Sections]
\family{D11.08.sections}\label{d11:08:sections}
For $E\in\Sigma\otimes\mathcal T$, its sections
$E_x=\{y:(x,y)\in E\}$ and $E^y=\{x:(x,y)\in E\}$ belong to
$\mathcal T$ and $\Sigma$, respectively, for every $x,y$.
Every section of an ordinary-product-measurable real, complex or
extended-real function is measurable. No measures are required.
\end{theorem}
\begin{proof}
For fixed $x$, taking the section commutes with complements and countable
unions. The sets whose $x$-sections lie in $\mathcal T$ therefore form a
sigma-algebra. It contains every rectangle, whose section is either $B$
or the empty set, hence contains the product sigma-algebra. The other
coordinate is identical. If $f$ is measurable and $H$ is any Borel set
in its target, $(f_x)^{-1}(H)=(f^{-1}(H))_x$ is measurable by the set
assertion. This proves the function assertion, in both coordinates.
\end{proof}

\begin{proposition}[Pairs and fixed-coordinate inclusions]
\family{D11.08.measurable_pair}\label{d11:08:measurable-pair}
For a measurable space $(Z,\Upsilon)$, a map
$h(z)=(u(z),v(z))$ into $(X\times Y,\Sigma\otimes\mathcal T)$ is
measurable if and only if both components are measurable. The maps
$y\mapsto(x,y)$ and $x\mapsto(x,y)$ with the other coordinate fixed are
measurable.
\end{proposition}
\begin{proof}
If $h$ is measurable, compose with the measurable projections. Conversely
$h^{-1}(A\times B)=u^{-1}(A)\cap v^{-1}(B)$ is measurable; the generator
test then proves measurability of $h$. A fixed component is a constant
map, whose inverse images are either the entire domain or empty, so the
last assertion follows.
\end{proof}

\begin{theorem}[One-sided measurability of section integrals]
\family{D11.08.section_integrals_one_sided}\label{d11:08:section-integrals-one-sided}
Let $\nu$ be a sigma-finite positive measure on $(Y,\mathcal T)$, and
let $(X,\Sigma)$ be only a measurable space. For every
$E\in\Sigma\otimes\mathcal T$, the map $x\mapsto\nu(E_x)$ is measurable.
More generally, for ordinary-product-measurable $f:X\times Y\to\Eplus$,
the map $x\mapsto\int f(x,y)\,d\nu(y)$ is measurable. There is no measure
hypothesis on $X$.
\end{theorem}
\begin{proof}
First suppose $\nu(Y)<\infty$. Let $\mathcal D$ consist of product
measurable $E$ for which $x\mapsto\nu(E_x)$ is measurable. It contains
$X\times Y$. If $E\subseteq F$ belong to it, finite total mass gives
$\nu((F\setminus E)_x)=\nu(F_x)-\nu(E_x)$, a difference of finite
measurable functions. If $E_j$ are disjoint members, the corresponding
map for their union is $\sum_j\nu((E_j)_x)$, hence measurable.
Thus $\mathcal D$ is a Dynkin class. It contains the rectangles because
their maps are $\nu(B)\one_A$. Rectangles form a pi-system generating the
product, so the pi--lambda theorem gives all product measurable sets.
For sigma-finite $\nu$ choose a disjoint measurable partition $(D_j)$ of
$Y$ with $\nu(D_j)<\infty$ (empty members are harmless). The finite
measures $\nu_j(B)=\nu(B\cap D_j)$ give measurable maps by the finite
case, and $\nu(E_x)=\sum_j\nu_j(E_x)$ is measurable.
Finally choose nonnegative simple $s_n\uparrow f$ on the product.
For a simple function its section-integral map is a finite sum of
coefficients times the just-proved measurable set maps. Monotone
convergence in $Y$ identifies the pointwise increasing limit of these
maps with the section integral of $f$, proving its measurability.
\end{proof}

\begin{construction}[The sigma-finite product measure]
\family{D11.08.product_measure}\label{d11:08:product-measure}
If $\mu$ and $\nu$ are sigma-finite positive measures on the two spaces,
there is a unique measure $\eta=\mu\times\nu$ on
$\Sigma\otimes\mathcal T$ such that
\[
 \eta(A\times B)=\mu(A)\nu(B),\qquad0\cdot\infty=0.
\]
It is sigma-finite, and for every product measurable $E$,
\[
 \eta(E)=\int_X\nu(E_x)\,d\mu(x)
        =\int_Y\mu(E^y)\,d\nu(y).
\]
Both integrands are measurable.
\end{construction}
\begin{proof}
The preceding theorem makes
$\eta(E)=\int_X\nu(E_x)\,d\mu$ meaningful. Its empty-set value is zero.
For disjoint $E_j$, the sections are disjoint for each $x$, so countable
additivity of $\nu$ followed by the nonnegative-series theorem gives
$\eta(\bigcup_jE_j)=\sum_j\eta(E_j)$. Thus $\eta$ is a measure.
For a rectangle the integrand is $\nu(B)\one_A$. When $\nu(B)<\infty$
the simple integral gives its prescribed value. When $\nu(B)=\infty$,
write this integrand as the increasing limit of $k\one_A$; the integral
is zero if $\mu(A)=0$ and infinite if $\mu(A)>0$, again the prescribed
product. Choose increasing finite-measure exhaustions $A_n\uparrow X$
and $B_n\uparrow Y$. Their rectangles exhaust the product and have
finite $\eta$-measure, proving sigma-finiteness.
Rectangles are a pi-system, and these same finite rectangles provide
common finite exhaustions for any two measures with the prescribed
values. The uniqueness theorem for measures on a generating pi-system
therefore proves uniqueness on the product sigma-algebra. Reversing the
construction, using the sigma-finiteness of $\mu$, yields a measure
$E\mapsto\int_Y\mu(E^y)\,d\nu$ with the same rectangle values. Uniqueness
identifies it with $\eta$, proving the other formula.
\end{proof}

\begin{theorem}[Tonelli]
\family{D11.08.tonelli}\label{d11:08:tonelli}
For sigma-finite positive $\mu,\nu$ and ordinary-product-measurable
$f:X\times Y\to\Eplus$, both section-integral maps are measurable and
\[
 \int f\,d(\mu\times\nu)
 =\int_X\left(\int_Y f(x,y)\,d\nu(y)\right)d\mu(x)
 =\int_Y\left(\int_X f(x,y)\,d\mu(x)\right)d\nu(y).
\]
Every integral here is a nonnegative extended integral.
\end{theorem}
\begin{proof}
For indicators these are the two construction formulas just proved.
Finite nonnegative linearity gives the identities for simple functions.
For general $f$ take $s_n\uparrow f$ simple. Monotone convergence on the
product gives convergence of their product integrals. For each fixed
parameter, monotone convergence gives increasing convergence of the
section integrals; those maps are measurable by the one-sided theorem.
A second use of monotone convergence in the outer factor gives exactly
each asserted equality, including infinite values.
\end{proof}

\begin{theorem}[Fubini for absolutely integrable scalar functions]
\family{D11.08.fubini}\label{d11:08:fubini}
Let $\mu,\nu$ be sigma-finite and let scalar $f$ be ordinary-product
measurable with $\int|f|\,d(\mu\times\nu)<\infty$.
Then $f_x$ is $\nu$-integrable for $\mu$-almost every $x$, and $f^y$ is
$\mu$-integrable for $\nu$-almost every $y$. Assign value zero to each
section-integral map where the section is not integrable. These maps are
measurable and integrable, and
\[
 \int f\,d(\mu\times\nu)
 =\int_X\left(\int_Y f(x,y)\,d\nu(y)\right)d\mu(x)
 =\int_Y\left(\int_X f(x,y)\,d\mu(x)\right)d\nu(y)
\]
with that zero convention at exceptional parameters.
\end{theorem}
\begin{proof}
Tonelli applied to $|f|$ shows that
$a(x)=\int_Y|f_x|\,d\nu$ is measurable with finite integral. Hence
$M=\{a=\infty\}$ is measurable and $\mu$-null. Every section is
measurable; outside $M$ it is integrable. For real $f$, the maps
$b_\pm(x)=\int_Y(f_x)^\pm\,d\nu$ are measurable, bounded by $a$, and
have finite integrals by Tonelli. Replace them by zero on $M$ and take
their difference. This gives exactly the zero-extended section integral,
which is measurable. Its modulus is at most $a$ off $M$, so it is
integrable. Linearity, null invariance and Tonelli for $f^+$ and $f^-$
give its integral as $\int f^+-\int f^-=\int f$.
For complex $f$ apply this argument to its two coordinates; their
section integrals are finite off $M$ and combine to the required map.
Its modulus is at most $a$ by the scalar integral bound, and the
coordinate equalities give the complex identity. Interchanging the
factors proves the other order by the same argument.
\end{proof}

\begin{theorem}[Null sections]
\family{D11.08.null_sections}\label{d11:08:null-sections}
For sigma-finite $\mu,\nu$ and $E\in\Sigma\otimes\mathcal T$,
$\eta(E)=0$ if and only if $\nu(E_x)=0$ for $\mu$-almost every $x$,
and if and only if $\mu(E^y)=0$ for $\nu$-almost every $y$.
\end{theorem}
\begin{proof}
The construction formula writes $\eta(E)$ as the nonnegative integral
of either measurable section-measure map. The zero-integral criterion
is precisely the claimed almost-everywhere vanishing in each case.
\end{proof}

\begin{theorem}[Tonelli and Fubini after completion]
\family{D11.08.completed_product}\label{d11:08:completed-product}
Let sigma-finite $\mu,\nu$ be given on $(X,\Sigma),(Y,\mathcal T)$,
without assuming completeness. Write $\bar\mu,\bar\nu$ for their
completions and $\bar\eta$ for the completion of
$\eta=\mu\times\nu$ on the ordinary product.
If $f:X\times Y\to\Eplus$ is $\bar\eta$-measurable, then $f_x$ is
$\bar\nu$-measurable for $\bar\mu$-almost every $x$ and $f^y$ is
$\bar\mu$-measurable for $\bar\nu$-almost every $y$. Give a section
integral value zero if the section is not measurable. The resulting maps
are measurable for the completed factors and
\[
 \int f\,d\bar\eta
 =\int_X\left(\int_Y f_x\,d\bar\nu\right)d\bar\mu
 =\int_Y\left(\int_X f^y\,d\bar\mu\right)d\bar\nu.
\]
If instead scalar $f$ is $\bar\eta$-integrable, its sections are
integrable almost everywhere, the section integrals extended by zero
where not integrable are integrable for the completed factors, and the
same equalities hold for scalar integrals.
\end{theorem}
\begin{proof}
We first record an integration fact about completion. If $h$ is
measurable for an original sigma-algebra, its integral against the
completed measure equals its original integral: the measures agree on
each level set of a simple $h$, and increasing simple approximations
give the claim for nonnegative $h$. Positive and negative parts, or
coordinates, give it for integrable scalar $h$.

The completion-representative theorem
\cref{d11:04:completion-representative} supplies an ordinary-product
measurable $g$ and an ordinary-product measurable $\eta$-null set $N$
such that $f=g$ off $N$. In the nonnegative case choose $g\geq0$:
the representative construction permits setting it to zero on $N$.
In the scalar case take finite scalar coordinate representatives, again
setting them to zero on $N$. Null sections give measurable null sets
$M\in\Sigma$, $L\in\mathcal T$ such that
$\nu(N_x)=0$ off $M$ and $\mu(N^y)=0$ off $L$.
For $x\notin M$, the ordinary measurable $g_x$ agrees with $f_x$ off
the measurable $\nu$-null set $N_x$. Completeness of $\bar\nu$ makes
$f_x$ measurable, and null invariance gives equality of their completed
section integrals whenever nonnegative, or whenever integrable.
The same assertion holds for the other coordinate off $L$.

In the nonnegative case the zero-extended actual section-integral map
for $f$ agrees off $M$ with the ordinary measurable map
$x\mapsto\int g_x\,d\nu$. The parameters with nonmeasurable sections
form a subset of $M$, hence a measurable null set for $\bar\mu$.
An arbitrary change on $M$ preserves measurability for that complete
factor. Thus the prescribed map is $\bar\mu$-measurable, has the same
integral as the map for $g$, and Tonelli identifies this with
$\int g\,d\eta$. This last integral equals $\int f\,d\bar\eta$, by
the completion integration fact and null invariance. Reversing factors
proves the other iterated equality.

For scalar $f$, the same completion fact and null invariance show that
$g$ is $\eta$-integrable. Ordinary Fubini supplies a measurable null set
of additional parameters outside which $g_x$ is integrable. Outside
its union with $M$, $f_x$ is completed-measurable and integrable with
the same section integral. Its actual failures of integrability form a
subset of this completed null set. The prescribed zero-extended map
therefore equals the ordinary Fubini map almost everywhere, is
completion-measurable and integrable, and has the same integral.
Ordinary Fubini and null invariance prove the scalar formula. The
argument in the other coordinate is identical. This explains both why
factor completions are present and why the exceptional qualification
cannot generally be removed.
\end{proof}

\begin{theorem}[Finite products and permutations]
\family{D11.08.finite_product_permutations}\label{d11:08:finite-product-permutations}
For $n\geq1$ sigma-finite spaces, the sigma-algebra generated by their
measurable rectangles and its rectangle-value product measure do not
depend on parenthesization under canonical tuple identifications.
Every permutation of factors induces a measurable measure-preserving
bijection onto the permuted product. Nonnegative and absolutely
integrable multiple integrals can consequently be evaluated in every
order, with the appropriate almost-everywhere qualifications for
integrable sections.
\end{theorem}
\begin{proof}
For three factors both parenthesizations have the sigma-algebra generated
by $A\times B\times C$. To verify this without assuming that a general
set in $\Sigma_1\otimes\Sigma_2$ is a rectangle, fix $C$ and observe that
the class of $D$ for which $D\times C$ belongs to the three-coordinate
rectangle sigma-algebra is a sigma-algebra containing the two-coordinate
rectangles; then vary $C$ through its measurable sets. Equivalently use
the minimal sigma-algebra making all three projections measurable, which
the paired-projection universal property identifies with either
parenthesization. Repetition proves the claim for any finite number.
Each iterated measure has the same value $\prod_i\mu_i(A_i)$ on a
rectangle. Associativity of nonnegative extended multiplication with
the zero convention holds even when one factor is zero and another
infinite: both parenthesizations are then zero. Products of the finite
factor exhaustions form common finite-measure rectangle exhaustions.
Uniqueness on the generating pi-system identifies the measures.
A coordinate permutation and its inverse have measurable coordinate
maps, hence are measurable. The transported measure has the same
rectangle values as the permuted product, so the same uniqueness proves
measure preservation. Finally repeated two-factor Tonelli, or Fubini
after applying Tonelli to the absolute value, gives each integration
order. There are only finitely many factors at each step, so unions of
the required exceptional null sets retain measure zero. For $n=1$
the assertions reduce to the identity map and the original integral.
\end{proof}

\begin{theorem}[Euclidean products and their completions]
\family{D11.08.euclidean_product_completion}\label{d11:08:euclidean-product-completion}
Under the coordinate identification, for $m,n\geq1$,
\[
 \BB(\R^m)\otimes\BB(\R^n)=\BB(\R^{m+n}),\qquad
 \lambda_m^B\times\lambda_n^B=\lambda_{m+n}^B.
\]
The completion is $(\LL_{m+n},\lambda_{m+n})$.
Completing the ordinary product on $\LL_m\otimes\LL_n$ of the completed
factor Lebesgue measures gives this same completed measure, although
$\LL_m\otimes\LL_n$ itself need not be $\LL_{m+n}$.
\end{theorem}
\begin{proof}
Coordinate projections are continuous, so inverse images of Borel sets
are Borel and every Borel rectangle is Borel in the joint space.
Conversely every rational open box in $\R^{m+n}$ is a rectangle of
rational open boxes in the two factors. Such boxes form a countable
basis, proving the equality of sigma-algebras. On bounded half-open
boxes the product measure has the product of all side lengths, so the
Euclidean uniqueness theorem identifies it with $\lambda_{m+n}^B$.
Its completion was already proved to be $\lambda_{m+n}$.

If $A\in\LL_m$ and $B\in\LL_n$, Borel approximation gives Borel
$A_0,B_0$ and Borel null $N_m,N_n$ with
$A\mathbin\triangle A_0\subseteq N_m$ and
$B\mathbin\triangle B_0\subseteq N_n$. Then
\[
 (A\times B)\mathbin\triangle(A_0\times B_0)
 \subseteq(N_m\times\R^n)\cup(\R^m\times N_n).
\]
The right side is a Borel product-null set, since its rectangle values
are zero, including $0\cdot\infty$. Consequently every completed-factor
rectangle belongs to $\LL_{m+n}$, and its measure equals
$\lambda_m(A)\lambda_n(B)$. The restriction of $\lambda_{m+n}$ to
$\LL_m\otimes\LL_n$ is thus the unique ordinary product measure.
We have inclusions
\[
 \BB(\R^{m+n})\subseteq\LL_m\otimes\LL_n\subseteq\LL_{m+n}.
\]
Completing the middle space adds only Lebesgue measurable sets, because
the right space is complete with the same restricted measure. Conversely
every Lebesgue measurable set is a Borel set modified inside a Borel
null set, both of which are already in the middle sigma-algebra.
Therefore its completion is all of $\LL_{m+n}$ with the stated measure.
\end{proof}

\begin{example}[Complete factors need not have a complete ordinary product]
\family{D11.08.completion_sections_failure}\label{d11:08:completion-sections-failure}
Choose a non-Lebesgue-measurable $V\subseteq[0,1]$ and put
$E=\{0\}\times V$. It lies inside the measurable product-null rectangle
$\{0\}\times[0,1]$, hence in the completed product. But its section at
$x=0$ is $V$, so the sections theorem rules out
$E\in\LL_1\otimes\LL_1$. Thus the ordinary product of two complete
Lebesgue spaces is not complete, and completed-product measurable sets
need not have every section measurable.
\end{example}

\begin{example}[A non-sigma-finite failure of Tonelli interchange]
\family{D11.08.nonsigma_finite_tonelli_failure}\label{d11:08:nonsigma-finite-tonelli-failure}
On $[0,1]$ with its Borel sigma-algebra let $\lambda$ be Borel Lebesgue
measure and $c$ counting measure. The diagonal is closed in the square,
so $f(x,y)=\one_{\{x=y\}}$ is ordinary-product Borel measurable. Every
$y$-section is a singleton of counting mass $1$, while every $x$-section
is a singleton of Lebesgue mass zero. Therefore
\[
 \int\left(\int f(x,y)\,dc(y)\right)d\lambda(x)=1,
 \qquad
 \int\left(\int f(x,y)\,d\lambda(x)\right)dc(y)=0.
\]
The measure $c$ is not sigma-finite: every finite-mass set is finite,
and a countable union of finite sets cannot cover the uncountable
interval. This is a failure of the iterated equality without the
product-theorem hypotheses, even for a nonnegative function.
\end{example}

\begin{example}[A scalar interchange without absolute integrability]
\family{D11.08.signed_fubini_failure}\label{d11:08:signed-fubini-failure}
On $\N\times\N$, here $\N=\{1,2,\ldots\}$, use counting measures and set
\[
 f(m,n)=\begin{cases}1&m=n,\\-1&m=n+1,\\0&\text{otherwise.}\end{cases}
\]
Every section is finitely supported. Summing first in $n$, the row at
$m=1$ has sum $1$ and all other rows have sum zero. The outer sum is
$1$. Summing first in $m$, every column has one $1$ and one $-1$, so the
outer sum is zero. The countable product sigma-algebra is the full power
set (every singleton is a rectangle), and $\sum_{m,n}|f(m,n)|=\infty$.
Thus no scalar product integral is licensed by absolute integrability,
and Fubini's hypothesis is essential here.
\end{example}

\begin{proposition}[Counting-factor interchange on arbitrary spaces]
\family{D11.08.sum_integral_counting}\label{d11:08:sum-integral-counting}
For any positive measure $\mu$ and nonnegative measurable $f_n$,
$\int\sum_nf_n\,d\mu=\sum_n\int f_n\,d\mu$.
For scalar $f_n$ with $\sum_n\int|f_n|<\infty$, the corresponding scalar
interchange holds and the series is absolutely convergent almost
everywhere. No sigma-finiteness of $\mu$ is required in this counting
factor case.
\end{proposition}
\begin{proof}
The nonnegative assertion is exactly the nonnegative-series theorem,
whose proof used increasing finite sums on $X$, not a product measure.
The scalar assertion is \cref{d11:06:absolute-series-interchange}.
Since integration against counting measure on $\N$ is summation, these
are the claimed counting-factor formulas. This proof does not assert
uniqueness of arbitrary non-sigma-finite product extensions.
\end{proof}

\begin{theorem}[Subgraphs and Cavalieri's formula]
\family{D11.08.subgraph_cavalieri}\label{d11:08:subgraph-cavalieri}
For sigma-finite $\mu$ and measurable $f:X\to\Eplus$, the strict
subgraph $G_f=\{(x,t):0<t<f(x)\}$ belongs to
$\Sigma\otimes\BB(\R)$ and
\[
 (\mu\times\lambda_1^B)(G_f)=\int f\,d\mu.
\]
For Lebesgue measurable $E\subseteq\R^{m+n}$,
\[
 \lambda_{m+n}(E)=\int_{\R^m}\lambda_n(E_x)\,d\lambda_m(x),
\]
where sections and the zero convention at nonmeasurable sections have
the completed-product almost-everywhere meaning.
\end{theorem}
\begin{proof}
The identity
$G_f=\bigcup_{q\in\Q,\ q>0}\{f>q\}\times(0,q)$ proves ordinary-product
measurability: for $0<t<f(x)$ choose rational $q$ strictly between $t$
and $f(x)$. Each section is $(0,f(x))$, with length $f(x)$, including
length infinity. The set construction formula gives the first equality.
The same formula holds with the completed real factor, since the
completed-factor rectangle product restricts to the same Borel product.
For Borel $E$ the second equality is the section formula for the Borel
Euclidean product. For a general Lebesgue measurable $E$, apply the
completed-product theorem to $\one_E$ and the Euclidean completion
identification. Almost every section is Lebesgue measurable and its
indicator integral is its measure, giving the stated formula.
\end{proof}

\begin{theorem}[Tensor integrals with prescribed rectangle values]
\family{D11.08.tensor_integrals}\label{d11:08:tensor-integrals}
Let $\mu,\nu$ be arbitrary positive measures and suppose a measure
$\eta$ on $\Sigma\otimes\mathcal T$ satisfies
$\eta(A\times B)=\mu(A)\nu(B)$ with the zero-infinity convention.
For nonnegative measurable $u$ on $X$ and $v$ on $Y$, the function
$h(x,y)=u(x)v(y)$ is measurable and
\[
 \int h\,d\eta=\left(\int u\,d\mu\right)
                 \left(\int v\,d\nu\right).
\]
For scalar integrable $u,v$, $h$ is $\eta$-integrable,
$\int|h|\,d\eta=(\int|u|\,d\mu)(\int|v|\,d\nu)$, and
$\int h\,d\eta=(\int u\,d\mu)(\int v\,d\nu)$.
When the factors are sigma-finite $\eta$ is the unique product above.
\end{theorem}
\begin{proof}
Coordinate measurability and measurable multiplication prove the first
measurability assertion, including the nonnegative extended product
with value zero at $(0,\infty)$ or $(\infty,0)$.
Choose increasing finite-valued simple $u_k\uparrow u$ and
$v_k\uparrow v$. Then $u_k(x)v_k(y)\uparrow u(x)v(y)$: if one limit is
zero every corresponding term is zero; if both are positive the usual
monotone product limit holds, also when a limit is infinite.
On the disjoint rectangular level partition, the rectangle-value
formula and finite distributivity give
\[
 \int u_k(x)v_k(y)\,d\eta
 =\left(\int u_k\,d\mu\right)\left(\int v_k\,d\nu\right).
\]
Monotone convergence on each of the three spaces proves the desired
identity. In taking the limit of the right side, if one limiting
integral is zero all of that factor's approximating integrals are zero,
so every product and the limit are zero. If both limits are positive,
eventual positive lower bounds in both factors justify the ordinary
extended monotone product limit. This checks the only zero-infinity
boundary without using a general Tonelli theorem.
Apply the nonnegative formula to $|u|,|v|$ to obtain the finite
absolute integral of a scalar tensor. For real tensors, expand each
factor into its positive and negative parts. Each of the four
nonnegative products has finite integral, and the nonnegative formula
and scalar linearity give the product of the two real integrals.
For complex tensors expand into real and imaginary coordinates and
apply this real conclusion to the four coordinate products. Finite
linearity gives the asserted complex product. Uniqueness in the
sigma-finite case has already been proved.
\end{proof}
\source{COH 5.1--5.2; FOL 2.5; AXL 5A--5B; ROY 20.1--20.2.}

% END chapters/08


% BEGIN chapters/09
\chapter{Transport and change of variables}\label{ch:d11-09}

Transport is first a statement about measurable maps. Euclidean Jacobians
are proved separately, including the passage from Borel to completed
Lebesgue measure.

\begin{definition}[Pushforward]\family{D11.09.pushforward}\label{d11:09:pushforward}
For a measurable $u:(X,\SSS,\mu)\to(Y,\mathcal T)$ define
$u_*\mu(B)=\mu(u^{-1}B)$, $B\in\mathcal T$. It is a positive measure
and $u_*\mu(Y)=\mu(X)$, without finiteness assumptions.
\end{definition}
\begin{proof}
Inverse images preserve the empty set and disjoint countable unions.
Apply the measure axioms of $\mu$. The total-mass identity follows from
$u^{-1}Y=X$.
\end{proof}

\begin{theorem}[Nonnegative transport]\family{D11.09.transport_nonnegative}\label{d11:09:transport-nonnegative}
For every measurable $h:Y\to[0,\infty]$,
$\int_Yh\,d(u_*\mu)=\int_Xh\circ u\,d\mu$ as extended nonnegative
integrals.
\end{theorem}
\begin{proof}
For an indicator the formula is the definition of pushforward. Add it
over the disjoint levels of a nonnegative simple function. For general
$h$ choose simple $h_k\uparrow h$; their compositions are measurable and
increase to $h\circ u$. MCT on both spaces gives the identity.
\end{proof}

\begin{theorem}[Scalar transport]\family{D11.09.transport_integrable}\label{d11:09:transport-integrable}
For measurable scalar $h$, $h\in L^1(u_*\mu)$ iff $h\circ u\in L^1(\mu)$.
In that case $\int h\,d(u_*\mu)=\int h\circ u\,d\mu$.
\end{theorem}
\begin{proof}
Apply nonnegative transport to $|h|$ to obtain the integrability equivalence.
Then apply it to positive and negative real parts and imaginary parts and
combine their finite integrals. No subtraction of infinite integrals is used.
\end{proof}

\begin{proposition}[Composition and induced isometries]\family{D11.09.pushforward_composition}\label{d11:09:pushforward-composition}
For measurable $u,v$, $(v\circ u)_*\mu=v_*(u_*\mu)$. If $u_*\mu=\nu$,
composition preserves nonnegative and integrable scalar integrals and gives
a linear isometry $L^p(\nu)\to L^p(\mu)$, $1\leq p\leq\infty$.
It need not be onto.
\end{proposition}
\begin{proof}
Evaluate both pushforwards on $B$: their value is
$\mu(u^{-1}(v^{-1}B))$. Transport gives the integral identities and,
applied to $|h|^p$, the finite-$p$ norm identity. For infinity the level
sets have the same measures, so their essential bounds agree.
Inverse images of $\nu$-null exceptional sets are $\mu$-null, proving
well-definedness on quotients and linearity. For non-surjectivity map a
two-point probability space with masses $1/2,1/2$ to one point. Its pullbacks are the
constant classes, whereas the two-point space has nonconstant classes.
\end{proof}

\begin{proposition}[A supplied density]\family{D11.09.pushforward_density_bridge}\label{d11:09:pushforward-density-bridge}
If $u_*\mu=h\nu$ for measurable $h\geq0$, then
$\int g(u(x))\,d\mu=\int gh\,d\nu$ for measurable $g\geq0$.
If $h$ is finite a.e. and $g:Y\to\mathbb K$ is measurable, then
$g\circ u$ is integrable iff
$\int|g|h\,d\nu<\infty$, and the scalar identity holds then.
\end{proposition}
\begin{proof}
Combine transport with \cref{d11:15:density-change-integral}. In the scalar
case first apply both nonnegative identities to $|g|$, and only then
combine the finite scalar integrals. This statement assumes the density;
its existence requires the separate RN hypotheses.
\end{proof}

\begin{theorem}[Linear Jacobian]\family{D11.09.linear_jacobian}\label{d11:09:linear-jacobian}
For an invertible real linear $T:\R^n\to\R^n$, every Lebesgue measurable
$E$ has measurable image and $m_n(T E)=|\det T|m_n(E)$. For nonnegative
measurable $h$, and for integrable scalar $h$,
\[
 \int h(Tx)|\det T|\,dx=\int h(y)\,dy.
\]
Scalar $h$ is integrable iff $h\circ T$ is integrable.
\end{theorem}
\begin{proof}
Coordinate permutations preserve the integral by sigma-finite Fubini.
A nonzero coordinate scaling has factor equal to its absolute scalar by
the one-dimensional interval formula and nonnegative transport. A shear
$x_i\mapsto x_i+c x_j$, $i\ne j$, preserves the integral: fix all
coordinates except $x_i$, apply one-dimensional translation invariance,
and then Tonelli. Gaussian elimination factors every invertible matrix
into these maps. Multiplying their factors gives $|\det T|$, so the
nonnegative formula holds first for Borel functions. Taking indicators
gives the Borel set formula.

A linear map is Lipschitz and sends null sets to outer-null sets: cover a
null set by cubes with total volume below $\eta$; the image of a cube of
side $s$ lies in a cube of side $C s$, where $C$ depends only on the
matrix and dimension. Thus image outer measure is at most $C^n\eta$.
Apply this also to $T^{-1}$. Write a completed measurable set as a Borel
set modulo a subset of a Borel null set. Its image is then measurable
with the same Borel formula. Nonnegative simple approximation proves the
completed integral formula. Its application to $|h|$ proves the scalar
integrability equivalence; real and imaginary parts give the scalar identity.
\end{proof}

\begin{theorem}[The completed nonlinear Jacobian]\family{D11.09.c1_jacobian}\label{d11:09:c1-jacobian}
Let $U\subset\R^n$ be open and let $G:U\to\R^n$ be injective and $C^1$,
with $J(x)=|\det DG(x)|>0$ everywhere. Then $V=G(U)$ is open and $G^{-1}$
is $C^1$. Every Lebesgue measurable $E\subset U$ has measurable image and
\[
 m(GE)=\int_EJ(x)\,dx.                                      \tag{9.1}
\]
For nonnegative measurable $h$ on $V$,
$\int_Vh=\int_Uh(Gx)J(x)\,dx$.
For scalar measurable $h$, integrability on $V$ is equivalent to absolute
integrability of the weighted composition on $U$, and the same identity
then holds.
\end{theorem}
\begin{proof}
The inverse function theorem gives local $C^1$ inverses. Injectivity makes
them agree, so $V$ is open and the global inverse is $C^1$.
Here are quantitative volume bounds that prove (9.1).

Use the sup norm temporarily, with its induced matrix norm. If $H$ is
$C^1$ near a convex cube $Q$ and $\|DH-I\|_\infty\leq\eps$ there,
the remainder $R=H-\id$ is $\eps$-Lipschitz on $Q$ by integrating its
derivative along segments. For side length $s$, the range of each
coordinate of $H$ has width at most $s+\eps s$. Consequently
\[
 m(HQ)\leq(1+\eps)^n m(Q).                                \tag{9.2}
\]
This factor tends to one; a fixed dimensional Lipschitz bound would not
suffice. At a point $c$, put $A=DG(c)$. Continuity of $DG$ and boundedness
of $A^{-1}$ give, on a sufficiently small surrounding cube,
$\|A^{-1}DG(x)-I\|_\infty\leq\eps$. The linear theorem and (9.2)
therefore give
\[
 m(GQ)\leq J(c)(1+\eps)^n m(Q)                             \tag{9.3}
\]
for any cube in that neighborhood. Euclidean and sup norms are equivalent
in finite dimension, so the stated $C^1$ assumption supplies exactly this
continuity, with no extra differentiability hypothesis.

We claim that on an open $W\subset U$ with $J\leq a<\infty$,
$m(GW)\leq a\,m(W)$. To see this, take the dyadic grid and call a closed
cube acceptable if it has side at most one, lies in $W$, and satisfies
the derivative estimate above with its center as $c$. Every point of $W$
belongs to arbitrarily small acceptable cubes: continuity is uniform near
that point and the inverses of nearby derivative matrices are bounded.
Choose the maximal acceptable dyadic cubes (equivalently stop the dyadic
subdivision the first time a cube qualifies). Their interiors are disjoint
and their closed union covers $W$; the countable grid faces account for
points belonging to more than one. Local Lipschitz continuity of $G$ sends
these faces to null sets. Indeed a compact cube inside $U$ has bounded
derivative, hence a Lipschitz constant; the null-image cube-cover argument
in the linear proof applies to any Lipschitz map. A countable such local
cover of $U$ handles all faces. Summing (9.3) gives
$m(GW)\leq a(1+\eps)^n m(W)$. Let $\eps\downarrow0$; when $m(W)=\infty$
the bound is automatic.

If $E$ is Borel with $J\leq a$ on $E$ and $m(E)<\infty$, continuity of
$J$ and outer regularity provide an open
$W\supset E$, $W\subset U\cap\{J<a+\delta\}$, with
$m(W)\leq m(E)+\delta$. Thus
$m(GE)\leq(a+\delta)(m(E)+\delta)$, and $\delta\downarrow0$ gives
$m(GE)\leq a m(E)$. For infinite $m(E)$ this upper bound is automatic
when $a>0$. Apply the same result to $G^{-1}$, whose determinant is
$1/J(G^{-1}y)$: if $J\geq b>0$ on $E$, then
$m(E)\leq b^{-1}m(GE)$, hence $m(GE)\geq b m(E)$.
All these sets are Borel because $G$ is a homeomorphism.

Fix $t>1$ and partition a Borel $E$ into
$E_k=E\cap\{t^k\leq J<t^{k+1}\}$, $k\in\Z$.
The upper and lower bounds yield
\[
 t^k m(E_k)\leq m(GE_k)\leq t^{k+1}m(E_k).
\]
The sets $GE_k$ are disjoint by injectivity. Summing and comparing the
same sums with $\int_EJ$ gives
$t^{-1}\int_EJ\leq m(GE)\leq t\int_EJ$, with extended nonnegative
values. Let $t\downarrow1$ to prove (9.1) for Borel sets.

Finally both $G$ and $G^{-1}$ are locally Lipschitz on countably many
compact cubes in their open domains. The cube-cover argument therefore
shows that each sends null subsets to outer-null subsets. Decomposing a
Lebesgue measurable $E$ into a Borel representative modulo a subset of a
Borel null set proves measurability of $GE$ and extends (9.1) to the
completion. The same argument for the inverse makes $h\circ G$
Lebesgue measurable for every Lebesgue measurable $h$. Indicators, simple
functions and MCT give the nonnegative formula. Apply it to $|h|$ before
forming any scalar integral to obtain the final assertion.
\end{proof}

\begin{theorem}[Polar measure]\family{D11.09.polar_measure}\label{d11:09:polar-measure}
For $n\geq2$ there is a unique finite Borel measure $\sigma$ on $S^{n-1}$
such that the map $x\mapsto(|x|,x/|x|)$ sends $m_n$ off zero to
$r^{n-1}\,dr\otimes\sigma$. Hence
\[
 \int_{\R^n}h(x)\,dx=\int_0^\infty\int_{S^{n-1}}
 h(r\theta)\,d\sigma(\theta)\,r^{n-1}\,dr
\]
for nonnegative Borel $h$ and integrable scalar $h$.
Moreover $\sigma(S^{n-1})=n m_n(B(0,1))$; completion gives the same
formula with a.e. sections.
\end{theorem}
\begin{proof}
The radial-direction map is a homeomorphism off zero. Define
$\sigma(A)=n m_n\{r\theta:0<r\leq1,\ \theta\in A\}$.
This is a finite Borel measure since the indicated cones are the inverse
images of the Borel rectangles $(0,1]\times A$. Linear scaling gives
the cone of radius $b$ measure $b^n\sigma(A)/n$. Subtract the nested finite
cones of radii $a,b$ to obtain for $0\leq a<b<\infty$
$m\{r\theta:a<r\leq b,\theta\in A\}
=\sigma(A)\int_a^b r^{n-1}\,dr$.
These rectangles generate the product Borel sigma-algebra (both spaces
are second countable). Uniqueness of sigma-finite measures, using radial
annuli and finiteness of the sphere measure, identifies the pushforward
with that product. It also proves uniqueness of $\sigma$ by $(0,1]\times A$.
The origin and the unit sphere have zero $m_n$-measure: for the sphere,
bound by annuli whose volumes tend to zero. Thus the total sphere measure
is the displayed ball constant. Transport and Tonelli give the formula.
The homeomorphism preserves null sets relative to these equal pushforward
measures; the completed-product theorem then gives the completed version
with its a.e. section qualification, rather than an assertion about every
section of an arbitrary completed-measurable function.
\end{proof}

\begin{corollary}[Radial power tests]\family{D11.09.radial_integrability}\label{d11:09:radial-integrability}
For $R>0$ and real $\alpha$, $|x|^{-\alpha}$ is integrable on
$0<|x|<R$ exactly when $\alpha<n$, and on $|x|>R$ exactly when
$\alpha>n$. Both integrals diverge at $\alpha=n$.
\end{corollary}
\begin{proof}
Polar integration reduces the tests to
$\int_0^R r^{n-1-\alpha}\,dr$ and
$\int_R^\infty r^{n-1-\alpha}\,dr$, with a finite positive angular factor.
The elementary power-integral tests give the claims. For $n=1$ split the
line into its positive and negative halves, equivalent to counting measure
on the two-point sphere.
\end{proof}

\begin{theorem}[Monotone AC substitution]\family{D11.09.monotone_ac_substitution}\label{d11:09:monotone-ac-substitution}
Let $F:[a,b]\to\R$ be nondecreasing and absolutely continuous,
$A=F(a)$, $B=F(b)$. For Lebesgue measurable $h$ on $[A,B]$, define
$H=h(F)F'$ where $F'$ exists and is positive, and $H=0$ where it is zero
or undefined. Then $H$ is measurable. For $h\geq0$ the identity
$\int_a^bH=\int_A^Bh$ holds with extended values; for integrable scalar
$h$ it holds absolutely and as a scalar identity. Use $0\cdot\infty=0$.
\end{theorem}
\begin{proof}
If $A=B$, $F$ is constant and $F'=0$, so all claims are immediate.
Otherwise give $[a,b]$ the measure $F'\,dx$. For $t\in[A,B]$, the set
$\{x:F(x)\leq t\}$ is an initial interval ending at the last point of
the level $t$, with an immaterial choice of a single endpoint. Continuity
and the AC fundamental theorem imply its $F'\,dx$ mass is $t-A$.
Thus the pushforward of $F'\,dx$ by $F$ is Lebesgue measure on $[A,B]$,
first on half-lines and then on Borel sets by finite-measure uniqueness.
Transport proves substitution for Borel functions.

A Lebesgue measurable $h$ agrees off a Borel null set $N$ with a Borel
function $h_0$. The just-proved pushforward identity gives
$\int_{F^{-1}N}F'=0$. Since $F'\geq0$, the measurable set
$F^{-1}N\cap\{F'>0\}$ is Lebesgue-null. The weighted compositions agree
outside this null set and the derivative-exceptional null set; on
$\{F'=0\}$ both are defined as zero. Completeness proves measurability
of $H$ even if the unweighted composition is nonmeasurable. The identities
then pass from $h_0$ to $h$, and applying the nonnegative identity to
$|h|$ proves absolute integrability before scalar combination.
\end{proof}

\begin{example}[Unweighted composition can fail]\family{D11.09.composition_measurability_failure}\label{d11:09:composition-measurability-failure}
There is a strictly increasing continuous AC $F$ and a null measurable
target set $E$ for which $F^{-1}E$ is not Lebesgue measurable.
\end{example}
\begin{proof}
Choose a closed nowhere dense $K\subset[0,1]$ of positive measure
as in \cref{d11:03:cantor-null-uncountable}, and
put $F(x)=\int_0^x\one_{[0,1]\setminus K}(t)\,dt$. It is AC and
strictly increasing, because every nonempty open interval meets the open
complement of $K$ in an interval of positive length. On each complementary
interval $(a_j,b_j)$ it is a translation with slope one. Their disjoint
image intervals have total length $m([0,1]\setminus K)=F(1)-F(0)$.
They exhaust the target except for $F(K)$ and possibly their endpoints;
therefore $F(K)$ is compact and null. Choose a nonmeasurable subset
$A\subset K$, using \cref{d11:03:positive-set-nonmeasurable-subset}.
Then $E=F(A)\subset F(K)$ is measurable and null, whereas injectivity
gives $F^{-1}E=A$. Taking $h=\one_E$ proves the composition failure.
\end{proof}

\Needspace{18\baselineskip}
\begin{theorem}[von Neumann's composition criterion]\family{D11.09.von_neumann_composition}\label{d11:09:von-neumann-composition}
Let $h:[a,b]\to[c,d]$ be continuous, strictly increasing and onto, with
finite $a<b$, $c<d$. Every Lebesgue measurable real $g$ on $[c,d]$
has measurable $g\circ h$ iff $h^{-1}$ is absolutely continuous.
The complex version follows componentwise.
\end{theorem}
\begin{proof}
Testing indicator functions shows that the first condition is equivalent
to measurability of $h^{-1}(E)$ for every measurable $E\subset[c,d]$.
Here $h^{-1}(E)$ is also the image of $E$ under the continuous monotone
function $h^{-1}:[c,d]\to[a,b]$. The null-image/measurable-image
characterization in \cref{d11:12:monotone-null-image-characterization}
makes this exactly absolute continuity of $h^{-1}$. Conversely measurable
inverse images of all measurable target sets imply measurability of every
composition by the definition of a measurable scalar function. Restricting
or extending $g$ by zero relates this interval formulation to the formulation
with $g$ defined on all of $\R$.
\end{proof}
\source{FOL 2.7 and 10.1; COH 6.1; STE Chapter 3 Exercises 20--21; ROY 6.6. The nonlinear theorem above is the full completed-measure theorem, not a Riemann substitution special case.}

% END chapters/09


% BEGIN chapters/10
\chapter{Scalar \texorpdfstring{$L^p$}{Lp} spaces}\label{ch:d11-10}

Throughout this chapter $(X,\SSS,\mu)$ is a positive measure space and
$\mathbb K$ is $\R$ or $\C$. Null sets mean ordinary measurable null sets.
Neither completeness nor sigma-finiteness is implicit.

\begin{definition}[Raw spaces and essential bounds]\family{D11.10.lp_definition}\label{d11:10:lp-definition}
For $1\leq p<\infty$ put
$\|f\|_p=(\int|f|^p\,d\mu)^{1/p}$ and let $L^p_{\rm raw}$ consist of
measurable scalar $f$ with finite norm. Put
\[
 \|f\|_\infty=\inf\{M\geq0:\mu(\{|f|>M\})=0\},\qquad
 \inf\varnothing=\infty.
\]
Finite essential norm is attained a.e., and an essentially bounded function
has a bounded measurable representative.
\end{definition}
\begin{proof}
If the infimum is $a<\infty$, every $a+1/k$ is an essential bound: the
upward-closed set of bounds has an element below $a+1/k$. Outside the
countable union of the corresponding null sets, $|f|\leq a$. Set $f$ to
zero on $\{|f|>a\}$ to get the bounded representative. This is a
modification on a measurable null set, so no completeness is needed.
\end{proof}

\begin{theorem}[H\"older]\family{D11.10.holder}\label{d11:10:holder}
For $p,q\in[1,\infty]$ with $1/p+1/q=1$, $f\in L^p_{\rm raw}$ and
$g\in L^q_{\rm raw}$ have integrable product and
$\int|fg|\leq\|f\|_p\|g\|_q$. When $1<p<\infty$ and both norms are
nonzero, equality holds exactly when $|f|^p$ and $|g|^q$ are positive
constant multiples a.e. A zero norm gives equality trivially.
\end{theorem}
\begin{proof}
For $a,b\geq0$, convexity or differentiation of
$a^p/p-ab+b^q/q$ in $a$ gives Young's inequality
$ab\leq a^p/p+b^q/q$, with equality exactly when $a^p=b^q$.
Normalize by the two nonzero norms and integrate. Equality of integrals
forces the nonnegative Young difference to vanish a.e., which proves the
characterization. Zero norms mean a.e. zero functions. At $p=1,q=\infty$
use the attained bound $|g|\leq\|g\|_\infty$ outside one null set;
the reversed endpoint is identical.
\end{proof}

\begin{corollary}[Cauchy--Schwarz]\family{D11.10.cauchy_schwarz}\label{d11:10:cauchy-schwarz}
For square-integrable scalar $f,g$, $fg$ is integrable and
$\int|fg|\leq\|f\|_2\|g\|_2$, hence
$|\int f\overline g|\leq\|f\|_2\|g\|_2$.
\end{corollary}
\begin{proof}
Take $p=q=2$ in H\"older and apply the scalar integral triangle inequality.
In the real case conjugation has no effect.
\end{proof}

\begin{theorem}[Minkowski and the norm identities]\family{D11.10.minkowski}\label{d11:10:minkowski}
For $1\leq p\leq\infty$ and $f,g\in L^p_{\rm raw}$,
$f+g\in L^p_{\rm raw}$ and $\|f+g\|_p\leq\|f\|_p+\|g\|_p$.
Also $\|cf\|_p=|c|\|f\|_p$ and $\|f\|_p=0$ iff $f=0$ a.e.
\end{theorem}
\begin{proof}
For finite $p$, $|f+g|^p\leq2^{p-1}(|f|^p+|g|^p)$ first proves
integrability. For $p>1$, H\"older applied to
$|f||f+g|^{p-1}$ and $|g||f+g|^{p-1}$ gives
$\|f+g\|_p^p\leq(\|f\|_p+\|g\|_p)\|f+g\|_p^{p-1}$.
Divide when the left norm is nonzero. For $p=1$ integrate the pointwise
triangle inequality; for infinity use simultaneous essential bounds.
Homogeneity follows directly, treating $c=0$ separately. A nonnegative
function has zero integral exactly when it is zero a.e.; essential norm zero
has the same conclusion by the thresholds $1/k$.
\end{proof}

\begin{definition}[The quotient norm]\family{D11.10.quotient_construction}\label{d11:10:quotient-construction}
The relation $f\sim g$ iff $\mu(\{f\ne g\})=0$ is an equivalence relation
compatible with addition and scalar multiplication on $L^p_{\rm raw}$.
The quotient $L^p(\mu)=L^p_{\rm raw}/\!\sim$ has the well-defined norm
$\|[f]\|_p=\|f\|_p$.
\end{definition}
\begin{proof}
Reflexivity and symmetry are immediate, and transitivity follows since the
new exceptional set is contained in the union of two null sets. The same
union argument proves compatibility of operations. Integrals and essential
bounds are unchanged under a.e. equality, so the norm is well-defined.
The preceding theorem supplies all norm axioms and shows that zero norm
is exactly the zero class.
\end{proof}

\begin{theorem}[Completeness]\family{D11.10.completeness}\label{d11:10:completeness}
Every Cauchy sequence in $L^p(\mu)$, $1\leq p\leq\infty$, has a unique
norm limit in that space.
\end{theorem}
\begin{proof}
For finite $p$, choose a subsequence represented by $f_k$ with
$\|f_{k+1}-f_k\|_p\leq2^{-k}$. Minkowski bounds the norms of the finite
sums $S_N=\sum_{k=1}^N|f_{k+1}-f_k|$ by one. MCT gives
$S=\sum_{k\geq1}|f_{k+1}-f_k|\in L^p$, finite outside a measurable null
set. There the series $f_1+\sum(f_{k+1}-f_k)$ converges; set its value to
zero on the null set. This gives a measurable $f$ and, by the same finite
sum bounds followed by Fatou,
$\|f-f_k\|_p\leq\sum_{j\geq k}2^{-j}$. Thus the subsequence, and by
Cauchyness the original sequence, converges to $[f]$.
For infinity take a subsequence with essential differences at most $2^{-k}$,
and choose its first representative bounded everywhere.
Outside a single measurable null set the representatives form a uniformly
Cauchy sequence of scalar functions. Its bounded limit there, extended by
zero on that set, is measurable and is the essential-uniform limit.
The triangle inequality again handles the original sequence. Two limits
have zero mutual distance, proving uniqueness.
\end{proof}

\begin{proposition}[Chebyshev]\family{D11.10.chebyshev}\label{d11:10:chebyshev}
For every measurable scalar $f$, real $p>0$, and $\alpha>0$,
$\alpha^p\mu(\{|f|>\alpha\})\leq\int|f|^p\,d\mu$, allowing infinity.
\end{proposition}
\begin{proof}
Integrate $\alpha^p\one_{\{|f|>\alpha\}}\leq|f|^p$.
\end{proof}

\begin{lemma}[Sigma-finite support]\family{D11.10.sigmafinite_support}\label{d11:10:sigmafinite-support}
For $1\leq p<\infty$ every raw $L^p$ function has reference-sigma-finite
nonzero support: $\{f\ne0\}=\bigcup_{k\geq1}\{|f|>1/k\}$ and each
set has measure at most $k^p\|f\|_p^p<\infty$.
\end{lemma}
\begin{proof}
The union identity is pointwise; the bound is Chebyshev. This is a property
of finite-$p$ functions, not of arbitrary finite target measures.
\end{proof}

\begin{proposition}[Norm convergence and pointwise subsequences]\family{D11.10.lp_convergence_measure}\label{d11:10:lp-convergence-measure}
If $\|f_k-f\|_p\to0$ with finite $p$, then $f_k\to f$ in measure and
some subsequence converges a.e. If $\|f_k-f\|_\infty\to0$, the whole
sequence converges uniformly on one measurable conull set.
\end{proposition}
\begin{proof}
Chebyshev proves convergence in measure. Choose $k_j$ so that
$\mu\{|f_{k_j}-f|>2^{-j}\}<2^{-j}$. The measure of the union of these
sets for $j\geq J$ tends to zero; their limsup is null. Outside it the
subsequence converges pointwise. For infinity take the union, over all $k$
and rational bounds larger than $\|f_k-f\|_\infty$, of the exceptional
null sets. Off this union the pointwise bounds imply uniform convergence.
\end{proof}

\begin{proposition}[Dominated norm convergence]\family{D11.10.dominated_lp}\label{d11:10:dominated-lp}
For finite $p$, let $f_k,f$ be measurable scalar functions and $g$ a
nonnegative measurable $L^p$ function. If $f_k\to f$ a.e. and
$|f_k|\leq g$ a.e. for each $k$, then
$f\in L^p$ and $\|f_k-f\|_p\to0$. The conclusion also holds with
convergence in measure in place of a.e. convergence.
\end{proposition}
\begin{proof}
On a common conull set $|f|\leq g$, and
$|f_k-f|^p\leq2^pg^p$. DCT proves the first assertion. For the second,
every subsequence has a further a.e. convergent subsequence by the
summable-exceptional-set construction. Along each such subsequence the
first assertion applies. If the full norms did not tend to zero, a
subsequence bounded away from zero would contradict this conclusion.
\end{proof}

\begin{proposition}[Finite measure inclusions]\family{D11.10.finite_measure_inclusion}\label{d11:10:finite-measure-inclusion}
If $\mu(X)<\infty$ and $1\leq p<q\leq\infty$, then
$L^q\subset L^p$ and
$\|f\|_p\leq\mu(X)^{1/p-1/q}\|f\|_q$.
The zero-measure space is separately the zero quotient.
\end{proposition}
\begin{proof}
For finite $q$, apply H\"older to $|f|^p\cdot1$ with exponents
$q/p$ and $q/(q-p)$ and take $p$th roots. For infinite $q$, integrate
the essential bound $|f|^p\leq\|f\|_\infty^p$. If $\mu(X)=0$ every
class is zero and no ambiguous power convention is required.
\end{proof}

\begin{proposition}[Elementary interpolation]\family{D11.10.elementary_norm_interpolation}\label{d11:10:elementary-norm-interpolation}
For $1\leq p<q<r\leq\infty$ and
$1/q=\theta/p+(1-\theta)/r$, $0<\theta<1$,
$\|f\|_q\leq\|f\|_p^\theta\|f\|_r^{1-\theta}$ whenever the two
right norms are finite.
\end{proposition}
\begin{proof}
For finite $r$, factor $|f|^q=|f|^{\theta q}|f|^{(1-\theta)q}$ and
apply H\"older with exponents $p/(\theta q)$ and
$r/((1-\theta)q)$. They exceed one and their reciprocals sum to one.
For $r=\infty$, use $\theta q=p$ and bound the remaining factor by its
essential supremum before integration. Zero norms are dealt with by a.e.
zero representatives.
\end{proof}

\begin{proposition}[Counting measure inclusions]\family{D11.10.sequence_space_inclusion}\label{d11:10:sequence-space-inclusion}
For every set $I$ with counting measure and $1\leq p<q\leq\infty$,
$\ell^p(I)\subset\ell^q(I)$ and $\|a\|_q\leq\|a\|_p$.
Every finite-$p$ element has at most countable nonzero support.
\end{proposition}
\begin{proof}
Every coordinate satisfies $|a_i|\leq\|a\|_p$; thus
$\sum|a_i|^q\leq\|a\|_p^{q-p}\sum|a_i|^p$ for finite $q$.
The supremum bound proves the infinite endpoint. Each threshold
$\{|a_i|>1/k\}$ is finite by Chebyshev, so their union is countable.
Integrals against counting measure are the suprema of finite subsums,
which justifies these sums for arbitrary $I$.
\end{proof}

\begin{example}[Inclusions can fail in both directions]\family{D11.10.inclusion_failure}\label{d11:10:inclusion-failure}
For $1\leq p<q<\infty$ on $(0,\infty)$, choose
$1/q\leq a<1/p$ and $1/q<b\leq1/p$. Then
$x^{-a}\one_{(0,1)}\in L^p\setminus L^q$ and
$x^{-b}\one_{(1,\infty)}\in L^q\setminus L^p$.
\end{example}
\begin{proof}
The integrals $\int_0^1x^{-s}\,dx$ and $\int_1^\infty x^{-s}\,dx$
are finite exactly for $s<1$ and $s>1$, respectively. At $s=1$ the
logarithmic integral diverges. Substitute $s=ap,aq,bp,bq$.
\end{proof}

\begin{theorem}[Distribution functions and layer cake]\family{D11.10.distribution_layer_cake}\label{d11:10:distribution-layer-cake}
For measurable scalar $f$, $\lambda_f(t)=\mu\{|f|>t\}$, $t>0$, is
decreasing and right-continuous, and for every real $p>0$,
\[
 \int|f|^p\,d\mu=p\int_0^\infty t^{p-1}\lambda_f(t)\,dt.
\]
Both sides may be infinite, and $\mu$ may be nonsigma-finite.
\end{theorem}
\begin{proof}
Monotonicity is immediate. For $t_k\downarrow t$ the level sets increase
to $\{|f|>t\}$, so continuity from below proves right continuity, including
infinite values. For a nonnegative simple $s=\sum a_j\one_{E_j}$ on
disjoint sets, the identity follows by integrating
$p t^{p-1}\sum_j\mu(E_j)\one_{(0,a_j)}(t)$; this also holds for infinite
level measures with the nonnegative conventions. Take $s_k\uparrow|f|$.
Their strict level sets increase to those of $|f|$ at every $t$, and MCT
on each side proves the formula. This avoids an unjustified product
Tonelli theorem over an arbitrary nonsigma-finite factor.
\end{proof}

\begin{proposition}[Weak level-set control]\family{D11.10.weak_endpoint_bridge}\label{d11:10:weak-endpoint-bridge}
For $p>0$ put $W_p(f)=\sup_{\alpha>0}\alpha\mu\{|f|>\alpha\}^{1/p}$.
Chebyshev gives $W_p(f)\leq\|f\|_p$ when the latter is finite.
A weak $(1,1)$ estimate controls $W_1$, not necessarily the $L^1$ norm.
\end{proposition}
\begin{proof}
Take $p$th roots in Chebyshev and then the supremum. The example in
\cref{d11:13:maximal-not-l1} proves that the stronger norm conclusion does
not follow.
\end{proof}

\begin{theorem}[Minkowski's integral inequality]\family{D11.10.minkowski_integral}\label{d11:10:minkowski-integral}
Let $\mu,\nu$ be sigma-finite and $F:X\times Y\to\mathbb K$ product
measurable. For $1\leq p<\infty$, if
$A=\int_Y\|F(\cdot,y)\|_p\,d\nu(y)<\infty$, its sections are
absolutely integrable in $y$ for a.e. $x$, their integral $g$ is measurable
after zero extension on the exceptional set, and $\|g\|_p\leq A$.
For nonnegative $F$ the inequality holds with extended sides as well.
\end{theorem}
\begin{proof}
Tonelli makes the section norms measurable. Put
$u(x)=\int|F(x,y)|\,d\nu(y)$. For $p=1$ Tonelli proves the claim.
For $p>1$ exhaust $X$ by finite-measure sets $D_j$, and set
$u_{j,N}=\one_{D_j}\min(u,N)$. Tonelli and H\"older on $X$ give
\[
 \|u_{j,N}\|_p^p\leq\int u\,u_{j,N}^{p-1}\,d\mu
 \leq\int_Y\|F(\cdot,y)\|_p
       \|u_{j,N}^{p-1}\|_q\,d\nu
 =A\|u_{j,N}\|_p^{p-1}.
\]
The truncated norm is finite, so division and then MCT give $\|u\|_p\leq A$.
Thus $u$ is finite a.e. Measurable real/imaginary positive-section integrals
define $g$ there, and $|g|\leq u$ proves the conclusion. If $A=\infty$
the extended inequality is automatic; otherwise the same argument applies.
No dual representation theorem was used.
\end{proof}

\begin{theorem}[Jensen]\family{D11.10.jensen}\label{d11:10:jensen}
If $\mu(X)=1$, $f$ is real integrable and $\Phi:\R\to\R$ is convex,
then $\Phi(\int f)\leq\int\Phi(f)$ whenever the right side is integrable,
or as an extended inequality when $\Phi\geq0$. Normalizing Lebesgue
measure gives the finite-interval average version.
\end{theorem}
\begin{proof}
At $a=\int f$, secant-slope monotonicity for a convex function shows that
every left secant slope is at most every right secant slope. Their finite
one-sided limiting bounds admit a real number $c$ between them.
Then $\Phi(t)\geq\Phi(a)+c(t-a)$ for all $t$. Convex finite functions
on $\R$ are continuous (the same local slope bounds prove local
Lipschitz continuity), so $\Phi(f)$ is measurable. Integrate the supporting
line; its linear term has integral zero. In the nonnegative case the
inequality is valid even if the integral is infinite.
\end{proof}

\begin{example}[Two endpoint boundaries]\family{D11.10.c0p_boundary}\label{d11:10:c0p-boundary}
For $0<p<1$ on two counting points, the two distinct point indicators have
norms one, whereas their sum has the displayed $p$-expression $2^{1/p}>2$.
That expression is not a norm. On one point of mass infinity every $L^1$
function is zero, while the constant-one class has ordinary essential norm
one. Hence arbitrary-measure $L^\infty$ norm cannot be recovered by pairing
with the $L^1$ unit ball.
\end{example}

\begin{proposition}[Cohn's endpoint convention]\family{D11.10.cohn_endpoint_translation}\label{d11:10:cohn-endpoint-translation}
On a sigma-finite space, locally null sets are null. Thus bounded measurable
functions modulo local a.e. equality give the same normed $L^\infty$ space
as the ordinary essential-bounded quotient. This identification is not
asserted on a general measure space.
\end{proposition}
\begin{proof}
Locally null means that intersection with each finite-measure measurable
set has measure zero. A countable such cover of $X$ makes the original
set null. The converse is immediate. Every ordinary class has a bounded
representative by the definition above; for bounded representatives the
same null sets give both the equivalence relation and the essential norm.
The infinite-atom example shows why sigma-finiteness cannot be dropped.
\end{proof}
\source{FOL 6.1--6.3; COH 3.1--3.2; AXL 7; ROY 7. Jensen is used only in its scalar form.}

% END chapters/10


% BEGIN chapters/11
\chapter{Approximation, translation and convolution}\label{ch:d11-11}

Finite-$p$ approximation begins on an arbitrary measure space. Topology
enters only when simple functions are replaced by continuous ones.

\begin{theorem}[Simple functions of finite support]\family{D11.11.finite_p_simple_density}\label{d11:11:finite-p-simple-density}
For arbitrary positive $\mu$, scalar $f\in L^p(\mu)$, $1\leq p<\infty$,
and $\eps>0$, there is a finite-valued measurable simple $s$, zero outside
a finite-measure set, with $\|f-s\|_p<\eps$.
\end{theorem}
\begin{proof}
First replace $f$ by $f\one_{\{1/k<|f|\leq k\}}$. Its error tends to zero
by DCT, and its support has finite measure by Chebyshev. On this bounded
range approximate real and imaginary parts by a finite grid of mesh
$\delta$, assigning zero off the support. The pointwise error on the
support is at most $\sqrt2\delta$, so its $L^p$ norm tends to zero as
$\delta\downarrow0$. Choose the two errors successively smaller than
$\eps/2$. The real case needs only one grid.
\end{proof}

\begin{proposition}[Simple approximation at infinity]\family{D11.11.bounded_simple_infinity}\label{d11:11:bounded-simple-infinity}
Every ordinary $L^\infty$ class is approximable in essential norm by
bounded simple functions. Their supports need not have finite measure.
\end{proposition}
\begin{proof}
Choose the bounded representative from \cref{d11:10:lp-definition}.
Partition its bounded scalar range into finitely many squares (intervals
over $\R$) of diameter below $\eps$, and assign a point of each cell.
The measurable simple function has pointwise error below $\eps$.
There was no restriction on the measure of its nonzero levels.
\end{proof}

\begin{theorem}[Rectangle step functions]\family{D11.11.step_density}\label{d11:11:step-density}
For $n\geq1$, $1\leq p<\infty$, finite linear combinations of indicators
of bounded half-open rectangles are dense in $L^p(\R^n,m_n)$.
The analogous finite-interval step functions are dense on an interval.
\end{theorem}
\begin{proof}
It suffices to approximate $\one_E$ for $m(E)<\infty$. Euclidean
regularity gives compact $K\subset E$ and open $U\supset E$ with
$m(U\setminus K)<\eta$. Cover $K$ by finitely many sufficiently small
half-open grid rectangles whose closures lie in $U$; a grid fine enough
exists by the positive distance of $K$ from $U^c$. Their finite union $A$
has $K\subset A\subset U$, so
$\|\one_E-\one_A\|_p\leq\eta^{1/p}$. Disjointizing along their
finitely many faces gives a finite rectangle combination. Apply this to
the finitely many levels of a simple approximant and use Minkowski.
Intersection with the interval gives the one-dimensional restriction.
\end{proof}

\begin{theorem}[Continuous approximation in Euclidean space]\family{D11.11.cc_density_euclidean}\label{d11:11:cc-density-euclidean}
For $1\leq p<\infty$, $C_c(\R^n,\mathbb K)$ is dense in $L^p(m_n)$.
\end{theorem}
\begin{proof}
For a bounded rectangle, replace each coordinate interval indicator by a
continuous piecewise-linear cutoff which is one on the interval except
within distance $\delta$ of its endpoints and zero beyond distance
$\delta$. The product is compactly supported, lies in $[0,1]$, and differs
from the rectangle indicator only in a set whose measure tends to zero.
Its $L^p$ error therefore tends to zero. Apply this to the finitely many
rectangles in the preceding theorem and add with their coefficients.
\end{proof}

\begin{theorem}[The Radon-space bridge]\family{D11.11.radon_cc_density_bridge}\label{d11:11:radon-cc-density-bridge}
Let $X$ be locally compact Hausdorff and $\mu$ Borel, finite on compact
sets, outer regular on Borel sets and inner regular on open sets.
Then $C_c(X,\mathbb K)$ is dense in ordinary $L^p(\mu)$ for finite $p$.
Ordinary completion gives the same classes isometrically.
\end{theorem}
\begin{proof}
For a Borel $E$ of finite measure, choose open $U\supset E$ of finite
measure with $\mu(U\setminus E)<\eta$. Inner regularity of $U$ gives
compact $K\subset U$ with $\mu(U\setminus K)<\eta$.
The LCH compact-open cutoff lemma, \cref{d11:18:compact-cutoff}, gives
$0\leq\phi\leq1$, $\phi=1$
on $K$, $\phi\in C_c(U)$. Then
$\int|\phi-\one_E|^p\leq\mu(U\setminus K)+\mu(U\setminus E)<2\eta$.
The finite-support simple density theorem finishes. The cutoff lemma is
the exact LCH result supplied in the topology/regularity development; no
normality of all of $X$ is assumed. For completion, approximate each
measurable function by simple functions, replace their completed-measurable
levels by Borel levels modulo null sets, and take an a.e. pointwise limit.
This supplies a Borel representative and preserves the norm. The claim
does not concern arbitrary larger Carath\'eodory domains.
\end{proof}

\begin{theorem}[Translations]\family{D11.11.translation_continuity}\label{d11:11:translation-continuity}
For $1\leq p<\infty$, $\tau_hf(x)=f(x-h)$ defines an isometry of
$L^p(\R^n)$ and $\|\tau_hf-\tau_kf\|_p\to0$ as $h\to k$.
\end{theorem}
\begin{proof}
Translation invariance of completed Lebesgue measure proves the isometry
and independence of representatives. For $g\in C_c$, uniform continuity
gives uniform convergence $\tau_hg\to g$ as $h\to0$, and for $|h|\leq1$
all supports lie in a fixed compact set. Thus convergence also holds in
$L^p$. Approximate $f$ by such a $g$; the error is bounded by
$2\|f-g\|_p+\|\tau_hg-g\|_p$. The group identity reduces $h\to k$ to
this case.
\end{proof}

\begin{theorem}[Young's convolution inequality]\family{D11.11.convolution_young}\label{d11:11:convolution-young}
For $k\in L^1(\R^n)$ and $f\in L^p(\R^n)$, $1\leq p\leq\infty$,
$k*f(x)=\int k(y)f(x-y)\,dy$ exists absolutely a.e., has a measurable
$L^p$ representative, and $\|k*f\|_p\leq\|k\|_1\|f\|_p$.
Its class is independent of representatives. On $L^1$, convolution is
bilinear, commutative and associative as classes.
\end{theorem}
\begin{proof}
Choose Borel representatives, possible by the completion comparison.
The function $(x,y)\mapsto k(y)f(x-y)$ is Borel. For finite $p$ its
$x$-section norm is $|k(y)|\|f\|_p$, so Minkowski's integral inequality
proves the assertion. For $p=\infty$ use a bounded representative of $f$:
absolute integrals are bounded by $\|f\|_\infty\|k\|_1$ at every $x$.
Changing $f$ on a null set changes each translated section only on a null
set; changing $k$ changes its integrand only on a null set. Tonelli, or
these translated-section observations, proves representative independence
a.e. for all choices. The substitution $y\mapsto x-y$ gives commutativity
where absolute integrals exist. For associativity, the integral of the
absolute triple expression on $\R^{3n}$ is
$\|f\|_1\|g\|_1\|k\|_1$ by translations and Tonelli. Fubini then
permits either grouping and the same linear substitution. Bilinearity
holds on the common conull set of absolute existence.
\end{proof}

\begin{theorem}[Approximate identities in finite norm]\family{D11.11.approximate_identity_lp}\label{d11:11:approximate-identity-lp}
Let $\phi\in L^1$, $a=\int\phi$, and
$\phi_t(x)=t^{-n}\phi(x/t)$. For $f\in L^p$, $1\leq p<\infty$,
$\|f*\phi_t-af\|_p\to0$ as $t\downarrow0$. No positivity or symmetry
of $\phi$ is required.
\end{theorem}
\begin{proof}
Substitution writes the difference as
$\int\phi(z)(\tau_{tz}f-f)\,dz$. Minkowski bounds its norm by
$\int|\phi(z)|\|\tau_{tz}f-f\|_p\,dz$.
For fixed $z$ the inner norm tends to zero by translation continuity,
and it is at most $2\|f\|_p$. DCT proves the result. Section measurability
and existence follow from the preceding convolution theorem, so the
integral manipulations are identities of a.e. classes.
\end{proof}

\begin{theorem}[Uniform approximate identities]\family{D11.11.approximate_identity_uniform}\label{d11:11:approximate-identity-uniform}
With the same kernels, bounded uniformly continuous $f$ satisfies
$\sup_x|f*\phi_t(x)-af(x)|\to0$. If $f$ is bounded measurable and
continuous on an open $U$, convergence is uniform on every compact
$K\subset U$.
\end{theorem}
\begin{proof}
For the first assertion split
$\int|\phi(z)||f(x-tz)-f(x)|\,dz$ at $|z|=R$. Its tail is at most
$2\|f\|_\infty\int_{|z|>R}|\phi|$, uniformly in $x,t$; on the bounded
part uniform continuity makes the difference uniformly small for small $t$.
For compact $K\subset U$, choose a compact neighborhood $K_1\subset U$.
The restriction of $f$ there is uniformly continuous, and for $|z|\leq R$
and small $t$, $x-tz\in K_1$ for all $x\in K$. The identical splitting
proves the local statement.
\end{proof}

\begin{theorem}[Values at Lebesgue points]\family{D11.11.approximate_identity_lebesgue_points}\label{d11:11:approximate-identity-lebesgue-points}
If $f\in L^p$, $1\leq p\leq\infty$, and $\phi$ is bounded, measurable,
compactly supported with integral one, then $f*\phi_t(x)\to f(x)$ at
every Lebesgue point of the fixed representative. This holds a.e. and
at each of its continuity points.
\end{theorem}
\begin{proof}
Let $\supp\phi\subset B(0,R)$. Local integrability of $f$ follows from
H\"older on compact sets. At a fixed finite-valued point the difference is
bounded by
\[
 \|\phi\|_\infty t^{-n}
 \int_{B(x,tR)}|f(y)-f(x)|\,dy.
\]
This tends to zero at a Lebesgue point. The same integral proves absolute
existence at every such point (indeed at every point for these compact
bounded kernels and locally integrable representatives). Differentiation
gives the a.e. assertion; ordinary continuity implies the Lebesgue-point
condition by taking the supremum on a small ball. This is not a claim of
$L^\infty$ norm convergence.
\end{proof}

\begin{theorem}[Smooth approximation and cutoffs]\family{D11.11.smooth_density}\label{d11:11:smooth-density}
For finite $p$, $C_c^\infty(\R^n)$ is dense in $L^p$.
For compact $K\subset U$ open there is $\psi\in C_c^\infty(U)$ with
$0\leq\psi\leq1$ and $\psi=1$ on $K$.
\end{theorem}
\begin{proof}
The function $b(x)=\exp(-1/(1-|x|^2))$ for $|x|<1$, zero otherwise,
is smooth: each derivative near the boundary is a polynomial in
$(1-|x|^2)^{-1}$ times that exponential and tends to zero. Normalize it
to a nonnegative $\phi$ of integral one. For $g\in C_c$, $g*\phi_t$ has
compact support and is smooth. Differentiation under the integral in the
form $\int g(y)\phi_t(x-y)\,dy$ is justified by compact support and
bounded derivatives, with the difference quotients dominated on a fixed
compact integration region. Approximate $f$ by $g$ and then choose $t$
using the approximate-identity theorem.
For the cutoff choose $\delta>0$ so the closed $3\delta$-neighborhood
of $K$ is compact in $U$. Convolve the indicator of its open
$2\delta$-neighborhood with $\phi_t$, $t<\delta$. The result is between
zero and one, equals one on $K$, is smooth by the same differentiation
argument, and has compact support inside $U$.
\end{proof}

\begin{theorem}[Separability]\family{D11.11.finite_p_separability}\label{d11:11:finite-p-separability}
If $\mu$ is sigma-finite and $\SSS$ countably generated, then $L^p(\mu)$
is separable for $1\leq p<\infty$. In particular this holds for completed
Lebesgue measure on $\R^n$.
\end{theorem}
\begin{proof}
Adjoin a countable finite-measure cover to the generators and form the
countable algebra $\mathcal A$ they generate. On each finite-measure piece
$D$, every set in $\SSS$ is approximable in symmetric-difference measure
by restrictions of members of $\mathcal A$: the approximable sets contain
the algebra, are closed under complements in $D$, finite unions, and
countable unions by first cutting off a finite union using continuity from
below. They therefore form a sigma-algebra. Thus any finite-measure set,
first cut off to a finite union of the cover, is approximable by members
of this countable algebra having finite measure. Rational real or complex
coefficients on finitely many such sets give a countable collection.
Simple-function density and $\|\one_E-\one_A\|_p^p=\mu(E\triangle A)$
show it is dense. Euclidean Borel sets are generated by rational boxes;
the completion comparison supplies Borel representatives for all Lebesgue
$L^p$ classes, so separability passes to that completion.
\end{proof}

\begin{example}[Failures at infinity]\family{D11.11.infinity_density_failure}\label{d11:11:infinity-density-failure}
On $[0,1]$, $\one_{[0,1/2]}$ has distance at least $1/2$ from every
continuous real function in essential norm. On $\R$,
$\|\tau_h\one_{[0,1]}-\one_{[0,1]}\|_\infty=1$ for every
$0<|h|<1$.
\end{example}
\begin{proof}
If the first distance were below $1/2$, some $\delta>0$ would give
$g\geq1/2+\delta$ a.e. to the left of $1/2$ and
$g\leq1/2-\delta$ a.e. to its right. Continuity extends these inequalities
to the common endpoint, contradiction. For the second claim the translated
interval and the original have a symmetric difference of positive measure;
the absolute difference equals one there and is everywhere at most one.
\end{proof}
\source{FOL 6.1 and 8.2; COH 3.2 and 7.4. Fourier transforms and distribution theory are not inferred from these convolution statements.}

% END chapters/11


% BEGIN chapters/12
\chapter{Variation, absolute continuity and Stieltjes measures}\label{ch:d11-12}

The proofs in this chapter have a deliberate order. Interval variation
constructs Stieltjes measures algebraically. Direct RN and differentiation,
proved independently in Chapters~\ref{ch:d11-15} and \ref{ch:d11-13}, then
give the AC fundamental theorem and the general BV differentiation theorem.
Thus the fundamental theorem is not used to prove its own hypotheses.

\begin{definition}[Variation]\family{D11.12.bv_definition}\label{d11:12:bv-definition}
For $F:[a,b]\to\mathbb K$, $a<b$, let
\[
 \Var_{[a,b]}F=\sup_{a=x_0<\cdots<x_k=b}
       \sum_{j=1}^k|F(x_j)-F(x_{j-1})|.
\]
Put $\Var_{[a,a]}F=0$ and $V_F(x)=\Var_{[a,x]}F$.
The function is BV when its variation is finite. For $a\leq c\leq b$,
$\Var_{[a,b]}F=\Var_{[a,c]}F+\Var_{[c,b]}F$.
\end{definition}
\begin{proof}
Insert $c$ into any partition; the triangle inequality cannot decrease
its sum, which gives one bound. Concatenate partitions on the two
subintervals within any chosen error of their suprema for the reverse
bound. This also proves the formula for an infinite supremum and for
degenerate subintervals.
\end{proof}

\begin{theorem}[Jordan decomposition of a function]\family{D11.12.jordan_functions}\label{d11:12:jordan-functions}
A real $F$ is BV iff it is the difference of two bounded nondecreasing
functions. Explicitly, for BV $F$,
\[
 F(x)=F(a)+P(x)-N(x),\quad
 P=\tfrac12(V_F+F-F(a)),\quad N=\tfrac12(V_F-F+F(a)).
\]
Both $P,N$ are nonnegative, nondecreasing and zero at $a$.
\end{theorem}
\begin{proof}
For $x<y$, additivity gives
$V_F(y)-V_F(x)=\Var_{[x,y]}F\geq|F(y)-F(x)|$.
Thus both displayed increments are nonnegative. They are bounded by the
finite variation and the boundedness $|F(x)|\leq|F(a)|+V_F(b)$.
Conversely the sum of the variations of two bounded nondecreasing
functions bounds the variation of their difference, by the triangle
inequality on each partition.
\end{proof}

\begin{proposition}[One-sided limits]\family{D11.12.bv_one_sided}\label{d11:12:bv-one-sided}
BV functions have finite left and right limits at interior points and the
corresponding one-sided endpoint limits. Their discontinuities form an
at most countable set. Globally BV functions have limits at both infinities,
and replacing $F(x)$ by $F(x+)$ changes at most countably many values.
\end{proposition}
\begin{proof}
A bounded nondecreasing real function has left limit the supremum of its
values to the left and right limit the corresponding infimum to the
right. At a discontinuity the open interval between these limits has
positive length. Such intervals for distinct points are disjoint, so
assigning a rational to each proves countability. The value at the point
lies between the limits. Jordan decomposition proves the assertions for
real BV functions; real and imaginary parts prove them for complex ones.
For real $F$ on $\R$, put $V(x)=\sup_{t<x}\Var_{[t,x]}F$, with the supremum also
including zero. Appending partitions proves
$V(y)-V(x)\geq|F(y)-F(x)|$ for $x<y$; hence $(V+F)/2$ and $(V-F)/2$
are bounded nondecreasing components. This construction does not assume
an endpoint limit in advance. Boundedness gives their limits at infinity.
For complex $F$, apply this real argument to its real and imaginary parts. The same
countability proof applies globally, or on the countable compact exhaustion.
\end{proof}

\begin{theorem}[Normalized Stieltjes correspondence]\family{D11.12.normalized_stieltjes}\label{d11:12:normalized-stieltjes}
Let $NBV(\R,\mathbb K)$ be globally BV functions, right-continuous and
zero at $-\infty$. They correspond bijectively to finite scalar Borel
measures through
\[
 F(x)=\nu(( -\infty,x]),\qquad \nu((a,b])=F(b)-F(a).
\]
In particular $\nu(\{x\})=F(x)-F(x-)$, and the cumulative distribution
of $|\nu|$ is $V_F(x)=\Var_{(-\infty,x]}F$.
\end{theorem}
\begin{proof}
For bounded nondecreasing right-continuous real $F$ zero at $-\infty$,
the interval premeasure $F(b)-F(a)$ extends uniquely to a finite Borel
measure by the interval extension theorem. Right continuity is precisely
the endpoint continuity used in that premeasure theorem. Its cumulative
values follow by sending $a\to-\infty$.

For general real $F$, its variation $V_F$ is right-continuous. Here is the
point needed to verify this, rather than merely assume it. Fix $x<z$ and
a partition of $[x,z]$ whose sum is within $\eta$ of its variation; let
$t>x$ be its first point. Inserting an almost maximizing partition of
$[x,y]$, $x<y<t$, and then $t$, shows
\[
 \Var_{[x,y]}F\leq\eta+|F(t)-F(x)|-|F(t)-F(y)|
 \leq\eta+|F(y)-F(x)|.
\]
Let $y\downarrow x$ and then $\eta\downarrow0$. Additivity of variation
proves the claim. The same reasoning applies to complex $F$.
We also need $V_F(-\infty)=0$. Let $T=\Var_\R F<\infty$ and
$\eta>0$. Choose a finite partition $t_0<\cdots<t_k$ whose increment
sum is greater than $T-\eta$. Append this partition to any finite
partition lying to the left of and ending at $t_0$. The resulting sum
is at most $T$, so the earlier sum is less than $\eta$. Taking the
supremum of earlier sums gives $V_F(t_0)\leq\eta$; the initial
zero value at $-\infty$ gives the same supremum, since $F(s)\to0$ as
$s\to-\infty$. For $y\leq t_0$, monotonicity gives
$0\leq V_F(y)\leq\eta$. This proves the required limit.
For real $F$, the functions $(V_F+F)/2,(V_F-F)/2$ are therefore bounded,
nondecreasing, right-continuous and zero at $-\infty$; apply the positive
construction and subtract. For complex $F$, do so to its real and imaginary
parts and combine. Uniqueness follows on half-open intervals by finite
measure uniqueness, componentwise.

Conversely the cumulative function of a finite scalar measure is
right-continuous and zero at $-\infty$ by variation continuity, and every
partition sum is bounded by $|\nu|(\R)$. To identify variation exactly,
first on $(a,b]$ the suprema of sums on interval partitions are
$\Var_{[a,b]}F$. They are at most $|\nu|((a,b])$.
For the reverse inequality approximate any finite measurable partition
of $(a,b]$ in $|\nu|$-measure by sets in its interval algebra, then
disjointize that approximation. Such approximation holds for every finite
measure: the approximable sets contain the generating algebra and are
closed under complements and countable unions, using finite unions first
and continuity from below. On the interval algebra all pieces refine to
one endpoint partition, whose sum is at most $\Var_{[a,b]}F$.
The errors in their $\nu$ values are bounded by their $|\nu|$ errors,
so the same upper bound holds for any measurable partition. Thus
$|\nu|((a,b])=V_F(b)-V_F(a)$. Passing $a\to-\infty$ proves the
cumulative identity; measure uniqueness proves it on all Borel sets.
The atom formula is continuity from below in $a\uparrow x$.
For a compact interval use $F-F(a)$ on $[a,b]$, zero to its left and
constant $F(b)-F(a)$ to its right whenever the generator is continuous
at the endpoints; for arbitrary BV first take the indicated right-limit
representative. This specifies the endpoint normalization, rather than
equating arbitrary a.e. representatives with the same Stieltjes function.
\end{proof}

\begin{definition}[Absolute continuity]\family{D11.12.ac_definition}\label{d11:12:ac-definition}
A scalar $F$ on $[a,b]$, $a<b$ finite, is absolutely continuous if for
every $\eps>0$ some $\delta>0$ makes
\[
 \sum_j(b_j-a_j)<\delta\quad\Longrightarrow\quad
 \sum_j|F(b_j)-F(a_j)|<\eps
\]
for every finite family of pairwise disjoint open subintervals.
Shared endpoints are permitted. On $\R$ we use the same uniform condition
for all finite families of bounded intervals, not merely local AC.
\end{definition}

\begin{proposition}[Basic AC properties]\family{D11.12.ac_basic}\label{d11:12:ac-basic}
AC functions on compact intervals are uniformly continuous and BV, and
their variation functions are AC. Finite linear combinations and products
are AC. Every Lipschitz function is AC; the converse fails, as shown by
$\sqrt x$ on $[0,1]$.
\end{proposition}
\begin{proof}
Use one interval in the definition for uniform continuity. To bound
variation, fix $\delta$ for error one and split $[a,b]$ into $M$ intervals
each of length less than $\delta$. Refine any partition by their endpoints.
Within each piece all its subintervals have total length below $\delta$,
so their increment sum is below one. Thus variation is at most $M$.
For disjoint $(a_j,b_j)$, additivity says
$\sum_j(V_F(b_j)-V_F(a_j))=\sum_j\Var_{[a_j,b_j]}F$.
Choose finite partitions within all these intervals approximating the
variations to arbitrary total error; AC bounds the combined increment
sum when the combined lengths are small. Let the error decrease to zero,
using $\eps/2$ in the AC condition to retain a strict $\eps$ bound.
This proves AC of $V_F$.
Linearity follows from the triangle inequality with a common small
$\delta$. Since compact-interval AC functions are bounded, the identity
$FG(v)-FG(u)=F(v)(G(v)-G(u))+G(u)(F(v)-F(u))$ proves AC of products.
A Lipschitz constant $L$ bounds the increment sum by $L$ times the length
sum. Finally $\sqrt v-\sqrt u=\int_u^v(2\sqrt t)^{-1}\,dt$ for
$0\leq u<v\leq1$, by elementary calculus and the limit at zero.
The density is integrable, so small-set control of its integral proves AC.
The quotients $(\sqrt x-0)/x=x^{-1/2}$ exclude Lipschitz continuity.
\end{proof}

\begin{theorem}[Primitives]\family{D11.12.ac_primitive}\label{d11:12:ac-primitive}
For $f\in L^1([a,b],\mathbb K)$ and $c\in\mathbb K$, let
$F(x)=c+\int_a^xf$. Then $F$ is AC, $F'=f$ a.e., and
$\Var_{[a,b]}F=\int_a^b|f|$. Its values are independent of the
representative of $f$.
\end{theorem}
\begin{proof}
For disjoint intervals the increment sum is at most the integral of
$|f|$ over their union; absolute continuity of the integral proves AC.
At a Lebesgue point of the extension of $f$ by zero, the one-sided interval
averages tend to $f(x)$ by regular shrinking, giving both derivative
quotients. Differentiation supplies these points a.e.
Every partition sum is at most $\int|f|$, giving one variation bound.
Choose interval step functions $s_k\to f$ in $L^1$. Their primitives
$S_k$ are polygonal and have variation $\int|s_k|$, since refining by
the finitely many step endpoints attains it. The preceding upper bound
applied to the primitive of $f-s_k$ gives
$\Var(F-S_k)\leq\int|f-s_k|$. The triangle inequality for variation
therefore yields
$\Var F\geq\int|s_k|-\int|f-s_k|\to\int|f|$.
Equal a.e. densities have identical integrals on every initial interval,
proving the last assertion.
\end{proof}

\begin{theorem}[AC and the Stieltjes measure]\family{D11.12.ac_stieltjes_equivalence}\label{d11:12:ac-stieltjes-equivalence}
For $F\in NBV(\R,\mathbb K)$, $F$ is AC on $\R$ iff
$\nu_F\ll m$. In this case $F'\in L^1(\R)$ and
$F(x)=\int_{-\infty}^xF'(t)\,dt$ for every $x$.
The compact-interval version uses the normalized constant extension.
\end{theorem}
\begin{proof}
If $F$ is AC, the preceding variation argument, valid also on $\R$ for
globally BV $F$, shows $V_F$ is AC. Given a Borel null set $E$, cover it
by an open set of total length below the AC $\delta$ for $V_F$. Its
components are countably many bounded disjoint intervals $(a_j,b_j)$;
finite subcollections and then positive additivity give
$|\nu_F|(E)\leq\sum_j(V_F(b_j)-V_F(a_j))\leq\eps$.
Continuity of $V_F$ makes these the masses of the open intervals, so no
endpoint atoms enter. Let $\eps\downarrow0$.
Conversely scalar RN relative to sigma-finite Lebesgue measure gives
$\nu_F=fm$ with $f\in L^1(\R)$, and the cumulative identity gives
$F(x)=\int_{-\infty}^xf$. Small-set integral control proves uniform AC
on $\R$, and the primitive differentiation proof gives $F'=f$ a.e.
For compact intervals, an AC function is continuous at both endpoints;
its stated constant extension is globally BV and satisfies the same AC
condition. This proves that version without any invocation of general
BV differentiation.
\end{proof}

\begin{theorem}[The AC fundamental theorem]\family{D11.12.ac_ftc}\label{d11:12:ac-ftc}
For scalar $F$ on finite $[a,b]$, $a<b$, the following are equivalent:
$F$ is AC; there is $f\in L^1$ with
$F(x)=F(a)+\int_a^xf$ for every $x$; $F$ is differentiable a.e.,
$F'\in L^1$ and the same formula holds with $F'$. The density is $F'$
a.e.
\end{theorem}
\begin{proof}
For AC $F$, use its normalized constant extension. The AC-to-Stieltjes
argument and direct RN just proved give the integral representation.
The primitive theorem identifies its derivative and proves the converse.
The third condition implies the second by taking $f=F'$; the second
implies the third by primitive differentiation. Thus all formulas hold
at every point, not only a.e. This proof has not assumed the general
monotone or BV differentiation theorem.
\end{proof}

\begin{theorem}[Differentiating a monotone function]\family{D11.12.monotone_differentiation}\label{d11:12:monotone-differentiation}
A nondecreasing real $F$ on $[a,b]$ has finite derivative a.e. in $(a,b)$.
After zero extension on the exceptional set it is measurable,
nonnegative, integrable, and $\int_a^bF'\leq F(b)-F(a)$.
\end{theorem}
\begin{proof}
Replace the interior values by right limits and extend constantly beyond
the endpoints using $F(a+)$ and $F(b-)$. Subtract $F(a+)$ to obtain
$H\in NBV$ with a finite positive Stieltjes measure of mass
$F(b-)-F(a+)\leq F(b)-F(a)$.
Measure differentiation gives, outside one null set, the density $g$ as
the limit of its ratios on both $(x,x+h]$ and $(x-h,x]$.
These are nice families, and their ratios are precisely the two
difference quotients of $H$. Thus $H'=g$ a.e. and
$\int g\leq\nu_H(\R)$.
It remains to justify returning to the original point values. Exclude its
countably many discontinuities. At a remaining $x$, $F(x)=F(x+)$.
For $h>0$ and fixed $0<\eta<1$,
\[
 F(x+(1-\eta)h+)\leq F(x+h)\leq F(x+h+).
\]
The right-limit function is differentiable at $x$, so division by $h$
bounds the lower and upper derivative limits by $(1-\eta)g(x)$ and
$g(x)$. Let $\eta\downarrow0$. For $h<0$ use
$F(x+(1+\eta)h+)\leq F(x+h)\leq F(x+h+)$, and reverse the inequalities
on division by $h$. This proves the derivative of the original $F$ is
$g$, without the false inference that a.e. equality preserves derivatives.
Derivative measurability follows from measurable rational difference
quotients of the Borel monotone function; their limsup equals the derivative
where it exists. The density bound proves integrability and the inequality.
\end{proof}

\begin{theorem}[Differentiating BV functions]\family{D11.12.bv_derivative}\label{d11:12:bv-derivative}
For scalar BV $F$ on $[a,b]$, $F'$ exists finitely a.e., is integrable,
and $\int_a^b|F'|\leq\Var_{[a,b]}F$.
\end{theorem}
\begin{proof}
Real Jordan components and then real/imaginary components give derivative
existence and integrability. For the sharp complex inequality take the
right-limit representative, constant extensions and subtraction of the
left constant as above, obtaining $H\in NBV$ with
$\Var_\R H\leq\Var_{[a,b]}F$: approximate each partition of the right
limits by partitions at nearby original points, and pass to limits.
Its scalar Stieltjes measure has that variation. Decompose it as
$\nu_s+fm$. Measure differentiation gives $H'=f$ a.e., and the
monotone squeeze argument applied to each component shows that $F'=H'$
a.e. Finally
$\int_a^b|f|\leq|\nu_H|(\R)=\Var_\R H\leq\Var_{[a,b]}F$.
The variation inequality uses the measure density and loses no factor
from summing real and imaginary estimates. Reconstruction is not implied;
the singular example below shows why.
\end{proof}

\begin{proposition}[Integration by parts for AC functions]\family{D11.12.ac_integration_parts}\label{d11:12:ac-integration-parts}
For AC scalar $F,G$ on $[a,b]$,
\[
 \int_a^bF'G+\int_a^bFG'=F(b)G(b)-F(a)G(a).
\]
The products in the integrals are bilinear; conjugating $G$ gives the
sesquilinear version.
\end{proposition}
\begin{proof}
Both functions are bounded and their derivatives are $L^1$, so the
integrands are integrable. The product is AC by the basic properties, and
at the common differentiability points its derivative is $F'G+FG'$ by
the scalar product rule. Apply the AC fundamental theorem to the product.
\end{proof}

\begin{proposition}[Stieltjes integration by parts]\family{D11.12.stieltjes_integration_parts}\label{d11:12:stieltjes-integration-parts}
For $F,G\in NBV$, if at least one is continuous, then for finite $a<b$,
\[
 \int_{(a,b]}F\,dG+\int_{(a,b]}G\,dF=F(b)G(b)-F(a)G(a).
\]
\end{proposition}
\begin{proof}
The integrands are bounded Borel and the measures finite. Set
$\alpha=dF|_{(a,b]}$, $\beta=dG|_{(a,b]}$.
Since $F(t)=F(a)+\alpha((a,t])$, Fubini on the finite positive
components of the scalar measures expresses the first integral as
$F(a)(G(b)-G(a))+\int\one_{\{s\leq t\}}\,d\alpha(s)d\beta(t)$.
The other gives the reversed triangle with coefficient
$G(a)(F(b)-F(a))$. The sum of triangle indicators is one on the square
plus one on its diagonal. If one generator is continuous its measure is
atomless, so the diagonal has zero product variation mass by Tonelli:
integrate its singleton sections against the other variation. Thus the
double integrals sum to $(F(b)-F(a))(G(b)-G(a))$. Expanding proves the
formula, with its specified half-open endpoints. Without continuity the
diagonal term is generally nonzero.
\end{proof}

\begin{theorem}[Atomic, AC and singular continuous parts]\family{D11.12.stieltjes_three_parts}\label{d11:12:stieltjes-three-parts}
A finite scalar Borel measure on $\R$ has a unique pairwise singular
decomposition $\nu=\nu_d+\nu_{ac}+\nu_{sc}$, where
\[
 \nu_d=\sum_{x\in D}\nu(\{x\})\delta_x,\qquad
 \sum_{x\in D}|\nu(\{x\})|<\infty,
\]
$\nu_{ac}\ll m$ and $\nu_{sc}$ is atomless singular. Their normalized
BV functions have the same decomposition: the discrete part is a
convergent pure-jump function, the AC part is AC, and the singular
continuous part is continuous with derivative zero a.e. The pure-jump
part also has derivative zero a.e.
\end{theorem}
\begin{proof}
For each $k$ only finitely many points have $|\nu|(\{x\})\geq1/k$,
so the set $D$ of nonzero atoms is countable. Finite partial sums of their
absolute masses are bounded by $|\nu|(\R)$, giving the summability and
variation convergence of the discrete series. Subtract it; the remainder
has no atoms. Its Lebesgue decomposition gives an AC part, which has no
atoms because singletons are Lebesgue-null, and hence an atomless singular
part. The discrete part is supported on $D$, the AC part ignores every
null set, and the singular continuous part ignores $D$ and is supported
on a null set, proving pairwise singularity. Atomic values fix the first
part, and Lebesgue-decomposition uniqueness fixes the other two.
Cumulative sums give the function decomposition, with uniform convergence
of the jump series from absolute summability. A cumulative function is
continuous exactly when its finite scalar measure has no atoms, by
one-sided measure continuity. The AC correspondence handles the middle
part; measure differentiation applied to each singular measure shows its
right and left quotients converge to zero a.e. This proves the derivative
statements, rather than interchanging derivatives of the jump series.
\end{proof}

\begin{example}[The singular continuous function]\family{D11.12.singular_continuous_counterexample}\label{d11:12:singular-continuous-counterexample}
There is continuous nondecreasing $C:[0,1]\to[0,1]$ with endpoints
$0,1$, variation one and $C'=0$ a.e. It is not AC, and its normalized
extension has a nonzero singular continuous probability measure.
\end{example}
\begin{proof}
On the middle-thirds Cantor set, replace the ternary digits $0,2$ by the
binary digits $0,1$; at the double expansions the endpoint values agree
with the constant value on the removed intervening interval. Extend
constantly on every removed interval. Equivalently interpolate the stage
$k$ endpoint values and keep earlier gaps constant. These continuous
nondecreasing stage functions differ uniformly by at most $2^{-k}$ after
stage $k$, so converge to continuous nondecreasing $C$ with the stated
endpoints. The Cantor set has measure zero, since its stage $k$ cover has
length $(2/3)^k$. On its complement $C$ is locally constant, so $C'=0$
a.e. Variation equals the endpoint increase, one. The AC fundamental
theorem rules out AC. Its Stieltjes measure is atomless by continuity,
supported on the Cantor set since each complementary interval has zero
mass, and has total mass one. This also proves strictness in the monotone
derivative inequality.
\end{proof}

\begin{example}[Continuous need not be BV]\family{D11.12.continuous_not_bv}\label{d11:12:continuous-not-bv}
The function $F(0)=0$, $F(x)=x\cos(\pi/(2x))$ for $0<x\leq1$ is
continuous and has infinite variation.
\end{example}
\begin{proof}
Its absolute value is at most $x$, proving continuity at zero. At
$x_n=1/(2n)$ its values are $(-1)^n/(2n)$. A partition containing the
first $N$ such points has variation at least
$\sum_{n=1}^{N-1}(1/(2n)+1/(2n+2))$, which diverges. Partitions
containing all $1,1/2,\ldots,1/(2N)$ refine these and have at least the
same sums. AC would imply BV, so continuity does not imply AC either.
\end{proof}

\begin{example}[AC composition boundary]\family{D11.12.ac_composition_boundary}\label{d11:12:ac-composition-boundary}
On $[-1,1]$, $F(t)=\sqrt[3]t$ and
$G(0)=0$, $G(x)=x^2\cos(\pi/(2x))$ for $x\ne0$ are AC, but
$F\circ G$ is not BV. A Lipschitz outer function composed with an AC
inner function is always AC.
\end{example}
\begin{proof}
The derivative $(1/3)|t|^{-2/3}$ is integrable and its primitive, including
the value at zero, is $\sqrt[3]t$, so the primitive theorem proves AC.
For $x\ne0$, $G'(x)=2x\cos(\pi/(2x))+(\pi/2)\sin(\pi/(2x))$;
also $G'(0)=0$. The derivative is bounded, so the mean value theorem
makes $G$ Lipschitz, hence AC. At $x_n=1/(2n)$ the composite values are
$(-1)^n(2n)^{-2/3}$. Their successive absolute differences have divergent
sum, proving infinite variation. If $H$ has Lipschitz constant $L$, the
increment sum for $H\circ G$ is at most $L$ times that for $G$ on the
same intervals; the AC definition proves the positive assertion.
\end{proof}

\begin{theorem}[Null images and AC for continuous monotone maps]\family{D11.12.monotone_null_image_characterization}\label{d11:12:monotone-null-image-characterization}
For continuous nondecreasing $F:[a,b]\to\R$, $a<b$ finite, the following
are equivalent: $F$ is AC; $F$ sends every Lebesgue-null subset to a null
set; $F$ sends every Lebesgue measurable subset to a measurable set.
Constant functions are included; measurability here is completed Lebesgue
measurability.
\end{theorem}
\begin{proof}
AC implies the null-image condition. Cover a null set, apart from the
two endpoints, by disjoint open intervals of total length below the AC
$\delta$. Their image intervals cover its image and have total length
at most $\eps$, by finite subcollections and passage to the countable
sum. The two endpoint images are harmless. A measurable subset of the
compact domain is a countable union of compact sets modulo a null set,
by inner regularity. Its compact images are compact; if null sets have
null images, its full image is therefore measurable.

Conversely assume measurable images. A null $N$ then has measurable
image. If that image had positive measure, choose a nonmeasurable subset
$B\subset F(N)$ using \cref{d11:03:positive-set-nonmeasurable-subset}.
The set $N\cap F^{-1}(B)$ is measurable by completeness, being a subset
of $N$, but its image is exactly $B$, contradiction. Thus images of null
sets are null. This step uses no unproved measurability of $F^{-1}(B)$.

To prove the remaining implication, subtract $F(a)$, and take the finite
Stieltjes measure on $[a,b]$ with no atom at $a$. For nonconstant $F$
define on $(F(a),F(b)]$
$q(t)=\inf\{x\in[a,b]:F(x)\geq t\}$. Continuity and monotonicity give
$\{t:q(t)\leq x\}=\{t:t\leq F(x)\}$ and $F(q(t))=t$.
Thus $q$ is Borel and its pushforward of Lebesgue measure is exactly that
Stieltjes measure, by the displayed half-line identity and finite-measure
uniqueness. For a Borel null $E\subset[a,b]$,
$q^{-1}(E)\subset F(E)$, which is null by hypothesis. Its pushforward
mass is zero. The Stieltjes measure is consequently absolutely continuous
relative to Lebesgue measure; the already proved AC correspondence makes
$F$ AC. If $F$ is constant, AC follows directly. This proof avoids any
analytic-set theorem or assumed unrestricted composition rule.
\end{proof}

\begin{proposition}[Rectifiable curves]\family{D11.12.rectifiable_curve_bridge}\label{d11:12:rectifiable-curve-bridge}
For continuous $\gamma:[a,b]\to\R^d$, its length is the supremum of
polygonal sums $\sum|\gamma(t_j)-\gamma(t_{j-1})|$. It is finite iff
every coordinate is BV. If the coordinates are AC, then
$\operatorname{length}(\gamma)=\int_a^b|\gamma'(t)|\,dt$.
In this AC case, for continuous $h$ on its image, the arc-length integral along this
parametrized curve is $\int_a^bh(\gamma(t))|\gamma'(t)|\,dt$;
it counts traversal multiplicity.
\end{proposition}
\begin{proof}
Each coordinate variation is at most the length, and Euclidean increments
are at most the sum of absolute coordinate increments. Taking suprema
proves the equivalence and the converse bound. In the AC case the vector
primitive formula gives the upper length bound by $\int|\gamma'|$.
Approximate $\gamma'$ in vector $L^1$ by finite interval step functions
$v_k$ (approximate coordinates and take a common partition). Their vector
primitives are polygonal with length $\int|v_k|$. The length of the
difference primitive is at most $\int|\gamma'-v_k|$, so the reverse
bound follows by the triangle inequality for polygonal sums and passage
to the limit. The same argument on every subinterval proves that the
length function is $L(t)=\int_a^t|\gamma'|$. By definition the
parametrized arc-length integral is the Stieltjes integral of $h\circ\gamma$
against this length function; its density yields the displayed formula.
The function is bounded continuous and the length finite, so the integral
exists. Repeated traversals increase $L$ each time; replacing this by
integration on the geometric image requires additional injectivity data.
\end{proof}
\source{FOL 3.5; COH 4.4 and 6.3; ROY Chapter 6; STE 3.3--3.4. Endpoint choices and right-continuous normalization are part of the statements.}

% END chapters/12


% BEGIN chapters/13
\chapter{Differentiation of functions and measures}\label{ch:d11-13}

This chapter proves the differentiation statements from a finite covering
argument. The function theorem does not use RN. The measure theorem uses
the direct RN construction of Chapter~\ref{ch:d11-15}, whose proof does not
use any differentiation statement.

\begin{definition}[Averages, maximal function and Lebesgue points]\family{D11.13.averages_maximal_definition}\label{d11:13:averages-maximal-definition}
For $f\in L^1_{\rm loc}(\R^n,\mathbb K)$ put
\[
 A_rf(x)=\frac1{m(B(x,r))}\int_{B(x,r)}f(y)\,dy,
 \qquad Mf(x)=\sup_{r>0}A_r|f|(x).
\]
Local integrability means absolute integrability on every compact set.
A finite-valued representative has a Lebesgue point at $x$ when
$A_r(|f-f(x)|)(x)\to0$ as $r\downarrow0$.
\end{definition}

\begin{lemma}[Finite ball selection]\family{D11.13.finite_ball_covering}\label{d11:13:finite-ball-covering}
From a finite family of Euclidean balls select pairwise disjoint balls
$B_{j_1},\ldots,B_{j_k}$ such that the original union is contained in
$\bigcup_{\ell=1}^k3B_{j_\ell}$. Its measure is therefore at most
$3^n\sum_\ell m(B_{j_\ell})$.
\end{lemma}
\begin{proof}
Choose a ball of largest radius, keep it, and discard every ball intersecting
it. Repeat until no balls remain. A discarded ball of radius $s$ meets a
kept ball of radius $r\geq s$; any of its points lies within $r+2s\leq3r$
of the kept center. Thus the triple dilates cover. Scaling of Lebesgue
measure gives the estimate. Disjoint means disjoint open balls; possible
tangencies cause no problem.
\end{proof}

\begin{theorem}[Weak maximal estimate]\family{D11.13.maximal_weak11}\label{d11:13:maximal-weak11}
For $f\in L^1(\R^n)$, $Mf$ is measurable and
$m\{Mf>\alpha\}\leq3^n\alpha^{-1}\|f\|_1$ for $\alpha>0$.
In particular it is finite a.e.
\end{theorem}
\begin{proof}
For fixed $r$, the integral over $B(x,r)$ is continuous in $x$.
Indeed nearby balls lie in a fixed compact region, their symmetric
difference has measure tending to zero, and absolute continuity of the
integral of $|f|$ on that region makes the difference of the integrals tend
to zero. Hence $\{Mf>\alpha\}$ is open as the union over all $r$ of
open sets. For a compact subset $K$ of this open set choose at each
$x\in K$ a ball centered at $x$ whose average exceeds $\alpha$.
A finite subcover and the selection lemma give
\[
 m(K)\leq3^n\sum_\ell m(B_{j_\ell})
 <\frac{3^n}{\alpha}\sum_\ell\int_{B_{j_\ell}}|f|
 \leq\frac{3^n}{\alpha}\|f\|_1.
\]
Inner regularity of open Lebesgue sets proves the estimate. Finally
$\{Mf=\infty\}\subset\{Mf>k\}$ for every $k$, whose measures tend to
zero; thus the infinite-value set is null.
\end{proof}

\begin{theorem}[Lebesgue differentiation]\family{D11.13.lebesgue_differentiation}\label{d11:13:lebesgue-differentiation}
For every $f\in L^1_{\rm loc}(\R^n)$ almost every point is a Lebesgue
point of its fixed measurable representative. Consequently $A_rf(x)\to f(x)$
a.e. The conclusion also holds for balls containing $x$ whose diameters
tend to zero.
\end{theorem}
\begin{proof}
First suppose $f\in L^1$ and define
$Df(x)=\limsup_{r\downarrow0}A_r(|f-f(x)|)(x)$.
For fixed $r$, the absolute average defining $Df$ is measurable in $x$:
choose a Borel representative and apply Tonelli's measurability assertion
to $\one_{\{|x-y|<r\}}|f(y)-f(x)|$ on the sigma-finite Euclidean
product. Changing back to the selected Lebesgue representative changes
this function of $x$ only on a measurable null set. For fixed $x$ the
average is continuous in $r>0$
by absolute continuity of the local integral and continuity of ball volume;
thus rational radii suffice. For $g\in C_c$, continuity gives $Dg=0$.
The triangle inequality gives
$Df(x)\leq M(f-g)(x)+|f(x)-g(x)|$.
Chebyshev and the weak estimate imply
\[
 m\{Df>\eps\}\leq \frac{2(3^n+1)}{\eps}\|f-g\|_1.
\]
The $C_c$ density theorem makes the right side arbitrarily small. Thus
$Df=0$ outside a null set, by taking $\eps=1/k$.
For locally integrable $f$, apply this argument to
$f\one_{B(0,j+1)}$. At $x\in B(0,j)$ all sufficiently small balls give
the original averages. Take the countable union of the exceptional sets.
The absolute-average condition implies convergence of scalar averages.
If $B(z,r)$ contains $x$, it lies in $B(x,2r)$ and has $2^{-n}$ of the
latter's measure. The same absolute-average comparison proves the last
assertion, using one exceptional set independent of the balls chosen.
\end{proof}

\begin{corollary}[Density points]\family{D11.13.density_points}\label{d11:13:density-points}
For every Lebesgue measurable $E\subset\R^n$, almost every $x\in E$
has density one, and almost every $x\notin E$ density zero, in centered
balls. There is no assumption $m(E)<\infty$.
\end{corollary}
\begin{proof}
The indicator is locally integrable; apply differentiation to $\one_E$.
\end{proof}

\begin{proposition}[Regular shrinking families]\family{D11.13.regular_shrinking}\label{d11:13:regular-shrinking}
At a Lebesgue point $x$, suppose measurable $E_i\subset B_i$, the balls
$B_i$ contain $x$, their radii tend to zero, and
$m(E_i)\geq c\,m(B_i)>0$ for fixed $c>0$. Then
$m(E_i)^{-1}\int_{E_i}|f-f(x)|\to0$, hence the averages of $f$ tend to
$f(x)$. Cubes containing $x$ with shrinking side lengths qualify.
\end{proposition}
\begin{proof}
If $B_i=B(z_i,r_i)$, then $E_i\subset B(x,2r_i)$, and the relevant
average is at most $2^n/c$ times the absolute average over that centered
ball. For a cube of side $s$ containing $x$, use the ball centered at the
cube center of radius $\sqrt n\,s$ (increasing it slightly if necessary);
the ratio of its volume to the cube volume is a dimensional constant.
The comparison makes no claim for rectangles of unbounded eccentricity.
\end{proof}

\begin{lemma}[A weak estimate for finite measures]\label{d11:13:measure-weak-bound}
For a finite positive Borel measure $\lambda$, define
$M\lambda(x)=\sup_{r>0}\lambda(B(x,r))/m(B(x,r))$.
Then $m\{M\lambda>\alpha\}\leq3^n\lambda(\R^n)/\alpha$.
\end{lemma}
\begin{proof}
For fixed radius, $x\mapsto\lambda(B(x,r))$ is lower semicontinuous:
if $x_k\to x$, every point of $B(x,r)$ belongs to all sufficiently late
$B(x_k,r)$, and Fatou proves the assertion. Thus the level set of the
supremum is open. Repeat the finite compact-cover proof above, replacing
the disjoint integrals by disjoint $\lambda$ masses. Lebesgue inner
regularity gives the estimate.
\end{proof}

\begin{theorem}[Differentiating a finite measure]\family{D11.13.measure_differentiation}\label{d11:13:measure-differentiation}
Let $\nu$ be a finite positive Borel measure on $\R^n$ and write
$\nu=g m_n|_{\BB}+\nu_s$ for its Lebesgue decomposition on Borel sets.
For almost every $x$, the ratios $\nu(Q)/m(Q)$ over all nondegenerate
closed axis-parallel cubes containing $x$ and shrinking to it converge to
$g(x)$. The singular ratios converge to zero.
\end{theorem}
\begin{proof}
We first prove the singular assertion for centered open balls. Choose a
Borel null support $S$ of $\nu_s$. The finite Borel regularity theorem
(\cref{d11:18:metric-closed-regular}) gives for each $\eta>0$ an open
$O\supset\R^n\setminus S$ with $\nu_s(O)<\eta$. At each $x\notin S$,
all sufficiently small balls lie in $O$. Therefore, for every $\alpha>0$,
\[
 \left\{x\notin S:\limsup_{r\downarrow0}
 \frac{\nu_s(B(x,r))}{m(B(x,r))}>\alpha\right\}
 \subset\{M(\nu_s|_O)>\alpha\}.
\]
The weak measure estimate bounds its outer Lebesgue measure by
$3^n\eta/\alpha$. Let $\eta\downarrow0$ and then take rational
$\alpha>0$. Since $S$ is Lebesgue-null, the singular centered ratios
vanish a.e. A closed cube of side $s$ containing $x$ lies in
$B(x,2\sqrt n\,s)$, whose volume is a fixed multiple of $s^n$.
Positivity therefore transfers this conclusion to all such cubes at once.
For the AC part, $g\in L^1$ and the regular-shrinking theorem gives its
cube averages at every Lebesgue point. Combine the two conull sets.
All decompositions initially have Borel domain; the density may then be
represented as a Lebesgue measurable function for differentiation.
\end{proof}

\begin{theorem}[Locally finite scalar measures and nice families]\family{D11.13.regular_measure_differentiation}\label{d11:13:regular-measure-differentiation}
Let $\nu$ on Borel $\R^n$ be positive, extended signed omitting one
infinity, or finite complex. Suppose $|\nu|$ is compact-finite and outer
regular. Write $\nu=\nu_s+f m_n|_{\BB}$, with $f\in L^1_{\rm loc}$.
There is one Lebesgue-null set $N$ such that for every $x\notin N$ and
every Borel family $E_r\subset B(x,r)$ with
$m(E_r)\geq c\,m(B(x,r))>0$ for some fixed $c>0$,
\[
 \frac{\nu(E_r)}{m(E_r)}\longrightarrow f(x),\qquad
 \frac{\nu_s(E_r)}{m(E_r)}\longrightarrow0 \quad(r\downarrow0).
\]
The null set is independent of the chosen family.
\end{theorem}
\begin{proof}
Compact-finiteness makes $|\nu|$ sigma-finite by the compact balls
$\overline B(0,j)$. The positive or signed RN decomposition is therefore
available; its AC density is locally integrable since
$\int_K|f|\leq|\nu|(K)$ for compact $K$. Restrict the singular variation
to each compact ball. It is a finite positive Borel measure, to which the
singular part of the preceding proof applies. At points in the interior
of that ball, small averaging balls see the unrestricted measure. The
countable union of the exceptional sets consequently gives
$|\nu_s|(B(x,r))/m(B(x,r))\to0$ a.e. This also handles real and imaginary
components, since their variations are dominated by $|\nu_s|$.
At the same points require the Lebesgue-point property of $f$.
For every indicated family,
\[
 \frac{|\nu_s(E_r)|}{m(E_r)}\leq c^{-1}
 \frac{|\nu_s|(B(x,r))}{m(B(x,r))},\qquad
 \frac{\int_{E_r}|f-f(x)|}{m(E_r)}\leq c^{-1}A_r(|f-f(x)|)(x).
\]
Both bounds tend to zero independently of the family. Small balls have
compact closures, so all numerators are finite. For signed targets no
undefined extended subtraction occurs on them.
\end{proof}

\begin{corollary}[Approximate continuity]\family{D11.13.approximate_continuity}\label{d11:13:approximate-continuity}
At a Lebesgue point of $f$, for every $\eps>0$,
$m\{y\in B(x,r):|f(y)-f(x)|>\eps\}/m(B(x,r))\to0$.
Thus a locally integrable measurable function is approximately continuous
in this density sense a.e.
\end{corollary}
\begin{proof}
The fraction is at most $\eps^{-1}A_r(|f-f(x)|)(x)$ by Chebyshev on
the ball. Apply the Lebesgue-point condition and differentiation.
\end{proof}

\begin{example}[The maximal function need not be integrable]\family{D11.13.maximal_not_l1}\label{d11:13:maximal-not-l1}
For $f=\one_{[0,1]}$ on $\R$ and $x\geq2$, the centered interval of
radius $x$ contains $(0,1)$, so $Mf(x)\geq1/(2x)$.
Thus $f\in L^1$ but $Mf\notin L^1$, since the logarithmic tail diverges.
\end{example}
\source{STE Chapter 3; FOL 3.4; COH 6.2. Finite metric Borel regularity is the independent topological-measure input in the singular proof.}

% END chapters/13


% BEGIN chapters/14
\chapter{Signed and complex measures}\label{ch:d11-14}

The constructions here precede Radon--Nikodym. In particular, variation is
constructed from partitions, rather than defined by a density whose existence
would already require that theorem. All scalar integrals in this chapter are
absolutely integrable integrals.

\begin{definition}[Scalar measures]\family{D11.14.signed_complex_definition}\label{d11:14:signed-complex-definition}
A finite signed measure on $(X,\SSS)$ is a countably additive map
$\nu:\SSS\to\R$ with $\nu(\varnothing)=0$; a complex measure has values in
$\C$. Countable additivity includes convergence of the resulting scalar series
for every disjoint sequence. An extended signed measure instead has values in
$[-\infty,\infty]$, omits at least one infinity, and satisfies countable
additivity with defined extended sums.
\end{definition}

\begin{lemma}[Extraction of a negative set]\label{d11:14:negative-extraction}
If $\nu$ is a finite signed measure and $E\in\SSS$, there is a negative set
$N\subset E$ with $\nu(N)\leq\nu(E)$. Here negative means that every
measurable subset has nonpositive measure.
\end{lemma}
\begin{proof}
Start with $E_0=E$. If $E_{k-1}$ is not negative, put
$s_k=\sup\{\nu(A):A\subset E_{k-1},\ A\in\SSS\}>0$ and choose
$A_k\subset E_{k-1}$ with
$\nu(A_k)>\frac12\min(s_k,1)$. Remove it: $E_k=E_{k-1}\setminus A_k$.
Stop if a negative remainder is reached. Otherwise the $A_k$ are disjoint;
countable additivity says that the series of their positive measures converges
to the finite number $\nu(\bigcup_k A_k)$. Consequently $s_k\to0$.
For $N=E\setminus\bigcup_k A_k$, every measurable $B\subset N$ satisfies
$\nu(B)\leq s_k$ for every $k$, hence $\nu(B)\leq0$.
Also $\nu(N)=\nu(E)-\sum_k\nu(A_k)\leq\nu(E)$.
\end{proof}

\begin{theorem}[Hahn decomposition]\family{D11.14.hahn}\label{d11:14:hahn}
Every finite signed measure has a measurable partition $X=P\sqcup N$ into a
positive and a negative set. Any two such positive sets differ by a set on
which the measure vanishes on every measurable subset, hence by a
$|\nu|$-null set.
\end{theorem}
\begin{proof}
Let $m$ be the infimum of $\nu(A)$ over negative sets $A$ (including the empty
set). Choose negative sets $N_k$ whose values approach $m$, or tend to
$-\infty$ if $m=-\infty$. A countable union of negative sets is negative:
disjointize the union, intersect an arbitrary subset with the disjoint pieces,
and use countable additivity. Thus $N=\bigcup N_k$ is negative and
$\nu(N)\leq\nu(N_k)$. Since $\nu(N)$ is finite, $m>-\infty$ and
$\nu(N)=m$. If $P=X\setminus N$ had a subset $E$ of negative measure,
the extraction lemma would give a negative $A\subset E$ of negative measure;
then $N\sqcup A$ would contradict minimality. This proves existence.
For two partitions, every subset of $P\cap N'$ is both positive and negative,
and has measure zero. The same holds for $N\cap P'$. The partition definition
of variation below therefore gives zero variation on their union.
\end{proof}

\begin{theorem}[Jordan decomposition]\family{D11.14.jordan}\label{d11:14:jordan}
There are unique mutually singular finite positive measures
$\nu^+,\nu^-$ with $\nu=\nu^+-\nu^-$. For a Hahn partition,
$\nu^+(E)=\nu(E\cap P)$ and $\nu^-(E)=-\nu(E\cap N)$.
If $\nu=\alpha-\beta$ for finite positive measures, then
$\nu^+\leq\alpha$ and $\nu^-\leq\beta$. Moreover
$|\nu|=\nu^++\nu^-$.
\end{theorem}
\begin{proof}
The displayed restrictions are positive measures; their total masses are
finite because $\nu(P),\nu(N)$ are finite real numbers. They are concentrated
on disjoint sets. For any alternative positive difference and measurable $E$,
$\nu^+(E)=\nu(E\cap P)\leq\alpha(E\cap P)\leq\alpha(E)$;
the negative inequality follows identically. If the alternative parts are
singular, take their supporting partition and apply these inequalities in both
directions, or evaluate on the common refinement of the two partitions.
This proves uniqueness. On any partition of $E$, the triangle inequality
bounds the sum of absolute values by $\nu^+(E)+\nu^-(E)$; the partition
$E\cap P,E\cap N$ attains that bound. This proves the variation identity.
\end{proof}

\begin{theorem}[Partition variation]\family{D11.14.total_variation}\label{d11:14:total-variation}
For a signed or complex measure define
\[
 |\nu|(E)=\sup_{E=\bigsqcup_{j=1}^kE_j}\sum_{j=1}^k|\nu(E_j)|.
\]
It is a finite positive measure, dominates $|\nu(E)|$, and is zero on $E$
exactly when $\nu(A)=0$ for every measurable $A\subset E$.
\end{theorem}
\begin{proof}
Write $\nu=\nu_1+i\nu_2$. Each component is a finite signed measure, so
Jordan decomposition bounds every partition sum by
$\lambda(E)=\nu_1^+(E)+\nu_1^-(E)+\nu_2^+(E)+\nu_2^-(E)$.
Thus the supremum is finite. Refining partitions shows finite additivity:
one inequality comes from intersecting a partition with two disjoint sets;
the other comes from combining partitions arbitrarily close to their suprema.
For disjoint $A_k$, put $A=\bigcup A_k$. Finite additivity gives
$|\nu|(A)\geq\sum_k|\nu|(A_k)$. Conversely any finite partition $(E_j)$ of
$A$ satisfies
\[
 \sum_j|\nu(E_j)|\leq\sum_j\sum_k|\nu(E_j\cap A_k)|
 \leq\sum_k|\nu|(A_k).
\]
The first series is absolutely convergent because it is dominated by
$\lambda(A)$; countable additivity of $\nu$ justifies its sum. Take the
supremum. The null-set assertion follows by using the two-piece partition
$A,E\setminus A$ and, in the converse direction, every partition of $E$.
\end{proof}

\begin{corollary}[Finite variation is automatic]\family{D11.14.complex_finite_variation}\label{d11:14:complex-finite-variation}
Every countably additive real- or complex-valued measure on a sigma-algebra
has finite total variation. In the complex case it is a linear combination
$\alpha_1-\alpha_2+i(\alpha_3-\alpha_4)$ of four finite positive measures.
\end{corollary}
\begin{proof}
The real and imaginary components satisfy the definition of a finite signed
measure, including finiteness on every set. Apply Hahn and Jordan to each.
The finite bound $|\nu|\leq\sum_{j=1}^4\alpha_j$ in the preceding proof
establishes the assertion without a density or a uniform-boundedness theorem.
\end{proof}

\begin{proposition}[Variation inequalities]\family{D11.14.variation_inequalities}\label{d11:14:variation-inequalities}
For scalar $c$ and finite scalar measures,
$|c\nu|=|c||\nu|$, $|\nu+\eta|(E)\leq|\nu|(E)+|\eta|(E)$, and
$|(\nu|_A)|(E)=|\nu|(E\cap A)$, where $\nu|_A(E)=\nu(E\cap A)$.
Consequently $\|\nu\|_{\rm var}=|\nu|(X)$ is a norm.
\end{proposition}
\begin{proof}
Scalar multiplication and the triangle inequality apply to every partition
sum. For restriction, intersect partitions with $A$ for one inequality, and
adjoin $E\setminus A$ to partitions of $E\cap A$ for the reverse one.
Zero norm implies $\nu(E)=0$ for every $E$; the other norm axioms are the
displayed inequalities.
\end{proof}

\begin{definition}[Mutual singularity]\family{D11.14.mutual_singularity}\label{d11:14:mutual-singularity}
Write $\nu\perp\eta$ if some measurable $A$ has
$|\nu|(X\setminus A)=0$ and $|\eta|(A)=0$; for a positive measure its
variation means itself. Equivalently the two measures are concentrated on
disjoint measurable parts. Taking the complementary part proves symmetry.
\end{definition}

\begin{definition}[Scalar integration]\family{D11.14.scalar_integration}\label{d11:14:scalar-integration}
A measurable scalar $f$ is integrable against $\nu$ when
$\int|f|\,d|\nu|<\infty$. Define its integral by Jordan differences and,
for complex $\nu$, real and imaginary differences. This definition is linear,
independent of the positive decomposition, and gives
$\int\one_E\,d\nu=\nu(E)$.
\end{definition}
\begin{proof}
Each positive component is bounded by $|\nu|$ (apply partitions to its real
or imaginary restriction), so all component integrals exist absolutely.
For simple $f=\sum c_j\one_{E_j}$ the definition is exactly
$\sum c_j\nu(E_j)$, independent of decomposition and partition. Approximate
an integrable $f$ in $L^1(|\nu|)$ by simple functions. The simple-function
estimate $|\sum c_j\nu(E_j)|\leq\sum|c_j||\nu|(E_j)$ makes their integrals
Cauchy and agrees with each component's limit. Thus independence and
linearity pass to the limit.
\end{proof}

\begin{proposition}[Integral bounds and dominated convergence]\family{D11.14.integral_variation_bound}\label{d11:14:integral-variation-bound}
If $f$ is $\nu$-integrable, then
$|\int_E f\,d\nu|\leq\int_E|f|\,d|\nu|$. If $f_k,f$ are measurable,
$f_k\to f$ $|\nu|$-a.e.
and $|f_k|\leq g\in L^1(|\nu|)$, then $f$ is integrable and their scalar
integrals converge.
\end{proposition}
\begin{proof}
The simple bound just proved passes to $L^1$ limits, including restriction
to $E$. Positive-measure dominated convergence gives
$\int|f_k-f|\,d|\nu|\to0$; the bound applied to the difference concludes.
\end{proof}

\begin{theorem}[Variation of a density]\family{D11.14.density_measure_variation}\label{d11:14:density-measure-variation}
For any positive $\mu$ and $h\in L^1(\mu,\mathbb K)$,
$\nu(E)=\int_Eh\,d\mu$ is a finite scalar measure and
$|\nu|(E)=\int_E|h|\,d\mu$. For real $h$, its Jordan parts have densities
$h^+,h^-$.
\end{theorem}
\begin{proof}
Dominated convergence on disjoint partial unions proves countable
additivity. The triangle inequality on partitions proves the upper bound.
For simple $s$ its disjoint level sets show $|s\mu|(E)=\int_E|s|\,d\mu$.
Take simple $s_k\to h$ in $L^1(\mu)$. The upper bound applied to $h-s_k$
and the variation triangle inequality give
\[
 \big||h\mu|(E)-|s_k\mu|(E)\big|\leq\int_E|h-s_k|\,d\mu.
\]
Also $\int_E|s_k|\to\int_E|h|$, proving equality. In the real case the
positive and negative supports are disjoint, so Jordan uniqueness identifies
the two densities.
\end{proof}

\begin{example}[Zero value versus null set]\family{D11.14.zero_set_not_null}\label{d11:14:zero-set-not-null}
For $\nu=\delta_a-\delta_b$ on $\{a,b\}$, $\nu(X)=0$ whereas
$|\nu|(X)=2$ and $\nu(\{a\})=1$. Thus cancellation of a set's value is
not the null-set condition needed for scalar integration.
\end{example}

\begin{proposition}[Normalized distribution functions]\family{D11.14.normalized_measure_function_bridge}\label{d11:14:normalized-measure-function-bridge}
For a finite scalar Borel measure on $\R$,
$F_\nu(x)=\nu(( -\infty,x])$ is right-continuous, globally BV and zero
at $-\infty$. Its interval differences recover $\nu$ and its variation
function recovers $|\nu|$. Conversely every normalized right-continuous BV
function gives exactly one such measure.
\end{proposition}
\begin{proof}
Continuity of finite measures from above and below, applied with the
variation bound, gives right continuity and the limit at $-\infty$.
Partition sums of $F_\nu$ are bounded by $|\nu|(\R)$. The converse and the
exact variation assertion are proved in
\cref{d11:12:normalized-stieltjes}. That proof uses only the partition
variation constructed above and interval-measure extension, not RN. These
normalizations are essential: arbitrary changes on a countable set can
alter BV point values and atoms.
\end{proof}

\begin{proposition}[Extended signed measures]\family{D11.14.extended_signed_variant}\label{d11:14:extended-signed-variant}
An extended signed measure omitting one infinity has a Hahn partition and
Jordan measures $\nu^+,\nu^-$, at least one of whose total masses is finite.
Its variation is $\nu^++\nu^-$, possibly infinite. Its selected scalar
integration domain is $L^1(|\nu|)$. Sigma-finiteness means sigma-finiteness
of $|\nu|$.
\end{proposition}
\begin{proof}
Suppose $-\infty$ is omitted; reverse signs otherwise. Negative sets have
finite values. The negative-set extraction proof remains valid: its removed
positive pieces cannot have infinite total value when the starting set has
negative finite value, because countable additivity would force that starting
value to be nonnegative infinite or force a forbidden negative-infinite
remainder. More explicitly, a set of negative finite value contains no
positive-infinite subset, since its complement would then have to have
value $-\infty$. The extraction argument therefore applies on that set.
Take the infimum over negative sets as in the Hahn proof. Their countable
union is negative by additivity and cannot have value $-\infty$; hence the
infimum is finite and attained. Its complement is positive by extraction
from any negative-valued subset. Restricting to the Hahn sets gives the
Jordan measures; the negative part has finite total mass. The same
partition proof gives their sum as variation. Uniqueness follows by
restriction to common Hahn refinements, where no infinity is subtracted
from itself. On the stated integrability domain all component integrals
are finite, so the preceding integration construction applies unchanged.
\end{proof}
\source{COH 4.1; AXL 9.2--9.3; FOL 3.1. The complex variation construction here is the direct partition construction.}

% END chapters/14


% BEGIN chapters/15
\chapter{Absolute continuity and Radon--Nikodym}\label{ch:d11-15}

The main density proof uses Hahn decomposition and a maximal integral.
It uses neither Hilbert-space representation nor $L^p$ duality. In particular
the arbitrary-measure duality theorem later in the book can use this chapter
without a dependency cycle.

\begin{definition}[Absolute continuity]\family{D11.15.absolute_continuity}\label{d11:15:absolute-continuity}
For positive measures on the same sigma-algebra, $\nu\ll\mu$ means
$\mu(E)=0\Rightarrow\nu(E)=0$ for every measurable $E$. For finite scalar
$\nu$ it means $|\nu|\ll\mu$. Equivalently, $\mu(E)=0$ implies
$\nu(E)=0$ for every measurable $E$: every measurable subset of a
$\mu$-null set is again $\mu$-null, so the partition criterion for variation
proves the reverse implication.
\end{definition}

\begin{proposition}[Small-set characterization]\family{D11.15.absolute_continuity_epsilon_delta}\label{d11:15:absolute-continuity-epsilon-delta}
For arbitrary positive $\mu$ and finite scalar $\nu$, $\nu\ll\mu$ iff for
every $\eps>0$ some $\delta>0$ satisfies
$\mu(E)<\delta\Rightarrow|\nu|(E)<\eps$. Replacing $|\nu|(E)$ by
$|\nu(E)|$ gives an equivalent criterion.
\end{proposition}
\begin{proof}
If the first criterion fails, choose $E_k$ with $\mu(E_k)<2^{-k}$ but
$|\nu|(E_k)\geq\eps$. The decreasing sets $H_j=\bigcup_{k\geq j}E_k$
have $\mu(\bigcap_jH_j)=0$, whereas finiteness and continuity from above
give $|\nu|(\bigcap_jH_j)\geq\eps$, a contradiction. Either small-set
criterion implies the null-set definition by testing a fixed null set with
arbitrarily small $\delta$. The variation criterion implies the scalar-value
criterion by $|\nu(E)|\leq|\nu|(E)$, completing all implications.
\end{proof}

\begin{theorem}[Finite positive RN theorem]\family{D11.15.rn_positive_finite}\label{d11:15:rn-positive-finite}
For finite positive measures $\nu\ll\mu$ there is a nonnegative measurable
$h\in L^1(\mu)$ with $\nu(E)=\int_Eh\,d\mu$ for every measurable $E$.
It is unique $\mu$-a.e., and can be chosen finite everywhere.
\end{theorem}
\begin{proof}
If $\mu(X)=0$, then $\nu=0$ and take $h=0$. Otherwise let $\mathcal H$
consist of measurable $g\geq0$ satisfying $\int_Eg\,d\mu\leq\nu(E)$ for
every $E$. It contains zero and is closed under maxima: partition $E$ where
one of two functions is larger. Let $s=\sup_{g\in\mathcal H}\int g\,d\mu$
and choose $g_k$ approaching $s$. Successive maxima form an increasing
sequence in $\mathcal H$; its limit $h$ belongs to $\mathcal H$ by MCT and
has integral $s\leq\nu(X)$. Thus $h$ is finite a.e. and
$\rho(E)=\nu(E)-\int_Eh\,d\mu$ is a finite positive measure with
$\rho\ll\mu$.

If $\rho(X)>0$, choose $c>0$ with $c\mu(X)<\rho(X)$. Apply finite Hahn
decomposition to $\rho-c\mu$. On its positive part $P$ every subset $A$
satisfies $\rho(A)\geq c\mu(A)$. Moreover
$\rho(P)-c\mu(P)\geq\rho(X)-c\mu(X)>0$, so $\rho(P)>0$ and
$\mu(P)>0$ by absolute continuity. Then $h+c\one_P\in\mathcal H$:
the extra integral on $E\cap P$ is at most $\rho(E\cap P)\leq\rho(E)$.
Its total integral exceeds $s$, contradiction. Hence $\rho=0$.
If $h,k$ are densities, the sets $\{h\geq k+1/j\}$ have measure zero:
integrating their nonnegative difference gives zero, and both total
integrals are finite. Their union and the reversed sets prove uniqueness.
Replacing the measurable null set where $h=\infty$ by zero gives the stated
representative without changing integrals.
\end{proof}

\begin{theorem}[Sigma-finite positive RN theorem]\family{D11.15.rn_positive_sigmafinite}\label{d11:15:rn-positive-sigmafinite}
If both positive $\mu,\nu$ are sigma-finite and $\nu\ll\mu$, there is a
measurable nonnegative density $h$, finite $\mu$-a.e., unique $\mu$-a.e.
It need not have finite total integral.
\end{theorem}
\begin{proof}
Intersect countable finite-measure covers for both measures and disjointize
the resulting cover to obtain $X=\bigsqcup A_k$ with both masses finite.
Apply the finite theorem to their restrictions, and define $h=h_k$ on
$A_k$. MCT and countable additivity give the identity on every $E$.
For uniqueness restrict any second density to each $A_k$, where its integral
is $\nu(A_k)<\infty$, and apply finite uniqueness. The countable union of
exceptional null sets is null. Modify it to select a finite representative.
\end{proof}

\begin{theorem}[Finite scalar RN theorem]\family{D11.15.rn_signed_complex}\label{d11:15:rn-signed-complex}
If $\mu$ is sigma-finite and a finite signed or complex $\nu$ satisfies
$\nu\ll\mu$, there is a unique $h\in L^1(\mu,\mathbb K)$ with
$\nu(E)=\int_Eh\,d\mu$, and $|\nu|(E)=\int_E|h|\,d\mu$.
\end{theorem}
\begin{proof}
The four positive Jordan components of the real and imaginary parts are
dominated by $|\nu|$, hence absolutely continuous. Positive RN supplies
their densities, each with finite integral. Combine them with their scalar
coefficients to obtain $h\in L^1$. Variation of a density gives the asserted
identity. For uniqueness, $\int_E(h-k)=0$ for every $E$ implies zero
variation of the density measure, hence $\int|h-k|=0$.
\end{proof}

\begin{lemma}[AC and singular implies zero]\family{D11.15.ac_singular_zero}\label{d11:15:ac-singular-zero}
A finite scalar measure both absolutely continuous and singular relative to
$\mu$ is zero. Thus finite AC/singular decompositions, if they exist, are
unique.
\end{lemma}
\begin{proof}
On a $\mu$-null supporting set its variation is zero by absolute continuity;
on the complement it is zero by singularity. For uniqueness, the difference
of the two AC parts is absolutely continuous and equals the difference of
the two singular parts. The union of their null supporting sets supports
that difference, so it is zero.
\end{proof}

\begin{theorem}[Finite-target Lebesgue decomposition]\family{D11.15.lebesgue_decomposition_finite_target}\label{d11:15:lebesgue-decomposition-finite-target}
For any positive reference measure $\mu$, and every finite scalar measure
$\nu$, there are unique finite scalar measures
$\nu=\nu_a+\nu_s$, $\nu_a\ll\mu$, $\nu_s\perp\mu$.
Positive $\nu$ has positive parts. No sigma-finiteness of $\mu$, nor a
density for $\nu_a$, is asserted.
\end{theorem}
\begin{proof}
Put $b=\sup\{|\nu|(B):\mu(B)=0\}$. Choose null $B_k$ with variations
approaching $b$ and let $B=\bigcup B_k$. Finiteness gives
$|\nu|(B)=b$. Set $\nu_s(E)=\nu(E\cap B)$ and
$\nu_a(E)=\nu(E\setminus B)$. The first is singular. If $\mu(E)=0$,
then $B\cup E$ is null and
$b\geq|\nu|(B\cup E)=b+|\nu|(E\setminus B)$, so the latter variation
is zero. Thus $\nu_a\ll\mu$. Restrictions retain positivity.
Uniqueness is the preceding lemma. This construction never deduces a
reference-sigma-finite support from the finiteness of the target.
\end{proof}

\begin{theorem}[Sigma-finite Lebesgue decomposition]\family{D11.15.lebesgue_decomposition_sigmafinite}\label{d11:15:lebesgue-decomposition-sigmafinite}
For sigma-finite positive $\mu,\nu$ there are unique positive sigma-finite
$\nu_a\ll\mu$, $\nu_s\perp\mu$ with $\nu=\nu_a+\nu_s$;
$\nu_a=h\mu$ for a nonnegative measurable density unique $\mu$-a.e.
For finite scalar $\nu$ the analogous density is in $L^1(\mu)$.
\end{theorem}
\begin{proof}
The measure $\rho=\mu+\nu$ is sigma-finite, using intersections of the two
finite covers. RN gives $a=d\mu/d\rho$, $b=d\nu/d\rho$; uniqueness gives
$a+b=1$ $\rho$-a.e. Choose representatives with this identity outside one
null set, and put $Z=\{a=0\}$, assigning the exceptional set to $Z$.
Then $\mu(Z)=0$. On $Z^c$ put $h=b/a$, and put $h=0$ on $Z$.
The elementary density-change identity proved below gives
$\int_Eh\,d\mu=\int_{E\setminus Z}b\,d\rho=\nu(E\setminus Z)$.
Thus restrictions to $Z,Z^c$ give the parts. They are sigma-finite since
they are bounded by $\nu$.

For any two decompositions take the union $Z'$ of their null supports.
Their AC parts vanish on $Z'$, while on its complement both equal the
original measure. This proves equality without subtracting infinities.
Density uniqueness is RN uniqueness. Finite scalar targets can instead use
the preceding maximal-null-set construction and scalar RN on its AC part;
the variation formula makes the density $L^1$.
\end{proof}

\begin{proposition}[Changing the measure in an integral]\family{D11.15.density_change_integral}\label{d11:15:density-change-integral}
For positive $\mu$ and measurable $h\geq0$, define $\nu=h\mu$.
Then $\int g\,d\nu=\int gh\,d\mu$ for every measurable $g\geq0$,
with $0\cdot\infty=0$. If $h$ is finite a.e., a scalar $g$ is
$\nu$-integrable exactly when $\int|g|h\,d\mu<\infty$, and then the
same scalar identity holds. For finite scalar $\nu=k\mu$ it holds whenever
$\int|g||k|\,d\mu<\infty$.
\end{proposition}
\begin{proof}
For indicators this is the definition, for nonnegative simple functions it
is finite additivity, and for nonnegative functions it follows by MCT.
Apply that identity to $|g|$ before subtracting real positive/negative parts
and combining imaginary parts. In the scalar-density case
$|\nu|=|k|\mu$; prove the formula for simple $g$, then approximate in
$L^1(|k|\mu)$ and use the integral variation bound. This proof requires
no RN theorem and may therefore be used in its decomposition construction.
\end{proof}

\begin{proposition}[RN chain and reciprocal rules]\family{D11.15.rn_chain_rule}\label{d11:15:rn-chain-rule}
If $\nu\ll\mu\ll\lambda$ are sigma-finite positive measures, their
nonnegative densities satisfy
$d\nu/d\lambda=(d\nu/d\mu)(d\mu/d\lambda)$ $\lambda$-a.e., with
measurable zero conventions. If $\lambda\ll\mu$ too, the reciprocal
densities have product one both $\lambda$- and $\mu$-a.e.
\end{proposition}
\begin{proof}
Writing $h=d\nu/d\mu$, $k=d\mu/d\lambda$, the nonnegative
density-change formula gives $\nu(E)=\int_Ehk\,d\lambda$.
Take finite representatives, setting the product zero where $k=0$; any
$\mu$-exceptional set where $h$ was changed has $\int k\,d\lambda=0$,
so this convention changes no integral. Uniqueness proves the chain rule.
Apply it to $\lambda\ll\mu\ll\lambda$; the density of a measure with
respect to itself is one. Equivalent measures have the same null sets, so
the reciprocal identity holds for both.
\end{proof}

\begin{theorem}[Polar decomposition]\family{D11.15.polar_decomposition}\label{d11:15:polar-decomposition}
A finite complex measure has a measurable density $u$ relative to $|\nu|$
with $|u|=1$ $|\nu|$-a.e., unique a.e. One may choose $|u|=1$ everywhere.
For real $\nu$ one may choose $u\in\{-1,1\}$.
\end{theorem}
\begin{proof}
The finite measure $|\nu|$ dominates $\nu$, so scalar RN supplies $u$.
Variation of a density says $|\nu|(E)=\int_E|u|\,d|\nu|$, and density
uniqueness gives $|u|=1$ a.e. Change $u$ to one on its exceptional
measurable null set. Real Jordan decomposition gives $u=1$ on the positive
Hahn part and $-1$ on the negative part, with an arbitrary value on null
overlaps. RN uniqueness gives the claimed uniqueness.
\end{proof}

\begin{example}[Failure without reference sigma-finiteness]\family{D11.15.rn_non_sigmafinite_failure}\label{d11:15:rn-non-sigmafinite-failure}
On Borel $[0,1]$, let $\mu$ be counting measure and $\nu$ Lebesgue
measure. Only the empty set is $\mu$-null, so $\nu\ll\mu$.
A nonnegative density would satisfy $h(x)=\int_{\{x\}}h\,d\mu=0$ at
every point, contrary to $\nu([0,1])=1$. On a single point, the measures
$\mu(\{a\})=\infty$, $\nu(\{a\})=1$ give an even smaller obstruction:
a nonnegative function has integral zero if its value is zero and infinity
otherwise. Also no reference-sigma-finite set can carry the target mass.
These examples do not contradict the arbitrary-reference finite-target
decomposition, which asserts parts but no density.
\end{example}

\begin{proposition}[The conditioning interface]\family{D11.15.conditioning_route}\label{d11:15:conditioning-route}
On a probability space $(\Omega,\SSS,P)$, for a sub-sigma-algebra
$\mathcal T$ and integrable real $X$, there is a $\mathcal T$-measurable
$Y\in L^1(P)$ with $\int_B Y\,dP=\int_BX\,dP$ for every $B\in\mathcal T$,
unique $P$-a.e. This interface asserts no kernels, disintegration or
martingale theorem.
\end{proposition}
\begin{proof}
On $\mathcal T$ the map $B\mapsto\int_BX\,dP$ is a finite signed measure
absolutely continuous relative to the finite restriction of $P$. Scalar RN
gives exactly $Y$ and its uniqueness. Integrability is with the restricted
measure, hence also with $P$ for this $\mathcal T$-measurable function.
\end{proof}

\begin{corollary}[Sigma-finite extended-signed version]\family{D11.15.lebesgue_rn_sigmafinite_signed}\label{d11:15:lebesgue-rn-sigmafinite-signed}
Let $\mu$ be positive sigma-finite and let an extended signed $\nu$ omit
one infinity and have sigma-finite $|\nu|$. There are unique signed parts
$\nu=\nu_a+\nu_s$, with sigma-finite variations, $\nu_a\ll\mu$ and
$\nu_s\perp\mu$. There is a real measurable $f$, finite a.e., such that
\[
 \nu_a(E)=\int_E f^+\,d\mu-\int_E f^-\,d\mu,
\]
where one of the two total integrals is finite. The density is unique a.e.
On $E$ with $|\nu_a|(E)<\infty$, $f\one_E\in L^1(\mu)$.
If $\nu\ll\mu$ then $\nu_s=0$.
\end{corollary}
\begin{proof}
Decompose the positive and negative Jordan measures by the positive
sigma-finite theorem. Each decomposition is by restriction to a common
$\mu$-null set and its complement: take the union of the two null supports.
Their differences therefore inherit the finite opposing Jordan mass from
$\nu$. This makes every displayed signed sum defined. RN on the AC
positive parts gives densities $h_+,h_-$. These parts are mutually singular,
since the original Jordan parts were, so $h_+h_-=0$ a.e. Put
$f=h_+-h_-$ after removing their null infinite-value sets. Then
$f^\pm=h_\pm$ a.e. and the asserted integrals and local $L^1$ statement
follow. Uniqueness follows on the complement of a common null singular
support, and for densities on a common countable cover finite for the
variations. No step subtracts two infinite values. Absolute continuity
annihilates the singular restriction.
\end{proof}
\source{FOL 3.2; COH 4.2; AXL 9.35--9.41. The finite-target arbitrary-reference argument is kept separate from density existence.}

% END chapters/15


% BEGIN chapters/16
\chapter{Riemann and Lebesgue integration}\label{ch:d11-16}

Riemann integrability concerns a particular bounded function at every
point. Lebesgue integrals depend only on a.e. classes. The precise
relationship is governed by oscillation.

\begin{theorem}[Lebesgue's Riemann criterion]\family{D11.16.riemann_criterion}\label{d11:16:riemann-criterion}
For bounded real $f$ on a finite interval $[a,b]$, $a<b$, Riemann
integrability is equivalent to its discontinuity set being Lebesgue-null.
In particular every such Riemann-integrable function is Lebesgue measurable.
Continuity and discontinuity at the endpoints use the relative topology.
\end{theorem}
\begin{proof}
Let $\omega_f(x)$ be the infimum over relative neighborhoods of $x$ of
the supremum of $|f(u)-f(v)|$ for $u,v$ in the neighborhood. It is zero
exactly at continuity points. The sets $D_k=\{\omega_f\geq1/k\}$ are
closed: if a neighborhood has oscillation below $1/k$, the same is true
at all its points in a sufficiently small neighborhood.

Suppose $f$ is Riemann integrable. Given $\eta>0$, take a partition with
upper minus lower sum below $\eta$. Apart from its finitely many endpoints,
every point of $D_k$ lies in a partition interval whose oscillation is
at least $1/k$. The sum of lengths of these intervals is at most $k\eta$.
They cover $D_k$ except for the finite endpoints, so $m(D_k)\leq k\eta$.
Let $\eta\downarrow0$ and take the union over $k$.

Conversely suppose each $D_k$ is null. Write $|f|\leq M$. Given
$\eps>0$, choose $k$ with $(b-a)/k<\eps/2$, and a relatively open
$O\supset D_k$ with total length less than $\eps/(4M)$ if $M>0$.
The case $M=0$ is immediate. Cover the compact set $[a,b]\setminus O$
by finitely many relative open neighborhoods on each of which the
oscillation is below $1/k$. Together with $O$ these form a finite open
cover of $[a,b]$. Its Lebesgue number exists by compactness: otherwise
subintervals of diameter tending to zero not contained in any member
would have a limiting point whose neighborhood is a cover member,
contradiction. Choose a partition with mesh below that number. Each
closed partition interval is either contained in $O$, or in one of the
small-oscillation neighborhoods. The first type has total length at most
$m(O)$, and the second contributes at most $(b-a)/k$ to upper minus
lower sum. The total is below $2M m(O)+(b-a)/k<\eps$.
This is Darboux's integrability criterion.

Finally, for Riemann-integrable $f$ choose partitions with upper minus
lower sums tending to zero and refine successively. Their lower and upper
step functions $l_j,u_j$ satisfy $l_j\uparrow l$, $u_j\downarrow u$,
and $l\leq f\leq u$, with $\int(u_j-l_j)\to0$.
Thus $l=u$ a.e. The bounded $f$ differs from the measurable $l$ only
on a subset of a measurable null set; completeness gives its Lebesgue
measurability. This last argument also handles arbitrary values of $f$
on its null discontinuity set.
\end{proof}

\begin{theorem}[Equality of proper integrals]\family{D11.16.proper_integral_equality}\label{d11:16:proper-integral-equality}
A bounded real Riemann-integrable function on finite $[a,b]$ is in
$L^1$ and its Riemann and Lebesgue integrals agree. The same holds for
complex functions whose two components are Riemann integrable, and for
bounded Riemann-integrable scalar functions on a compact Euclidean rectangle.
\end{theorem}
\begin{proof}
Use the measurable lower and upper step functions in the preceding proof.
Their integrals are exactly the Darboux sums, since individual partition
endpoints are null. They squeeze the Lebesgue integral of $f$, which
exists absolutely because the interval has finite measure and $f$ is
bounded. Their gap tends to zero, and their common limit is the Riemann
integral. Apply this to real and imaginary components for complex $f$.
On a rectangle, use rectangular partitions and their step functions;
faces have $n$-dimensional measure zero by the product theorem. Refining
partitions with Darboux gap tending to zero gives the same measurable
lower/upper limits and the identical squeeze. This uses the ordinary
rectangle Riemann convention, without an unproved substitution theorem.
\end{proof}

\begin{theorem}[Absolutely convergent improper integrals]\family{D11.16.absolute_improper_bridge}\label{d11:16:absolute-improper-bridge}
Let $f:[a,\infty)\to\mathbb K$ be Riemann integrable on every compact
$[a,R]$. Then
$\int_{[a,\infty)}|f|\,dm=\lim_{R\to\infty}\int_a^R|f|$ as an
extended nonnegative value. Absolute improper convergence is equivalent
to $f\in L^1$, and in that case the improper scalar value is its Lebesgue
integral. Finite singular endpoints and two-sided intervals require the
corresponding separate one-sided absolute limits.
\end{theorem}
\begin{proof}
The absolute value is Riemann integrable on each compact interval: its
oscillation is at most that of the scalar function, componentwise if
needed. Proper equality applies, and the countably many compact
restrictions make $f$ measurable on the half-line. MCT on expanding
intervals proves the displayed absolute-integral identity. If it is
finite, DCT applied to $f\one_{[a,R]}$, dominated by $|f|$, gives scalar
integral convergence and proper equality identifies each term. Conversely
absolute improper convergence makes the displayed integral finite.
For a finite singular endpoint use a sequence of truncated compact
intervals increasing to the one-sided domain; for two ends do so on each
side separately. The same MCT/DCT proof works. A principal value alone
does not provide either absolute limit.
\end{proof}

\begin{example}[Conditional improper convergence]\family{D11.16.conditional_improper_counterexample}\label{d11:16:conditional-improper-counterexample}
Let $f=0$ below one and $f(x)=(-1)^n/n$ on $[n,n+1)$, $n\geq1$.
Its improper Riemann integral on $[0,\infty)$ exists, whereas
$\int|f|=\infty$ and both its positive and negative integrals are infinite.
\end{example}
\begin{proof}
Every compact restriction is a finite step function, hence Riemann
integrable. At integer endpoints the integrals are partial alternating
harmonic sums, which converge. The remainder on a final partial interval
has absolute value at most $1/\lfloor R\rfloor$, so arbitrary endpoints
have the same limit. The absolute integral is $\sum_{n\geq1}1/n$.
The even and odd subseries each diverge, proving both signed parts
infinite. Thus there is no ordinary signed Lebesgue integral of $f$.
\end{proof}

\begin{example}[A Lebesgue-integrable function which is not Riemann integrable]\family{D11.16.lebesgue_not_riemann}\label{d11:16:lebesgue-not-riemann}
On $[0,1]$, $f=\one_\Q$ is Borel and zero a.e., so its Lebesgue
integral is zero. Every relative neighborhood contains both rational and
irrational points, so $f$ is discontinuous everywhere and is not Riemann
integrable. A.e. equality therefore preserves the Lebesgue integral but
not Riemann integrability of representatives.
\end{example}

\begin{remark}[Gauge integrals: the boundary of this comparison]\family{D11.16.gauge_integral_boundary}\label{d11:16:gauge-integral-boundary}
Cohn's Appendix H introduces Henstock--Kurzweil and McShane integrals
with tagged partitions and gauges. The present comparison proves the
proper and absolute-improper Riemann relationships above; it does not
assert a general gauge-integral equivalence or an existence theorem.
A future such statement must fix the interval, scalar field and precise
gauge and tag conventions. This is a scope boundary, not an invoked
mathematical premise.
\end{remark}
\source{AXL 3; ROY 5; COH 2.5 and Appendix H; FOL 2.3; STE 2.1.}

% END chapters/16


% BEGIN chapters/17
\chapter{Uniform integrability and Vitali convergence}\label{ch:d11-17}

Throughout this chapter $(X,\Sigma,\mu)$ is a positive measure space. Functions are measurable scalar representatives; integral conditions descend to their almost-everywhere classes. We distinguish three different controls. A supremum over an empty family of nonnegative quantities is zero. No atomlessness or finiteness of $\mu(X)$ is implicit.

\section{Three controls and their distinctions}

\begin{definition}[Uniform absolute continuity of integrals (small-set uniform integrability)]\family{D11.17.uniform_absolute_continuity}\label{d11:17:uniform-absolute-continuity}
A family $\mathcal F\subset L^1(\mu)$ has \emph{small-set uniform integrability}, denoted $\mathrm{UI}_{\rm small}$, if for every $\eps>0$ there is $\delta>0$ such that
\[ A\in\Sigma,\quad \mu(A)<\delta\quad\Longrightarrow\quad
 \sup_{f\in\mathcal F}\int_A|f|\dd\mu<\eps.\]
The functions in the family are individually integrable; a uniform bound on their norms is a separate assertion.
\end{definition}
\source{ROY 11.3 definition; ROY 5.1 definition}

\begin{definition}[Uniform control outside sets of finite measure]\family{D11.17.finite_measure_tails}\label{d11:17:finite-measure-tails}
The family has \emph{finite-measure spatial tails} if for every $\eps>0$ some measurable $E$ with $\mu(E)<\infty$ satisfies
\[\sup_{f\in\mathcal F}\int_{X\setminus E}|f|\dd\mu<\eps.\]
For an $L^p$ family we apply this condition to $|f|^p$. Royden calls this tightness of integrable functions. The set $E$ need not be compact; compact tightness of measures is a different condition.
\end{definition}
\source{ROY 11.3 definition; FOL 6.1 Exercise15(iii)}

\begin{definition}[Uniform large-value tail criterion]\family{D11.17.amplitude_tails}\label{d11:17:amplitude-tails}
Write $\mathrm{UI}_{\rm tail}$ for
\[\lim_{R\to\infty}\sup_{f\in\mathcal F}\int_{\{|f|>R\}}|f|\dd\mu=0.\]
Uniform $L^1$ boundedness means $\sup_{f\in\mathcal F}\|f\|_1<\infty$. Neither condition is included by definition in $\mathrm{UI}_{\rm small}$.
\end{definition}
\source{FOL 6.1 Exercise15 and 3.2 Exercise11; ROY 11.3 Proposition14 and definitions}

\begin{proposition}[Small sets, norm bounds and amplitude tails]\family{D11.17.definition_comparison}\label{d11:17:definition-comparison}
On any positive measure space, $\mathrm{UI}_{\rm tail}$ implies $\mathrm{UI}_{\rm small}$. Small-set control together with a uniform $L^1$ bound implies tail control. If $\mu(X)<\infty$, tail control also implies a uniform $L^1$ bound.
\end{proposition}
\begin{proof}
Choose $R$ with the tail integrals below $\eps/2$. Then
\[\int_A|f|\le R\mu(A)+\int_{\{|f|>R\}}|f|,\]
so $\delta=\eps/(2R)$ works when $R>0$; we may enlarge $R$ to be positive. Conversely let $\|f\|_1\le M$ uniformly. Chebyshev gives $\mu\{|f|>R\}\le M/R$. Choose the small-set $\delta$ first, then $R>M/\delta$. If $M=0$, all relevant integrals vanish. Finally on a finite space the same splitting with $A=X$ gives $\|f\|_1\le R\mu(X)+\eps/2$. These arguments include the zero measure and the empty family.
\end{proof}
\source{ROY 5.1 and 11.3 small-set definitions; FOL 6.1 Exercise15}

\begin{proposition}[Translate Folland's signed-integral UI convention]\family{D11.17.folland_ui_convention}\label{d11:17:folland-ui-convention}
For real or complex $L^1$ families, uniform small-set control of $\left|\int_A f\dd\mu\right|$ is equivalent to $\mathrm{UI}_{\rm small}$.
\end{proposition}
\begin{proof}
Integral-of-absolute-value control implies the displayed control by the triangle inequality. For the converse, split a given $A$ into the measurable subsets where $\Re f$ is nonnegative and negative. The integral of $|\Re f|$ over $A$ is the difference of the real parts of the integrals of $f$ over these two subsets, hence is at most twice the common signed-integral bound. Repeating for $\Im f$ and using $|f|\le|\Re f|+|\Im f|$ gives a factor at most four. Apply the assumed condition with $\eps/4$; real functions need only two pieces. The subsets may depend on $f$, which is permitted by the uniform quantifiers.
\end{proof}
\source{FOL 3.2 Exercise11; ROY 11.3 definition}

\begin{example}[Small-set UI is vacuous on an indivisible atom]\family{D11.17.atomic_ui_counterexample}\label{d11:17:atomic-ui-counterexample}
On the one-point probability space, $f_n=n$ has small-set control but unbounded $L^1$ norms and no amplitude-tail control.
\end{example}
\begin{proof}
Every measurable set of measure less than one is empty, so it witnesses any small-set bound. But $\|f_n\|_1=n$, and for every $R$ the supremum of $\int_{\{|f_n|>R\}}|f_n|$ is infinite. Thus atomlessness cannot be silently used in a comparison of definitions.
\end{proof}
\source{ROY 11.3 definition and general-space warning}

\section{Uniform controls and convergence}

\begin{proposition}[Finite and dominated families satisfy both controls]\family{D11.17.dominated_and_finite_families}\label{d11:17:dominated-and-finite-families}
Finite subsets of $L^1$ satisfy both small-set control and finite-measure spatial tails. So does every family dominated almost everywhere by one $g\in L^1$, $g\ge0$.
\end{proposition}
\begin{proof}
For one integrable $g$, choose a bounded simple $s$ of finite-measure support with $0\le s\le g$ and $\int(g-s)<\eps/2$, by simple approximation and monotone convergence. Then $\int_Ag\le\|s\|_\infty\mu(A)+\eps/2$, and outside the support of $s$ the integral is below $\eps/2$. These give both controls. For a finite family take the minimum of its finitely many small-set thresholds and the union of its finitely many finite-measure supports. For a dominated family use $\int_A|f|\le\int_Ag$ and the same inequality on complements. Representatives need only satisfy the domination almost everywhere.
\end{proof}
\source{ROY 5.1 Proposition3 and example; ROY 11.3 Proposition14 and following paragraph; FOL 3.2 Exercise11(a)}

\begin{proposition}[L1 convergence implies the uniform controls]\family{D11.17.norm_convergence_controls}\label{d11:17:norm-convergence-controls}
If $f_n\to f$ in $L^1$, the family $\{f_n\}$ has both controls and bounded norms. More generally, if $f_n\to f$ in $L^p$, $1\le p<\infty$, then the family $\{|f_n|^p\}$ has both controls and bounded $L^1$ norms.
\end{proposition}
\begin{proof}
For $L^1$, choose $N$ with $\|f_n-f\|_1<\eps/3$ for $n\ge N$. Apply the preceding proposition to $f$ and the finite initial segment, with error $\eps/3$. The triangle inequality gives the required bounds uniformly for the tail as well. Norm boundedness follows from convergence. For general $p$ use
\[|a|^p\le2^{p-1}(|a-b|^p+|b|^p).\]
Choose $N$ with $\|f_n-f\|_p^p$ sufficiently small, then use the individual controls for $|f|^p$ and the finite initial segment. The same bound on a set and its complement proves the claim; convergence already bounds $\|f_n\|_p^p$.
\end{proof}
\source{FOL 3.2 Exercise11(b); FOL 6.1 Exercise15 necessity}

\begin{theorem}[Vitali criterion for Lp Cauchy sequences]\family{D11.17.vitali_lp_cauchy}\label{d11:17:vitali-lp-cauchy}
Let $1\le p<\infty$ and $f_n\in L^p(\mu)$. Then $(f_n)$ is norm-Cauchy if and only if it is Cauchy globally in measure, $\{|f_n|^p\}$ has small-set control, and this family has finite-measure spatial tails.
\end{theorem}
\begin{proof}
If the sequence is norm-Cauchy, completeness \cref{d11:10:completeness} gives a norm limit. The preceding proposition supplies both controls. Chebyshev applied to $|f_n-f_m|^p$ proves the Cauchy-in-measure condition.

Conversely fix a desired integral error $\eps>0$, and put $C=2^{p-1}$. Choose finite-measure $E$ such that every outside integral of $|f_n|^p$ is below $\eps/(6C)$. The integral of $|f_n-f_m|^p$ outside $E$ is then below $\eps/3$. Choose $\delta>0$ so that each integral of $|f_n|^p$ on a set of measure less than $\delta$ is below $\eps/(6C)$. Choose $\eta>0$ with $\eta^p\mu(E)<\eps/3$; if $\mu(E)=0$ any positive $\eta$ works. For large $n,m$, the measurable set $A=E\cap\{|f_n-f_m|>\eta\}$ has measure less than $\delta$. Its difference integral is below $\eps/3$ by the elementary power inequality. On $E\setminus A$ it is at most $\eta^p\mu(E)$. Adding the three estimates proves norm-Cauchyness. No finiteness of $\mu(X)$ was used.
\end{proof}
\source{FOL 6.1 Exercise15}

\begin{corollary}[Vitali convergence from convergence in measure]\family{D11.17.vitali_lp_convergence}\label{d11:17:vitali-lp-convergence}
For measurable $f$, $f_n\in L^p$ converges to $f$ in $L^p$ if and only if it converges globally in measure to $f$ and the two uniform controls in the preceding theorem hold. The finite-measure spatial-tail condition is automatic when $\mu(X)<\infty$.
\end{corollary}
\begin{proof}
Norm convergence implies convergence in measure by \cref{d11:07:norm-to-measure}, and the controls were proved above. Conversely convergence in measure implies Cauchyness in measure by the triangle inequality for exceptional sets. The theorem gives an $L^p$ limit $g$. This sequence also converges in measure to $g$. Uniqueness of a measurable in-measure limit, \cref{d11:07:measure-subsequence-uniqueness}, gives $f=g$ almost everywhere. Thus the given measurable $f$ has the required integrable class and the norm difference tends to zero. When the whole space has finite measure, take $E=X$.
\end{proof}
\source{FOL 6.1 Exercise15 with Lp completeness and in-measure uniqueness; ROY 11.3 Vitali comparison}

\begin{proposition}[Preserve the distinct pointwise Vitali statements]\family{D11.17.royden_vitali_versions}\label{d11:17:royden-vitali-versions}
For Lebesgue measurable $E\subset\R$ of finite measure, a small-set uniformly integrable sequence converging pointwise to a finite measurable $f$ has $f\in L^1$ and convergent integrals. On arbitrary Lebesgue $E$, add finite-measure spatial tails. Royden's general-measure variant, with these two controls and pointwise convergence, retains the stated extra assumption $f\in L^1$.
\end{proposition}
\begin{proof}
On finite Lebesgue $E$, partition $E$ into finitely many measurable sets of measure below a small-set threshold for error one. Such a partition follows by first cutting off a small tail and then partitioning a bounded interval into sufficiently short intervals. The uniform control bounds all $\int_E|f_n|$ by the number of pieces. Fatou gives $f\in L^1$. Finite-space pointwise convergence implies convergence in measure by \cref{d11:07:finite-ae-to-measure}; Vitali convergence then applies.

On arbitrary Lebesgue $E$, choose a finite-measure set controlling the outside integrals. Fatou controls the outside integral of $|f|$. The finite-domain argument applies on that set and makes its $L^1$ difference small; the two outside integrals make the remaining difference small. For the stated general-measure variant, $f$ is already integrable. On each finite-measure truncation, Fatou also transfers small-set control to $f$; Egorov \cref{d11:07:egorov} and the split into a small exceptional set and a uniform-convergence set prove $L^1$ convergence there. Spatial tails and Fatou finish the same estimate. We have kept, rather than erased, this source's integrable-limit hypothesis.
\end{proof}
\source{ROY 5.1 Vitali and General Vitali; ROY 11.3 Vitali and following explanation}

\begin{example}[L1 boundedness does not imply uniform integrability]\family{D11.17.spikes_counterexample}\label{d11:17:spikes-counterexample}
On $(0,1)$, $f_n=n\one_{(0,1/n)}$ is $L^1$ bounded and tends to zero almost everywhere and in measure, but fails small-set control and $L^1$ convergence.
\end{example}
\begin{proof}
Its norm is one. Every fixed positive point eventually lies outside its support, and the support measure is $1/n$. Nevertheless its entire integral is concentrated on that support, giving integral one on sets whose measures tend to zero.
\end{proof}
\source{STE Chapter2 integration concentration example; FOL 6.2 Exercise22(b)}

\begin{example}[Small-set control alone does not control mass at infinity]\family{D11.17.escaping_mass_counterexample}\label{d11:17:escaping-mass-counterexample}
On $\R$, $f_n=\one_{[n,n+1]}$ has small-set control and tends pointwise to zero, but has neither uniform finite-measure spatial tails nor $L^1$ or global in-measure convergence to zero.
\end{example}
\begin{proof}
For every $A$, $\int_Af_n\le m(A)$, while each norm and each measure $m\{|f_n|>1/2\}$ equals one. If a finite-measure $E$ captured more than half the integral of every $f_n$, the disjoint intervals would force $m(E)=\infty$. Thus tails fail, despite pointwise convergence and local convergence in measure.
\end{proof}
\source{ROY 5.1 infinite-domain counterexample}

% END chapters/17


% BEGIN chapters/18
\chapter{Radon regularity, tightness, and products}\label{ch:d11-18}

All locally compact spaces in this chapter are Hausdorff. Compact means topologically compact; finite measure does not mean compact support. The topological inputs are the compact-neighborhood, shrinking, and cutoff consequences of the declared topology foundation. Measure-theoretic regularity is proved below and is never inferred from its name.

\section{Definitions and the topological interface}

\begin{definition}[Borel and Baire sigma-algebras on LCH spaces]\family{D11.18.borel_baire}\label{d11:18:borel-baire}
For a topological space $X$, $\BB(X)$ is the sigma-algebra generated by the open sets. On an LCH space put
\[\operatorname{Ba}(X)=\sigma\{f^{-1}(A):f\in C_c(X,\R),\ A\in\BB(\R)\}.\]
Every continuous real function is Borel measurable, so this $C_c$-generated Baire sigma-algebra is contained in $\BB(X)$. Equality needs an additional argument.
\end{definition}
\source{FOL Chapter7 introduction and 7.1 Baire discussion}

\begin{definition}[Outer, inner, compact-finite and local-finite predicates]\family{D11.18.regularity_predicates}\label{d11:18:regularity-predicates}
A positive Borel measure is \emph{outer regular on $E$} if
$\mu(E)=\inf_{U\supset E,\ U\text{ open}}\mu(U)$, and \emph{compact-inner regular on $E$} if
$\mu(E)=\sup_{K\subset E,\ K\text{ compact}}\mu(K)$.
It is compact-finite if every compact set has finite measure, and locally finite if every point has an open neighborhood of finite measure. The collection of sets on which a regularity assertion holds is part of the assertion.
\end{definition}
\source{FOL 7.1 regularity definitions; COH 7.2 regularity definitions}

\begin{remark}[Fix Radon predicates and translate source terminology]\family{D11.18.radon_and_source_names}\label{d11:18:radon-and-source-names}
Our positive \emph{Radon} measure on LCH $X$ is a Borel measure which is compact-finite, outer regular on Borel sets, and compact-inner regular on open sets. Cohn's regular Borel measure has these properties. Folland calls this class Radon; his stronger ``regular'' includes compact-inner regularity on every Borel set. The words do not change the predicates. The pinned Mathlib predicate \texttt{Measure.Regular} has our three Radon predicates; its separate all-measurable-set inner regularity is stronger.
\end{remark}
\source{FOL 7.1 definitions; COH 7.2 definitions}

\begin{proposition}[Local finiteness and compact finiteness on LCH spaces]\family{D11.18.lch_local_finite_bridge}\label{d11:18:lch-local-finite-bridge}
For a positive Borel measure on an LCH space, compact finiteness and local finiteness are equivalent.
\end{proposition}
\begin{proof}
Every point has an open neighborhood with compact closure, so compact finiteness bounds that neighborhood's mass. Conversely cover a compact $K$ by the finite-measure neighborhoods supplied by local finiteness. A finite subcover bounds $\mu(K)$ by a finite sum. This proves no outer or inner regularity assertion.
\end{proof}
\source{COH 7.1 locally compact preliminaries and 7.2 compact finiteness}

\begin{lemma}[Selected compact-support cutoff interface]\family{D11.18.compact_cutoff}\label{d11:18:compact-cutoff}
If $K\subset U\subset X$, $K$ is compact and $U$ open in an LCH space, there is $\phi\in C_c(X,\R)$ with $0\le\phi\le1$, $\phi=1$ on $K$, and $\supp\phi\subset U$.
\end{lemma}
\begin{proof}
The topology foundation supplies, for each $x\in K$, precompact open neighborhoods with closure inside $U$. A finite union gives an open $V$ with $K\subset V\subset\overline V\subset U$ and $\overline V$ compact. A further shrinking gives $W$ with $K\subset W\subset\overline W\subset V$. On the compact Hausdorff space $\overline V$, the Urysohn separation theorem applied to $K$ and $\overline V\setminus W$ gives a function equal to one on $K$ and zero off $W$. Extend it by zero outside $\overline V$. It vanishes in a neighborhood of the boundary, so the extension is continuous and supported in $\overline W\subset U$. This is the exact separation consequence being used from topology.
\end{proof}
\source{COH 7.1 Proposition7.1.9; FOL 7.1 Proposition7.1 proof}

\begin{theorem}[Baire equals Borel under second countability]\family{D11.18.baire_borel_second_countable}\label{d11:18:baire-borel-second-countable}
In a second-countable LCH space every compact set is a $G_\delta$, and $\operatorname{Ba}(X)=\BB(X)$.
\end{theorem}
\begin{proof}
Precompact open sets form a base: shrink a neighborhood of a point within any prescribed open set. Second countability lets us choose a countable such base. An open $U$ is covered by the closures of countably many base elements whose closures lie in $U$; hence it is sigma-compact. For compact $K$, choose from the countable base every element whose closure misses $K$. Their union is $X\setminus K$; writing those compact closures as $L_n$, $K=\bigcap_n(X\setminus L_n)$, a $G_\delta$ (the finite or empty list is harmless).

For each relatively compact base element $V$ and each smaller base element with closure $K\subset V$, choose a cutoff equal to one on $K$ and supported in $V$. This is a countable family. Each open $U$ is the union of the positivity sets of the chosen cutoffs supported in base elements contained in $U$, because shrinking supplies one positive at every point of $U$. Those positivity sets belong to $\operatorname{Ba}(X)$. Thus every open belongs to it, proving equality.
\end{proof}
\source{FOL 7.1 Exercises4-5}

\section{What regularity actually gives}

\begin{theorem}[Radon measures are inner regular on sigma-finite Borel sets]\family{D11.18.radon_sigmafinite_inner}\label{d11:18:radon-sigmafinite-inner}
If $\mu$ is Radon and a Borel $E$ is a countable union of finite-measure Borel sets, then $\mu(E)=\sup_{K\subset E,\ K\text{ compact}}\mu(K)$.
\end{theorem}
\begin{proof}
The precise ingredients are \cref{d11:18:radon-and-source-names}. First suppose $\mu(E)<\infty$. Choose open $U\supset E$ with $\mu(U\setminus E)<\eps$ and $\mu(U)<\infty$. Outer regularity gives open $V\supset U\setminus E$ with $\mu(V)<2\eps$. Inner regularity on $U$ gives compact $H\subset U$ with $\mu(U\setminus H)<\eps$. Then $K=H\setminus V$ is compact, lies in $E$, and $\mu(E\setminus K)<3\eps$. For general $E$, replace its finite-measure covering by an increasing finite-measure covering $E_n$. If $a<\mu(E)$, continuity from below gives $n$ with $\mu(E_n)>a$, and the finite case supplies compact $K\subset E_n$ with $\mu(K)>a$. This also covers infinite mass. Zero mass and the empty set are immediate.
\end{proof}
\source{FOL 7.2 Proposition7.5; COH 7.2 Proposition7.2.6}

\begin{corollary}[Sigma-finite and sigma-compact regularity specializations]\family{D11.18.sigmafinite_radon_regular}\label{d11:18:sigmafinite-radon-regular}
A sigma-finite Radon measure is compact-inner and outer regular on every Borel set. A Radon measure on a sigma-compact LCH space is sigma-finite and has the same conclusion.
\end{corollary}
\begin{proof}
The precise ingredients are \cref{d11:18:radon-sigmafinite-inner}. Intersect a sigma-finite covering with the Borel set and apply the preceding theorem; outer regularity is in the definition. A countable compact covering is a finite-measure covering by compact finiteness.
\end{proof}
\source{FOL 7.2 Corollary7.6}

\begin{theorem}[Compact-finite Borel measures on second-countable LCH spaces]\family{D11.18.second_countable_borel_regular}\label{d11:18:second-countable-borel-regular}
Suppose every open set of LCH $X$ is sigma-compact. Every compact-finite positive Borel measure on $X$ is Radon and is compact-inner and outer regular on all Borel sets. In the second-countable case it is sigma-finite.
\end{theorem}
\begin{proof}
The precise ingredients are \cref{d11:18:baire-borel-second-countable}, \cref{d11:18:sigmafinite-radon-regular}. Every open $U$ is a countable union of compact sets; their finite unions increase to $U$, so it is compact-inner regular. In particular $X$ is sigma-compact. We justify outer regularity rather than assume it. Cover $X$ by precompact open sets $O_n$ with $\overline O_n\subset O_{n+1}$. Construct them recursively: include the previous compact closure together with the next compact of an exhaustion in finitely many precompact neighborhoods, and take their union. On any precompact open $O$, the restricted measure is finite. Every relatively open subset of $O$ is a union of countably many compact subsets of $X$ contained in it, hence can be approximated from within by closed subsets of $O$.

For a finite measure on any space, let $\mathcal D$ consist of sets which can, for each $\eps>0$, be placed between a closed set and an open set with difference measure below $\eps$. If all opens have closed inner approximations, they belong to $\mathcal D$; complements preserve $\mathcal D$. For $A=\bigcup_i A_i$, choose open supersets with errors summing to less than $\eps/2$, and choose $N$ so that $\mu(A\setminus\bigcup_{i\le N}A_i)<\eps/4$ by finite-measure continuity. A finite union of the corresponding closed inner approximations has remaining error below $\eps/2$. Thus $\mathcal D$ is a sigma-algebra containing the opens. Apply this argument in each $O_{n+1}$ to the Borel set $E\cap\overline O_n$, choosing a relative open superset with error $\eps2^{-n-1}$. These relative opens are open in $X$, and their union covers $E$ with total error below $\eps$. This proves global outer regularity. We now have a Radon measure, and $X$ is sigma-compact, so the preceding corollary gives all-Borel compact-inner regularity. Second countability supplies the sigma-compact-open hypothesis by the proof above.
\end{proof}
\source{FOL 7.2 Theorem7.8; COH 7.2 Propositions7.2.3 and7.2.5}

\begin{theorem}[Metric Borel regularity with finite balls]\family{D11.18.metric_closed_regular}\label{d11:18:metric-closed-regular}
Let $X$ be metric and let every finite-radius ball have finite Borel measure. For Borel $E$ and $\eps>0$, there are open $U\supset E$ and closed $F\subset E$ with $\mu(U\setminus E)<\eps$ and $\mu(E\setminus F)<\eps$. The closed $F$ need not be compact.
\end{theorem}
\begin{proof}
For a finite Borel measure on a metric space, every closed $F$ has the decreasing open neighborhoods $\{d(x,F)<1/n\}$, with intersection $F$. Continuity from above gives open approximations. Complements give closed inner approximations of opens. The finite-measure sigma-algebra argument in the preceding proof gives both approximations for all Borel sets.

If $X$ is nonempty fix $x_0$ and write $B_n=\{d(x,x_0)\le n\}$, $O_n=\{d(x,x_0)<n+1\}$. Restrict the measure to the finite-measure metric space $O_n$. For $E\cap B_n$ choose a relatively open superset $V_n\subset O_n$ with error below $\eps2^{-n-1}$. It is open in $X$. Their union covers $E$ and its excess over $E$ has measure below $\eps$. Applying this outer approximation to $X\setminus E$ and taking a complement gives the closed inner approximation. Empty $X$ requires no construction. No completeness or compactness of closed balls was asserted.
\end{proof}
\source{STE Chapter6 §1.2 Proposition1.3}

\begin{theorem}[Ulam regularity and tightness of finite Polish-space measures]\family{D11.18.polish_tightness_regular}\label{d11:18:polish-tightness-regular}
A finite Borel measure on a separable complete metric space is compactly tight. Moreover every Borel $E$ has compact $K\subset E$ and open $U\supset E$ with $\mu(E\setminus K)<\eps$ and $\mu(U\setminus E)<\eps$.
\end{theorem}
\begin{proof}
The precise ingredients are \cref{d11:18:metric-closed-regular}. For each $n$, the balls of radius $2^{-n}$ about a countable dense set cover $X$. Finiteness and continuity from below permit a finite subfamily whose union misses mass less than $\eps2^{-n-2}$. Let $A_n$ be the union of their closed balls, and put $K_0=\bigcap_n A_n$. This set is closed and totally bounded: at every scale it lies in a finite union of balls of that scale (enlarge radii slightly to use open balls). It is complete as a closed subset, hence compact by the complete, totally bounded compactness theorem from topology. The union bound gives $\mu(X\setminus K_0)<\eps/2$. Apply metric closed regularity to obtain closed $F\subset E$ missing mass below $\eps/2$ and open $U\supset E$ with the required outer error. Then $K=F\cap K_0$ is compact and has the asserted inner error. A zero measure and empty space obey the same estimates.
\end{proof}
\source{ROY 13.6 Lemma28 and Ulam regularity theorem}

\begin{corollary}[Finite Radon measures are compact tight]\family{D11.18.finite_radon_tightness}\label{d11:18:finite-radon-tightness}
If $\mu$ is a finite Radon measure on LCH $X$, every $\eps>0$ admits compact $K$ with $\mu(X\setminus K)<\eps$. For a family $\mathcal P$ of finite measures, \emph{uniform compact tightness} requires one such $K$ for all its members.
\end{corollary}
\begin{proof}
The precise ingredients are \cref{d11:18:radon-and-source-names}. Apply compact-inner regularity to the open set $X$ and its finite mass. The family assertion is a definition, not a consequence: the probabilities $\delta_n$ on $\R$ are individually Radon and tight, whereas no compact contains every integer, so no compact captures all their masses with error below one.
\end{proof}
\source{FOL 7.2 Proposition7.5; ROY 13.6 Lemma28 for comparison}

\begin{theorem}[Support from null open sets]\family{D11.18.radon_support}\label{d11:18:radon-support}
Let $N$ be the union of all open $\mu$-null sets for Radon $\mu$. Then $N$ is open and null, and $\supp\mu=X\setminus N$ is closed. A point $x$ belongs to the support exactly when every $f\in C_c(X,[0,1])$ with $f(x)>0$ has positive integral.
\end{theorem}
\begin{proof}
Any compact subset of $N$ has a finite cover by null opens, so has zero mass. Inner regularity on the open $N$ gives $\mu(N)=0$, even if its defining union is uncountable. If $x\in\supp\mu$ and $f(x)>0$, continuity gives an open neighborhood on which $f>c>0$. That neighborhood has positive measure, hence so does the integral. If $x\in N$, use a compact cutoff inside a null neighborhood, equal to one at $x$; it has zero integral. This proves both implications without countability of a base.
\end{proof}
\source{FOL 7.1 Exercise2}

\section{Products without hidden countability}

\begin{proposition}[Borel product versus product Borel sigma-algebras]\family{D11.18.product_borel_bridge}\label{d11:18:product-borel-bridge}
For LCH $X,Y$, $\BB(X)\otimes\BB(Y)\subset\BB(X\times Y)$. If both spaces are second countable, equality holds. In that case the ordinary product of compact-finite Borel measures is Radon.
\end{proposition}
\begin{proof}
The precise ingredients are \cref{d11:18:second-countable-borel-regular}, \cref{d11:08:product-measure}. Coordinate inverse images of Borel sets are Borel, proving the inclusion. A countable product base writes every open set as a countable union of open rectangles, proving equality in the second-countable case. Each compact in the product lies in the product of its compact projections, so the ordinary product is compact-finite; its construction is unambiguous since the factors are sigma-finite by a countable precompact cover. The product is second-countable LCH, so the preceding Borel regularity theorem applies.
\end{proof}
\source{FOL 7.4 Theorem7.20}

\begin{lemma}[Finite separate-variable approximation of a compact-support product function]\family{D11.18.cc_tensor_approximation}\label{d11:18:cc-tensor-approximation}
For $f\in C_c(X\times Y,\mathbb K)$ and $\eps>0$, a finite sum $\sum_jg_j(x)h_j(y)$ with $g_j\in C_c(X,\mathbb K)$ and $h_j\in C_c(Y,\mathbb K)$ approximates $f$ uniformly within $\eps$. All approximants can have support in one fixed precompact product neighborhood. In particular $f$ is product-Borel measurable.
\end{lemma}
\begin{proof}
Let $K,L$ be the compact projections of $\supp f$, and choose precompact open neighborhoods $O,P$ of them. On $\overline O\times\overline P$, continuity and the finite-cover form of the tube lemma imply the following: for each $y\in\overline P$ and each $\eta>0$ there is a neighborhood $V_y$ such that $\sup_{x\in\overline O}|f(x,z)-f(x,y)|<\eta$ for $z\in V_y\cap\overline P$. Choose finitely many such neighborhoods covering $L$, all contained in $P$, and a compact-support partition $h_j$ subordinate to them with sum one on $L$ and at most one everywhere, as constructed in \cref{d11:19:finite-compact-partition}. Choose a cutoff $a=1$ on $K$, supported in $O$, and put $g_j(x)=a(x)f(x,y_j)$. On $\overline O\times P$, the error between $f\sum_jh_j$ and $\sum_jg_jh_j$ is at most $\eta$: where $a\ne1$, $f(x,z)=0$ for all $z$, since $x\notin K$, and the same estimate still holds. Also $f\sum_jh_j=f$, since $f=0$ off $X\times L$. Outside $O\times P$ both expressions vanish. Take $\eta<\eps$. A sequence of such sums converges pointwise to $f$ and consists of product-measurable functions. The finite partition is purely topological and uses no measure representation.
\end{proof}
\source{FOL 7.4 Proposition7.21 and start of7.22}

\begin{construction}[Radon product from compact-support iterated integrals]\family{D11.18.radon_product}\label{d11:18:radon-product}
For arbitrary positive Radon $\mu,\nu$ on LCH $X,Y$, every $f\in C_c(X\times Y)$ is product measurable and supported in a finite-measure rectangle. Its iterated integrals agree. There is a unique Radon measure $\mu\otimes_R\nu$ on $\BB(X\times Y)$ satisfying
\[\int f\dd(\mu\otimes_R\nu)=\int_X\int_Y f(x,y)\dd\nu(y)\dd\mu(x).\]
\end{construction}
\begin{proof}
The precise ingredients are \cref{d11:18:cc-tensor-approximation}. The preceding lemma proves measurability. Compact projection supports $K,L$ have finite masses. Apply finite-measure Fubini, \cref{d11:08:fubini}, to the restrictions to $K,L$, extending the resulting section integrals by zero. Alternatively uniform tensor approximation shows directly that these section integrals are continuous with compact support: all tensor approximants have support in fixed compact neighborhoods $K'\times L'$, and their uniform error changes a section integral by at most that error times $\nu(L')$. Thus the iterated integrals exist and agree. They define a positive linear functional on real $C_c(X\times Y)$. The positive Riesz--Markov theorem \cref{d11:19:positive-riesz-markov}, whose proof uses no product measure, supplies the unique Radon measure. Complex functions follow by real and imaginary parts. This forward reference is an acyclic dependency, not an assumption of a product Fubini theorem on all product Borel sets.
\end{proof}
\source{FOL 7.4 Propositions7.21-7.22 and definition}

\begin{theorem}[Agreement of sigma-finite Radon and ordinary products]\family{D11.18.sigmafinite_radon_product}\label{d11:18:sigmafinite-radon-product}
If the two Radon factors are sigma-finite, then for every $E\in\BB(X\times Y)$ the section masses are Borel measurable and
\[(\mu\otimes_R\nu)(E)=\int_X\nu(E_x)\dd\mu(x)=\int_Y\mu(E^y)\dd\nu(y).\]
On $\BB(X)\otimes\BB(Y)$ this is the ordinary product.
\end{theorem}
\begin{proof}
The precise ingredients are \cref{d11:18:radon-product}, \cref{d11:19:recover-open-compact-mass}. First observe that for open rectangles,
\[(\mu\otimes_R\nu)(U\times V)=\mu(U)\nu(V).
\]
Indeed open-mass recovery identifies the left side with the supremum of the defining integrals of $0\le f\le1$ supported there. Each such integral is at most the product of open masses. Product tests $f(x,y)=a(x)b(y)$, with compact cutoffs supported in $U,V$, approach that product from below by open-mass recovery for each factor. This also handles zero factors under the convention $0\cdot\infty=0$. For sigma-finite Radon factors, a finite-measure Borel exhaustion can be enlarged by outer regularity to a countable finite-measure open cover. Its open rectangles give a countable finite-mass cover of the product. The pi-system of open rectangles generates $\BB(X)\otimes\BB(Y)$, so sigma-finite measure uniqueness shows that the restriction of the Radon product is the ordinary product. Thus product-measurable integrals may already be computed there by ordinary Fubini.

We now give the additional open-set step for the larger product Borel sigma-algebra, since open product sets need not be countable unions of rectangles. Suppose first both factors are finite. For open $W$, $x\mapsto\nu(W_x)$ is lower semicontinuous: a compact $L\subset W_x$ stays inside $W_z$ for all $z$ in a neighborhood of $x$, by a finite rectangle cover; compact-inner regularity gives the assertion. Further,
\[\int_X\nu(W_x)\dd\mu=\sup\left\{\int f\dd(\mu\otimes_R\nu):f\in C_c,\ 0\le f\le1,\ \supp f\subset W\right\}.
\]
The inequality $\ge$ is immediate. To prove $\le$, approximate the bounded Borel function $v(x)=\nu(W_x)$ from below in integral by a finite simple function $\sum_ja_j\one_{E_j}$ with pairwise disjoint Borel level strata $E_j\subset\{v>a_j\}$ and $a_j>0$. Such strict lower steps have arbitrarily small integral error; $v\le\nu(Y)$ and $\mu(X)<\infty$. Finite Radon compact inner regularity gives disjoint compact $K_j\subset E_j$ with arbitrarily small further loss in $\sum_ja_j\mu(K_j)$. At each point $x\in K_j$ choose $h\in C_c(Y,[0,1])$ supported in $W_x$ with integral exceeding $a_j$. A neighborhood $V_x$ of $x$ times $\supp h$ remains in $W$, by the compact tube argument. Shrink these neighborhoods and choose a finite cover of $K_j$ by smaller neighborhoods whose compact closures lie in the corresponding $V_x$. Cutoffs equal to one on those closures and supported in $V_x$, multiplied by the chosen $h$, are $C_c(X\times Y,[0,1])$ tests supported in $W$. Their finite pointwise maximum $f$ has section integral greater than $a_j$ at every point of $K_j$. Since the $K_j$ are disjoint, its iterated integral is at least $\sum_ja_j\mu(K_j)$. Let both approximation errors decrease to zero. This proves the reverse inequality. The supremum equals the Radon product's mass of $W$ by open inner regularity and cutoffs.

The class of Borel sets for which the section functions are Borel and the identity holds contains the opens. With finite factor masses and finite product mass it is closed under complements and disjoint countable unions, by subtraction and monotone convergence. The pi-lambda theorem, applied to the pi-system of opens, gives all product Borel sets. Interchanging factors gives the second identity.

\enlargethispage{2\baselineskip}
For sigma-finite factors take increasing finite-measure open exhaustions $A_n,B_n$, obtained from the finite open covers in the first paragraph. Restricting a Radon measure to any finite-measure Borel set gives a finite Radon measure: compact inner approximation holds on every restricted Borel set by \cref{d11:18:radon-sigmafinite-inner}; applying it to complements in the finite restricted measure gives outer approximation by open complements of compact sets. The same observation applies to the Radon product restricted to $A_n\times B_n$, whose mass is finite by open rectangle equality. For a $C_c$ test, the integral against this restriction is computed on the ordinary product sigma-algebra established in the first paragraph, since the test is product measurable by \cref{d11:18:cc-tensor-approximation}. Finite Fubini makes it equal to the functional for the two finite restricted factors. Riesz--Markov uniqueness therefore identifies the restricted Radon product with their Radon product on every Borel set. Apply the finite-factor result and let $n$ increase, using monotone convergence on both section integrals. This proves all the claimed identities. No identification of the two Borel sigma-algebras was used.
\end{proof}
\source{FOL 7.4 Theorem7.26}

% END chapters/18


% BEGIN chapters/19
\chapter{Positive and signed measure representation}\label{ch:d11-19}

This chapter constructs measures from functionals rather than treating representation as a foundation. All spaces are LCH, all scalar functionals are scalar-linear, and the complex integral pairing has no hidden conjugation. The construction below precedes and supports the Radon product and Haar arguments.

\section{Continuous functions and finite partitions}

\begin{definition}[Compact support and vanishing at infinity]\family{D11.19.cc_czero}\label{d11:19:cc-czero}
The space $C_c(X,\mathbb K)$ consists of continuous functions whose support $\overline{\{f\ne0\}}$ is compact. The space $C_0(X,\mathbb K)$ consists of continuous functions for which $\{|f|\ge\eps\}$ is compact for every $\eps>0$. Such functions are bounded; the uniform norm makes $C_0$ a normed space. Moreover $C_c\subset C_0$ is uniformly dense.
\end{definition}
\begin{proof}
The precise ingredients are \cref{d11:18:compact-cutoff}. For $f\in C_0$, the compact set $\{|f|\ge1\}$ has bounded continuous image; off it $|f|<1$, proving boundedness. Given $\eps>0$, choose a cutoff $\phi=1$ on $\{|f|\ge\eps\}$. Then $\phi f\in C_c$ and $\|f-\phi f\|_\infty\le\eps$. Compact support immediately implies the $C_0$ condition. Uniform limits preserve continuity and the vanishing condition: a level set of the limit lies as a closed subset inside a compact half-level set of a sufficiently close approximant. Thus $C_0$ is also complete, by scalar completeness.
\end{proof}
\source{FOL 7.1 Cc and 7.3 recalled C0 interface; COH 7.1 support and compact-support spaces}

\begin{lemma}[Finite compact-support decomposition subordinate to an open cover]\family{D11.19.finite_compact_partition}\label{d11:19:finite-compact-partition}
If compact $K$ is covered by finitely many opens $U_1,\ldots,U_m$, there are $\phi_i\in C_c(X,[0,1])$, supported in $U_i$, whose sum is one on $K$ and at most one everywhere. Consequently any $f\in C_c$ supported in their union splits as $f=\sum_i f_i$ with $\supp f_i\subset U_i$; if $f\ge0$ the pieces can be nonnegative.
\end{lemma}
\begin{proof}
A finite shrinking of the cover of $K$ gives compact $K_i\subset U_i$ covering $K$; it follows by first taking finitely many precompact neighborhoods with closures in the prescribed opens and then grouping the closures. Choose $h_i\in C_c(X,[0,1])$ equal to one on $K_i$, supported in $U_i$, using \cref{d11:18:compact-cutoff}. Put $\phi_i=h_i\prod_{j<i}(1-h_j)$. Then $\sum_i\phi_i=1-\prod_i(1-h_i)$, which lies in $[0,1]$ and equals one on $K$. Each support is contained in $\supp h_i$. For $K=\supp f$, set $f_i=f\phi_i$. Empty $K$ gives zero pieces.
\end{proof}
\source{COH 7.1 Propositions7.1.9-7.1.12}

\begin{lemma}[Positivity gives compact-support norm bounds]\family{D11.19.positive_functional_bound}\label{d11:19:positive-functional-bound}
If $I:C_c(X,\R)\to\R$ is positive linear, then for every compact $K$ there is finite $C_K$ with $|I(f)|\le C_K\|f\|_\infty$ whenever $\supp f\subset K$. A complex-linear functional taking nonnegative real functions to nonnegative real values has the analogous bound.
\end{lemma}
\begin{proof}
The precise ingredients are \cref{d11:18:compact-cutoff}. Choose $h\in C_c$ with $h\ge1$ on $K$, $h\ge0$. The pointwise inequalities $-\|f\|_\infty h\le f\le\|f\|_\infty h$ give the real estimate with $C_K=I(h)$. In the complex case, real functions have real images by their positive-negative decomposition. For any $f$ choose a scalar $z$ of modulus one with $|I(f)|=\Re(zI(f))=I(\Re(zf))\le\|f\|_\infty I(h)$. This also proves the sharp local bound without a factor two.
\end{proof}
\source{FOL 7.1 Proposition7.1}

\section{The positive construction}

\begin{construction}[Open content of a positive Cc functional defines a Radon measure]\family{D11.19.positive_content_measure}\label{d11:19:positive-content-measure}
For positive real $I$ put
\[c(U)=\sup\{I(f):f\in C_c,\ 0\le f\le1,\ \supp f\subset U\},\qquad
\mu^*(A)=\inf_{U\supset A,\ U\text{ open}}c(U).
\]
Then $\mu^*$ is an outer measure, all Borel sets are Caratheodory measurable, and its Borel restriction is Radon.
\end{construction}
\begin{proof}
The precise ingredients are \cref{d11:19:finite-compact-partition}, \cref{d11:19:positive-functional-bound}. Monotonicity and the empty-set value are immediate. If $U\subset\bigcup_n U_n$ are open, every admissible $f$ for $U$ has compact support covered by finitely many $U_n$. The finite partition splits it into $f_n\ge0$ with sum $f$, each at most one and supported in the corresponding $U_n$. Thus $c(U)\le\sum_n c(U_n)$. Choosing open supersets of arbitrary $A_n$ with summed errors below $\eps$ proves countable subadditivity of $\mu^*$; cases with an infinite right side need no approximation. Notice $\mu^*(U)=c(U)$ for open $U$, by monotonicity.

We prove an open $G$ measurable. It suffices to test $A$ with $\mu^*(A)<\infty$. Choose open $V\supset A$ with $c(V)<\mu^*(A)+\eps$. For $0\le f\le1$ supported in $V\cap G$, put $K=\supp f$. Every admissible $g$ for $V\setminus K$ has disjoint support from $f$, so $f+g\le1$ and
$c(V)\ge I(f)+c(V\setminus K)\ge I(f)+\mu^*(A\setminus G)$.
Take the supremum over $f$ to obtain
$\mu^*(A)+\eps\ge c(V\cap G)+\mu^*(A\setminus G)\ge\mu^*(A\cap G)+\mu^*(A\setminus G)$.
Let $\eps\downarrow0$. The reverse inequality is subadditivity; if $\mu^*(A)=\infty$ it is automatic in the extended order. Caratheodory's theorem \cref{d11:02:caratheodory-theorem} gives a measure containing the Borel sets.

For compact $K$, a cutoff $h\ge1$ on $K$ and the open set $\{h>1/2\}$ give $\mu(K)\le2I(h)<\infty$. More generally, if $h\ge0$ and $h\ge1$ on $A$, then $\mu^*(A)\le I(h)$: use $\{h>1-\eta\}$ and $(1-\eta)f\le h$, then let $\eta\downarrow0$. If $0\le f\le1$ is supported in compact $K$, then $I(f)\le c(U)$ for every open $U\supset K$, hence $I(f)\le\mu(K)$. It follows that $c(U)\le\sup_{K\subset U}\mu(K)$; the reverse inequality is monotonicity. This proves open inner regularity. Outer regularity is built into the definition and $\mu(U)=c(U)$.
\end{proof}
\source{COH 7.2 Proposition7.2.9, Lemma7.2.10, Proposition7.2.11}

\begin{lemma}[Identify the constructed measure integral with the functional]\family{D11.19.positive_content_integral}\label{d11:19:positive-content-integral}
For the preceding constructed measure, $I(f)=\int f\dd\mu$ for all real $f\in C_c$. If $f\ge0$, $K=\supp f$, and $\eta>0$, the two quantities lie between the same level sums with difference at most $\eta\mu(K)$.
\end{lemma}
\begin{proof}
The precise ingredients are \cref{d11:19:positive-content-measure}. Choose an integer $N$ with $N\eta\ge\|f\|_\infty$, and set $f_j=\min(\eta,\max(f-(j-1)\eta,0))$, $1\le j\le N$. These continuous compact-support functions sum to $f$. Put $K_0=K$ and $K_j=\{f\ge j\eta\}$ for $j\ge1$. Each $K_j$ is compact. The bounds
\[\eta\one_{K_j}\le f_j\le\eta\one_{K_{j-1}}\]
give the same bounds on their integrals. They also give them on $I(f_j)$: the lower bound follows from the $h\ge1$ outer-measure estimate in the preceding proof; the upper bound follows from $f_j/\eta\le1$ supported in $K_{j-1}$. Thus both $I(f)$ and $\int f$ lie between $\eta\sum_{j=1}^N\mu(K_j)$ and $\eta\sum_{j=1}^N\mu(K_{j-1})$. Their difference is at most $\eta\mu(K)<\infty$. Let $\eta\downarrow0$, then split a real $f$ into $f^+-f^-$.
\end{proof}
\source{COH 7.2 Lemma7.2.10 and Proposition7.2.11}

\begin{theorem}[Riesz-Markov representation for positive Cc functionals]\family{D11.19.positive_riesz_markov}\label{d11:19:positive-riesz-markov}
Every positive linear $I:C_c(X,\R)\to\R$ has a unique positive Radon representing measure $\mu$ on $\BB(X)$ with $I(f)=\int f\dd\mu$. Every positive Radon measure induces such a functional. The complex-linear positive version has the same positive measure.
\end{theorem}
\begin{proof}
The precise ingredients are \cref{d11:19:positive-content-measure}, \cref{d11:19:positive-content-integral}. Existence and the integral identity were constructed in the preceding two results. Conversely a continuous compact-support function is bounded on a finite-measure compact set, so its integral is finite; linearity and positivity follow from integration. For any Radon representing measure, cutoffs and open inner regularity give
$\mu(U)=\sup\{\int f\dd\mu:0\le f\le1,\supp f\subset U\}$.
Indeed every compact $K\subset U$ has a cutoff one on it, and every such test is bounded by $\one_U$. Thus $I$ determines all open masses. Outer regularity determines all Borel masses, proving uniqueness. In the complex version restrict to real functions, represent there, and use $I(u+iv)=I(u)+iI(v)$. No sequential-continuity hypothesis is needed.
\end{proof}
\source{COH 7.2 Theorem7.2.8 with its complete proof; FOL 7.1 Theorem7.2}

\begin{corollary}[Recover open and compact masses from a positive functional]\family{D11.19.recover_open_compact_mass}\label{d11:19:recover-open-compact-mass}
For this representation,
\[\mu(U)=\sup_{0\le f\le1,\ \supp f\subset U}I(f),\qquad
\mu(K)=\inf_{g\in C_c,\ g\ge0,\ g|_K\ge1}I(g).
\]
The first formula is for open $U$, the second for compact $K$.
\end{corollary}
\begin{proof}
The precise ingredients are \cref{d11:19:positive-riesz-markov}, \cref{d11:18:compact-cutoff}. The first was proved in uniqueness. Every $g$ in the second formula has integral at least $\mu(K)$. Conversely choose open $U\supset K$ with $\mu(U)<\mu(K)+\eps$, and a cutoff equal to one on $K$ supported in $U$; its integral is at most $\mu(U)$. Compact finiteness ensures the approximation is meaningful. Let $\eps\downarrow0$.
\end{proof}
\source{FOL 7.1 formulas7.3-7.4 and uniqueness proof; COH 7.2 Lemma7.2.7}

\section{Finite measures and the full scalar dual}

\begin{theorem}[Finite mass characterizes bounded C0 extension]\family{D11.19.finite_positive_czero}\label{d11:19:finite-positive-czero}
The positive integration functional on $C_c$ extends boundedly and positively to $C_0$ exactly when $\mu(X)<\infty$. The extension is unique, is ordinary integration on $C_0$, and has norm $\mu(X)$.
\end{theorem}
\begin{proof}
The precise ingredients are \cref{d11:19:cc-czero}, \cref{d11:19:recover-open-compact-mass}. For finite $\mu$, $|\int f|\le\mu(X)\|f\|_\infty$ proves boundedness; uniform approximation by $C_c$ proves uniqueness and positivity of the extension. The open-mass formula for $X$ gives norm at least $\mu(X)$ by tests of norm at most one. Conversely any bounded extension bounds all those tests by its norm, so the same formula forces $\mu(X)<\infty$. The zero measure gives the zero functional.
\end{proof}
\source{FOL 7.3 opening discussion}

\begin{lemma}[Positive decomposition of real C0 functionals]\family{D11.19.functional_jordan}\label{d11:19:functional-jordan}
For a bounded real-linear $I:C_0(X,\R)\to\R$, define, for $f\ge0$,
\[I^+(f)=\sup_{0\le g\le f,\ g\in C_0}I(g).
\]
It extends to a bounded positive linear functional, and $I^-=I^+-I$ is bounded positive. Hence $I=I^+-I^-$.
\end{lemma}
\begin{proof}
The supremum is finite between zero and $\|I\|\|f\|_\infty$. It is positively homogeneous. For $0\le g\le f_1+f_2$, put $g_1=\min(g,f_1)$ and $g_2=g-g_1$; they lie in $C_0$, with $0\le g_i\le f_i$. This proves subadditivity of the displayed supremum. Conversely sums of independently near-maximizing $g_i$ prove the reverse inequality, so it is additive on the positive cone. Define $I^+(u-v)=I^+(u)-I^+(v)$ for nonnegative $u,v$. Additivity makes this independent of the decomposition, and gives linearity. For arbitrary $f$, positivity gives $|I^+(f)|\le I^+(|f|)\le\|I\|\|f\|_\infty$, proving boundedness. Finally $I^+(f)\ge I(f)$ for $f\ge0$, so $I^-$ is positive and bounded.
\end{proof}
\source{FOL 7.3 Lemma7.15}

\begin{definition}[Finite signed/complex Radon measures and variation]\family{D11.19.complex_radon_variation}\label{d11:19:complex-radon-variation}
A finite signed Borel measure is signed Radon if its Jordan parts are positive Radon. A finite complex Borel measure is complex Radon if its real and imaginary parts are signed Radon. Equivalently its total variation is Radon. The norm is $\|\nu\|=|\nu|(X)$, zero only for the zero measure.
\end{definition}
\begin{proof}
The precise ingredients are \cref{d11:18:radon-sigmafinite-inner}. We justify the equivalence, which is not merely terminology. A positive finite measure dominated by a finite Radon measure is Radon: use all-Borel compact inner and outer approximations for the dominating measure, and bound the errors by domination. If the four Jordan parts are Radon, their finite sum is Radon and dominates $|\nu|$ (use the finite-partition definition of variation). Hence $|\nu|$ is Radon. Conversely $|\Re\nu|,|\Im\nu|$ and their Jordan parts are dominated by $|\nu|$ and are Radon. Variation zero implies $|\nu(E)|=0$ for every Borel $E$. This uses the already proved signed/complex variation calculus, not a density theorem.
\end{proof}
\source{FOL 7.3 definitions and Proposition7.16}

\begin{lemma}[Continuous compact-support functions detect total variation]\family{D11.19.variation_test_function}\label{d11:19:variation-test-function}
For finite signed or complex Radon $\nu$ and $\eps>0$, there is $f\in C_c(X,\mathbb K)$ with $\|f\|_\infty\le1$ and $|\int f\dd\nu|\ge|\nu|(X)-\eps$.
\end{lemma}
\begin{proof}
The precise ingredients are \cref{d11:19:complex-radon-variation}, \cref{d11:18:radon-sigmafinite-inner}, \cref{d11:18:compact-cutoff}. Choose a finite Borel partition $A_1,\ldots,A_m$ with $\sum_j|\nu(A_j)|>|\nu|(X)-\eps/3$, by the definition of total variation. By finite-Radon compact inner regularity choose disjoint compact $K_j\subset A_j$ with $\sum_j|\nu|(A_j\setminus K_j)<\eps/6$. Disjoint finitely many compact sets have disjoint open neighborhoods in a Hausdorff space (separate each pair by neighborhoods and intersect the finite choices). Choose compact cutoffs $\phi_j=1$ on $K_j$ supported in these disjoint neighborhoods. Let $z_j$ be the scalar of modulus one with $z_j\nu(K_j)=|\nu(K_j)|$; choose any such scalar when the mass is zero. Then $f=\sum_j z_j\phi_j$ has norm at most one, and on $K=\bigcup K_j$ its integral is $\sum_j|\nu(K_j)|$. The remaining integral has absolute value at most $|\nu|(X\setminus K)<\eps/6$. Comparing the compact and original partition sums loses another $\eps/6$. The resulting lower bound exceeds $|\nu|(X)-\eps$. Empty or zero pieces cause no problem. No polar density or RN theorem was used.
\end{proof}
\source{COH 7.3 proof of Theorem7.3.6}

\begin{theorem}[Finite Radon measures represent the bounded dual of C0]\family{D11.19.czero_duality}\label{d11:19:czero-duality}
The map $\nu\mapsto T_\nu$, $T_\nu(f)=\int f\dd\nu$, is a scalar-linear isometric bijection from finite signed Radon measures onto $C_0(X,\R)^*$, and from finite complex Radon measures onto $C_0(X,\C)^*$. Its norm is $|\nu|(X)$.
\end{theorem}
\begin{proof}
The precise ingredients are \cref{d11:19:functional-jordan}, \cref{d11:19:positive-riesz-markov}, \cref{d11:19:finite-positive-czero}, \cref{d11:19:variation-test-function}. The variation inequality gives a bounded functional with norm at most the variation mass; the preceding test-function lemma gives equality. For real $I$, split $I=I^+-I^-$ by the functional Jordan lemma, restrict each positive part to $C_c$, and represent by positive Radon measures. Their masses are finite by boundedness and the finite-positive theorem. Their difference represents $I$ first on $C_c$, then on $C_0$ by density. In the complex case restrict $I$ to real functions and apply the real result separately to $\Re I$ and $\Im I$. Their representing signed measures $\alpha,\beta$ give $\nu=\alpha+i\beta$; complex linearity supplies the identity for all complex functions. Variation regularity was proved above. If two representations agree, their difference integrates every $C_c$ test to zero, and the test-function lemma makes its variation zero. This proves uniqueness and injectivity, and linearity follows directly from integration.
\end{proof}
\source{COH 7.3 Theorem7.3.6 and Lemma7.3.7; FOL 7.3 Theorem7.17}

\section{Compact specialization and a literal source boundary}

\begin{corollary}[Compact-space C(K) specialization]\family{D11.19.compact_representation}\label{d11:19:compact-representation}
On compact Hausdorff $X$, $C_c=C_0=C(X)$. Positive real functionals have unique finite positive regular Borel measures; bounded scalar functionals have unique finite signed/complex Radon measures, with norm the variation mass.
\end{corollary}
\begin{proof}
The precise ingredients are \cref{d11:19:positive-riesz-markov}, \cref{d11:19:czero-duality}, \cref{d11:18:radon-sigmafinite-inner}. All continuous functions on compact $X$ are bounded and compactly supported, and their levels are compact. A positive functional satisfies $|I(f)|\le I(1)\|f\|_\infty$, hence is bounded. Apply positive representation and the full $C_0$ duality. Since a finite Radon measure is compact-inner regular on all Borel sets, the asserted regularity follows.
\end{proof}
\source{FOL 7.3 Corollary7.18; STE Chapter6 Exercise9 compact interval}

\begin{corollary}[Compact-interval positive-functional variant]\family{D11.19.stein_exercise9_interface}\label{d11:19:stein-exercise9-interface}
For $a\le b$, a positive $I:C([a,b],\R)\to\R$ is integration against a unique finite positive Borel measure on $[a,b]$, and conversely. Its norm is $I(1)$; no extra sequential-continuity assumption is needed.
\end{corollary}
\begin{proof}
The precise ingredients are \cref{d11:19:compact-representation}, \cref{d11:18:second-countable-borel-regular}. Apply compact representation to the compact interval. Conversely every finite Borel measure there is regular by the metric/second-countable regularity theorem and defines the positive functional. The order bound by $I(1)$ proves boundedness directly. The case $a=b$ is a one-point space and obeys the same formula.
\end{proof}
\source{STE Chapter 6 Exercise 9, physical PDF p.335, is this valid compact-interval interface; the proof above also makes its boundedness explicit.}

\begin{example}[Bounded-support C0 on arbitrary metrics is not the canonical C0 theorem]\family{D11.19.stein_problem5_caution}\label{d11:19:stein-problem5-caution}
In the discrete metric $d(m,n)=1$ on $X=\N$, let $\mu\{n\}=2^{-n}$. All balls have finite measure and $X$ is itself a closed ball. The continuous function $f(n)=2^n$ is supported in that closed ball but has infinite integral.
\end{example}
\begin{proof}
The measure is a finite Borel measure of total mass one. Every function is continuous in the discrete topology, every subset is bounded, and a closed ball of radius one is $X$. But $\int f\dd\mu=\sum_{n\ge1}2^n2^{-n}=\infty$. Thus integration is not a finite-valued linear functional on the literal continuous-closed-ball-support space in STE Chapter 6 Problem 5 on arbitrary metric spaces. Any sequential-continuity premise on a proposed finite-valued functional cannot repair this failed converse for all such measures. Compact support, a bounded-function restriction, or other appropriate hypotheses are required; this problem is source context, not a supporting generic representation theorem.
\end{proof}
\source{STE Chapter 6 Problem 5, physical PDF p.342. This counterexample concerns its literal general-metric support formulation; it does not affect Exercise 9 or the LCH compact-support theorem proved here.}

% END chapters/19


% BEGIN chapters/20
\chapter{Duality and a minimal Hilbert-space bridge}\label{ch:d11-20}

The scalar field is $\mathbb K=\R$ or $\C$. The $L^p$ spaces use ordinary almost-everywhere equality unless a local-null quotient is explicitly displayed. The bilinear pairing is $\int fg$; the Hilbert inner product is linear in its first argument. The Hilbert proof uses only completeness and inner-product algebra. The general $L^p$ dual proof uses the independently established RN theorem and creates its own sigma-finite support.

\section{Integral pairings and their norms}

\begin{definition}[Scalar Lp-Lq integral pairing]\family{D11.20.integral_pairing}\label{d11:20:integral-pairing}
For conjugate $1\le p,q\le\infty$, put $T_g(f)=\int fg\dd\mu$ for $f\in L^p$, $g\in L^q$. This is well defined on a.e. classes, scalar-linear in both $f$ and $g$, and $\|T_g\|\le\|g\|_q$.
\end{definition}
\begin{proof}
Holder's inequality \cref{d11:10:holder} makes the product integrable and bounds its integral. Changing either representative on a null set changes the product only on a null set. Integral linearity proves the remaining assertions. Zero functions give zero maps.
\end{proof}
\source{FOL 6.2 opening; AXL 9B Theorem9.42 pairing convention}

\begin{theorem}[Norm equality for finite dual exponent]\family{D11.20.pairing_norm_finite_q}\label{d11:20:pairing-norm-finite-q}
On any positive measure space, if $1\le q<\infty$ and $p$ is conjugate, $\|T_g\|=\|g\|_q$. In particular $L^q$ injects isometrically into $(L^p)^*$, including $L^1\hookrightarrow(L^\infty)^*$.
\end{theorem}
\begin{proof}
The precise ingredients are \cref{d11:20:integral-pairing}. The upper bound is Holder. If $q=1$, use $f=\overline g/|g|$ on $\{g\ne0\}$ and zero elsewhere; this is a measurable function of modulus at most one and gives $\int fg=\|g\|_1$. If $1<q<\infty$ and $g\ne0$, set $f=(\overline g/|g|)|g|^{q-1}/\|g\|_q^{q-1}$. Since $(q-1)p=q$, its $L^p$ norm is one and its pairing is $\|g\|_q$. For real scalars use the real sign. If $g=0$ equality is immediate. A zero functional therefore forces the $L^q$ class to be zero.
\end{proof}
\source{FOL 6.2 Proposition6.13}

\begin{theorem}[Endpoint L-infinity norm equality under semifiniteness]\family{D11.20.pairing_norm_endpoint}\label{d11:20:pairing-norm-endpoint}
If $\mu$ is semifinite, in the sense that $\mu(E)=\sup\{\mu(A):A\subset E,\ A\text{ measurable},\ \mu(A)<\infty\}$, then every ordinary $g\in L^\infty$ satisfies $\|T_g\|_{(L^1)^*}=\|g\|_\infty$. Sigma-finite measures are semifinite.
\end{theorem}
\begin{proof}
The precise ingredients are \cref{d11:20:integral-pairing}. For $0<c<\|g\|_\infty$, the set $\{|g|>c\}$ has positive measure. Semifiniteness supplies a subset $A$ of finite positive measure. The test $f=(\overline g/|g|)\one_A/\mu(A)$ has norm one and gives pairing at least $c$. Let $c$ increase to the essential supremum; zero norm is harmless. For sigma-finite $\mu$, intersections with a countable finite-measure exhaustion prove the displayed supremum property by continuity from below. Without semifiniteness a nonnull set can be invisible to every $L^1$ function, which explains the hypothesis.
\end{proof}
\source{FOL 6.2 Proposition6.13}

\section{A route independent of RN}

\begin{proposition}[L2 is a Hilbert space without a separability assumption]\family{D11.20.l2_hilbert_interface}\label{d11:20:l2-hilbert-interface}
For arbitrary positive $\mu$, $\ip f g=\int f\overline g\dd\mu$ makes ordinary $L^2(\mu)$ a Hilbert space, linear in the first slot. No separability or sigma-finiteness is assumed.
\end{proposition}
\begin{proof}
Holder at exponent two gives existence and Cauchy--Schwarz; integral algebra gives sesquilinearity and conjugate symmetry. The identity $\ip f f=\|f\|_2^2$ and the a.e. quotient give positive definiteness. Completeness is \cref{d11:10:completeness}. These are exactly the inner-product and completeness inputs needed below.
\end{proof}
\source{STE Chapter4 §1 and Chapter6 §2 abstract L2; AXL 8B L2 self-duality discussion}

\begin{theorem}[Projection onto a closed Hilbert subspace]\family{D11.20.orthogonal_projection}\label{d11:20:orthogonal-projection}
For a closed linear subspace $V$ of a real or complex Hilbert space $H$, every $x$ has a unique minimizing $v=P_Vx\in V$. The residual lies in $V^\perp$. The map $P_V$ is scalar-linear, idempotent, contractive, and gives the unique decomposition $H=V\oplus V^\perp$.
\end{theorem}
\begin{proof}
Let $d=\inf_{v\in V}\|x-v\|$ and choose $v_n$ with $\|x-v_n\|\to d$. The parallelogram identity and $(v_n+v_m)/2\in V$ give
\[\|v_n-v_m\|^2\le2\|x-v_n\|^2+2\|x-v_m\|^2-4d^2\longrightarrow0.
\]
Completeness and closedness give $v_n\to v\in V$ with $\|x-v\|=d$. For $u\in V$ and real $t$, minimality of $v+tu$ implies $\|x-v-tu\|^2\ge\|x-v\|^2$, so the linear coefficient $\Re\ip{x-v}u$ vanishes. Replace $u$ by $iu$ in the complex case to obtain its imaginary part. Thus the residual is orthogonal to $V$. Conversely an orthogonal residual gives $\|x-w\|^2=\|x-v\|^2+\|v-w\|^2$ for all $w\in V$, proving uniqueness. Sums and scalar multiples of the decompositions remain decompositions, so uniqueness proves linearity. For $x\in V$ its projection is itself, giving idempotence. Pythagoras gives $\|P_Vx\|\le\|x\|$. This proof has used no measure representation.
\end{proof}
\source{AXL 8B Theorem8.28 and orthogonal decomposition8.43; STE Chapter4 §4}

\begin{theorem}[Bounded Hilbert functionals have unique representing vectors]\family{D11.20.hilbert_riesz}\label{d11:20:hilbert-riesz}
For a bounded scalar-linear $\ell:H\to\mathbb K$ there is a unique $h\in H$ with $\ell(x)=\ip x h$ for every $x$, and $\|h\|=\|\ell\|$.
\end{theorem}
\begin{proof}
The precise ingredients are \cref{d11:20:orthogonal-projection}. For $\ell=0$ take $h=0$. Otherwise $V=\ker\ell$ is closed and proper. Choose $x\notin V$ and normalize its nonzero perpendicular component to a unit $u\in V^\perp$. Then $\ell(u)\ne0$, since $u\in V$ would force $u=0$. Every $x$ has $x-\ell(x)u/\ell(u)\in V$, so taking its inner product with $u$ gives $\ip x u=\ell(x)/\ell(u)$. Put $h=\overline{\ell(u)}u$; because the second slot is conjugate-linear, $\ip x h=\ell(u)\ip x u=\ell(x)$. Cauchy--Schwarz gives $\|\ell\|\le\|h\|$, and testing $h/\|h\|$ gives the reverse inequality. If two vectors represent $\ell$, test their difference against itself to prove equality.
\end{proof}
\source{AXL 8B Theorem8.47; STE Chapter4 §5.1 Theorem5.3}

\begin{corollary}[Relate Hilbert and bilinear integral representatives]\family{D11.20.hilbert_bilinear_bridge}\label{d11:20:hilbert-bilinear-bridge}
For complex $L^2$, the Hilbert vector $h$ represents the bilinear density $g=\overline h$: $\ell(f)=\int f\overline h=\int fg$. Both norms equal $\|\ell\|$. The Hilbert-vector assignment is conjugate-linear in $\ell$; the density assignment is scalar-linear.
\end{corollary}
\begin{proof}
The precise ingredients are \cref{d11:20:l2-hilbert-interface}, \cref{d11:20:hilbert-riesz}. Complex conjugation preserves the $L^2$ norm and is well defined on a.e. classes. The identity follows from the inner-product definition. Multiplication of a functional by $a$ requires its Hilbert vector to be multiplied by $\overline a$, while its bilinear density is multiplied by $a$. Additivity and uniqueness prove both assertions.
\end{proof}
\source{AXL 8B L2 discussion following8.47; FOL 6.2 opening convention note}

\section{RN duality and the support it actually requires}

\begin{lemma}[The local sigma-finite scalar duality construction]\family{D11.20.lp_dual_sigmafinite_local}\label{d11:20:lp-dual-sigmafinite-local}
For sigma-finite $\mu$, $1\le p<\infty$, and bounded scalar-linear $\ell\in(L^p)^*$, there is a unique $g\in L^q$, $q$ conjugate, with $\ell(f)=\int fg$ and $\|g\|_q=\|\ell\|$.
\end{lemma}
\begin{proof}
The precise ingredients are \cref{d11:20:pairing-norm-finite-q}, \cref{d11:20:pairing-norm-endpoint}. First assume $\mu(X)<\infty$. Define $\nu(A)=\ell(\one_A)$. For disjoint $A_n$, the indicators of the finite unions converge in $L^p$ to that of their union, so $\nu$ is countably additive. For any finite measurable partition choose phases $z_j$ with $z_j\nu(A_j)=|\nu(A_j)|$. Then
$\sum_j|\nu(A_j)|=\ell(\sum_jz_j\one_{A_j})\le\|\ell\|\mu(X)^{1/p}$ in absolute value. Thus $\nu$ is a finite signed/complex measure and vanishes on $\mu$-null sets. The finite-reference RN theorem \cref{d11:15:rn-signed-complex} gives an integrable density $g$, and the desired identity holds for bounded simple functions.

If $p>1$, the bounded test $f_n=(\overline g/|g|)|g|^{q-1}\one_{\{1/n\le|g|\le n\}}$ is in $L^p$. Its simple approximations give $S_n=\int_{\{1/n\le|g|\le n\}}|g|^q\le\|\ell\|S_n^{1/p}$. Hence $S_n^{1/q}\le\|\ell\|$, with the zero case read directly. Monotone convergence gives $g\in L^q$ with norm at most $\|\ell\|$. If $p=1$, a positive-measure set with $|g|>\|\ell\|+\eta$ and its bounded phase-indicator test contradict the same functional bound, so $\|g\|_\infty\le\|\ell\|$. Simple density \cref{d11:11:finite-p-simple-density} and Holder extend the identity to every $f$. The pairing-norm results give equality and uniqueness on the finite space.

For general sigma-finite $\mu$, take increasing $E_n$ of finite measure covering $X$. Restrict $\ell$ by extension by zero to each $E_n$. The densities agree a.e. on overlaps by the finite-case uniqueness. Change them on the countable union of exceptional null sets and glue to measurable $g$. For finite $q$, monotone convergence gives $\|g\|_q\le\|\ell\|$; for $q=\infty$, the countable exceptional union gives the same bound. Since $f\one_{E_n}\to f$ in $L^p$, boundedness and Holder give the full identity. The norm and uniqueness follow from the pairing-norm results (sigma-finiteness is used at the endpoint). This argument does not reduce an arbitrary reference measure to a sigma-finite one merely because $\nu$ is finite.
\end{proof}
\source{FOL 6.2 proof of Theorem6.15; AXL 9B finite measure part of9.42; COH 4.5 proof of4.5.1}

\begin{lemma}[An arbitrary-measure finite-p functional has a visible sigma-finite support]\family{D11.20.functional_visible_sigmafinite_support}\label{d11:20:functional-visible-sigmafinite-support}
For arbitrary $\mu$, $1<p<\infty$, and bounded $\ell\in(L^p)^*$, there is a measurable sigma-finite set $F$ with $\ell(f)=\ell(f\one_F)$ for every $f\in L^p$.
\end{lemma}
\begin{proof}
The precise ingredients are \cref{d11:20:lp-dual-sigmafinite-local}. For sigma-finite measurable $E$ let $a(E)$ be the norm of the restriction of $\ell$ to functions supported in $E$, and let $M=\sup_E a(E)\le\|\ell\|$. If $M=0$, every $L^p$ function has sigma-finite nonzero support by \cref{d11:10:sigmafinite-support}, so $\ell=0$ and $F=\varnothing$ works. Otherwise choose $E_n$ with $a(E_n)\to M$ and let $F=\bigcup_nE_n$. Then $a(F)=M$. The preceding local lemma represents the restriction on $F$ by $g_F\in L^q(F)$ of norm $M$.

If sigma-finite $A\supset F$, its density $g_A$ agrees with $g_F$ a.e. on $F$ and has norm at most $M$. Since $q<\infty$,
$\int_A|g_A|^q=\int_F|g_F|^q+\int_{A\setminus F}|g_A|^q\le M^q$,
so the last integral is zero. For any $f\in L^p$, take $A=F\cup\{f\ne0\}$, which is sigma-finite by that same support lemma. The local representation on $A$ now gives $\ell(f)=\ell(f\one_F)$. This support is manufactured from $\ell$, not inferred from the finiteness of a target measure.
\end{proof}
\source{FOL 6.2 arbitrary-measure part of6.15; AXL 9B arbitrary-measure part of9.42; COH 4.5 arbitrary-measure part of4.5.1}

\begin{theorem}[Lp dual representation for arbitrary positive measures and 1<p<infinity]\family{D11.20.lp_dual_arbitrary_measure}\label{d11:20:lp-dual-arbitrary-measure}
For every positive $\mu$, $1<p<\infty$, and $q=p/(p-1)$, the map $g\mapsto T_g$ is a scalar-linear isometric bijection $L^q(\mu)\to(L^p(\mu))^*$.
\end{theorem}
\begin{proof}
The precise ingredients are \cref{d11:20:functional-visible-sigmafinite-support}, \cref{d11:20:pairing-norm-finite-q}. Take the functional-visible sigma-finite $F$ and its density $g_F$ from the preceding proof. Extend $g_F$ by zero outside $F$. For every $f$, support reduction gives $\ell(f)=\ell(f\one_F)=\int fg$. The finite-$q$ pairing-norm theorem gives norm equality and injectivity on the original arbitrary measure space. It also gives a.e. uniqueness. Thus no sigma-finiteness, semifiniteness, or separability is missing from the stated theorem.
\end{proof}
\source{AXL 9B Theorem9.42 and arbitrary-measure extension; FOL 6.2 Theorem6.15; COH 4.5 Theorem4.5.1}

\begin{corollary}[L1 dual representation under sigma-finiteness]\family{D11.20.lone_dual_sigmafinite}\label{d11:20:lone-dual-sigmafinite}
For sigma-finite positive $\mu$, ordinary $L^\infty(\mu)$ is scalar-linearly and isometrically $(L^1(\mu))^*$ under $\int fg$.
\end{corollary}
\begin{proof}
The precise ingredients are \cref{d11:20:lp-dual-sigmafinite-local}. This is the $p=1$ case of the local RN lemma. Sigma-finiteness supplies both gluing and the endpoint norm/uniqueness argument. The proof does not assert this ordinary-quotient statement for arbitrary measures.
\end{proof}
\source{FOL 6.2 Theorem6.15; AXL 9B Theorem9.42 endpoint statement; COH 4.5 Theorem4.5.1}

\section{The genuinely different local-null endpoint}

\begin{definition}[Cohn's separate locally-null L-infinity variant]\family{D11.20.local_null_variant}\label{d11:20:local-null-variant}
A measurable $N$ is \emph{locally null} if $\mu(N\cap E)=0$ for every measurable finite-measure $E$. Let $L^\infty_{\rm loc}$ be bounded measurable functions modulo locally-null equality, with
\[\|[g]\|_{\rm loc}=\inf\{C\ge0:\{|g|>C\}\text{ is locally null}\}.
\]
This quotient need not be ordinary a.e. $L^\infty$.
\end{definition}
\begin{proof}
Locally-null sets form a sigma-ideal, because their intersections with each finite-measure set are a countable union of null sets. The triangle inequality for level sets proves a seminorm before quotienting. If its value is zero, $\{g\ne0\}=\bigcup_n\{|g|>1/n\}$ is locally null, proving definiteness on the quotient. A local-essentially bounded measurable representative can always be made pointwise bounded by changing it on its exceptional locally-null level set, which explains the bounded-representative formulation.
\end{proof}
\source{COH 3.3 locally-null definition and quotient; FOL 6.2 Exercise23}

\begin{proposition}[Identify local-null and ordinary-null quotients in sigma-finite settings]\family{D11.20.local_null_sigmafinite_bridge}\label{d11:20:local-null-sigmafinite-bridge}
Null sets are locally null. If $\mu$ is sigma-finite the converse holds, and the local-null quotient identifies isometrically with ordinary $L^\infty$ after bounded representatives are chosen.
\end{proposition}
\begin{proof}
The precise ingredients are \cref{d11:20:local-null-variant}. Intersect a locally-null $N$ with each member of a countable finite-measure cover. Their countable union is $N$ and is null. Consequently the exceptional-level predicates in the two norm definitions coincide. Ordinary essentially bounded functions have pointwise bounded representatives by replacing their null exceptional set by zero.
\end{proof}
\source{COH 3.3 preliminary observations and quotient}

\begin{theorem}[Cohn's nonsigma-finite L1 dual theorem in its own measure framework]\family{D11.20.local_null_endpoint_specialization}\label{d11:20:local-null-endpoint-specialization}
Let $\mu$ be Radon on LCH $X$, put $\mu^*(A)=\inf_{U\supset A}\mu(U)$ over opens, let $\Sigma^*$ be the Caratheodory sigma-algebra and $\mu_1=\mu^*|_{\Sigma^*}$. Then
\[L^\infty_{\rm loc}(X,\Sigma^*,\mu_1)\longrightarrow L^1(X,\Sigma^*,\mu_1)^*,\quad [g]\longmapsto T_g
\]
is an isometric bijection, without sigma-finiteness.
\end{theorem}
\begin{proof}
The precise ingredients are \cref{d11:20:lp-dual-sigmafinite-local}, \cref{d11:20:local-null-variant}, \cref{d11:10:sigmafinite-support}. We prove the localization and gluing, including their measurability content. The construction of \cref{d11:19:positive-content-measure} applied to $f\mapsto\int f\dd\mu$ gives this same $\mu^*$ by open-mass recovery. It is complete on $\Sigma^*$. Every finite-measure $A\in\Sigma^*$ is compact-inner regular: choose finite-mass open $U\supset A$ with small excess, then an open set covering $U\setminus A$ with small outer mass and remove it from a compact inner approximation to $U$, exactly as in \cref{d11:18:radon-sigmafinite-inner}. This argument uses outer measure for $U\setminus A$ and so also works when $A$ is not Borel.

A subset $B\subset X$ lies in $\Sigma^*$ if and only if $B\cap K\in\Sigma^*$ for every compact $K$. For the nontrivial direction, a finite-mass open $U$ has compact $K_n\subset U$ whose union misses a null set. Completeness gives $B\cap U\in\Sigma^*$. For arbitrary $A$ of finite outer mass choose finite-mass opens $U\supset A$ with masses decreasing to $\mu^*(A)$; their measurable splits by $B\cap U$ give the Caratheodory inequality for $A$. Infinite-outer-mass test sets impose no further inequality, so $B\in\Sigma^*$.

Every positive-mass compact $K$ contains a compact $K_0$ of the same mass such that each open set meeting $K_0$ meets it in positive mass. Indeed remove from $K$ the union of opens $V$ with $\mu_1(V\cap K)=0$. Every compact subset of that union intersected with $K$ has zero mass by a finite subcover; finite-set compact inner regularity makes the removed part null. The remaining closed subset $K_0$ has the required property.

Use Zorn's lemma to choose a maximal disjoint family $\mathcal C$ of positive-mass compacts with that property. Unions of chains are admissible families, so maximality is justified. If $A\in\Sigma^*$ has finite mass, choose a finite-mass open $U\supset A$. Every $K\in\mathcal C$ meeting $A$ has $\mu_1(U\cap K)>0$. Disjointness and finiteness of $\mu_1(U)$ imply only countably many such $K$: those with mass greater than $1/n$ form a finite family for every $n$. Remove their countable union from $A$. A remaining positive-mass part would contain a positive compact and then a full-support compact of the preceding type, disjoint from all of $\mathcal C$, contradicting maximality. Thus $A$ is covered modulo null sets by those countably many members.

A subset $B$ whose intersection with every $K\in\mathcal C$ is measurable is measurable: test on an arbitrary compact $L$; it meets only countably many members of $\mathcal C$ and its remaining part is null by the preceding paragraph. Completeness and the compact-testing criterion finish the proof. Consequently a function is $\Sigma^*$-measurable if its restrictions to all $K\in\mathcal C$ are measurable (test inverse images; the zero value outside the family is harmless).

Given bounded $\ell$ on $L^1$, restrict it to each finite-measure $K\in\mathcal C$ by extension by zero. The finite-measure endpoint theorem supplies a measurable $g_K$ of pointwise modulus at most $\|\ell\|$ representing that restriction; replace null exceptional values by zero. Choose these representatives, define $g=g_K$ on each disjoint $K$, and zero elsewhere. The preceding measurability criterion proves $g$ measurable, despite an uncountable family. Each $L^1$ function has sigma-finite nonzero support; applying the finite-set decomposition countably often shows it vanishes a.e. off a countable union of members of $\mathcal C$. Finite partial sums of its restrictions converge in $L^1$. The two bounded functionals $\ell$ and $T_g$ agree on those partial sums, hence everywhere.

Finally locally-null changes do not affect any $L^1$ pairing, since an $L^1$ function has sigma-finite support. If $C$ is above the local norm, the exceptional level is locally null and $|\int fg|\le C\|f\|_1$. Conversely a level $\{|g|>c\}$ which is not locally null has positive intersection with some finite-measure set. Its phase-indicator test of norm one has pairing at least $c$. Hence $\|T_g\|=\|[g]\|_{\rm loc}$ and the map is injective on precisely this quotient. These arguments establish the full asserted isometric bijection; neither the Borel domain nor the ordinary quotient was silently substituted.
\end{proof}
\source{COH 7.5 introduction and Theorem7.5.4}

\section{Endpoint failures}

\begin{example}[Nonsigma-finite L1 duality can fail by measurability]\family{D11.20.endpoint_surjectivity_failure}\label{d11:20:endpoint-surjectivity-failure}
On $\R$ with its countable-or-cocountable sigma-algebra and restricted counting measure, $\ell(f)=\sum_{x>0}f(x)$ is a bounded functional on $L^1$ with no measurable bounded representing density.
\end{example}
\begin{proof}
For any integrable $f$, each $\{|f|>1/n\}$ is finite, so its nonzero support is countable and its integral is the absolutely convergent sum on that support. Thus $\ell$ is linear and $|\ell(f)|\le\|f\|_1$. If a measurable $g$ represented it, testing $\one_{\{x\}}$ would give $g(x)=1$ for $x>0$ and zero otherwise. The positive half-line is neither countable nor cocountable, so this $g$ is not measurable. This defeats surjectivity even with a semifinite reference measure.
\end{proof}
\source{COH 4.5 Example4.5.2; FOL 6.2 endpoint example}

\begin{example}[The dual of L-infinity is generally larger than L1]\family{D11.20.linfinity_dual_boundary}\label{d11:20:linfinity-dual-boundary}
On counting-measure $\N$ there is a positive bounded functional on $\ell^\infty$ with $\ell(1)=1$ and $\ell(e_n)=0$ for every $n$. It has no $\ell^1$ density.
\end{example}
\begin{proof}
Extend the cofinite filter to a free ultrafilter $\mathcal U$ by the ultrafilter lemma. Every bounded real or complex sequence has a unique $\mathcal U$-limit: its values lie in a compact scalar disk, and closures of values on members of $\mathcal U$ have the finite-intersection property; compactness gives a point and disjoint neighborhoods give uniqueness. Define $\ell(a)$ to be that limit. Intersecting the ultrafilter neighborhoods for two sequences and using scalar continuity proves additivity and homogeneity. Real nonnegative sequences have nonnegative limits, and $|\ell(a)|\le\|a\|_\infty$. The constant one has limit one, whereas each singleton sequence vanishes on a cofinite set. If $\ell(a)=\sum_na_ng_n$ with $g\in\ell^1$, singleton tests force all $g_n=0$, contradicting $\ell(1)=1$. Thus the integral embedding into this endpoint dual is proper.
\end{proof}
\source{FOL 6.2 Exercise19 and preceding L-infinity discussion}

% END chapters/20


% BEGIN chapters/21
\chapter{Haar measure on general locally compact groups}\label{ch:d11-21}

The group need not be second countable, metrizable, or sigma-compact. We construct a compact content using a compact product, turn it into a Radon measure, and prove uniqueness by compactly supported kernels. All product integrations in that proof are justified on finite compact rectangles, after product measurability has been proved. The modular function is introduced only after uniqueness.

\section{The group and a compact-content construction}

\begin{definition}[Locally compact Hausdorff group and translations]\family{D11.21.lch_group_interface}\label{d11:21:lch-group-interface}
A topological group has continuous multiplication and inversion. Here a locally compact group is Hausdorff and locally compact. Left and right translations $L_a(x)=ax$, $R_a(x)=xa$, and inversion are homeomorphisms, hence preserve Borel measurability and compactness.
\end{definition}
\begin{proof}
Their inverses are $L_{a^{-1}}$, $R_{a^{-1}}$, and inversion itself; continuity follows from the group operations. Thus inverse images and images of opens, compacts, and Borel sets have the stated properties.
\end{proof}
\source{COH 9.1 definition and Proposition9.1.2; FOL 11.1 conventions}

\begin{definition}[Left and right Haar measures in a precise regularity class]\family{D11.21.haar_definition}\label{d11:21:haar-definition}
A left Haar measure is a nonzero positive Radon Borel measure $\mu$ with $\mu(aE)=\mu(E)$ for every $a$ and Borel $E$. A right Haar measure satisfies $\mu(Ea)=\mu(E)$. The Radon predicates are exactly those of \cref{d11:18:radon-and-source-names}; invariant set functions without that regularity are not included.
\end{definition}
\source{COH 9.2 definition; FOL 11.1 Haar definition}

\begin{lemma}[Finite compact covering ratios with fixed normalization]\family{D11.21.compact_covering_ratios}\label{d11:21:compact-covering-ratios}
Choose compact $K_0$ with nonempty interior. Let $N(K:U)$ be the least number of left translates of an identity neighborhood $U$ covering compact $K$, and set $h_U(K)=N(K:U)/N(K_0:U)$. Then $h_U(K_0)=1$, $0\le h_U(K)\le N(K:K_0)$, and $h_U$ is left invariant, monotone, and subadditive on unions. If $KU^{-1}\cap LU^{-1}=\varnothing$, then $h_U(K\cup L)=h_U(K)+h_U(L)$.
\end{lemma}
\begin{proof}
A precompact identity neighborhood supplies $K_0$. Every nonempty open set has finitely many left translates covering a given compact, so the covering numbers are finite nonnegative integers; the denominator is positive. Cover $K$ by $r=N(K:K_0)$ translates of $K_0$, and each translate by $N(K_0:U)$ translates of $U$, giving the bound. Translating covers, restricting covers, and taking unions prove the other inequalities. A translate $aU$ meeting $K$ has $a\in KU^{-1}$. Under the separation hypothesis no such translate can meet both $K$ and $L$, so any cover of their union divides into two covers and has at least the sum of the minimal numbers. The empty compact has number zero.
\end{proof}
\source{COH 9.2 Lemma9.2.3}

\begin{construction}[Obtain invariant additive compact content by a compact-product cluster point]\family{D11.21.compact_content_cluster}\label{d11:21:compact-content-cluster}
There is a finite nonnegative function $h$ on compact subsets with $h(\varnothing)=0$, $h(K_0)=1$, left invariance, monotonicity, finite subadditivity, and additivity on disjoint compact sets.
\end{construction}
\begin{proof}
The precise ingredients are \cref{d11:21:compact-covering-ratios}. Form the compact product $P=\prod_K[0,N(K:K_0)]$, indexed by all compact subsets. Compactness is precisely the Tychonoff compact-product theorem from the topology foundation; no sequential compactness is being inferred. For each identity neighborhood $V$, let $S_V$ be the closure in $P$ of $\{h_U:U\subset V,\ U\text{ an identity neighborhood}\}$. Any finitely many $V$ have a common smaller identity neighborhood, so the closed sets $S_V$ have the finite-intersection property. Choose $h\in\bigcap_VS_V$.

Every finite-coordinate equality or inequality valid for all $h_U$ is closed in $P$, so the preceding lemma transfers normalization, invariance, monotonicity, and subadditivity. If compact $K,L$ are disjoint, the compact $K^{-1}L$ misses the identity. Choose $V$ so small that $V^{-1}V\cap K^{-1}L=\varnothing$. Then $KV^{-1}\cap LV^{-1}=\varnothing$, and the preceding additivity equality holds for every $U\subset V$. It transfers to $S_V$ and hence to $h$. Empty sets cause no exception. This argument selects a cluster point by finite constraints, rather than assuming a countable neighborhood base or a convergent sequence.
\end{proof}
\source{COH 9.2 compact-product construction and Lemma9.2.4}

\begin{construction}[Convert invariant compact content into a nonzero invariant Radon measure]\family{D11.21.compact_content_radon_measure}\label{d11:21:compact-content-radon-measure}
Put $c(U)=\sup_{K\subset U,\ K\text{ compact}}h(K)$ and $\mu^*(A)=\inf_{U\supset A,\ U\text{ open}}c(U)$. Then $\mu^*$ is an outer measure, its Borel restriction is left invariant and Radon, and $\mu(K_0)\ge1$.
\end{construction}
\begin{proof}
The precise ingredients are \cref{d11:21:compact-content-cluster}. If $K\subset\bigcup_nU_n$ is compact, a finite shrinking of its finite open subcover gives compact pieces $K_i\subset U_i$ covering $K$. Subadditivity of $h$ gives $h(K)\le\sum_i h(K_i)\le\sum_n c(U_n)$. Taking a supremum proves open countable subadditivity, and open approximation with summable errors proves outer-measure subadditivity exactly as in the positive-functional construction. Also $\mu^*(U)=c(U)$ for opens.

For an open $G$, a finite-outer-mass test set $A$, and an open $V\supset A$ with $c(V)<\mu^*(A)+\eps$, fix compact $K\subset V\cap G$. For every compact $L\subset V\setminus K$, disjoint compact additivity gives $c(V)\ge h(K)+h(L)$. Hence $c(V)\ge h(K)+c(V\setminus K)\ge h(K)+\mu^*(A\setminus G)$. Supremizing over $K$ proves the Caratheodory inequality for $G$; let $\eps\downarrow0$. Infinite test masses impose no further inequality. Borel measurability follows from \cref{d11:02:caratheodory-theorem}.

For compact $K$ choose precompact open $W\supset K$. Then $\mu(K)\le c(W)\le h(\overline W)<\infty$. Conversely $h(K)\le c(U)$ for every open $U\supset K$, so $h(K)\le\mu(K)$. Thus $c(U)=\sup_{K\subset U}\mu(K)$, proving open inner regularity. Outer regularity is definitional. Translations permute the compact and open comparison sets and preserve $h$, hence preserve $c$, $\mu^*$, and the Borel measure. Finally $\mu(K_0)\ge h(K_0)=1$, so the measure is nonzero. Equality at $K_0$ has not been asserted; if desired divide by its finite positive mass afterward.
\end{proof}
\source{COH 9.2 completion of Theorem9.2.2}

\begin{theorem}[Existence of left Haar measure]\family{D11.21.haar_existence}\label{d11:21:haar-existence}
Every LCH group has a left Haar measure, without any countability assumption.
\end{theorem}
\begin{proof}
The precise ingredients are \cref{d11:21:compact-content-radon-measure}. The measure in the preceding construction is Borel, Radon, nonzero, and left invariant, exactly the four properties in the definition. Its construction used finite compact coverings, the compact product theorem, outer-measure extension, and separation of compact sets. No Haar existence theorem or functional-analysis compactness principle was taken as background.
\end{proof}
\source{COH 9.2 Theorem9.2.2 complete construction; FOL 11.1 Theorem11.8 alternative functional route}

\begin{lemma}[Haar measure is positive on nonempty opens]\family{D11.21.haar_positive_open}\label{d11:21:haar-positive-open}
A Haar measure assigns positive mass to every nonempty open set. Every nonzero nonnegative $f\in C_c(G)$ has positive integral, and the measure has full support.
\end{lemma}
\begin{proof}
Since $\mu\ne0$, $\mu(G)>0$; open inner regularity supplies a positive-mass compact $K$. If a nonempty open $U$ had mass zero, finitely many left translates of $U$ would cover $K$, forcing its mass to vanish. If $f(x)>0$, continuity gives a nonempty open set on which $f>c>0$, so its integral is positive. The support characterization \cref{d11:18:radon-support} now gives full support.
\end{proof}
\source{COH 9.2 Lemma9.2.5; FOL 11.1 Proposition11.4(c)}

\begin{proposition}[Measure invariance and translated integrals]\family{D11.21.haar_translation_integral}\label{d11:21:haar-translation-integral}
For left Haar $\mu$, $\int f(a^{-1}x)\dd\mu(x)=\int f\dd\mu$ for nonnegative Borel or integrable scalar $f$. Conversely this identity for all real $C_c$ tests, for every $a$, makes a nonzero Radon measure left Haar.
\end{proposition}
\begin{proof}
Apply set invariance to indicator functions, then simple approximation and monotone convergence, and finally integrable positive-negative or real-imaginary decomposition. Conversely the pushforward by $L_a$ is Radon because $L_a$ is a homeomorphism. Its $C_c$ integrals equal those of $\mu$, so positive Riesz--Markov uniqueness \cref{d11:19:positive-riesz-markov} makes the measures equal. This is precisely left invariance.
\end{proof}
\source{COH 9.2 translated integration identity; FOL 11.1 Proposition11.4(b)}

\section{Uniqueness on compact rectangles}

\begin{lemma}[Compact-support translate integrals are continuous]\family{D11.21.compact_translate_integral_continuity}\label{d11:21:compact-translate-integral-continuity}
If $\nu$ is Radon and $g\in C_c(G)$, then $B(x)=\int g(tx)\dd\nu(t)$ is continuous. For Haar $\nu$ and $0\ne g\ge0$, it is finite and positive everywhere.
\end{lemma}
\begin{proof}
The precise ingredients are \cref{d11:21:haar-positive-open}. Near a fixed $x_0$, restrict $x$ to a precompact neighborhood $O$. The $t$-supports lie in $\supp(g)\,\overline O^{-1}$, a compact set of finite measure. Continuity on a compact set, followed by a finite tube cover, makes $g(tx)$ converge uniformly in $t$ to $g(tx_0)$ as $x\to x_0$. Multiplication by that finite compact mass proves continuity of the integrals, in the neighborhood definition of continuity rather than a merely sequential test. Each translated function is a nonzero nonnegative $C_c$ function, so positivity follows from the preceding Haar lemma. Compact finiteness gives finiteness.
\end{proof}
\source{COH 9.1 Corollary9.1.6 as used in9.2.6; FOL 11.1 Proposition11.2}

\begin{lemma}[A compact double-integral identity fixes ratios of Haar integrals]\family{D11.21.compact_kernel_ratio_identity}\label{d11:21:compact-kernel-ratio-identity}
For left Haar $\mu,\nu$ and fixed $0\ne g\in C_c(G,\R)$, $g\ge0$, set $B(x)=\int g(tx)\dd\nu(t)$. Then for every real $f\in C_c(G)$,
\[\int f\dd\mu=\left(\int g\dd\mu\right)\int\frac{f(y^{-1})}{B(y^{-1})}\dd\nu(y).
\]
\end{lemma}
\begin{proof}
The precise ingredients are \cref{d11:21:compact-translate-integral-continuity}, \cref{d11:21:haar-translation-integral}. Set $K=\supp f$, $L=\supp g$, and $h(x,y)=f(x)g(yx)/B(x)$. The preceding lemma makes it continuous; it is supported in $K\times LK^{-1}$, hence compactly supported. The two transformed kernels $h(y^{-1}x,y)$ and $h(y^{-1},xy)$ have supports respectively in $L\times LK^{-1}$ and $L\times K^{-1}$. By \cref{d11:18:cc-tensor-approximation} all three kernels are measurable for $\BB(G)\otimes\BB(G)$. Restrict $\mu$ to the compact set $K\cup L$ and $\nu$ to $LK^{-1}\cup K^{-1}$ whenever invoking Fubini. These are finite measures and dominate every displayed kernel support. Thus each exchange of integrals is the finite-measure theorem \cref{d11:08:fubini}; no global sigma-finiteness or equality of the two product Borel sigma-algebras is used.

Writing the integrals over $G$ by zero extension, and using left invariance for the changes $x\mapsto y^{-1}x$ and $y\mapsto xy$, gives
\begin{align*}
\int f\dd\mu
 &=\int\!\int h(x,y)\dd\nu(y)\dd\mu(x)\\
 &=\int\!\int h(y^{-1}x,y)\dd\mu(x)\dd\nu(y)\\
 &=\int\!\int h(y^{-1}x,y)\dd\nu(y)\dd\mu(x)\\
 &=\int\!\int h(y^{-1},xy)\dd\nu(y)\dd\mu(x)\\
 &=\int\!\int \frac{f(y^{-1})g(x)}{B(y^{-1})}\dd\mu(x)\dd\nu(y).
\end{align*}
The first equality uses the definition of $B$, and the final equality yields the formula. Translation invariance is applied to the original integrals over $G$, not incorrectly to the restricted measures. No right invariance or modular formula has entered.
\end{proof}
\source{COH 9.2 proof of Theorem9.2.6}

\begin{theorem}[Uniqueness of Haar measure up to a positive scalar]\family{D11.21.haar_uniqueness}\label{d11:21:haar-uniqueness}
Two left Haar measures $\mu,\nu$ have $\nu=c\mu$ on all Borel sets for a unique finite real $c>0$. The same holds for right Haar measures.
\end{theorem}
\begin{proof}
The precise ingredients are \cref{d11:21:compact-kernel-ratio-identity}, \cref{d11:21:haar-inversion}. Apply the kernel identity first to $(\mu,\nu)$ and then to $(\nu,\nu)$, with the same $g$. Since both integrals of $g$ are finite and positive, it gives $\int f\dd\nu=c\int f\dd\mu$ for every real $C_c$ test, where $c=\int g\dd\nu/\int g\dd\mu>0$. Radon uniqueness \cref{d11:19:positive-riesz-markov} gives Borel equality. Testing $g$ proves uniqueness of the scalar. Inversion changes right Haar into left Haar by the next elementary pushforward result, giving the right version. Thus the uniqueness class and scalar finiteness are both explicit.
\end{proof}
\source{COH 9.2 Theorem9.2.6 and 9.3 Corollary9.3.2; FOL 11.1 Theorem11.9}

\begin{proposition}[Inversion exchanges left and right Haar measures]\family{D11.21.haar_inversion}\label{d11:21:haar-inversion}
For Radon $\mu$, $\check\mu(E)=\mu(E^{-1})$ is Radon and
$\int f\dd\check\mu=\int f(x^{-1})\dd\mu(x)$ for nonnegative Borel or integrable scalar $f$. Left Haar corresponds under inversion to right Haar.
\end{proposition}
\begin{proof}
Inversion is a homeomorphism, so its pushforward preserves compact finiteness and both regularity comparison formulas. The integral formula follows for indicators, then by simple approximation and decomposition. Since $(Ea)^{-1}=a^{-1}E^{-1}$, left invariance of $\mu$ gives right invariance of $\check\mu$. Inversion is an involution, proving the converse and preserving nonzero mass. This proof does not use Haar uniqueness, so the preceding right-uniqueness reference is acyclic.
\end{proof}
\source{COH 9.3 Proposition9.3.1; FOL 11.1 Proposition11.4(a)}

\begin{theorem}[Finite total Haar mass characterizes compact groups]\family{D11.21.haar_finite_iff_compact}\label{d11:21:haar-finite-iff-compact}
A left Haar measure has finite total mass exactly when $G$ is compact. Thus a compact Hausdorff group has a unique left Haar probability measure.
\end{theorem}
\begin{proof}
The precise ingredients are \cref{d11:21:haar-positive-open}, \cref{d11:21:haar-uniqueness}. Compact finiteness proves one direction. Conversely choose a compact identity neighborhood $K$ with nonempty interior and positive finite mass. If $G$ is not compact, choose recursively $a_n$ outside the finite compact union $\bigcup_{j<n}a_jKK^{-1}$. Then the translates $a_nK$ are pairwise disjoint: an intersection would put $a_n$ in one of these forbidden sets. Their masses are the same positive number, forcing $\mu(G)=\infty$. Thus a finite Haar measure has compact group, and its total mass is positive. Divide by that mass to obtain probability; uniqueness up to scale makes it the unique one.
\end{proof}
\source{COH 9.3 Proposition9.3.3 and following normalization; FOL 11.1 Proposition11.4(d)}

\section{The modular convention and examples}

\begin{definition}[Modular function with an explicit right-set convention]\family{D11.21.modular_function}\label{d11:21:modular-function}
For left Haar $\mu$, define $\Delta(a)>0$ by $\mu(Ea)=\Delta(a)\mu(E)$ for every Borel $E$. This convention uses right translation of sets. The scalar exists, is finite and positive, and is independent of normalization of $\mu$.
\end{definition}
\begin{proof}
The precise ingredients are \cref{d11:21:haar-uniqueness}. The measure $E\mapsto\mu(Ea)$ is another nonzero left Haar measure, because right translation commutes with left translation and is a homeomorphism. Uniqueness supplies its scalar; scaling $\mu$ cancels from the equality. This introduction comes after, and was not used in, the uniqueness proof.
\end{proof}
\source{COH 9.3 modular-function definition; FOL 11.1 modular discussion}

\begin{theorem}[Continuity and homomorphism property of the modular function]\family{D11.21.modular_properties}\label{d11:21:modular-properties}
The map $\Delta:G\to(0,\infty)$ is a continuous homomorphism, $\Delta(e)=1$ and $\Delta(a^{-1})=\Delta(a)^{-1}$. Moreover
$\int f(xa^{-1})\dd\mu(x)=\Delta(a)\int f\dd\mu$
for nonnegative Borel or integrable scalar $f$.
\end{theorem}
\begin{proof}
The precise ingredients are \cref{d11:21:modular-function}, \cref{d11:21:compact-translate-integral-continuity}. Use a compact Borel set with finite positive mass in the set-translation formula twice to get $\Delta(ab)=\Delta(a)\Delta(b)$, and use $e$ and inverses for the remaining algebra. Indicator functions, simple approximation, and decomposition give the integral identity. For fixed nonzero $g\ge0$ in $C_c$ it expresses $\Delta(a)$ as $\int g(xa^{-1})\dd\mu(x)/\int g\dd\mu$. Compact-translate continuity and inversion make this continuous.
\end{proof}
\source{COH 9.3 Proposition9.3.4 and equation2; FOL 11.1 Proposition11.10}

\begin{corollary}[Abelian and compact groups are unimodular]\family{D11.21.unimodular_cases}\label{d11:21:unimodular-cases}
A group is unimodular when $\Delta\equiv1$, equivalently its left Haar measures are also right Haar. Abelian and compact Hausdorff groups are unimodular, and compact Haar probability is invariant on both sides.
\end{corollary}
\begin{proof}
The precise ingredients are \cref{d11:21:modular-function}, \cref{d11:21:haar-finite-iff-compact}. The equivalence is exactly the right-set formula. In an abelian group right and left translations coincide. For compact $G$, substitute $E=G$ of finite positive mass into $\mu(Ga)=\Delta(a)\mu(G)$, giving $\Delta(a)=1$. Normalization preserves both invariances.
\end{proof}
\source{COH 9.3 unimodularity and Proposition9.3.5}

\begin{theorem}[The inversion measure has a modular density]\family{D11.21.left_right_density}\label{d11:21:left-right-density}
With this modular convention,
\[\check\mu(E)=\int_E\Delta(x^{-1})\dd\mu(x)\quad(E\text{ Borel}).\]
Thus left and right Haar measures have the same Borel null sets.
\end{theorem}
\begin{proof}
The precise ingredients are \cref{d11:21:modular-properties}, \cref{d11:21:haar-uniqueness}, \cref{d11:21:haar-inversion}, \cref{d11:18:radon-sigmafinite-inner}. Let $w(x)=\Delta(x^{-1})$ and $\rho(E)=\int_Ew\dd\mu$. We prove the needed weighted regularity explicitly. Boundedness of $w$ on compacts gives compact finiteness. For open $U$, the dyadic lower sums
$2^{-n}\sum_{j=1}^{n2^n}\one_{U\cap\{w>j2^{-n}\}}$
increase to $w\one_U$. Each level set is open, so its measure is the supremum of compact inner masses. Given a finite target below the integral of one such finite sum, choose compact subsets of its finitely many open levels meeting that target, and take their union $K\subset U$. Then $\int_Kw$ exceeds the target. This proves $\rho(U)=\sup_{K\subset U}\rho(K)$, including infinite values.

For Borel $E$ with $\rho(E)<\infty$, partition it into $E_k=E\cap\{2^k\le w<2^{k+1}\}$, $k\in\Z$. Then $\mu(E_k)\le2^{-k}\rho(E)<\infty$. Choose positive errors $\eta_k$ with $\sum_{k\in\Z}2^{k+2}\eta_k<\eps$. Outer regularity of $\mu$ gives open $V_k\supset E_k$ with $\mu(V_k\setminus E_k)<\eta_k$. Put $O_k=V_k\cap\{w<2^{k+2}\}$. It still contains $E_k$, and
$\rho(O_k\setminus E_k)\le2^{k+2}\eta_k$.
The open $O=\bigcup_kO_k$ contains $E$ and satisfies $\rho(O\setminus E)<\eps$. For infinite-$\rho$ Borel sets outer regularity is automatic by monotonicity. Thus $\rho$ is Radon, and it is nonzero since $w>0$ everywhere and $\mu\ne0$.

The right-translation integral formula gives $\rho(Ea)=\Delta(a)\int_E w(xa)\dd\mu(x)=\rho(E)$, since $w(xa)=\Delta(a)^{-1}w(x)$. Thus $\rho$ and $\check\mu$ are right Haar, so $\check\mu=c\rho$ for $c>0$. Choose a symmetric precompact identity neighborhood $V$ so small that $1-\eps<w<1+\eps$ on $V$. Its mass is finite and positive and $V^{-1}=V$, so
$\mu(V)=\check\mu(V)=c\int_Vw\dd\mu$.
Dividing gives $c(1-\eps)\le1\le c(1+\eps)$. Let $\eps\downarrow0$ to obtain $c=1$. The weight is everywhere positive finite, so $\rho$ and $\mu$ have the same null sets; right uniqueness extends this to any right Haar normalization.
\end{proof}
\source{COH 9.3 Proposition9.3.6 and Corollary9.3.7}

\begin{example}[Standard invariant measures and a nonunimodular example]\family{D11.21.haar_examples_and_nonunimodular}\label{d11:21:haar-examples-and-nonunimodular}
Lebesgue measure on additive $\R^d$, counting measure on any discrete group, and arc-length measure on the circle are two-sided Haar. On the affine group $(a,b)(c,d)=(ac,b+ad)$, $a,c>0$, left and right Haar measures are $da\,db/a^2$ and $da\,db/a$, respectively, and $\Delta(a,b)=1/a$.
\end{example}
\begin{proof}
The precise ingredients are \cref{d11:21:modular-function}, \cref{d11:18:second-countable-borel-regular}. Translation invariance and Euclidean regularity give the first example. In a discrete space every set is Borel and open and every compact is finite, so counting measure is compact-finite and both regularities follow directly by using finite subsets; translations are bijections. For the circle, push Lebesgue measure on $[0,2\pi)$ forward by $t\mapsto e^{it}$. Splitting an interval at the wrap point proves rotational invariance for Borel sets by ordinary translation invariance; finite metric Borel regularity supplies Radon regularity.

For the affine example the group operations and inversion are continuous on the open half-plane. Left translation by $(a,b)$ has Jacobian $a^2$ in $(c,d)$, while its new first coordinate is $ac$, so the density $1/c^2$ is invariant. Right translation by $(a,b)$ maps $(c,d)$ to $(ca,d+cb)$ with Jacobian $a$ and new first coordinate $ca$, so the density $1/c$ is right invariant. These positive continuous densities define Radon measures, locally finite on the half-plane. For left density $1/c^2$, right translation changes mass by the factor $a/(a^2)=1/a$, proving the displayed modular function and failure of unimodularity when $a\ne1$. Each fixed translation in these coordinates is an invertible linear map followed by an ordinary translation. The displayed mass computations therefore apply the proved Lebesgue linear Jacobian theorem \cref{d11:09:linear-jacobian} and Euclidean translation invariance, restricted to the open half-plane; they do not assume a new calculus-only Lebesgue substitution rule.
\end{proof}
\source{COH 9.2 Examples9.2.1 and Exercise5; 9.3 Exercise2; FOL 11.1 Haar examples}

% END chapters/21


% BEGIN chapters/22
\chapter{Hausdorff measure and elementary dimension}\label{ch:d11-22}

All covering sets below may be arbitrary subsets, so our statements concern outer measures until measurability is established. We use coefficient-one diameter powers. This choice is decisive at dimension zero and in exact Cantor masses; Euclidean volume constants are not imported from a different normalization or metric.

\section{Covers, metric outer measures, and zero dimension}

\begin{definition}[Diameter-cover content at scale delta]\family{D11.22.hausdorff_scale_content}\label{d11:22:hausdorff-scale-content}
For metric $X$, $s\ge0$, $\delta>0$, and $E\subset X$, define
\[H^s_\delta(E)=\inf\left\{\sum_j(\diam U_j)^s:E\subset\bigcup_jU_j,\quad \varnothing\ne U_j\subset X,\quad\diam U_j\le\delta\right\}.
\]
Covers are at most countable; the empty cover gives zero for $E=\varnothing$, and an empty admissible class has infimum infinity. When $s=0$ each nonempty covering set costs one, even a singleton. When $s>0$, a singleton costs zero.
\end{definition}
\source{FOL 11.2 Hausdorff definition and zero-dimensional exercise; STE Chapter7 §1 definition}

\begin{theorem}[Hausdorff outer measure as the small-scale limit]\family{D11.22.hausdorff_outer_measure}\label{d11:22:hausdorff-outer-measure}
The contents increase as $\delta$ decreases. Their supremum $H^s(E)=\sup_{\delta>0}H^s_\delta(E)$ is an outer measure. If $\dist(A,B)>0$, then $H^s(A\cup B)=H^s(A)+H^s(B)$.
\end{theorem}
\begin{proof}
Smaller diameter bounds give fewer covers. For fixed $\delta$, monotonicity and the empty-set value are immediate, and adjoining near-minimal covers with summable errors proves countable subadditivity; if any right-hand term is infinite no approximation is needed. Thus
$H^s_\delta(\bigcup_iE_i)\le\sum_iH^s_\delta(E_i)\le\sum_iH^s(E_i)$.
Take the supremum over $\delta$ to obtain outer-measure subadditivity. If $\delta<\dist(A,B)$, a covering set of diameter at most $\delta$ cannot meet both. Partitioning any cover accordingly gives $H^s_\delta(A\cup B)\ge H^s_\delta(A)+H^s_\delta(B)$. Let $\delta\downarrow0$; the increasing extended limits add, giving the asserted lower bound, and outer subadditivity gives the reverse bound.
\end{proof}
\source{FOL 11.2 Proposition11.17; STE Chapter7 §1 Properties1-3}

\begin{theorem}[Metric outer measures make Borel sets measurable]\family{D11.22.metric_outer_borel}\label{d11:22:metric-outer-borel}
Every outer measure $m^*$ on a metric space which is additive on positively separated sets makes all Borel sets Caratheodory measurable. In particular $H^s|_{\BB(X)}$ is a measure and its full Caratheodory extension is complete.
\end{theorem}
\begin{proof}
The precise ingredients are \cref{d11:22:hausdorff-outer-measure}. It suffices to prove every closed $F$ measurable. Test a set $A$ with $m^*(A)<\infty$; infinite-mass tests impose no additional inequality. Put $A_n=\{x\in A:d(x,F)\ge1/n\}$ and $C_n=A_{n+1}\setminus A_n$. The sets $A_n$ increase to $A\setminus F$. Although outer measures are not generally continuous from below, here we have the needed substitute. Within either parity, every finite collection of the $C_n$ is positively separated: the distance-to-$F$ coordinate lies in $[1/(n+1),1/n)$, and a skipped index leaves a positive gap. Repeated separated additivity therefore bounds each parity's sum of $m^*(C_n)$ by $m^*(A)$. Both series converge and their tails tend to zero. Since
$A\setminus F=A_n\cup\bigcup_{k\ge n}C_k$,
subadditivity gives $m^*(A\setminus F)\le m^*(A_n)+o(1)$. The separated sets $A_n$ and $A\cap F$ give $m^*(A)\ge m^*(A_n)+m^*(A\cap F)$. Combine and let $n\to\infty$ to obtain the Caratheodory inequality. If $F=\varnothing$, measurability is immediate. The Caratheodory measurable sets form a sigma-algebra and the restriction is complete by \cref{d11:02:caratheodory-theorem}, proving the result.
\end{proof}
\source{FOL 11.2 Proposition11.16; STE Chapter7 §1 Borel restriction}

\begin{theorem}[Zero-dimensional Hausdorff measure counts points]\family{D11.22.zero_dimension_counting}\label{d11:22:zero-dimension-counting}
With the stated convention, $H^0(E)$ is the number of points when $E$ is finite and is infinity when $E$ is infinite.
\end{theorem}
\begin{proof}
Singleton covers give the finite upper bound. For any $N$ distinct points in $E$, their finite set has a positive minimum pairwise distance. A cover at smaller diameter needs at least $N$ nonempty members. Thus $H^0(E)\ge N$. For finite $E$ choose $N=|E|$; for infinite $E$ let $N$ be arbitrary. The empty set has mass zero. This also identifies the exact place where a singleton-zero library convention would fail to match ours.
\end{proof}
\source{FOL 11.2 Exercise13; STE Chapter7 §1 Property6}

\section{Images and the critical exponent}

\begin{theorem}[Isometry invariance and similarity scaling]\family{D11.22.isometry_similarity}\label{d11:22:isometry-similarity}
An isometry preserves every $H^s$. A similarity with distance multiplier $r>0$ satisfies $H^s(f(E))=r^sH^s(E)$, including $s=0$.
\end{theorem}
\begin{proof}
Intersect an ambient cover of $E$ with $E$, discard empty pieces, and send it forward. Diameters multiply by $r$, and zero-dimensional costs remain one. This proves $H^s_{r\delta}(f(E))\le r^sH^s_\delta(E)$. Conversely intersect any cover in the target with $f(E)$ and pull back by the inverse similarity on its image. It gives the reverse inequality. Take small-scale limits. This works for embeddings into a larger space as well as surjective isometries.
\end{proof}
\source{FOL 11.2 Proposition11.18; STE Chapter7 §1 Property5}

\begin{proposition}[The s=0 and zero-Lipschitz cases]\family{D11.22.lipschitz_zero_cases}\label{d11:22:lipschitz-zero-cases}
If $f:E\to Y$ is $L$-Lipschitz with $L>0$, then $H^s(f(A))\le L^sH^s(A)$ for every $A\subset E$ and $s\ge0$. If $L=0$, all nonempty images are singletons: they have zero $H^s$ for $s>0$ and $H^0$ equal to one; the empty image has all masses zero.
\end{proposition}
\begin{proof}
The precise ingredients are \cref{d11:22:zero-dimension-counting}, \cref{d11:22:hausdorff-scale-content}. Intersect covers with $A$ and send them forward; their diameters are at most $L$ times the old diameters. For $s=0$ the number of nonempty image sets does not increase, so the same proof applies. If $L=0$, any two image points have zero distance and are equal; singleton covers and the preceding counting theorem give precisely the stated values. No value for $0^0$ has been silently imposed.
\end{proof}
\source{FOL 11.2 Proposition11.18 and Exercise13; STE Chapter7 Property6}

\begin{theorem}[Holder maps bound Hausdorff measures and dimension]\family{D11.22.holder_image_bound}\label{d11:22:holder-image-bound}
If $0<\alpha\le1$, $C>0$, and $d(f(x),f(y))\le C d(x,y)^\alpha$ on $E$, then, for $s>0$,
\[H^{s/\alpha}(f(E))\le C^{s/\alpha}H^s(E),\qquad
\dim_H f(E)\le\dim_H E/\alpha.
\]
\end{theorem}
\begin{proof}
The precise ingredients are \cref{d11:22:exponent-threshold}, \cref{d11:22:dimension-definition}. The image of an intersected $\delta$-cover has diameters at most $C\delta^\alpha$, and the cost with exponent $s/\alpha$ is bounded by $C^{s/\alpha}$ times its old cost. Infima and small-scale limits prove the outer-measure inequality without an image measurability assumption. If $t>\dim_HE/\alpha$, the dimension threshold proved next gives $H^{\alpha t}(E)=0$, and the inequality gives $H^t(f(E))=0$. Taking infima gives the dimension assertion; if the original dimension is infinity it is vacuous, and empty sets are immediate.
\end{proof}
\source{STE Chapter7 §2.1 Lemma2.2; FOL 11.2 Proposition11.18 Lipschitz case}

\begin{theorem}[Hausdorff measure has a critical exponent]\family{D11.22.exponent_threshold}\label{d11:22:exponent-threshold}
If $H^s(E)<\infty$ and $t>s\ge0$, then $H^t(E)=0$. Equivalently $H^t(E)>0$ forces $H^s(E)=\infty$ for $s<t$.
\end{theorem}
\begin{proof}
Every $\delta$-cover satisfies $\sum(\diam U)^t\le\delta^{t-s}\sum(\diam U)^s$, including singleton costs under the stipulated convention. Hence $H^t_\delta(E)\le\delta^{t-s}H^s_\delta(E)\le\delta^{t-s}H^s(E)$. Let $\delta\downarrow0$. The equivalent assertion is its contrapositive. No conclusion about the mass at the critical exponent follows.
\end{proof}
\source{FOL 11.2 Proposition11.19; STE Chapter7 §1 Property8}

\begin{definition}[Hausdorff dimension and endpoint conventions]\family{D11.22.dimension_definition}\label{d11:22:dimension-definition}
For nonempty $E$, define $\dim_HE=\inf\{s\ge0:H^s(E)=0\}$ with infimum infinity for the empty set of exponents, and define $\dim_H\varnothing=0$. Above this critical value all masses vanish; below it they are infinite. At dimension zero there is no lower nonnegative exponent.
\end{definition}
\begin{proof}
The precise ingredients are \cref{d11:22:exponent-threshold}. If $t$ exceeds the infimum, choose a zero exponent $s<t$ and use the threshold theorem. If $s$ is below the infimum and $H^s(E)$ were finite, every larger exponent would vanish, contradicting the position of the infimum. The assertions also cover infinite critical value. The critical mass itself may be zero, positive finite, or infinite; the definition does not decide it.
\end{proof}
\source{FOL 11.2 dimension discussion; STE Chapter7 §2 definition}

\begin{theorem}[Countable stability of Hausdorff dimension]\family{D11.22.countable_stability}\label{d11:22:countable-stability}
For a countable family, $\dim_H(\bigcup_nE_n)=\sup_n\dim_HE_n$, with the empty-family supremum zero. Dimension is monotone, and every countable set has dimension zero.
\end{theorem}
\begin{proof}
The precise ingredients are \cref{d11:22:dimension-definition}, \cref{d11:22:hausdorff-outer-measure}. Monotonicity of outer measure gives monotonicity of dimension and the lower bound. If $s$ exceeds the supremum, every $H^s(E_n)=0$, so outer subadditivity gives $H^s(\bigcup_nE_n)=0$. Let $s$ decrease to the supremum. If that supremum is infinite, the lower bound already proves equality. A singleton has zero $H^s$ for every $s>0$ by its diameter-zero cover; countable stability proves the last assertion, with the empty-set convention covering no points.
\end{proof}
\source{FOL 11.2 Exercise14; STE Chapter7 Exercise6}

\section{Euclidean normalization without an unsupported constant}

\begin{theorem}[Euclidean n-dimensional measure agrees up to a positive constant]\family{D11.22.euclidean_lebesgue_normalization}\label{d11:22:euclidean-lebesgue-normalization}
For each integer $n\ge1$ there is $c_n>0$ with $m_n(E)=c_nH^n(E)$ for Borel $E\subset\R^n$ in the Euclidean metric. Here $c_1=1$. For $n>1$ no explicit constant is asserted.
\end{theorem}
\begin{proof}
The precise ingredients are \cref{d11:22:metric-outer-borel}, \cref{d11:22:isometry-similarity}. For any cover of the unit cube, a set of diameter $D$ is contained in an axis-parallel box of side lengths at most $D$ (enlarge endpoints by arbitrarily small amounts if extrema are not attained). Lebesgue outer measure therefore gives $1\le\sum_j(\diam U_j)^n$, so $H^n([0,1]^n)\ge1$. A grid of $k^n$ cubes of side $1/k$ gives cost $n^{n/2}$ at diameter $\sqrt n/k$, so $H^n([0,1]^n)\le n^{n/2}$. Thus the Borel measure $H^n$ is nonzero and finite on bounded Borel sets, by covering bounded sets with finitely many unit cubes. It is translation invariant by isometry invariance. Euclidean translation-invariant measure uniqueness \cref{d11:03:euclidean-measure-uniqueness} gives $H^n=b_nm_n$ with finite $b_n>0$; put $c_n=b_n^{-1}$. When $n=1$ the two displayed cube bounds are both one, so $b_1=1$. This proof uses no isodiametric theorem, no gamma-function normalization, and no change from Euclidean to sup metric.
\end{proof}
\source{FOL 11.2 Proposition11.20 and normalization warning; STE Chapter7 §1 Properties6-7prime}

\begin{theorem}[Ambient Euclidean dimension bounds]\family{D11.22.euclidean_dimension}\label{d11:22:euclidean-dimension}
Every $E\subset\R^n$ has dimension at most $n$. Every nonempty open set has dimension $n$ and infinite $H^s$ for $0\le s<n$. A nondegenerate isometrically embedded $k$-cube has dimension $k$ and positive finite $H^k$. Also $H^s(\R^n)=0$ for $s>n$.
\end{theorem}
\begin{proof}
The precise ingredients are \cref{d11:22:euclidean-lebesgue-normalization}, \cref{d11:22:exponent-threshold}, \cref{d11:22:isometry-similarity}. The normalization theorem gives finite $H^n$ on bounded cubes; the threshold theorem gives zero $H^s$ there for $s>n$. A countable cube covering of $\R^n$ gives its zero $H^s$ mass and dimension at most $n$. Every nonempty open set contains a positive-side-length cube, which has positive finite $H^n$ by normalization and scaling. The threshold theorem forces all its lower-exponent masses infinite, proving its dimension and that of the open set. Apply the same argument in $\R^k$ and use isometric invariance for an embedded $k$-cube. Nondegeneracy is needed for its positive $k$-dimensional mass.
\end{proof}
\source{STE Chapter7 §1 observations and §2; FOL 11.2 Proposition11.20}

\begin{example}[Lower-dimensional Hausdorff measures need not be sigma-finite]\family{D11.22.hausdorff_nonsigmafinite}\label{d11:22:hausdorff-nonsigmafinite}
For $0\le s<n$, the Borel measure $H^s$ on $\R^n$ is not sigma-finite.
\end{example}
\begin{proof}
The precise ingredients are \cref{d11:22:exponent-threshold}, \cref{d11:22:euclidean-lebesgue-normalization}. Every finite-$H^s$ Borel set has zero $H^n$ by the threshold theorem, hence zero Lebesgue measure by normalization. A countable union of such sets cannot cover the unit cube of positive Lebesgue measure. Thus neither sigma-finiteness nor compact finiteness may be assumed automatically for a Hausdorff measure.
\end{proof}
\source{STE Chapter7 Exercise1}

\section{The Cantor lower estimate, including exact mass}

\begin{lemma}[The Cantor digit map is onto and has exact Holder constant one]\family{D11.22.cantor_holder_surjection}\label{d11:22:cantor-holder-surjection}
Let $C$ be the middle-thirds Cantor set and $\alpha=\log2/\log3$. Every $x\in C$ has a unique digit sequence $x=\sum_{n\ge1}2a_n3^{-n}$, $a_n\in\{0,1\}$. The map $F(x)=\sum_{n\ge1}a_n2^{-n}$ maps $C$ onto $[0,1]$ and satisfies $|F(y)-F(x)|\le|y-x|^\alpha$.
\end{lemma}
\begin{proof}
The nested construction intervals provide the digit representation. Two distinct allowed digit sequences with first difference at $k$ differ in value by at least $2\cdot3^{-k}-\sum_{j>k}2\cdot3^{-j}=3^{-k}>0$, proving uniqueness. Every binary expansion supplies an allowed ternary sequence, so $F$ is onto (binary nonuniqueness simply gives two different preimages at some dyadic values).

First prove $F(t)\le t^\alpha$ on $C$. The left-branch formulas $t=u/3$, $F(t)=F(u)/2$ preserve that inequality because $3^\alpha=2$. On the right branch $t=(2+u)/3$, $F(t)=(1+F(u))/2$, use
$1+u^\alpha\le(2+u)^\alpha$ for $0\le u\le1$.
Indeed the difference $(2+u)^\alpha-1-u^\alpha$ decreases to zero at $u=1$, as its derivative is negative on $(0,1]$. Induction proves the bound on finite digit expansions, and limits prove it on all of $C$. Reflection of digits gives $1-F(u)=F(1-u)\le(1-u)^\alpha$.

For $x<y$, remove their common digit prefix. Both coordinate and image differences acquire factors $3^{-k}$ and $2^{-k}=(3^{-k})^\alpha$, so it suffices to treat $x=u/3$ and $y=(2+v)/3$. Set $a=1-u$, $b=v$, and $t=a+b\in[0,2]$. The preceding bounds and concavity give
\[F(y)-F(x)\le\frac{a^\alpha+b^\alpha}{2}
\le\frac{2^{1-\alpha}t^\alpha}{2}
\le\frac{(1+t)^\alpha}{2}=|y-x|^\alpha.
\]
For the third inequality take $\alpha$-th roots: $2^{(1-\alpha)/\alpha}t\le1+t$ because $2^{1/\alpha}=3$ and $t\le2$. Equal points need no estimate. This proves the exact Holder constant one, not merely some unspecified constant.
\end{proof}
\source{STE Chapter7 Lemmas2.2-2.3 and Exercise7}

\begin{theorem}[Elementary middle-thirds Cantor dimension]\family{D11.22.cantor_dimension}\label{d11:22:cantor-dimension}
For the middle-thirds set, $0<H^\alpha(C)\le1$, $\dim_HC=\alpha$, and $m_1(C)=0$.
\end{theorem}
\begin{proof}
The precise ingredients are \cref{d11:22:cantor-holder-surjection}, \cref{d11:22:holder-image-bound}, \cref{d11:22:exponent-threshold}. Its level-$k$ construction uses $2^k$ closed intervals of diameter $3^{-k}$, giving total $\alpha$-cost $2^k3^{-k\alpha}=1$. Hence the upper bound. The Holder image inequality and $F(C)=[0,1]$ give $1=H^1([0,1])\le H^\alpha(C)$, in particular positivity. The threshold theorem now yields the dimension. The same level cover has total length $(2/3)^k\to0$, so Lebesgue outer measure of $C$ is zero. Since $C$ is closed it is Lebesgue measurable.
\end{proof}
\source{STE Chapter7 §2.1 Theorem2.1 and Lemmas2.2-2.3}

\begin{theorem}[Exact critical mass of the middle-thirds Cantor set]\family{D11.22.cantor_exact_mass}\label{d11:22:cantor-exact-mass}
In this coefficient-one diameter normalization, $H^\alpha(C)=1$.
\end{theorem}
\begin{proof}
The precise ingredients are \cref{d11:22:cantor-holder-surjection}, \cref{d11:22:cantor-dimension}, \cref{d11:22:euclidean-lebesgue-normalization}. For completeness spell out the lower covering estimate behind the preceding image inequality. Given any countable cover $C\subset\bigcup_jE_j$, the sets $F(E_j\cap C)$ cover $[0,1]$. Each is contained in an interval of length at most $\diam(E_j)^\alpha$ by the exact Holder estimate; an arbitrarily small enlargement is allowed when endpoints are not attained. Countable subadditivity of Lebesgue outer measure gives $1\le\sum_j\diam(E_j)^\alpha$. This uses outer measure and does not assume these arbitrary images are measurable. Take infima over every diameter-limited cover. The construction covers in the preceding proof give the reverse bound one at all small scales. Thus the exact mass is one.
\end{proof}
\source{STE Chapter7 Exercise7}

\begin{example}[Large-scale content is not Hausdorff measure]\family{D11.22.content_not_measure}\label{d11:22:content-not-measure}
The unrestricted one-dimensional diameter content of $[0,1]^2$ is at most $\sqrt2$, whereas its $H^1$ mass is infinity.
\end{example}
\begin{proof}
The precise ingredients are \cref{d11:22:euclidean-dimension}, \cref{d11:22:exponent-threshold}. The square itself is one cover member of diameter $\sqrt2$. Its positive finite $H^2$ mass and the exponent threshold give infinite $H^1$ mass. Hence removing the requirement of arbitrarily small covering diameters changes the quantity, even on a compact Euclidean set.
\end{proof}
\source{FOL 11.2 small-diameter discussion; STE Chapter7 §1 observations}

% END chapters/22


% BEGIN chapters/23
\chapter{The Daniell--Stone construction on a sigma-ring}\label{ch:d11-23}

Here $X$ is an arbitrary set, possibly empty. The functions in the lattice are finite real-valued functions. We use an undergraph construction and the already established one-dimensional measure theory, rather than claiming an integration-first foundational reconstruction. The sigma-ring generated by positive level sets is the exact uniqueness domain; adjoining a unit is an additional hypothesis, not an implicit reduction.

\section{The order hypotheses and the two generated domains}

\begin{definition}[Real vector lattices and Stone's condition]\family{D11.23.vector_lattice_stone}\label{d11:23:vector-lattice-stone}
A real vector lattice $V\subset\R^X$ is a vector subspace closed under pointwise minimum and maximum. It contains $f^+,f^-$ and $|f|$ for $f\in V$. Stone's condition is $f\wedge1\in V$ for every $f\in V$; this does not require $1\in V$.
\end{definition}
\begin{proof}
The positive and negative parts are $f\vee0$ and $(-f)\vee0$, and $|f|=f\vee(-f)$. Scaling Stone's condition shows $f\wedge b=b((f/b)\wedge1)\in V$ for every $b>0$. Thus all truncations used below belong to the original lattice. When $X=\varnothing$, $\R^X$ contains only the zero function and these statements still hold.
\end{proof}
\source{COH 7.7 definitions}

\begin{definition}[Positive order-continuous elementary integral]\family{D11.23.elementary_integral}\label{d11:23:elementary-integral}
An elementary integral is a real-linear $L:V\to\R$ which is positive on nonnegative functions and satisfies $L(f_n)\to0$ whenever $f_n\in V$ decreases pointwise to zero at every point. No norm or a.e. qualification occurs before the measure is built.
\end{definition}
\source{COH 7.7 elementary integral definition}

\begin{lemma}[Monotone and series properties inside the elementary lattice]\family{D11.23.elementary_monotone_properties}\label{d11:23:elementary-monotone-properties}
For nonnegative $f,f_n\in V$: if $f_n\uparrow f$, then $L(f_n)\uparrow L(f)$; if $f=\sum_nf_n$, then $L(f)=\sum_nL(f_n)$; and if $f\le\sum_nf_n$, then $L(f)\le\sum_nL(f_n)$.
\end{lemma}
\begin{proof}
The precise ingredients are \cref{d11:23:elementary-integral}. For the first assertion use order continuity on $f-f_n\downarrow0$ and linearity. Apply it to finite partial sums for the second assertion. For the inequality apply it to $f\wedge\sum_{i\le n}f_i\uparrow f$. Positivity bounds its $L$ value by the corresponding partial sum of $L(f_i)$. Extended infinite right sides are permitted; all individual values of $L$ remain finite.
\end{proof}
\source{COH 7.7 Lemma7.7.1}

\begin{definition}[Generated sigma-ring and generated measurable space]\family{D11.23.generated_ring_algebra}\label{d11:23:generated-ring-algebra}
Let $\mathcal F=\{\{f>b\}:f\in V,\ b>0\}$, let $\mathcal R_V$ be the sigma-ring it generates, and let $\mathcal A_V=\sigma(\mathcal F)$. A sigma-ring contains the empty set and is closed under differences and countable unions; it need not contain $X$. The sigma-algebra $\mathcal A_V$ is the smallest making every member of $V$ measurable.
\end{definition}
\begin{proof}
Existence of the two generated families follows by intersection of all families with the required closure properties. Positive thresholds and countable unions determine $\{f>0\}$ and every positive superlevel of $f$. Apply the same observation to $-f$ and take complements and countable limits to obtain negative thresholds as well. These sets generate the real Borel sigma-algebra, so every $f$ is measurable for $\mathcal A_V$. Conversely any sigma-algebra making every $f$ measurable contains $\mathcal F$.
\end{proof}
\source{COH 7.7 sigma-ring discussion and generated families}

\begin{lemma}[Extend a sigma-ring measure to its generated sigma-algebra]\family{D11.23.sigma_ring_extension}\label{d11:23:sigma-ring-extension}
For any sigma-ring $\mathcal R$ on $X$, $\sigma(\mathcal R)=\mathcal R\cup\{X\setminus A:A\in\mathcal R\}$. If $\mu_0$ is a measure on $\mathcal R$, one extension to this sigma-algebra is
\[\mu(A)=\begin{cases}\mu_0(A),&A\in\mathcal R,\\ \infty,&A\notin\mathcal R.\end{cases}\]
This is not a uniqueness statement outside $\mathcal R$.
\end{lemma}
\begin{proof}
The displayed family is closed under complements. For a countable union, either all members are in $\mathcal R$, or the complement of one member is in $\mathcal R$; in the latter case the complement of the union is obtained from that ring member by countable intersections and differences, so is in $\mathcal R$. Thus it is a sigma-algebra, proving the formula. If disjoint measurable $A_i$ have union in $\mathcal R$, each $A_i$ lies in $\mathcal R$: if its complement is in $\mathcal R$, subtract that complement from the union. Countable additivity then follows from $\mu_0$. If their union is outside $\mathcal R$, not all $A_i$ can belong to the ring. Some $A_i$ is outside and has assigned mass infinity, making both sides infinity. Two disjoint outside members cannot occur unless $X\in\mathcal R$, in which case no outside members exist. The empty set retains mass zero. These checks prove the extension for all cases.
\end{proof}
\source{COH 7.7 sigma-ring properties(a)-(c)}

\section{A content on vertical intervals}

\begin{lemma}[Elementary integral induces content on under-graph intervals]\family{D11.23.undergraph_content}\label{d11:23:undergraph-content}
For $f,g\in V$ let $[f,g)=\{(x,t):f(x)\le t<g(x)\}$. These sets form a semiring $\mathcal I$. For $f\le g$ define $J([f,g))=L(g-f)$. This is well defined and nonnegative; it is countably additive for disjoint unions which remain in $\mathcal I$, and satisfies $J(I)\le\sum_nJ(I_n)$ whenever $I\subset\bigcup_nI_n$.
\end{lemma}
\begin{proof}
The precise ingredients are \cref{d11:23:elementary-monotone-properties}. Replacing $f$ by $f\wedge g$ normalizes the endpoints without changing the interval. Equal intervals have equal pointwise section lengths, including zero for empty sections, so $g-f$ is representation independent after normalization. Intersections have endpoints $f_1\vee f_2$, $g_1\wedge g_2$. For normalized second endpoints, the difference is the disjoint union
\[[f_1,g_1)\setminus[f_2,g_2)=[f_1,g_1\wedge f_2)\ \dot\cup\ [f_1\vee g_2,g_1).
\]
Empty pieces are normalized as above. This proves the semiring closure and finite-difference property.

For each fixed $x$, a disjoint union or a cover of vertical intervals has respectively equal or dominating sums of their lengths, by the already established Lebesgue length on $\R$, \cref{d11:03:lebesgue-measure}. The nonnegative length functions $g-f,g_n-f_n$ belong to $V$. The preceding elementary monotone properties transfer those pointwise equalities or inequalities through $L$. This proves both countable statements, and in particular finite additivity. This use of earlier measure theory is explicit; we are not asserting an independent integration-first development.
\end{proof}
\source{COH 7.7 Lemmas7.7.5-7.7.6}

\begin{construction}[Covering by under-graph intervals gives an exact measure on their sigma-algebra]\family{D11.23.undergraph_outer_measure}\label{d11:23:undergraph-outer-measure}
For arbitrary $A\subset X\times\R$ put
\[\nu^*(A)=\inf\{\sum_nJ(I_n):A\subset\bigcup_nI_n,\ I_n\in\mathcal I\},\]
with value infinity if there is no countable cover. Then $\nu^*$ is an outer measure, every $I\in\mathcal I$ is Caratheodory measurable, and $\nu^*(I)=J(I)$. Its restriction to $\mathcal B=\sigma(\mathcal I)$ is a measure $\nu$.
\end{construction}
\begin{proof}
The precise ingredients are \cref{d11:23:undergraph-content}, \cref{d11:02:caratheodory-theorem}. This covering family need not contain $X\times\R$ or cover it countably. That causes no gap: the empty interval covers the empty set at cost zero; monotonicity follows by restricting the target set. To prove countable subadditivity, if the sum of outer costs is infinite there is nothing to prove. Otherwise each target $A_n$ has a countable cover of cost below $\nu^*(A_n)+\eps2^{-n}$ by the definition of a finite infimum. Adjoin all those covers and let $\eps\downarrow0$. This proves outer-measure subadditivity without a whole-space covering hypothesis.

Fix $I\in\mathcal I$ and test $A$ of finite outer cost. Choose a cover $I_n$ costing less than $\nu^*(A)+\eps$. The finite-difference property splits each $I_n$ into $I_n\cap I$ and at most two semiring pieces in $I_n\setminus I$; finite additivity splits its cost exactly. These pieces cover respectively $A\cap I$ and $A\setminus I$, hence
$\nu^*(A)+\eps\ge\nu^*(A\cap I)+\nu^*(A\setminus I)$.
Let $\eps\downarrow0$; infinite-cost tests are automatic. Thus $I$ is measurable. Its own cover gives $\nu^*(I)\le J(I)$, while the cover inequality of the preceding lemma gives the reverse bound for every countable cover. Caratheodory's theorem now supplies the measure on $\mathcal B$.
\end{proof}
\source{COH 7.7 Lemma7.7.7}

\section{Cylinders recover the required integrals}

\begin{construction}[Stone truncations produce the representing measure and identify its integrals]\family{D11.23.level_cylinder_representation}\label{d11:23:level-cylinder-representation}
For $f\in V$, $b>0$, the functions $u_n=\min(1,n(f-f\wedge b))$ lie in $V$ and increase to $\one_{\{f>b\}}$. For $c>0$, $\{f>b\}\times[0,c)\in\mathcal B$. The formula $\mu_0(A)=\nu(A\times[0,1))$ defines a measure on $\mathcal R_V$; its preceding sigma-ring extension to $\mathcal A_V$ integrates every $f\in V$ and has $\int f\dd\mu=L(f)$.
\end{construction}
\begin{proof}
The precise ingredients are \cref{d11:23:undergraph-outer-measure}, \cref{d11:23:sigma-ring-extension}, \cref{d11:05:monotone-convergence}. The truncation assertion follows from Stone's condition and positivity of $f-f\wedge b$. The sets $[0,cu_n)$ increase to $\{f>b\}\times[0,c)$, proving measurability. The sets $A$ for which $A\times[0,c)$ is in $\mathcal B$ form a sigma-ring, so contain $\mathcal R_V$. For $c=1$, disjoint cylinder unions prove countable additivity of $\mu_0$. Extend it by the sigma-ring lemma.

For a positive level set $E=\{f>b\}$, continuity from below and homogeneity give
\[\nu(E\times[0,c))=\lim_nL(cu_n)=c\lim_nL(u_n)=c\mu_0(E).
\]
All these values are finite: $u_n\le f^+/b$, so $\mu_0(E)\le L(f^+)/b$. For nested positive level sets, subtraction of these finite quantities proves the same cylinder formula for their differences.

Take $f\ge0$ in $V$. Define dyadic lower steps
$s_n=2^{-n}\sum_{k=1}^{n2^n}\one_{\{f>k2^{-n}\}}$.
They are nonnegative simple $\mathcal R_V$-supported functions and increase pointwise to $f$; at exact dyadic levels the strict inequality gives the lower adjacent step, which still increases to the level. Decompose each step into its finitely many disjoint bands between successive positive levels, with the final tail band. The cylinder formula for these differences gives
$\nu([0,s_n))=\int s_n\dd\mu$.
The sets $[0,s_n)$ increase to $[0,f)$, hence measure continuity and monotone convergence give
\[L(f)=J([0,f))=\nu([0,f))=\lim_n\nu([0,s_n))=\int f\dd\mu.
\]
This is finite, since $L$ is real-valued. Apply it to $f^+$ and $f^-$, both in $V$, to conclude $|f|$ is integrable and $\int f=L(f)$ for arbitrary $f$. The proof has not required the unit function or finite total mass.
\end{proof}
\source{COH 7.7 proof of Theorem7.7.4}

\begin{lemma}[Positive level sets localize uniqueness to the generated sigma-ring]\family{D11.23.localized_sigma_ring_uniqueness}\label{d11:23:localized-sigma-ring-uniqueness}
Any two measures on $\mathcal A_V$ representing $L$ agree on every set of $\mathcal R_V$.
\end{lemma}
\begin{proof}
The precise ingredients are \cref{d11:23:elementary-monotone-properties}, \cref{d11:05:monotone-convergence}. For a finite nonempty intersection $E=\bigcap_{i=1}^k\{f_i>b_i\}$, put
$g_n=\min_i\min(1,n(f_i-f_i\wedge b_i))$.
These belong to $V$, are nonnegative, and increase to $\one_E$. Monotone convergence makes both representing masses equal to $\lim_nL(g_n)$. Each positive level set has finite mass, since $\one_{\{f>b\}}\le f^+/b$ and $f^+$ is integrable. Thus the intersections have finite masses as well.

Fix one positive level set $H$. The finite intersections with $H$ form a pi-system containing $H$, on which the two finite restricted measures agree. The pi-lambda uniqueness theorem gives equality on its generated sigma-algebra on $H$. This sigma-algebra is exactly $\{A\cap H:A\in\mathcal R_V\}$: intersections and differences within $H$ stay in the ring, and conversely the class of $A$ whose trace lies in the generated sigma-algebra is a sigma-ring containing every positive level generator.

Every $A\in\mathcal R_V$ is contained in a countable union of positive level sets. Indeed the sets admitting such a cover form a sigma-ring containing the generators. Disjointify a covering $H_n$ by $B_n=H_n\setminus\bigcup_{j<n}H_j$. The sets $A\cap B_n$ are ring sets contained in $H_n$, so their masses agree by the localized finite result. Countable additivity gives equality on $A$. No uncountable sum or whole-space finite uniqueness argument has been substituted.
\end{proof}
\source{COH 7.7 uniqueness proof of7.7.4}

\begin{theorem}[Daniell-Stone existence and sigma-ring uniqueness]\family{D11.23.daniell_stone_representation}\label{d11:23:daniell-stone-representation}
For an elementary integral on a Stone vector lattice $V$ on any set $X$, there is a positive measure on $(X,\mathcal A_V)$ integrating all $f\in V$ with $L(f)=\int f\dd\mu$. Its restriction to $\mathcal R_V$ is unique among all representing measures; uniqueness on all of $\mathcal A_V$ requires a further hypothesis.
\end{theorem}
\begin{proof}
The precise ingredients are \cref{d11:23:level-cylinder-representation}, \cref{d11:23:localized-sigma-ring-uniqueness}. The undergraph outer-measure construction and level-cylinder construction give existence, finite absolute integrability, and the exact identity. The localized uniqueness lemma gives the stated uniqueness domain. If $X$ is empty all sets and functions used are empty or zero, and the zero measure realizes the same conclusion. For nonempty $X$, the construction still need not make $\mu(X)$ finite or determine it from $L$; that boundary is essential, as the counterexample below shows.
\end{proof}
\source{COH 7.7 Theorem7.7.4 and complete proof}

\begin{proposition}[Measure representation forces order continuity on integrable elementary functions]\family{D11.23.order_continuity_necessary}\label{d11:23:order-continuity-necessary}
If a measure makes all $f\in V$ integrable and $L(f)=\int f\dd\mu$, then $L$ is positive linear and order-continuous above at zero, hence an elementary integral.
\end{proposition}
\begin{proof}
Integral positivity and linearity give the first assertions. If $f_n\downarrow0$ pointwise, all $f_n$ are nonnegative and bounded by the integrable $f_1$. Dominated convergence \cref{d11:06:dominated-convergence} gives $L(f_n)\to0$. There is no appeal here to a norm continuity hypothesis in place of pointwise order continuity.
\end{proof}
\source{COH 7.7 definitions and Examples7.7.3}

\begin{corollary}[A constant unit yields finite total mass and full uniqueness]\family{D11.23.unit_implies_finite_unique}\label{d11:23:unit-implies-finite-unique}
If $1\in V$, then $X=\{1>1/2\}\in\mathcal R_V$, so $\mathcal R_V=\mathcal A_V$. The representing measure is unique and finite, with $\mu(X)=L(1)$; when $L(1)=1$ it is a probability. More generally $X\in\mathcal R_V$ suffices for full uniqueness.
\end{corollary}
\begin{proof}
The precise ingredients are \cref{d11:23:daniell-stone-representation}. A sigma-ring containing $X$ is a sigma-algebra, which proves equality of the generated domains. Apply sigma-ring uniqueness and integrate the unit for its mass. The broader full-uniqueness assertion uses only equality of the domains, not finiteness.
\end{proof}
\source{COH 7.7 Theorem7.7.4 and generated-family discussion}

\section{Examples, bridges, and two sharp failures}

\begin{theorem}[Compact and compact-support continuous-function examples]\family{D11.23.dini_elementary_examples}\label{d11:23:dini-elementary-examples}
On a compact Hausdorff space, nonnegative continuous $f_n\downarrow0$ converge uniformly to zero. Thus the classical Riemann integral on $C([a,b],\R)$ is elementary. On LCH $X$, $C_c(X,\R)$ is a Stone vector lattice and every positive linear functional on it is elementary.
\end{theorem}
\begin{proof}
For $\eps>0$ and each $x$ choose $n_x$ and an open neighborhood $U_x$ with $f_{n_x}<\eps$ on $U_x$. A finite subcover exists by compactness. For $N=\max n_x$ over that finite cover, monotonicity gives $0\le f_n<\eps$ everywhere for $n\ge N$. This proves Dini's assertion directly. A Riemann integral obeys $|L(f_n)|\le(b-a)\|f_n\|_\infty$, so it is order-continuous; positivity and linearity are calculus properties. On $C_c$, min, max, and Stone truncation preserve continuity and support. For a sequence decreasing to zero, every support lies in $\supp f_1$; apply the compact Dini argument there and the compact-support bound \cref{d11:19:positive-functional-bound}. This proves elementary order continuity without first invoking measure representation.
\end{proof}
\source{COH 7.7 Lemma7.7.2 and Examples7.7.3(a)-(b)}

\begin{corollary}[An integral-first construction agrees with Lebesgue integration on its initial lattice]\family{D11.23.riemann_lebesgue_bridge}\label{d11:23:riemann-lebesgue-bridge}
Daniell--Stone for the Riemann functional on $C([a,b],\R)$ yields the usual finite Borel Lebesgue measure on the interval. Their ordinary completions, and hence all integrals on their common completed space, agree.
\end{corollary}
\begin{proof}
The precise ingredients are \cref{d11:23:dini-elementary-examples}, \cref{d11:23:daniell-stone-representation}, \cref{d11:23:unit-implies-finite-unique}, \cref{d11:16:proper-integral-equality}. Positive level generators of continuous functions generate the interval Borel sigma-algebra (use the continuous distance cutoffs for the countable interval base). The unit belongs to the lattice, so the represented measure is finite and unique. Independently the usual interval Lebesgue measure integrates each continuous function to its Riemann integral, by the already established proper Riemann--Lebesgue bridge. Compact-interval representation \cref{d11:19:stein-exercise9-interface} identifies the two Borel measures. Completion adjoins precisely the subsets of their common Borel null sets, so the completed measures and their simple approximations, hence integrals, agree.
\end{proof}
\source{COH 7.7 Examples7.7.3(a), Theorem7.7.4; STE Chapter6 Exercise9}

\begin{example}[A positive functional need not be an elementary integral]\family{D11.23.positivity_without_order_continuity}\label{d11:23:positivity-without-order-continuity}
Let $V$ consist of continuous real functions on $[0,\infty)$ with a finite limit at infinity and set $L(f)=\lim_{x\to\infty}f(x)$. It is positive linear on a Stone vector lattice but not elementary and has no measure representation integrating all its functions.
\end{example}
\begin{proof}
The precise ingredients are \cref{d11:23:order-continuity-necessary}. Finite limits are preserved by real linear combinations, minima, maxima, and truncation by one, proving the lattice and Stone assertions. The limit functional is positive linear. Define $f_n$ to be zero before $n$, equal to $x-n$ on $[n,n+1)$, and one afterward. They are continuous, decrease pointwise to zero, and have limit at infinity one. Thus $L(f_n)=1$ for every $n$, contradicting order continuity. Necessity proved above rules out the asserted representation. Positivity alone is therefore insufficient.
\end{proof}
\source{COH 7.7 Example7.7.3(e)}

\begin{example}[Representing measures can differ outside the generated sigma-ring]\family{D11.23.sigma_algebra_nonuniqueness}\label{d11:23:sigma-algebra-nonuniqueness}
For $X=\R$, let $V$ consist of $C_c$ functions whose supports lie in $(0,\infty)$ and let $L$ be their Riemann integral. The measures $\mu_0(E)=m(E\cap(0,\infty))$ and $\mu_1(E)=\mu_0(E)+\delta_0(E)$ on $\mathcal A_V$ both represent $L$ but differ on $(-\infty,0]$.
\end{example}
\begin{proof}
The precise ingredients are \cref{d11:23:dini-elementary-examples}, \cref{d11:23:generated-ring-algebra}. The lattice is Stone and $L$ is elementary by the preceding compact-support example. Each generator is a Borel subset of $(0,\infty)$, so $\mathcal A_V$ consists of Borel sets and both measures restrict to it. Countably many positive compact bumps cover $(0,\infty)$ by positive level sets (take bumps equal to one on $[1/n,n]$); thus this half-line lies in $\mathcal R_V$, and its complement lies in $\mathcal A_V$. Every ring set is contained in $(0,\infty)$, since that containment is preserved by differences and countable unions. Consequently $(-\infty,0]$ lies outside the ring. All $f\in V$ vanish at zero and on the negative half-line, so both measures give the same integral $L(f)$. Their complement masses are zero and one. This proves actual nonuniqueness on the generated sigma-algebra while respecting uniqueness on its generated sigma-ring. The canonical extension constructed earlier may assign infinity there; it is one allowed extension, not a forced value.
\end{proof}
\source{COH 7.7 Exercise2}

% END chapters/23

\backmatter

% BEGIN source-register
\chapter*{Sources, comparison scope and cautions}

\addcontentsline{toc}{chapter}{Sources, comparison scope and cautions}

Book keys refer to the supplied, archived editions. Locations in the companion result ledger are physical PDF pages, not automatically printed book pages. Source comparisons identify adopted formulations and variants; they are not substitutes for the proofs in this manuscript.

\begin{description}[style=nextline,leftmargin=13mm]

\item[AXL] Sheldon Axler; Measure, Integration \& Real Analysis; 2020 open-access edition; supplied author PDF corrected 12 June 2026.

\item[ROY] Halsey L. Royden, Patrick M. Fitzpatrick; Real Analysis; Fifth edition, Pearson, copyright 2023.

\item[FOL] Gerald B. Folland; Real Analysis: Modern Techniques and Their Applications; Second edition, Wiley, 1999; reflowed electronic PDF.

\item[STE] Elias M. Stein, Rami Shakarchi; Real Analysis: Measure Theory, Integration, and Hilbert Spaces; Princeton Lectures in Analysis III; 2005 work, supplied electronic printing.

\item[COH] Donald L. Cohn; Measure Theory; Second edition, Birkhauser/Springer, 2013.

\end{description}

The full source-by-chapter disposition is retained in SOURCE\_CHAPTER\_MATRIX.json, and the precise adopted statement and physical-page mapping is retained in UNIFIED\_RESULT\_LEDGER.json. Those navigation records do not assert that every source proves every variant in a topic family.

\section*{Genuine differences retained}

Ordinary almost-everywhere and local-null quotients, compact-inner regularity on open sets and on all Borel sets, and sigma-ring and sigma-algebra uniqueness are distinguished in the relevant chapters. A finite target measure alone does not supply a Radon--Nikodym density with respect to an arbitrary positive reference measure. The stronger valid arbitrary-reference decomposition has a separate proof.

The bounded-support metric-space function class in Stein--Shakarchi, Chapter 6, Problem 5, is not identified with compact support or vanishing at infinity on a general metric space. The explicit bounded discrete counterexample and the adopted corrected representation framework appear in the representation chapter. This caution is not an alteration to an archived textbook PDF.

\section*{Scope routing}

General weak compactness, operator and spectral theory, interpolation, Fourier and distribution theory, probability and ergodic theory, and advanced geometric measure theory are routed to their own directions. Minimal Hilbert representation, scalar duality, Haar measure, Hausdorff dimension and Daniell--Stone are included only in the precise selected formulations proved here. No entire book is claimed to have been transcribed or formalized.

% END source-register


% BEGIN verification-register
\chapter*{Verification status and later formalization}
\addcontentsline{toc}{chapter}{Verification status and later formalization}
Independent cross-audits and mathematical integration checks are complete. The accompanying audit record binds their findings and resolutions to the frozen mathematical sources. This is manuscript acceptance; later Lean acceptance is separate.

Five isolated reuse statements have passed native Lean compilation, protected expected-statement comparison and kernel checking against Lean 4.35.0, the pinned Mathlib revision and MathGenome v0.5.0. Further declaration-type checks identify candidate interfaces. They are not a claim that the whole manuscript is already formalized.

Every independently adopted mathematical result must retain its actual hypotheses and conclusion under normal formalization. Equivalent proofs and correctly matched existing library results are permitted. The steps of one original proof route do not become additional independent acceptance obligations. The source/result index is navigation, and does not replace independent reading of the mathematical source.

% END verification-register

\end{document}
